%% =========================================================
%% Chapter: PEPS Fundamental Theorem
%% =========================================================

\chapter{PEPS Fundamental Theorem}%
\label{ch:peps_ft}

A \emph{projected entangled pair state} (PEPS) is built from a finite simple
graph $G=(V,E)$ by placing a local tensor $A_v$ at every vertex $v$ and
contracting the tensors along the edges. Each vertex carries one physical index
$\sigma(v)\in\{0,\ldots,d-1\}$, and each edge $e$ carries a virtual index of
bond dimension $D_e$ shared by its two endpoints. On a square lattice this is
the picture
% Formula: |Psi_A> = sum_sigma c_A(sigma)|sigma>, specialized here to a
% finite 3-by-3 square graph.  Ink-to-index: each dot is A_v, each thin
% in-plane segment contracts one shared virtual edge index, and each north
% stub is the physical index sigma(v).
% Contraction and boundary signature: all 12 graph edges are contracted; no
% virtual boundary leg is present and the nine physical legs remain open, so
% (V,P)=(0,9).  The central label identifies one representative local tensor.
% Source verdict: exact finite square-graph specialization of
% Papers/1804.04964/paper_normal.tex:961--979; the source theorem itself
% applies to arbitrary finite simple graphs.
\begin{center}
    \begin{tenkzlattice}[rows=3, cols=3, physical=up]
        \tnsite[label=$A_v$]{(2,2)}
    \end{tenkzlattice}
\end{center}
in which every black dot is a local tensor, the thin horizontal and vertical
bonds are the virtual edges, and the short thick stub at each vertex is its
physical leg.

The state generated by $A$ is the tensor whose coefficient at a physical
configuration $\sigma\colon V\to\{0,\ldots,d-1\}$ is obtained by summing the
product of the local tensors over all virtual configurations:
\[
    c_A(\sigma)=\sum_{\eta}\ \prod_{v\in V} A_v(\eta_v,\sigma(v)),
\]
where $\eta$ ranges over the assignments of a bond index $\eta(e)\in\{0,\ldots,
D_e-1\}$ to every edge and $\eta_v$ is the family of indices that $\eta$ places
on the edges incident to $v$. Closing every virtual bond while keeping the
physical legs open gives the full state tensor; fixing the legs to a
configuration $\sigma$ evaluates the scalar $c_A(\sigma)$. On the five-vertex
example graph this is
% Formula: c_A(sigma)=sum_eta prod_{i=1}^5 A_i(eta_i,sigma_i), for the
% source graph with cycle 1--2--3--4--5--1 and chord 2--5.
% Ink-to-index: the six thin joins are precisely the six summed virtual edge
% indices, while the five north legs are sigma_1,...,sigma_5.
% Contraction and boundary signature: every virtual index is contracted and
% the five physical indices remain visible, hence (V,P)=(0,5); fixing their
% labels evaluates the displayed coefficient.
% Source verdict: exact graph and contraction of
% Papers/1804.04964/paper_normal.tex:961--979, redrawn in house notation.
\begin{center}
    \begin{tenkzfree}
        \tnput[ports={0:virtual,110:virtual,north:physical}]
            {stateAone}{(-10mm,-5mm)}{A_1}
        \tnput[ports={180:virtual,55:virtual,141:virtual,north:physical}]
            {stateAtwo}{(8mm,-5mm)}{A_2}
        \tnput[ports={235:virtual,155:virtual,north:physical}]
            {stateAthree}{(16mm,6mm)}{A_3}
        \tnput[ports={-25:virtual,190:virtual,north:physical}]
            {stateAfour}{(4mm,15mm)}{A_4}
        \tnput[ports={10:virtual,-70:virtual,-39:virtual,north:physical}]
            {stateAfive}{(-10mm,10mm)}{A_5}
        \tnput[boundary, ports={south:physical}]{statePone}{(-10mm,3mm)}{}
        \tnput[boundary, ports={south:physical}]{statePtwo}{(8mm,3mm)}{}
        \tnput[boundary, ports={south:physical}]{statePthree}{(16mm,14mm)}{}
        \tnput[boundary, ports={south:physical}]{statePfour}{(4mm,23mm)}{}
        \tnput[boundary, ports={south:physical}]{statePfive}{(-10mm,18mm)}{}
        \tnjoin{stateAone.0}{stateAtwo.180}
        \tnjoin{stateAtwo.55}{stateAthree.235}
        \tnjoin{stateAthree.155}{stateAfour.-25}
        \tnjoin{stateAfour.190}{stateAfive.10}
        \tnjoin{stateAfive.-70}{stateAone.110}
        \tnjoin{stateAtwo.141}{stateAfive.-39}
        \tnjoin{stateAone.north}{statePone.south}
        \tnjoin{stateAtwo.north}{statePtwo.south}
        \tnjoin{stateAthree.north}{statePthree.south}
        \tnjoin{stateAfour.north}{statePfour.south}
        \tnjoin{stateAfive.north}{statePfive.south}
    \end{tenkzfree}
\end{center}

A PEPS is \emph{injective} when at every vertex the physical vectors
$A_v(\eta_v,\cdot)\in\C^d$, indexed by the virtual configurations $\eta_v$ on
the incident edges, are linearly independent; equivalently, the linear map they
define from the coefficient space on virtual configurations into $\C^d$ has
trivial kernel:
% Formula: A_v : tensor_{e incident to v} C^{D_e} -> C^d with ker(A_v)=0.
% Ink-to-index: the single west wire is the bundled family of all virtual
% indices incident to v, and the east wire is the physical output.
% Contraction and boundary signature: the map contracts no index internally;
% it has one bundled virtual input and one physical output, (V,P)=(1,1).
% Source verdict: exact arbitrary-degree map formulation of
% Papers/1804.04964/paper_normal.tex:961--981, not a square-degree-four star.
\begin{center}
    \begin{tenkzfree}
        \tnput[boundary, ports={east:virtual}, label pos=south]
            {introInjectiveDomain}{(-18mm,0)}
            {\displaystyle\bigotimes_{e\ni v}\C^{D_e}}
        \tnput[box, ports={west:virtual,east:physical}]
            {introInjectiveMap}{(0,0)}{A_v}
        \tnput[boundary, ports={west:physical}, label pos=south]
            {introInjectiveCodomain}{(18mm,0)}{\C^d}
        \tnjoin{introInjectiveDomain.east}{introInjectiveMap.west}
        \tnjoin{introInjectiveMap.east}{introInjectiveCodomain.west}
    \end{tenkzfree}
    \qquad $\ker A_v=\{0\}$.
\end{center}
It is \emph{normal} when blocking the tensors over suitable regions makes the
blocked tensors injective, following
\cite[Section~4, lines~1316--1318]{Molnar2018NormalPEPS}.

Two PEPS that generate the same state need not have equal tensors. The virtual
indices are internal, so on any edge $e$ one may insert an invertible matrix
$X_e$ on one endpoint and its inverse on the other without changing any
coefficient. Such a \emph{gauge} on the bonds is the only freedom, and the two
main results of the chapter make this precise.
\begin{itemize}
    \item For injective PEPS on a connected graph with positive bond
        dimensions, two tensors generate the same state if and only if they are
        related by a gauge on the edges, with the gauges unique modulo the
        balanced edge scalars of Section~\ref{sec:peps_balanced_edge_scalars}.
        A gauge always gives the same state
        (Theorem~\ref{thm:GaugeEquiv_sameState}); the converse is the
        positive-bond case of \cite[Theorem~2]{Molnar2018NormalPEPS}
        (Theorem~\ref{thm:fundamentalTheorem_PEPS}).
    \item For a normal PEPS on a connected graph the same conclusion holds once
        the tensors become injective after blocking suitable regions
        (Theorem~\ref{thm:fundamentalTheorem_normalPEPS_connected}; the general
        normal-PEPS theorem in \cite[Section~4]{Molnar2018NormalPEPS}, following
        its Theorem~3).
\end{itemize}

In the citations below, \cite[Theorem~3]{Molnar2018NormalPEPS} denotes only
the translationally invariant square-lattice theorem. The connected-graph
normal theorem and the normal-MPS corollaries are cited as results of
\cite[Section~4]{Molnar2018NormalPEPS}.

The chapter builds these results in stages.
Section~\ref{sec:peps_foundations} fixes the tensors, the generated state, the
bond gauges, and vertex injectivity, together with the local inverse used
throughout.
Section~\ref{sec:peps_edge_middle} reduces a single edge to a three-site
contraction by blocking its complementary middle region.
Section~\ref{sec:peps_edge_kernel_gauges} turns a virtual insertion on an edge
into a physical operation at an endpoint and extracts the per-edge gauge.
Section~\ref{sec:peps_ft} transports the edge gauges to all regions and proves
the Fundamental Theorem for injective PEPS.
Section~\ref{sec:peps_normal_union} shows that injective regions are closed
under unions, the step that lifts local injectivity to the larger regions a
normal PEPS requires.
Sections~\ref{sec:peps_normal_square} and~\ref{sec:peps_torus} supply the
square-lattice and torus geometries.
Section~\ref{sec:peps_normal_capstone} proves the Fundamental Theorem for
normal PEPS, Section~\ref{sec:peps_normal_mps} specializes it to matrix product
states on the closed chain, and Section~\ref{sec:peps_balanced_edge_scalars}
records the scalar freedom that remains in the edge gauges.

\input{chapter/ch24_peps_ft_foundations}
\input{chapter/ch24_peps_ft_edge_middle}
\section{Edge Kernels and Gauge Transport}\label{sec:peps_edge_kernel_gauges}
When the endpoint tensors are injective, a matrix inserted on the bond $e$ is
realized by a physical operation on a neighbouring physical leg, and conversely.
This sets up an isomorphism between the bond-$e$ matrix algebra of $A$ and that
of $B$, which by the Skolem--Noether theorem is conjugation by an invertible
matrix $X_e$, the edge gauge: for every physical configuration $\sigma$ and bond
matrix $M$,
$C_A(e,\sigma,M)=C_B(e,\sigma,X_eMX_e^{-1})$.
Collecting the $X_e$ over all edges and reconciling them on the shared bonds
produces the global gauge family that relates the two tensors.

\begin{definition}[Finite kernel descent datum]%
    \label{def:peps_finiteRegionKernelDescent}
    \lean{TNLean.PEPS.FiniteRegionKernelDescent}
    \leanok
    Let $V$ be a vertex type with decidable equality. A finite kernel descent
    datum is a predicate $K$ on finite subsets $S\subseteq V$ together with
    implications
    $j\in S\Longrightarrow K(S)\Longrightarrow K(S\setminus\{j\})$.
    In the edge-blocking proof, $K(S)$ is the assertion that the boundary
    coefficients $c_{\rho}$, with the virtual indices exposed at that stage,
    give zero after the tensors in $S$ have been contracted.
\end{definition}

\begin{lemma}[Finite descent of kernel conditions]%
    \label{lem:peps_finiteRegionKernelDescent}
    \lean{TNLean.PEPS.FiniteRegionKernelDescent.descend_to_empty}
    \lean{TNLean.PEPS.FiniteRegionKernelDescent.descend}
    \leanok
    \uses{def:peps_finiteRegionKernelDescent}
    For a finite kernel descent datum and any finite $R\subseteq V$,
    $K(R)\Longrightarrow K(\varnothing)$. More generally, if
    $K(\varnothing)\Longrightarrow Q$, then $K(R)\Longrightarrow Q$.
\end{lemma}
\begin{proof}\leanok
    Induct on the finite set $R$. If $R=\varnothing$, there is nothing to
    prove. Otherwise write $R=S\cup\{j\}$ with $j\notin S$. The deletion
    implication gives $K(R)\Longrightarrow K(S)$, and the induction hypothesis
    gives $K(S)\Longrightarrow K(\varnothing)$.
\end{proof}

\begin{definition}[Edge-middle kernel descent data]%
    \label{def:peps_edgeMiddleKernelDescentData}
    \lean{TNLean.PEPS.EdgeMiddleBoundaryLabel,
        TNLean.PEPS.EdgeMiddleKernelDescentData}
    \leanok
    \uses{def:peps_edgeMiddleTensorFamily,
        def:peps_finiteRegionKernelDescent}
    Let $e=(u,v)$ and $M_e=V\setminus\{u,v\}$. An edge-middle kernel descent
    datum is indexed by boundary labels $\rho$ consisting of the two residual
    endpoint configurations. It assigns to every finitely supported coefficient
    family $c=(c_{\rho})_{\rho}$ a finite kernel descent datum $K_c(S)$ such that
    \begin{align}
        \left(\forall\tau,\
            \sum_{\rho}c_{\rho}T_{A,M_e}^{\rho}(\tau)=0\right)
        &\Longrightarrow K_c(M_e),
        \notag\\
        K_c(\varnothing)
        &\Longrightarrow \forall\rho,\ c_{\rho}=0.
        \notag
    \end{align}
\end{definition}

\begin{theorem}[Kernel descent gives edge-middle injectivity]%
    \label{thm:peps_edgeMiddleTensorInjective_of_kernelDescent}
    \lean{TNLean.PEPS.EdgeMiddleKernelDescentData.edgeMiddleTensorInjective,
        TNLean.PEPS.IsVertexInjective.edgeBlockedThreeSiteInjective_of_kernelDescent}
    \leanok
    \uses{def:peps_edgeMiddleKernelDescentData,
        lem:peps_finiteRegionKernelDescent,
        thm:peps_edgeBlockedThreeSiteInjective_of_middle}
    Let $e=(u,v)$ and $M_e=V\setminus\{u,v\}$. If an edge-middle kernel descent
    datum is given, then $T_{A,M_e}$ is injective. Consequently, if $A$ is
    vertex-injective, the edge-blocked three-site chain at $e$ is injective.
\end{theorem}
\begin{proof}\leanok
    Let $c=(c_{\rho})_{\rho}$ be a finite linear relation among the middle
    tensor vectors:
    $\sum_{\rho}c_{\rho}T_{A,M_e}^{\rho}(\tau)=0$ for every $\tau$.
    The initial part of the descent datum gives $K_c(M_e)$. The finite descent
    lemma gives $K_c(M_e)\Longrightarrow K_c(\varnothing)$, and the terminal
    part gives $\forall\rho,\ c_{\rho}=0$. Thus every finite relation is
    trivial, which is $\inj(T_{A,M_e})$. Vertex injectivity at
    $u$ and $v$ gives $\inj(T_{A,u})$ and
    $\inj(T_{A,v})$; the edge-blocked three-site injectivity at
    $e$ follows.
\end{proof}

\begin{theorem}[Edge blocking preserves injectivity]%
    \label{thm:peps_edgeBlockedThreeSiteInjective}
    \lean{TNLean.PEPS.IsVertexInjective.edgeBlockedThreeSiteInjective}
    \leanok
    \uses{def:IsVertexInjective, def:peps_edgeBlockedThreeSiteInjective,
        thm:peps_edgeMiddleTensorInjective_of_kernelDescent,
        lem:peps_localCoeff_eq_zero_of_contract_zero,
        lem:peps_not_edgeMiddleTensorInjective_of_isEmpty}
    For a tensor map or blocked tensor family $T$, write
    $\inj(T)$ for the assertion that $T$ is injective. If $A$
    is vertex-injective and every bond dimension is positive, then for every
    edge $e=(u,v)$,
    $\inj(T_{A,u})$, $\inj(T_{A,V\setminus\{u,v\}})$, and
    $\inj(T_{A,v})$.
    Here $T_{A,u}$ and $T_{A,v}$ are the two endpoint tensor maps, and
    $T_{A,V\setminus\{u,v\}}$ is the blocked middle tensor family. This is the
    precise assertion after the edge-blocking diagram in
    \cite[Section~3]{Molnar2018NormalPEPS}. The positive-bond hypothesis is the
    source's standing assumption that injective PEPS have nonzero virtual bond
    spaces. Without it, an empty complement configuration space can make the
    middle tensor vanish identically;
    Lemmas~\ref{lem:peps_edgeMiddleTensorFamily_eq_zero_of_isEmpty}
    and~\ref{lem:peps_not_edgeMiddleTensorInjective_of_isEmpty} record
    this obstruction.
\end{theorem}
\begin{proof}\leanok
    \uses{thm:peps_edgeMiddleTensorInjective_of_kernelDescent,
        lem:peps_localCoeff_eq_zero_of_contract_zero}
    Fix $e=(u,v)$ and put $M_e=V\setminus\{u,v\}$. From vertex injectivity,
    construct the edge-middle descent datum required by
    Theorem~\ref{thm:peps_edgeMiddleTensorInjective_of_kernelDescent}: for each
    finite relation
    $\sum_{\rho}c_{\rho}T_{A,M_e}^{\rho}(\tau)=0$ for every $\tau$,
    the associated conditions $K_c(S)$ satisfy $K_c(M_e)$,
    $j\in S\Longrightarrow K_c(S)\Longrightarrow K_c(S\setminus\{j\})$, and
    $K_c(\varnothing)\Longrightarrow \forall\rho,\ c_{\rho}=0$. The descent
    implication is the local contraction identity: for $j\in S$, vertex
    injectivity at $j$ gives a left inverse $L_j$, and applying $L_j$ to the
    relation contracts away that tensor,
    $j\in S,\ K_c(S)\mathrel{\stackrel{L_j}{\Longrightarrow}}
    K_c(S\setminus\{j\})$.
    The contraction at $j$ uses the factorization
    \begin{align}
        \mc C_e
        &\cong
        \mc V_j
        \times
        \prod_{\substack{f\ne e\\ f\not\ni j}}
        \{0,\ldots,\dim A_f-1\},
        \notag
    \end{align}
    where $\mc C_e$ is the complement virtual-configuration space for
    the distinguished edge $e$ and $\mc V_j$ is the local virtual
    configuration space at $j$. The terminal step
    $K_c(\varnothing)\Longrightarrow \forall\rho,\ c_{\rho}=0$ isolates each
    boundary coefficient by surjectivity of the boundary labelling, which is
    where the positive bond dimensions are used to fill the interior virtual
    indices of a witness configuration.
    The kernel descent theorem gives $\inj(T_{A,M_e})$. The same
    theorem, together with vertex injectivity at $u$ and $v$, gives
    $\inj(T_{A,u})$, $\inj(T_{A,M_e})$, and
    $\inj(T_{A,v})$ as the edge-blocked three-site injectivity
    assertion.
\end{proof}

\begin{theorem}[Local virtual insertion realized physically on either endpoint]%
    \label{thm:peps_virtualInsertionPhysicalRealization}
    \lean{TNLean.PEPS.edgeVirtualInsertionPhysicalRealization}
    \leanok
    \uses{def:IsVertexInjective, def:peps_localIncidentMatrixOp,
        thm:peps_localIncidentMatrixOp_physicalRealization}
    Let $A$ be vertex-injective and let $e=(u,v)$ be an edge. For every
    matrix $M$ on the bond space of $e$, the one-leg virtual action induced
    by $M$ on the local tensor at $u$, and likewise at $v$, is realized by
    a physical operator on that endpoint's physical space.
    On a closed three-site chain, the corresponding state equality is
    % Formula (paper_normal.tex:303--352): inserting M on the (1,2) bond equals
    % applying O_1 to the physical leg of A_1 and also equals applying O_2 to the
    % physical leg of A_2.  Ink-to-index: the M panel has four virtual joins,
    % while each O panel has the three virtual joins of the closed chain; their
    % three north legs are the physical indices, with O_1 or O_2 inserted on the
    % indicated physical leg.
    % Boundary signature: every panel has (V,P)=(0,3).  This is the exact
    % closed-chain consequence drawn in Papers/1804.04964/paper_normal.tex:303--330,
    % not the boundary signature of the theorem's general local tensor.
    \begin{center}
        \tenkzkernel
        \begin{tenkz}[rows={ket}, cols=3, physical=up, boundary=periodic]
            \tn{A_1} & \tn{A_2} & \tn{A_3}
            \tn[skin=ring, species=operator, at=on bond-1-1-1-2 0.5]{M}
        \end{tenkz}
        $ = $
        \begin{tenkz}[rows={ket}, cols=3, physical=up, boundary=periodic]
            \tn[name=opSite]{A_1} & \tn{A_2} & \tn{A_3}
            \tn[at=1 n of opSite, skin=ring, species=operator,
                name=opRing, ports={270:physical, 90:physical}]{O_1}
            \tnwire{opSite.90}{opRing.270}
        \end{tenkz}
        $ = $
        \begin{tenkz}[rows={ket}, cols=3, physical=up, boundary=periodic]
            \tn{A_1} & \tn[name=opSite]{A_2} & \tn{A_3}
            \tn[at=1 n of opSite, skin=ring, species=operator,
                name=opRing, ports={270:physical, 90:physical}]{O_2}
            \tnwire{opSite.90}{opRing.270}
        \end{tenkz}
    \end{center}

    % This is the local endpoint part of Molnar2018NormalPEPS, Section 3. The corresponding
    % coefficient identity is
    % thm:peps_edgeInsertedCoeff_endpointPhysicalRealization.
\end{theorem}
\begin{proof}\leanok
    Apply the one-edge realization theorem to the two endpoint incidences of
    the distinguished edge.
\end{proof}

\begin{theorem}[Left-endpoint projected recovery]%
    \label{thm:peps_edgeLeftLocalVirtualOpOfPhysicalOp_eq_iff_projected_realization_eq}
    \lean{TNLean.PEPS.edgeLeftLocalVirtualOpOfPhysicalOp_eq_iff_projected_realization_eq}
    \leanok
    \uses{thm:peps_localVirtualOpOfPhysicalOp_eq_iff_projected_realization_eq,
        def:peps_localIncidentMatrixOp}
    Let $e=(u,v)$. A physical operator at the left endpoint pulls back to
    the one-edge virtual action of $M^{\mathsf T}$ if and only if its
    compressed physical action is the canonical realization of that
    one-edge virtual action.
\end{theorem}
\begin{proof}\leanok
    This is the projected local recovery equivalence applied to the left
    incident edge of $e$, with the transpose convention used by the
    left-endpoint coefficient identity.
\end{proof}

\begin{theorem}[Right-endpoint projected recovery]%
    \label{thm:peps_edgeRightLocalVirtualOpOfPhysicalOp_eq_iff_projected_realization_eq}
    \lean{TNLean.PEPS.edgeRightLocalVirtualOpOfPhysicalOp_eq_iff_projected_realization_eq}
    \leanok
    \uses{thm:peps_localVirtualOpOfPhysicalOp_eq_iff_projected_realization_eq,
        def:peps_localIncidentMatrixOp}
    Let $e=(u,v)$. A physical operator at the right endpoint pulls back to
    the one-edge virtual action of $M$ if and only if its compressed physical
    action is the canonical realization of that one-edge virtual action.
\end{theorem}
\begin{proof}\leanok
    This is the projected local recovery equivalence applied to the right
    incident edge of $e$.
\end{proof}

\begin{theorem}[Endpoint recovery of a common virtual insertion]%
    \label{thm:peps_edgeEndpointLocalVirtualOpOfPhysicalOp_eq_of_projected_realization_eq}
    \lean{TNLean.PEPS.edgeEndpointLocalVirtualOpOfPhysicalOp_eq_of_projected_realization_eq}
    \leanok
    \uses{thm:peps_edgeLeftLocalVirtualOpOfPhysicalOp_eq_iff_projected_realization_eq,
        thm:peps_edgeRightLocalVirtualOpOfPhysicalOp_eq_iff_projected_realization_eq}
    Let $e=(u,v)$. If
    $P_uO_1P_u=\Phi_uT_{u,e}(M^{\mathsf T})\Psi_u$ and
    $P_vO_2P_v=\Phi_vT_{v,e}(M)\Psi_v$, then
    $\Psi_uO_1\Phi_u=T_{u,e}(M^{\mathsf T})$ and
    $\Psi_vO_2\Phi_v=T_{v,e}(M)$.
    Here $T_{w,e}(N)$ denotes the virtual operator on the legs incident to $w$
    which acts by $N$ on the leg belonging to $e$ and by the identity on the
    other incident legs.
\end{theorem}
\begin{proof}\leanok
    The endpoint equivalences give
    \begin{align}
        P_uO_1P_u=\Phi_uT_{u,e}(M^{\mathsf T})\Psi_u
        &\iff
        \Psi_uO_1\Phi_u=T_{u,e}(M^{\mathsf T}),
        \notag\\
        P_vO_2P_v=\Phi_vT_{v,e}(M)\Psi_v
        &\iff
        \Psi_vO_2\Phi_v=T_{v,e}(M).
        \notag
    \end{align}
    The two hypotheses are the left sides of these equivalences.
\end{proof}

\begin{theorem}[Endpoint physical-to-virtual recovery under image preservation]%
    \label{thm:peps_edgePhysicalToVirtualInsertion_of_projected_realization_eq}
    \lean{TNLean.PEPS.edgePhysicalToVirtualInsertion_of_projected_realization_eq}
    \leanok
    \uses{thm:peps_edgeEndpointLocalVirtualOpOfPhysicalOp_eq_of_projected_realization_eq,
        thm:peps_localVirtualOpOfPhysicalOp_realizes_of_projector}
    Let $e=(u,v)$. Suppose that
    $P_uO_1P_u=\Phi_uT_{u,e}(M^{\mathsf T})\Psi_u$ and
    $P_vO_2P_v=\Phi_vT_{v,e}(M)\Psi_v$. Suppose also that
    $O_1$ maps the image of $\Phi_u$ into itself and $O_2$ maps the image of
    $\Phi_v$ into itself. Then
    \begin{align}
        O_1\Phi_u
        &=\Phi_uT_{u,e}(M^{\mathsf T}),
        \notag\\
        O_2\Phi_v
        &=\Phi_vT_{v,e}(M).
        \notag
    \end{align}
\end{theorem}
\begin{proof}\leanok
    The preceding endpoint theorem identifies the two virtual pullbacks with
    the one-edge actions of $M^{\mathsf T}$ and $M$. The pullback
    specification gives the projected physical actions, and the
    image-preservation hypotheses remove the endpoint projectors.
\end{proof}

\begin{theorem}[Right-endpoint physical realization of the inserted coefficient]%
    \label{thm:peps_edgeInsertedCoeff_rightPhysicalRealization}
    \lean{TNLean.PEPS.edgeInsertedCoeff_eq_sum_right_physicalRealization}
    \leanok
    \uses{def:peps_edgeInsertedCoeff,
        def:peps_edgeBoundaryRightIndexEquivInsertedBoundaryConfig,
        thm:peps_localTensorMap_localIncidentMatrixOp_single}
    Let $e=(u,v)$ and let $\sigma$ be a physical configuration. If a physical
    operator on $v$ realizes the one-edge virtual matrix action on the
    distinguished incident edge, then the inserted-edge coefficient at
    $\sigma$ is obtained by leaving the left endpoint and open middle
    contraction explicit and applying that physical operator to the right
    endpoint tensor.
\end{theorem}
\begin{proof}\leanok
    Expand the physical realization on a basis local configuration at the
    right endpoint. The resulting sum over the right bond index is then
    reindexed together with the ordinary edge-boundary data.
\end{proof}

\begin{theorem}[Left-endpoint physical realization of the inserted coefficient]%
    \label{thm:peps_edgeInsertedCoeff_leftPhysicalRealization}
    \lean{TNLean.PEPS.edgeInsertedCoeff_eq_sum_left_physicalRealization}
    \leanok
    \uses{def:peps_edgeInsertedCoeff,
        def:peps_edgeBoundaryLeftIndexEquivInsertedBoundaryConfig,
        thm:peps_localTensorMap_localIncidentMatrixOp_single}
    Let $e=(u,v)$ and let $\sigma$ be a physical configuration. If a physical
    operator on $u$ realizes the one-edge virtual action of $M^{\mathsf T}$ on
    the distinguished incident edge, then the inserted-edge coefficient for
    $M$ at $\sigma$ is obtained by applying that physical operator to the left
    endpoint tensor and leaving the open middle contraction and the right
    endpoint explicit.
\end{theorem}
\begin{proof}\leanok
    Expand the physical realization on a basis local configuration at the left
    endpoint. The transpose converts the ordinary right bond index and the new
    left bond index into the coefficient $M_{ij}$; the resulting finite sum
    is then reindexed with the ordinary edge-boundary data.
\end{proof}

\begin{theorem}[Endpoint physical realizations of the inserted coefficient]%
    \label{thm:peps_edgeInsertedCoeff_endpointPhysicalRealization}
    \lean{TNLean.PEPS.edgeInsertedCoeff_endpointPhysicalRealization}
    \leanok
    \uses{thm:peps_localIncidentMatrixOp_physicalRealization,
        thm:peps_edgeInsertedCoeff_leftPhysicalRealization,
        thm:peps_edgeInsertedCoeff_rightPhysicalRealization}
    If the PEPS tensor is vertex-injective, then for every physical
    configuration $\sigma$, each inserted edge matrix is realized by physical
    operators at the endpoint tensors: the left endpoint realizes
    $M^{\mathsf T}$, the right endpoint realizes $M$, and both realizations
    reproduce the same inserted-edge coefficient at $\sigma$ in the
    edge-centered three-site decomposition.
\end{theorem}
\begin{proof}\leanok
    Apply the local physical-realization theorem at the two endpoints and then
    use the two coefficient identities above.
\end{proof}

\begin{theorem}[Equal neighboring physical actions recover a virtual insertion]%
    \label{thm:peps_physical_to_virtual_insertion}
    \lean{TNLean.PEPS.physical_to_virtual_insertion}
    \leanok
    \uses{def:peps_edgeBlockedThreeSiteInjective, def:peps_edgeBoundaryConfig,
        def:peps_edgeLeftLocalConfig, def:peps_edgeRightLocalConfig,
        def:peps_edgeOpenMiddleWeight, def:peps_localTensorMap,
        def:peps_localIncidentMatrixOp}
    Let $e=(u,v)$. Assume that the edge-blocked three-site chain at $e$ is
    injective and that every virtual bond has positive dimension. Let $O_1,O_2$
    be physical operators at the endpoint tensors, and
    write $\Phi_u,\Phi_v$ for the two local tensor maps. Suppose that, for every
    physical configuration $\sigma$,
    \begin{align}
        &\sum_{\beta}
        O_1(A_u(\beta_u))_{\sigma_u}
        C_{\mathrm{mid}}(\sigma;\beta_{\mathrm{L}},\beta_{\mathrm{R}})
        A_v(\beta_v)_{\sigma_v}
        \notag\\
        &\qquad =
        \sum_{\beta}
        A_u(\beta_u)_{\sigma_u}
        C_{\mathrm{mid}}(\sigma;\beta_{\mathrm{L}},\beta_{\mathrm{R}})
        O_2(A_v(\beta_v))_{\sigma_v}.
        \notag
    \end{align}
    Then
    \begin{align}
        \exists M\in\MN{D(e)},\qquad
        O_1\Phi_u
        &=\Phi_uT_{u,e}(M^{\mathsf T}),
        \notag\\
        O_2\Phi_v
        &=\Phi_vT_{v,e}(M).
        \notag
    \end{align}
    The same implication is represented diagrammatically by the
    $O_1,O_2\mapsto W$ step:
    % Formula (paper_normal.tex:355--485): O_1 acting at B_1 and O_2 acting at
    % B_2 give the same three-site state only if both arise from one bond matrix W.
    % Ink-to-index: each panel is the same closed virtual three-cycle; the first two
    % panels place O_1 and O_2 on physical legs, and the last splits the (1,2) bond
    % by W.  The three north legs remain the physical indices in every panel.
    % Boundary signature: each panel has (V,P)=(0,3).  Source verdict: exact
    % implication assembled from paper_normal.tex:355--375 and 422--485.
    \begin{center}
        \tenkzkernel
        \begin{tenkz}[rows={ket}, cols=3, physical=up, boundary=periodic]
            \tn[name=opSite]{B_1} & \tn{B_2} & \tn{B_3}
            \tn[at=1 n of opSite, skin=ring, species=operator,
                name=opRing, ports={270:physical, 90:physical}]{O_1}
            \tnwire{opSite.90}{opRing.270}
        \end{tenkz}
        $ = $
        \begin{tenkz}[rows={ket}, cols=3, physical=up, boundary=periodic]
            \tn{B_1} & \tn[name=opSite]{B_2} & \tn{B_3}
            \tn[at=1 n of opSite, skin=ring, species=operator,
                name=opRing, ports={270:physical, 90:physical}]{O_2}
            \tnwire{opSite.90}{opRing.270}
        \end{tenkz}
        $ \Longrightarrow $
        \begin{tenkz}[rows={ket}, cols=3, physical=up, boundary=periodic]
            \tn{B_1} & \tn{B_2} & \tn{B_3}
            \tn[skin=ring, species=operator, at=on bond-1-1-1-2 0.5]{W}
        \end{tenkz}
    \end{center}
    This is the converse part of the insertion-algebra lemma in
    \cite[Section~3]{Molnar2018NormalPEPS}, where the recovered matrix is
    denoted $W$.
\end{theorem}
\begin{proof}\leanok
    Contracting the middle block out of the hypothesis leaves a bond-contracted
    identity between the two endpoint tensors, holding for every residual
    boundary configuration. Inverting the right endpoint of this identity writes
    the action of $O_1$ on a left tensor vector as a bond-indexed combination of
    left tensor vectors with the same residual; the combining coefficients are
    read at a fixed reference right residual and define $M$. Inverting the left
    endpoint writes the action of $O_2$ on a right tensor vector as a bond-indexed
    combination of right tensor vectors. The left inverse identifies the two
    coefficient families, which is the source step $V=W$, so the right-endpoint
    combination is also governed by $M$. The reference residuals exist because
    every virtual bond has positive dimension.
\end{proof}

\begin{theorem}[Injectivity of an inserted edge coefficient]%
    \label{thm:peps_edgeInsertedCoeff_injective}
    \lean{TNLean.PEPS.edgeInsertedCoeff_injective}
    \leanok
    \uses{def:peps_edgeBlockedThreeSiteInjective, def:peps_edgeInsertedCoeff,
        thm:peps_edgeInsertedCoeff_leftPhysicalRealization,
        thm:peps_edgeInsertedCoeff_rightPhysicalRealization,
        thm:peps_physical_to_virtual_insertion}
    Let the edge-blocked three-site chain of $B$ at $e$ be injective, and
    assume every virtual bond of $B$ has positive dimension. If two matrices
    $M,M'\in\MN{D_B(e)}$ give the same inserted coefficient at every physical
    configuration, $C_B(e,\sigma,M)=C_B(e,\sigma,M')$ for every $\sigma$, then
    $M=M'$.
\end{theorem}

\begin{proof}\leanok
    Realize $M$ on the right endpoint and $M'^{\mathsf T}$ on the left
    endpoint. The endpoint realizations and the equality
    $C_B(e,\sigma,M)=C_B(e,\sigma,M')$ give, for every physical configuration
    $\sigma$,
    \begin{align}
        &\sum_{\beta}
        O_1(B_u(\beta_u))_{\sigma_u}
        C_{\mathrm{mid}}^B(\sigma;\beta_{\mathrm{L}},\beta_{\mathrm{R}})
        (B_v(\beta_v))_{\sigma_v}
        \notag\\
        &\qquad =
        \sum_{\beta}
        (B_u(\beta_u))_{\sigma_u}
        C_{\mathrm{mid}}^B(\sigma;\beta_{\mathrm{L}},\beta_{\mathrm{R}})
        O_2(B_v(\beta_v))_{\sigma_v}.
        \notag
    \end{align}
    Theorem~\ref{thm:peps_physical_to_virtual_insertion} gives a single virtual
    matrix realizing both endpoint operators. Endpoint uniqueness then
    identifies this matrix with $M$ and with $M'$.
\end{proof}

\begin{definition}[Edge-blocked insertion algebra correspondence]%
    \label{def:peps_edgeBlockedInsertionAlgebraIsomorphism_property}
    \lean{TNLean.PEPS.IsEdgeBlockedInsertionAlgebraIsomorphism}
    \leanok
    \uses{def:peps_edgeInsertedCoeff}
    For two PEPS tensors $A$ and $B$ and an edge $e$, this property is
    \begin{align}
        \exists \Phi_e:\MN{D_A(e)}
        \simeq_{\mathrm{alg}}
        \MN{D_B(e)},\qquad
        \forall \sigma,\ \forall X\in\MN{D_A(e)},\quad
        C_A(e,\sigma,X)=C_B(e,\sigma,\Phi_e(X)).
        \notag
    \end{align}
\end{definition}

\begin{theorem}[Edge-blocked insertion algebra isomorphism]%
    \label{thm:peps_edgeBlockedInsertionAlgebraIsomorphism}
    \lean{TNLean.PEPS.isEdgeBlockedInsertionAlgebraIsomorphism}
    \leanok
    \uses{def:peps_edgeBlockedInsertionAlgebraIsomorphism_property,
        thm:peps_edgeBlockedThreeSiteInjective,
        thm:peps_virtualInsertionPhysicalRealization,
        thm:peps_edgeInsertedCoeff_endpointPhysicalRealization,
        thm:peps_edgeInsertedCoeff_injective,
        thm:peps_physical_to_virtual_insertion,
        thm:peps_sameState_edgeBlockedCoeff_eq, def:peps_edgeInsertedCoeff}
    Suppose the two edge-blocked three-site chains associated to $A$ and $B$ at
    the edge $e$ are injective, suppose every virtual bond of $A$ and $B$ has
    positive dimension, and suppose the two PEPS tensors define the same state.
    % Scope restriction: the unrestricted Section 3 statement requires deriving
    % positive bond dimensions from injectivity; see
    % docs/paper-gaps/peps_injective_ft_section3_route.tex.
    Then
    \begin{align}
        \exists \Phi_e:\MN{D_A(e)}
        \simeq_{\mathrm{alg}}
        \MN{D_B(e)},\qquad
        \forall \sigma,\ \forall X\in\MN{D_A(e)},\quad
        C_A(e,\sigma,X)=C_B(e,\sigma,\Phi_e(X)).
        \notag
    \end{align}
    This is the edge-blocked form of the algebra isomorphism $X\mapsto Y$ in
    \cite[Section~3]{Molnar2018NormalPEPS}:
    % Formula: C_A(e,X;sigma)=C_B(e,Y;sigma) for the corresponding three-site
    % blocked chains.  Ink-to-index: X and Y split the distinguished (1,2) bond;
    % the periodic return closes the remaining virtual contraction, and the three
    % upper legs are sigma_1,sigma_2,sigma_3.
    % Boundary signature: both panels have (V,P)=(0,3).  Source verdict: exact
    % reproduction of paper_normal.tex:255--274 and 565--580 in house style.
    \begin{center}
        \tenkzkernel
        \begin{tenkz}[rows={wire,wire}, cols=6, bonds=none, align=1]
            \tn[at=(1,2), name=a1, label pos=s,
                ports={180:virtual, 0:virtual, 90:physical:$\sigma_1$}]{A_1}
            \tn[at=(1,3), skin=ring, name=x, ports={180:virtual, 0:virtual}]{X}
            \tn[at=(1,4), name=a2, label pos=s,
                ports={180:virtual, 0:virtual, 90:physical:$\sigma_2$}]{A_2}
            \tn[at=(1,5), name=a3, label pos=s,
                ports={180:virtual, 0:virtual, 90:physical:$\sigma_3$}]{A_3}
            \tn[at=(2,1), skin=none, name=jw, ports={0:virtual, 90:virtual}]{}
            \tn[at=(2,6), skin=none, name=je, ports={90:virtual, 180:virtual}]{}
            \tnwire{a1.0}{x.180}
            \tnwire{x.0}{a2.180}
            \tnwire{a2.0}{a3.180}
            \tnwire[route=arc]{a3.0}{je.90}
            \tnwire{je.180}{jw.0}
            \tnwire[route=arc]{jw.90}{a1.180}
        \end{tenkz}
        $=$
        \begin{tenkz}[rows={wire,wire}, cols=6, bonds=none, align=1]
            \tn[at=(1,2), name=b1, label pos=s,
                ports={180:virtual, 0:virtual, 90:physical:$\sigma_1$}]{B_1}
            \tn[at=(1,3), skin=ring, name=y, ports={180:virtual, 0:virtual}]{Y}
            \tn[at=(1,4), name=b2, label pos=s,
                ports={180:virtual, 0:virtual, 90:physical:$\sigma_2$}]{B_2}
            \tn[at=(1,5), name=b3, label pos=s,
                ports={180:virtual, 0:virtual, 90:physical:$\sigma_3$}]{B_3}
            \tn[at=(2,1), skin=none, name=jw, ports={0:virtual, 90:virtual}]{}
            \tn[at=(2,6), skin=none, name=je, ports={90:virtual, 180:virtual}]{}
            \tnwire{b1.0}{y.180}
            \tnwire{y.0}{b2.180}
            \tnwire{b2.0}{b3.180}
            \tnwire[route=arc]{b3.0}{je.90}
            \tnwire{je.180}{jw.0}
            \tnwire[route=arc]{jw.90}{b1.180}
        \end{tenkz}
    \end{center}
\end{theorem}
\begin{proof}\leanok
    For a matrix $X\in\MN{D_A(e)}$, realize the insertion of $X$ at the two
    endpoints of the first edge-blocked chain by physical operators $O_1,O_2$.
    Write $C_{\mathrm{mid}}^A$ and $C_{\mathrm{mid}}^B$ for the open middle
    weights of the two edge-blocked chains. For every physical configuration
    $\sigma$, endpoint realization of $X$ and equality of the PEPS states give
    \begin{align}
        C_A(e,\sigma,X)
         & =
        \sum_{\beta}
        O_1(A_u(\beta_u))_{\sigma_u}
        C_{\mathrm{mid}}^A(\sigma;\beta_{\mathrm{L}},\beta_{\mathrm{R}})
                             (A_v(\beta_v))_{\sigma_v}
        \notag \\
         & =
        \sum_{\beta}
        O_1(B_u(\beta_u))_{\sigma_u}
        C_{\mathrm{mid}}^B(\sigma;\beta_{\mathrm{L}},\beta_{\mathrm{R}})
                             (B_v(\beta_v))_{\sigma_v},
        \notag
    \end{align}
    and
    \begin{align}
        C_A(e,\sigma,X)
         & =
        \sum_{\beta}
        (A_u(\beta_u))_{\sigma_u}
        C_{\mathrm{mid}}^A(\sigma;\beta_{\mathrm{L}},\beta_{\mathrm{R}})
        O_2(A_v(\beta_v))_{\sigma_v}
        \notag \\
         & =
        \sum_{\beta}
        (B_u(\beta_u))_{\sigma_u}
        C_{\mathrm{mid}}^B(\sigma;\beta_{\mathrm{L}},\beta_{\mathrm{R}})
        O_2(B_v(\beta_v))_{\sigma_v}.
        \notag
    \end{align}
    Therefore, for every physical configuration $\sigma$,
    \begin{align}
        \sum_{\beta}
        O_1(B_u(\beta_u))_{\sigma_u}
        C_{\mathrm{mid}}^B(\sigma;\beta_{\mathrm{L}},\beta_{\mathrm{R}})
                             (B_v(\beta_v))_{\sigma_v}
         & =
        \sum_{\beta}
        (B_u(\beta_u))_{\sigma_u}
        C_{\mathrm{mid}}^B(\sigma;\beta_{\mathrm{L}},\beta_{\mathrm{R}})
        O_2(B_v(\beta_v))_{\sigma_v}.
        \notag
    \end{align}
    The positive-bond physical-to-virtual recovery applied to the second
    edge-blocked chain gives a matrix $Y\in\MN{D_B(e)}$ such that
    $C_A(e,\sigma,X)=C_B(e,\sigma,Y)$ for every physical configuration
    $\sigma$.
    Thus $X\mapsto Y$ gives the required coefficient correspondence. Denote
    this correspondence by $T$. For every physical configuration $\sigma$,
    \begin{align}
        C_B(e,\sigma,T(X+X'))
         & =C_A(e,\sigma,X+X')
        \notag                               \\
         & =C_A(e,\sigma,X)+C_A(e,\sigma,X')
        \notag                               \\
         & =C_B(e,\sigma,T(X)+T(X')).
        \notag
    \end{align}
    Similarly, $T(cX)=cT(X)$. The multiplicative identity is read from
    the inserted two-endpoint realization:
    \begin{align}
        C_B(e,\sigma,T(XX'))
        =C_A(e,\sigma,XX')
        =C_B(e,\sigma,T(X)T(X')).
        \notag
    \end{align}
    The unit satisfies
    $C_B(e,\sigma,T(1))=C_A(e,\sigma,1)=C_B(e,\sigma,1)$.
    Theorem~\ref{thm:peps_edgeInsertedCoeff_injective} identifies the inserted
    matrices in these coefficient identities. Exchanging $A$ and $B$ gives the
    inverse correspondence, hence a $\C$-algebra isomorphism.
\end{proof}

\begin{theorem}[Matrix-algebra isomorphisms preserve size]%
    \label{thm:matrixAlgEquiv_fin_eq}
    \lean{MPSTensor.matrixAlgEquiv_fin_eq}
    \leanok
    If the full matrix algebras $\MN{D}$ and $\MN{E}$ are isomorphic as
    $\C$-algebras, then $D=E$.
\end{theorem}

\begin{proof}\leanok
    The isomorphism preserves complex vector-space dimension, and
    $\dim_{\C}\MN{D}=D^2$ and $\dim_{\C}\MN{E}=E^2$, so $D=E$.
\end{proof}

\begin{theorem}[Matrix-algebra isomorphisms are inner after identifying size]%
    \label{thm:matrixAlgEquiv_inner_of_fin}
    \lean{MPSTensor.matrixAlgEquiv_inner_of_fin}
    \leanok
    \uses{thm:matrixAlgEquiv_fin_eq, thm:skolem_noether}
    Let $\varphi:\MN{D}\to\MN{E}$ be a $\C$-algebra isomorphism. Then
    $D=E$; let $h$ denote this equality and
    $\iota_h:\MN{D}\to\MN{E}$ the canonical index-relabelling isomorphism.
    There is $X\in\GL_E(\C)$ with
    $\varphi(M)=X\iota_h(M)X^{-1}$ for all $M\in\MN{D}$.
\end{theorem}

\begin{proof}\leanok
    By Theorem~\ref{thm:matrixAlgEquiv_fin_eq}, $D=E$, so $\iota_h$ is defined.
    Then $\varphi\circ\iota_h^{-1}$ is a $\C$-algebra automorphism of
    $\MN{E}$, and Skolem--Noether
    (Theorem~\ref{thm:skolem_noether}) gives $X\in\GL_E(\C)$ with
    $(\varphi\circ\iota_h^{-1})(N)=XNX^{-1}$. Evaluating at
    $N=\iota_h(M)$ gives $\varphi(M)=X\iota_h(M)X^{-1}$.
\end{proof}

\begin{theorem}[Edge gauge from the insertion algebra isomorphism]%
    \label{thm:peps_edgeGaugeFromInsertionAlgebraIsomorphism}
    \lean{TNLean.PEPS.edgeGaugeFromInsertionAlgebraIsomorphism}
    \leanok
    \uses{thm:matrixAlgEquiv_inner_of_fin}
    Suppose the edge insertion correspondence is a $\C$-algebra isomorphism
    between the two full matrix algebras on the chosen bond. Then the two bond
    dimensions on that edge are equal. If $h_e:D_A(e)=D_B(e)$ is this
    equality, and if
    $\phi_{h_e}:\MN{D_A(e)}\to\MN{D_B(e)}$ is the induced algebra
    isomorphism, then, for an invertible edge matrix $Z_e$, the insertion
    correspondence is
    % Formula/source: Papers/1804.04964/paper_normal.tex, lines 254--279
    % and 576--582, identifies the bond insertion map as an algebra
    % isomorphism and writes its inner-conjugation form Y=ZXZ^{-1}.
    % Ink-to-index map: exact algebra is used here because X, phi_h_e(X),
    % and Y are bond endomorphisms, not local tensors with physical legs.
    % Contraction set: the two ordinary matrix products sum the two intermediate
    % bond indices; there are no physical or PEPS lattice contractions.
    % Type verdict: an endomorphism of the B-side edge bond space.
    % Boundary signature: not applicable to the displayed algebra; no tensor
    % boundary is asserted by this formula.
    $Y=Z_e\phi_{h_e}(X)Z_e^{-1}$.
    This is the finite-dimensional matrix-algebra step in
    \cite[Section~3]{Molnar2018NormalPEPS}.
\end{theorem}

\begin{proof}\leanok
    Apply Theorem~\ref{thm:matrixAlgEquiv_inner_of_fin} to the full matrix
    algebras determined by the two edge bond spaces.
\end{proof}

\begin{theorem}[Per-edge gauge family from insertion algebra isomorphisms]%
    \label{thm:peps_exists_edgeGaugeFamily}
    \lean{TNLean.PEPS.exists_edgeGaugeFamily}
    \leanok
    \uses{thm:peps_edgeBlockedThreeSiteInjective,
        thm:peps_edgeBlockedInsertionAlgebraIsomorphism,
        thm:peps_edgeGaugeFromInsertionAlgebraIsomorphism}
    Let $A$ and $B$ be vertex-injective PEPS tensors with the same PEPS state.
    Suppose the corresponding virtual bond dimensions are equal and every
    virtual bond of $A$ has positive dimension. Then there is a family
    $X_e\in\GL(D_A(e))$ such that, for every edge $e$, there is a
    $\C$-algebra isomorphism
    $\Phi_e:\MN{D_A(e)}\to\MN{D_B(e)}$ satisfying
    $C_A(e,\sigma,M)=C_B(e,\sigma,\Phi_e(M))$ for every physical configuration
    $\sigma$ and matrix $M$. Writing
    $\phi_{h_e}:\MN{D_A(e)}\to\MN{D_B(e)}$ for the algebra
    isomorphism induced by the dimension equality on $e$, one has
    $\Phi_e(M)=\phi_{h_e}(X_eMX_e^{-1})$.
\end{theorem}

\begin{proof}\leanok
    Vertex injectivity and positivity give edge-blocked three-site injectivity
    for $A$ and $B$ at each edge.
    Theorem~\ref{thm:peps_edgeBlockedInsertionAlgebraIsomorphism}
    gives the algebra isomorphism $\Phi_e$ preserving inserted coefficients.
    Theorem~\ref{thm:peps_edgeGaugeFromInsertionAlgebraIsomorphism} realizes
    $\Phi_e$ by conjugation, and the bond-dimension equality transports the
    conjugating matrix to the $A$-side bond.
\end{proof}

% Specializations of the per-edge gauge to the translated interior edges of the
% open square lattice, where the gauge comes from a one-edge blocking datum.

\begin{theorem}[Per-edge gauge at a translated horizontal interior edge]%
    \label{thm:peps_exists_regionEdgeGauge_normalSquareHorizontalTranslatedEdge}
    \lean{TNLean.PEPS.exists_regionEdgeGauge_normalSquareHorizontalTranslatedEdge}
    \leanok
    \uses{lem:peps_normalSquareHorizontalTranslatedEdge_blockingData}
    Let $A$ and $B$ be normal PEPS on the open $W\times H$ square lattice that
    satisfy the rectangular-injectivity hypotheses with union closure, have
    positive physical dimension $d>0$, have equal and everywhere-positive virtual
    bond dimensions, and define the same PEPS state. Fix a translated horizontal
    interior frame with two rows and two
    columns of margin, that is $2\le x_0$, $2\le y_0$, $x_0+8\le W$, and
    $y_0+8\le H$, and let $e$ be the distinguished horizontal edge of that frame.
    Then the bond dimensions agree on $e$, $D_A(e)=D_B(e)$, and there are an
    invertible edge matrix $Z\in\GL(D_B(e))$ and a forward map $\mathrm{fwd}$
    given by the conjugation
    $\mathrm{fwd}(M)=Z\phi_{h_e}(M)Z^{-1}$, where
    $\phi_{h_e}:\MN{D_A(e)}\to\MN{D_B(e)}$ is the reindexing induced by the
    bond-dimension equality on $e$. No single-vertex injectivity is used; the
    only injectivity hypotheses are the blocked-region injectivities of the
    interior blocking data and the red and host injectivities of $B$.
\end{theorem}

\begin{proof}\leanok
    The translated horizontal frame supplies, for both $A$ and $B$, the one-edge
    interior blocking datum. Because the red, blue, and complementary blocks are
    coordinate regions, independent of the tensor, the two data share them:
    $R_A=R_B$, $B_A=B_B$, and $C_A=C_B$.
    The distinguished edge $e$ is the unique crossing edge of the red and
    blue blocks. Applying the per-edge gauge construction from a one-edge
    blocking datum to the two blocking data yields the bond-dimension equality on
    $e$ and the conjugating gauge $Z$.
\end{proof}

\begin{theorem}[Per-edge gauge at a translated vertical interior edge]%
    \label{thm:peps_exists_regionEdgeGauge_normalSquareVerticalTranslatedEdge}
    \lean{TNLean.PEPS.exists_regionEdgeGauge_normalSquareVerticalTranslatedEdge}
    \leanok
    The coordinate-swap counterpart of
    Theorem~\ref{thm:peps_exists_regionEdgeGauge_normalSquareHorizontalTranslatedEdge}.
    Under the same hypotheses on $A$ and $B$, for a translated vertical interior
    frame with two rows and two columns of margin, that is $2\le x_0$,
    $2\le y_0$, $x_0+8\le W$, and $y_0+8\le H$, with distinguished vertical
    edge $e$, the bond dimensions agree on $e$, $D_A(e)=D_B(e)$, and there are an
    invertible edge matrix $Z\in\GL(D_B(e))$ and a forward map $\mathrm{fwd}$
    given by the conjugation
    $\mathrm{fwd}(M)=Z\phi_{h_e}(M)Z^{-1}$, with
    $\phi_{h_e}:\MN{D_A(e)}\to\MN{D_B(e)}$ the reindexing induced by the
    bond-dimension equality on $e$.
\end{theorem}

\begin{proof}
    \lean{TNLean.PEPS.Region_map_squareLatticeCoordinateSwap_verticalTranslatedEdgeRed}
    \lean{TNLean.PEPS.Region_map_squareLatticeCoordinateSwap_verticalTranslatedEdgeBlue}
    \leanok
    \uses{thm:peps_exists_regionEdgeGauge_normalSquareHorizontalTranslatedEdge}
    Interchange the two lattice coordinates. This graph isomorphism sends the
    vertical translated frame to the horizontal translated frame, transposes the
    rectangular injectivity hypotheses, and sends the distinguished vertical edge
    to the distinguished horizontal edge. Apply
    Theorem~\ref{thm:peps_exists_regionEdgeGauge_normalSquareHorizontalTranslatedEdge}
    to the transported tensors. Transporting its bond-dimension equality and
    conjugating gauge back along the inverse coordinate swap gives the required
    equality on $e$ and the matrix $Z$.
\end{proof}

\begin{definition}[Factorized local gauge datum]%
    \label{def:peps_hasFactorizedLocalGauge}
    \lean{TNLean.PEPS.HasFactorizedLocalGauge}
    \leanok
    \uses{def:peps_localTensorMap, def:peps_localLeftInverse}
    Fix $A$, $B$, a bond-dimension equality, and a vertex $v$. Let
    $\mathcal T^A_v$ be the local tensor map of $A$, let $L^A_v$ be the
    chosen left inverse of $\mathcal T^A_v$, and let
    $\Pi^A_v=\mathcal T^A_vL^A_v$. The factorized local gauge datum consists
    of the following two identities for local virtual configurations
    $\eta,\eta'$:
    \begin{align}
        \Pi^A_v B_v(\eta,-)
         & =B_v(\eta,-),
        \notag                                 \\
        L^A_vB_v(\eta,-)(\eta')
         & =\prod_{e\ni v}X_e(\eta_e,\eta'_e),
        \notag
    \end{align}
    with $X_e\in\GL(D_A(e))$ on each incident edge. This is the local output
    required from the blocked-middle reduction in
    \cite[Section~3]{Molnar2018NormalPEPS}.
\end{definition}

\begin{definition}[Blocked-middle local gauge formula]%
    \label{def:peps_blockedMiddleGaugeFormula}
    \lean{TNLean.PEPS.BlockedMiddleGaugeFormula}
    \leanok
    The blocked-middle local gauge formula at $v$ is the assertion that there
    are matrices $X_e\in\GL(D_A(e))$, one on each incident edge, such that for
    all local virtual configurations $\eta$ and physical indices $\sigma$,
    \begin{align}
        B_v(\eta,\sigma)
        =
        \sum_{\eta'}
        \left(\prod_{e\ni v}X_e(\eta_e,\eta'_e)\right)
        A_v(\eta',\sigma).
        \notag
    \end{align}
\end{definition}

\begin{theorem}[Blocked-middle formula implies factorized local gauge]%
    \label{thm:peps_hasFactorizedLocalGauge_of_blockedMiddleGaugeFormula}
    \lean{TNLean.PEPS.hasFactorizedLocalGauge_of_blockedMiddleGaugeFormula}
    \leanok
    \uses{def:peps_blockedMiddleGaugeFormula, def:peps_hasFactorizedLocalGauge}
    If the blocked-middle local gauge formula holds at $v$ with matrices
    $X_e$, then the factorized local gauge datum holds at $v$ with the same
    matrices.
    % Open: derive the displayed blocked-middle formula from equality of PEPS
    % states by the edge-centred three-site reduction.
\end{theorem}
\begin{proof}\leanok
    For fixed $\eta$, set
    $c_\eta(\eta')=\prod_{e\ni v}X_e(\eta_e,\eta'_e)$. The displayed
    blocked-middle formula is precisely
    $B_v(\eta,-)=\mathcal T^A_vc_\eta$. Hence
    \begin{align}
        \Pi^A_vB_v(\eta,-)
         & =\Pi^A_v\mathcal T^A_vc_\eta
        =\mathcal T^A_vc_\eta
        =B_v(\eta,-),
        \notag                               \\
        L^A_vB_v(\eta,-)(\eta')
         & =L^A_v\mathcal T^A_vc_\eta(\eta')
        =c_\eta(\eta')
        =\prod_{e\ni v}X_e(\eta_e,\eta'_e).
        \notag
    \end{align}
\end{proof}

\begin{theorem}[Blocked-middle formula gives a local gauge formula]%
    \label{thm:peps_localGauge_exists_of_blockedMiddleGaugeFormula}
    \lean{TNLean.PEPS.localGauge_exists_of_blockedMiddleGaugeFormula}
    \leanok
    \uses{def:peps_blockedMiddleGaugeFormula,
        thm:peps_hasFactorizedLocalGauge_of_blockedMiddleGaugeFormula,
        thm:localGauge_exists}
    Suppose that the edge-centred three-site reduction supplies invertible
    matrices $X_e$ on the incident edges of $v$ such that, for all local virtual
    configurations $\eta$ and physical indices $\sigma$,
    \begin{align}
        B_v(\eta,\sigma)
        =
        \sum_{\eta'}
        \left(\prod_{e\ni v}X_e(\eta_e,\eta'_e)\right)
        A_v(\eta',\sigma).
        \notag
    \end{align}
    Then the same family $X_e$ is a local gauge at $v$:
    % Formula: the displayed component identity above is the local gauge formula.
    % Ink-to-index/type verdict: its boundary is (V,P)=(deg(v),1).  No matching
    % standalone picture occurs in paper_normal.tex, and the old picture hard-coded
    % five virtual legs, so the valence-independent equation is authoritative.
    % Open: derive the blocked-middle formula from SameState
    % by the edge-centred three-site reduction.
\end{theorem}
\begin{proof}\leanok
    By the preceding theorem,
    $L^A_vB_v(\eta,-)(\eta')=\prod_{e\ni v}X_e(\eta_e,\eta'_e)$. Since
    $\Pi^A_vB_v(\eta,-)=B_v(\eta,-)$ and
    $\Pi^A_v=\mathcal T^A_vL^A_v$,
    \begin{align}
        B_v(\eta,\sigma)
         & =\sum_{\eta'}L^A_vB_v(\eta,-)(\eta')A_v(\eta',\sigma)
        \notag                                                   \\
         & =\sum_{\eta'}
        \left(\prod_{e\ni v}X_e(\eta_e,\eta'_e)\right)
        A_v(\eta',\sigma).
        \notag
    \end{align}
\end{proof}

\begin{theorem}[Absorbing edge gauges into the second PEPS]%
    \label{thm:peps_edge_gauge_absorption}
    \lean{TNLean.PEPS.edge_gauge_absorption}
    \leanok
    \uses{def:peps_absorbEdgeGauges, thm:peps_absorbEdgeGauges_bondDim,
        thm:peps_absorbEdgeGauges_component}
    Given edge gauges $X_e$, set $\widetilde B=B^X$. Then
    $D_{\widetilde B}=D_B$ and, for every vertex $v$,
    \begin{align}
        \widetilde B_v(\eta,\sigma)
        &=
        \sum_{\eta'}
        \left(\prod_{ie}M_{ie}(\eta(ie),\eta'(ie))\right)
        B_v(\eta',\sigma),
        \notag
    \end{align}
    where $M_{ie}=X_e$ or $M_{ie}=(X_e^{-1})^t$ according to the orientation of
    the endpoint $ie$.
    % Formula: the preceding component identity is gauge absorption at v.
    % Ink-to-index/type verdict: paper_normal.tex:1212--1314 gives only
    % vertex-specific pictures, each with boundary (V,P)=(deg(v),1); the old
    % degree-four star cannot represent the theorem on an arbitrary graph.
    After this absorption, arbitrary insertions of the same matrix on the same
    edge are compared in the first PEPS and the modified second PEPS.
\end{theorem}
\begin{proof}\leanok
    With $\widetilde B=B^X$, Theorem~\ref{thm:peps_absorbEdgeGauges_bondDim}
    gives $D_{\widetilde B}=D_B$, and
    Theorem~\ref{thm:peps_absorbEdgeGauges_component} gives
    \begin{align}
        \widetilde B_v(\eta,\sigma)
        &=
        \sum_{\eta'}
        \left(\prod_{ie}M_{ie}(\eta(ie),\eta'(ie))\right)
        B_v(\eta',\sigma).
        \notag
    \end{align}
\end{proof}

\begin{definition}[Post-absorption inserted-edge comparison]%
    \label{def:peps_postAbsorptionEdgeInsertionEquality}
    \lean{TNLean.PEPS.PostAbsorptionEdgeInsertionEquality}
    \leanok
    \uses{def:peps_edgeInsertedCoeff}
    For tensor families $A,\widetilde B$, the post-absorption comparison holds
    when $h:D_A=D_{\widetilde B}$ and the following equality holds. For each
    edge $e$, write
    $\phi_{h_e}:\MN{D_A(e)}
        \to\MN{D_{\widetilde B}(e)}$ for the matrix-algebra
    transport induced by $h_e:D_A(e)=D_{\widetilde B}(e)$. Then, for every
    matrix $X\in\MN{D_A(e)}$ and every physical configuration $\sigma$,
    $C_A(e,X;\sigma)
        =C_{\widetilde B}(e,\phi_{h_e}(X);\sigma)$.
\end{definition}

\begin{theorem}[Post-absorption edge insertion equality]%
    \label{thm:peps_postAbsorptionEdgeInsertionEquality}
    \lean{TNLean.PEPS.post_absorption_edge_insertion_equality}
    \leanok
    \uses{thm:peps_edge_gauge_absorption,
        def:peps_postAbsorptionEdgeInsertionEquality}
    Let $A$ and $B$ be vertex-injective PEPS with the same state, assume every
    bond dimension is positive, and assume the bond-dimension equality
    $h:D_A=D_B$. After the edge gauges have been
    absorbed into the second tensor family, there are gauges $Z_e$ and
    $\widetilde B=B^Z$ such that $D_{\widetilde B}=D_A$ and, for every edge
    $e$, write $\phi_{h_e}:\MN{D_A(e)}\to\MN{D_{\widetilde B}(e)}$ for the
    matrix-algebra transport induced by the corresponding equality
    $h_e:D_A(e)=D_{\widetilde B}(e)$. Then, for every matrix
    $X\in\MN{D_A(e)}$ and every physical configuration $\sigma$,
    $C_A(e,X;\sigma)
        =C_{\widetilde B}(e,\phi_{h_e}(X);\sigma)$.
    On the five-site graph used in the proof, this equality reads
    % Formula: C_A(e,X;sigma)=C_{\widetilde B}(e,phi_{h_e}(X);sigma), with
    % the transport identified on the common bond in the displayed house spelling.
    % Ink-to-index: each panel has the five physical tensors of eq:inj_equal_edge,
    % six graph bonds (the 5-cycle and the 2--5 chord), and X splits that chord.
    % All virtual indices are contracted; the five north legs are sigma_1,...,sigma_5.
    % Boundary signature: both source panels have (V,P)=(0,5), specializing the
    % general theorem's (0,|V(G)|) physical boundary.  Source verdict: exact graph
    % and insertion from paper_normal.tex:1037--1065.
    \begin{center}
        \tenkzkernel
        \begin{tenkz}[rows={wire}, cols=5, frame=circle, bonds=none]
            \tn[at=(1,1), name=edgeEqA1, label pos=145,
                ports={0:virtual, 180:virtual, 90:physical}]{A_1}
            \tn[at=(1,2), name=edgeEqA2, label pos=145,
                ports={0:virtual, 180:virtual, -72:virtual, 90:physical}]{A_2}
            \tn[at=(1,3), name=edgeEqA3, label pos=145,
                ports={0:virtual, 180:virtual, 90:physical}]{A_3}
            \tn[at=(1,4), name=edgeEqA4, label pos=145,
                ports={0:virtual, 180:virtual, 90:physical}]{A_4}
            \tn[at=(1,5), name=edgeEqA5, label pos=145,
                ports={0:virtual, 180:virtual, 252:virtual, 90:physical}]{A_5}
            \tnwire{edgeEqA1.0}{edgeEqA2.180}
            \tnwire{edgeEqA2.0}{edgeEqA3.180}
            \tnwire{edgeEqA3.0}{edgeEqA4.180}
            \tnwire{edgeEqA4.0}{edgeEqA5.180}
            \tnwire{edgeEqA5.0}{edgeEqA1.180}
            \tnwire[name=edgeEqAChord]{edgeEqA2.-72}{edgeEqA5.252}
            \tn[skin=ring, species=operator, at=on edgeEqAChord 0.5]{X}
        \end{tenkz}
        $ = $
        \begin{tenkz}[rows={wire}, cols=5, frame=circle, bonds=none]
            \tn[at=(1,1), name=edgeEqB1, label pos=145,
                ports={0:virtual, 180:virtual, 90:physical}]{\widetilde B_1}
            \tn[at=(1,2), name=edgeEqB2, label pos=145,
                ports={0:virtual, 180:virtual, -72:virtual, 90:physical}]{\widetilde B_2}
            \tn[at=(1,3), name=edgeEqB3, label pos=145,
                ports={0:virtual, 180:virtual, 90:physical}]{\widetilde B_3}
            \tn[at=(1,4), name=edgeEqB4, label pos=145,
                ports={0:virtual, 180:virtual, 90:physical}]{\widetilde B_4}
            \tn[at=(1,5), name=edgeEqB5, label pos=145,
                ports={0:virtual, 180:virtual, 252:virtual, 90:physical}]{\widetilde B_5}
            \tnwire{edgeEqB1.0}{edgeEqB2.180}
            \tnwire{edgeEqB2.0}{edgeEqB3.180}
            \tnwire{edgeEqB3.0}{edgeEqB4.180}
            \tnwire{edgeEqB4.0}{edgeEqB5.180}
            \tnwire{edgeEqB5.0}{edgeEqB1.180}
            \tnwire[name=edgeEqBChord]{edgeEqB2.-72}{edgeEqB5.252}
            \tn[skin=ring, species=operator, at=on edgeEqBChord 0.5]{X}
        \end{tenkz}
    \end{center}
    as in \cite[Section~3, lines 1037--1065]{Molnar2018NormalPEPS}. The
    positive-bond hypothesis is the source's standing assumption that injective
    PEPS have nonzero virtual bond spaces; it is needed because the edge gauges
    come from blocking the PEPS around each edge into a three-site injective
    chain.
\end{theorem}
\begin{proof}\leanok
    \uses{thm:peps_exists_edgeGaugeFamily, thm:peps_edge_gauge_absorption,
        def:peps_postAbsorptionEdgeInsertionEquality}
    Block the PEPS around each edge into a three-site injective chain and compare
    the matched matrix insertions: this assigns to every edge an invertible bond
    matrix $X_e$ realizing the insertion-algebra transport as conjugation. Absorb
    the transpose of $X_e$ into the second tensor family on that edge. The
    open-edge gauge action then conjugates each inserted matrix by the transpose
    of the absorbed gauge, while the insertion-algebra transport conjugates by
    $X_e$; the two conjugations agree because the absorbed gauge is the transpose
    of $X_e$, giving the post-absorption insertion equality.
\end{proof}

\begin{theorem}[Local gauge extraction]%
    \label{thm:localGauge_exists}
    \lean{TNLean.PEPS.localGauge_exists}
    \lean{TNLean.PEPS.localGauge_exists_of_factorizedLocalGauge}
    \leanok
    \uses{def:peps_hasFactorizedLocalGauge}
    Under the factorized local gauge hypothesis, the local gauge formula follows.
    What remains open is to upgrade equality of the edge-blocked coefficients to
    the blocked-middle formula using the corresponding three-site
    injective-chain argument. The physical-realization and physical-to-virtual
    steps in that argument are recorded separately as
    Theorems~\ref{thm:peps_virtualInsertionPhysicalRealization} and
    \ref{thm:peps_physical_to_virtual_insertion}.
    After Theorem~\ref{thm:peps_postAbsorptionEdgeInsertionEquality} supplies
    the factorized local gauge hypothesis, injectivity on the star neighborhood
    provides a left inverse that isolates the local edge gauges. Thus there are
    incident-edge matrices $X_e$ such that
    % Formula: local extraction yields the valence-independent component equation.
    % Ink-to-index/type verdict: both sides have boundary (V,P)=(deg(v),1).
    % No standalone source picture states this extraction, and the old two panels
    % had unequal virtual boundary types, so the exact equation replaces them.
    \begin{align}
        B_v(\eta,\sigma)
        &=
        \sum_{\eta'}
        \left(\prod_{e\ni v}X_e(\eta_e,\eta'_e)\right)
        A_v(\eta',\sigma).
        \notag
    \end{align}
\end{theorem}
\begin{proof}\leanok
    The factorized local gauge datum already supplies the incident-edge
    invertible matrices and the factorized local coefficient identity.
\end{proof}

\begin{theorem}[Gauge consistency]%
    \label{thm:gaugeConsistency}
    \lean{TNLean.PEPS.gaugeConsistency}
    \leanok
    \uses{thm:peps_exists_edgeGaugeFamily,
        thm:peps_postAbsorptionEdgeInsertionEquality,
        thm:localGauge_exists, def:gaugeVertex}
    Suppose $A$ and $B$ are vertex-injective PEPS tensors with the same PEPS
    state, that the corresponding virtual bond dimensions are equal, that
    every virtual bond of $A$ has positive dimension, and that the graph is
    connected. The connectivity assumption is necessary: on a disconnected
    graph the per-vertex scalars produced by the reduction below cannot be
    absorbed into edge gauges, since on the empty graph on two vertices the
    products of the vertex scalars can agree while no edge gauge relates the two
    tensors.
    Choose an edge $e=(u,v)$ and block all other vertices into one middle
    tensor. The three resulting tensors form an injective chain, so the
    one-dimensional injective comparison assigns an invertible matrix to the
    edge $e$. Repeating this for every edge and absorbing these matrices into
    the second tensor family gives the post-absorption insertion equality of
    \cite[Section~3]{Molnar2018NormalPEPS}: for every edge and every matrix,
    inserting that matrix on the edge in the first PEPS gives the same state as
    inserting it on the same edge in the modified second PEPS. Applying the
    local extraction at each vertex then gives one edge gauge family whose
    endpoint actions transform the first tensor family into the second:
    % Formula: one global family X supplies the local gauge equation at every v.
    % Ink-to-index/type verdict: the local component has boundary
    % (V,P)=(deg(v),1), while a two-vertex comparison has
    % (deg(u)+deg(w)-2,2).  The source gives the vertexwise equations in
    % paper_normal.tex:1212--1314, not the leg-suppressed scalar legacy picture.
    \begin{align}
        B_v(\eta,\sigma)
        &=
        (A^X)_v(\eta,\sigma)
        =
        \sum_{\eta'}
        \left(\prod_{e\ni v}X_{v,e}(\eta_e,\eta'_e)\right)
        A_v(\eta',\sigma).
        \notag
    \end{align}
\end{theorem}
\begin{proof}\leanok
    Absorbing the edge gauges into the second tensor family gives the
    post-absorption insertion equality. On each vertex with an incident edge,
    the one-vertex-versus-complement comparison supplies a nonzero scalar
    relating the first tensor to the absorbed second tensor. Because the state
    is nonzero, these per-vertex scalars multiply to one, so on the connected
    graph a spanning-tree construction realizes their reciprocals as the
    oriented incidence product of one edge-scalar family. Folding those scalars
    into the inverse of the absorbed edge gauges produces the global gauge. On a
    one-vertex graph there are no edges and the gauge is trivial.
\end{proof}

\begin{remark}
    The preceding identities are the algebraic content of the edge-gauge
    construction in \cite[Section~3]{Molnar2018NormalPEPS}. On an ordered edge
    $e=(u,w)$, the actions $X_e$ at $u$ and $(X_e^{-1})^\top$ at $w$ cancel
    upon contraction. Absorbing these matrices into $B$ gives
    $\widetilde B$, for which every inserted-edge coefficient agrees with the
    corresponding coefficient of $A$. Comparing one vertex with its injective
    complement then yields the local gauge equation. Repeating the argument on
    every edge produces a single compatible gauge family; the two-injective
    comparison below shows that the remaining three-leg operator is scalar.
\end{remark}

\begin{definition}[Left-identity three-leg residual form]%
    \label{def:peps_threeLegLeftIdentityForm}
    \lean{TNLean.PEPS.threeLegLeftIdentityForm}
    \leanok
    For three virtual leg spaces $\alpha,\beta,\gamma$ and an endomorphism
    $Z$ of $\beta\otimes\gamma$, define the induced endomorphism of
    $\alpha\otimes\beta\otimes\gamma$ by acting as the identity on the
    $\alpha$ leg and as $Z$ on the remaining two legs. This is the
    $\Id_\alpha\otimes Z$ form appearing after the first tensor is inverted in
    \cite[Section~3]{Molnar2018NormalPEPS}.
\end{definition}

\begin{definition}[Right-identity three-leg residual form]%
    \label{def:peps_threeLegRightIdentityForm}
    \lean{TNLean.PEPS.threeLegRightIdentityForm}
    \leanok
    For an endomorphism $U$ of $\alpha\otimes\beta$, define the induced
    endomorphism of $\alpha\otimes\beta\otimes\gamma$ by acting as $U$ on the
    first two legs and as the identity on the $\gamma$ leg. This is the
    $U\otimes\Id_\gamma$ form in \cite[Section~3]{Molnar2018NormalPEPS}.
\end{definition}

\begin{definition}[Middle-identity three-leg residual form]%
    \label{def:peps_threeLegMiddleIdentityForm}
    \lean{TNLean.PEPS.threeLegMiddleIdentityForm}
    \leanok
    For an endomorphism $W$ of $\alpha\otimes\gamma$, define the induced
    endomorphism of $\alpha\otimes\beta\otimes\gamma$ by acting as $W$ on the
    outer two legs and as the identity on the middle $\beta$ leg. This is the
    third residual form in \cite[Section~3]{Molnar2018NormalPEPS}.
\end{definition}

\begin{theorem}[Scalar reduction of the residual two-tensor gauges]%
    \label{thm:peps_twoInjectiveGaugeScalarReduction}
    \lean{TNLean.PEPS.threeLeg_residual_forms_scalar}
    \leanok
    \uses{def:peps_threeLegLeftIdentityForm,
        def:peps_threeLegRightIdentityForm,
        def:peps_threeLegMiddleIdentityForm}
    In the proof of the generalized two-injective-tensor comparison, applying
    the inverse of the second injective tensor leaves three possible virtual
    operators on the exposed legs of the first tensor:
    % Formula (paper_normal.tex:1157--1193): A_1 equals B_1 with a two-leg
    % residual Z, equals B_1 with U on the complementary pair, and equals B_1
    % with W on the first and third legs.  Ink-to-index: each operator box has
    % two incoming and two outgoing virtual ports; the untouched third index is
    % the remaining direct wire.  The north leg is the common physical index.
    % Boundary signature: every panel has (V,P)=(3,1).  Source verdict: exact
    % four-panel scalar-reduction identity in house style.
    \begin{center}
        \tenkzkernel
        \begin{tenkz}[rows={wire}, cols=1, bonds=none]
            \tn[name=scalarA, label pos=w,
                ports={35:virtual, 0:virtual, -35:virtual,
                       90:physical}]{A_1}
        \end{tenkz}
        $ = $
        \begin{tenkz}[rows={wire}, cols=2, bonds=none]
            \tn[name=scalarBZ, label pos=w,
                ports={35:virtual, 0:virtual, -35:virtual,
                       90:physical}]{B_1}
            \tn[at=0.4 s of 1.2 e of scalarBZ, name=scalarZ, skin=box,
                species=operator,
                ports={160:virtual, 200:virtual, 20:virtual,
                       -20:virtual}]{Z}
            \tnwire{scalarBZ.0}{scalarZ.160}
            \tnwire{scalarBZ.-35}{scalarZ.200}
        \end{tenkz}
        $ = $
        \begin{tenkz}[rows={wire}, cols=2, bonds=none]
            \tn[name=scalarBU, label pos=w,
                ports={35:virtual, 0:virtual, -35:virtual,
                       90:physical}]{B_1}
            \tn[at=0.4 n of 1.2 e of scalarBU, name=scalarU, skin=box,
                species=operator,
                ports={160:virtual, 200:virtual, 20:virtual,
                       -20:virtual}]{U}
            \tnwire{scalarBU.35}{scalarU.160}
            \tnwire{scalarBU.0}{scalarU.200}
        \end{tenkz}
        $ = $
        \begin{tenkz}[rows={wire}, cols=2, bonds=none]
            \tn[name=scalarBW, label pos=w,
                ports={35:virtual, 0:virtual, -35:virtual,
                       90:physical}]{B_1}
            \tn[at=0.7 n of 1.2 e of scalarBW, name=scalarW, skin=box,
                species=operator,
                ports={160:virtual, 200:virtual, 20:virtual,
                       -20:virtual}]{W}
            \tnwire{scalarBW.35}{scalarW.160}
            \tnwire{scalarBW.-35}{scalarW.200}
        \end{tenkz}
    \end{center}
    Inverting the first injective tensor then shows that the corresponding
    residual operator on the three exposed virtual legs has the three
    decompositions
    \begin{align}
        \sum_i \Id\otimes Z_i^{(1)}\otimes Z_i^{(2)}
        &=
        \sum_i U_i^{(1)}\otimes U_i^{(2)}\otimes\Id
        =
        \sum_i W_i^{(1)}\otimes\Id\otimes W_i^{(2)}.
        \notag
    \end{align}
    Equivalently, for nonzero virtual leg spaces $\alpha,\beta,\gamma$, if one
    endomorphism of $\alpha\otimes\beta\otimes\gamma$ lies simultaneously in
    the three subspaces $\Id_\alpha\otimes\End(\beta\otimes\gamma)$,
    $\End(\alpha\otimes\beta)\otimes\Id_\gamma$, and the analogous
    subspace with $\Id_\beta$ in the middle leg, then it is a scalar multiple of
    the identity. Thus the residual virtual operators are scalar multiples of
    the identity. This is the scalar step in the proof of
    \cite[Section~3]{Molnar2018NormalPEPS}.
\end{theorem}
\begin{proof}\leanok
    Write the common operator as $T$. The equality
    $T=\Id_\alpha\otimes Z=U\otimes\Id_\gamma$ gives
    $\delta_{aa'}Z_{bc,b'c'}=U_{ab,a'b'}\delta_{cc'}$ for all indices.
    If $c\neq c'$, the right hand side is zero; choosing $a=a'$ gives
    $Z_{bc,b'c'}=0$. Hence $Z$ is diagonal in the $\gamma$-index. Similarly,
    the equality $T=\Id_\alpha\otimes Z=W_{\alpha\gamma}\otimes\Id_\beta$
    gives
    $\delta_{aa'}Z_{bc,b'c'}=W_{ac,a'c'}\delta_{bb'}$,
    so $Z_{bc,b'c'}=0$ whenever $b\neq b'$. Thus $Z_{bc,b'c'}=0$ unless
    $b=b'$ and $c=c'$.

    It remains to compare the surviving diagonal entries. From
    $\delta_{aa'}Z_{bc,bc}=U_{ab,a'b}$, the value
    $Z_{bc,bc}$ is independent of $c$. From the middle-identity equality it is
    also independent of $b$. Therefore
    $Z=\lambda\Id_{\beta\otimes\gamma}$. Substitution into
    \begin{align}
        \Id_\alpha\otimes Z
        &=
        U\otimes\Id_\gamma
        =
        W_{\alpha\gamma}\otimes\Id_\beta
        \notag
    \end{align}
    gives $U=\lambda\Id_{\alpha\otimes\beta}$ and
    $W_{\alpha\gamma}=\lambda\Id_{\alpha\otimes\gamma}$.
\end{proof}

\begin{theorem}[Many-leg scalar extraction]%
    \label{thm:peps_piProductFormsScalar}
    \lean{TNLean.PEPS.piProduct_forms_scalar}
    \leanok
    Let a finite family of legs be given, each a nonzero finite index set, with
    invertible matrices $M_i$ and $N_i$ on each leg. If
    $\prod_i(M_i)_{p_iq_i}=\prod_i(N_i)_{p_iq_i}$ for every configuration pair
    $(p,q)$, then there are nonzero scalars $c_i$ with $M_i=c_iN_i$ for each
    leg and $\prod_i c_i=1$.
\end{theorem}
\begin{proof}\leanok
    Each $N_i$ is invertible, so each of its rows has a nonzero entry; pick a
    reference configuration $(p^0,q^0)$ on which every $(N_i)_{p^0_iq^0_i}$ is
    nonzero. The product identity at $(p^0,q^0)$ then makes every
    $(M_i)_{p^0_iq^0_i}$ nonzero as well. Fix a leg $i$ and a pair $(a,b)$, and
    apply the identity to the configuration that agrees with the reference
    except at leg $i$, where it takes $(a,b)$. Splitting off leg $i$ gives
    $(M_i)_{ab}K_i=(N_i)_{ab}L_i$, where $K_i,L_i$ are the products of the
    other legs at the reference, both nonzero. The reference identity at leg
    $i$ relates $K_i$ and $L_i$, and eliminating them yields
    $(M_i)_{ab}=c_i(N_i)_{ab}$ with
    $c_i=(M_i)_{p^0_iq^0_i}/(N_i)_{p^0_iq^0_i}$, independent of $(a,b)$.
    Taking the product of the scalars over all legs and using the reference
    identity gives $\prod_i c_i=1$.
\end{proof}

\begin{theorem}[Pointwise product cancellation for two injective blocks]%
    \label{thm:peps_twoBlockPointwiseProductCancellation}
    \lean{TNLean.PEPS.twoBlockReciprocalScalarProportional_of_pointwise_mul_eq}
    \leanok
    Let $A_1,A_2,B_1,B_2$ be tensors with matching external,
    shared-boundary, and physical index sets. Assume that the two external
    boundary spaces and the shared-boundary configuration space are nonempty,
    and that $A_1$ and $A_2$ are injective. Suppose that, for every pair of
    external indices, every pair of shared-boundary configurations, and every
    pair of physical indices,
    \begin{align}
        A_1(\eta_1,\mu,\sigma_1)A_2(\eta_2,\nu,\sigma_2)
        &=
        B_1(\eta_1,\mu,\sigma_1)B_2(\eta_2,\nu,\sigma_2).
        \notag
    \end{align}
    Then there exists $\lambda\in\C^\times$ such that
    $A_1=\lambda B_1$ and $A_2=\lambda^{-1}B_2$.
\end{theorem}
% This theorem is the rank-one cancellation step once
% \cite[Section~3]{Molnar2018NormalPEPS} has been reduced to pointwise product
% equality of separated coefficients.
\begin{proof}\leanok
    Choose external and shared-boundary indices, and use injectivity of
    $A_1$ and $A_2$ to choose physical indices for which the two chosen
    coefficients are nonzero. The product equality then shows that the
    corresponding coefficients of $B_1$ and $B_2$ are also nonzero. Define
    $\lambda=
        B_2(\eta_2^0,\nu^0,\sigma_2^0)/
        A_2(\eta_2^0,\nu^0,\sigma_2^0)$.
    For every $(\eta_1,\mu,\sigma_1)$, the product identity with
    $(\eta_2^0,\nu^0,\sigma_2^0)$ fixed gives
    \begin{align}
        A_1(\eta_1,\mu,\sigma_1)
        A_2(\eta_2^0,\nu^0,\sigma_2^0)
        &=
        B_1(\eta_1,\mu,\sigma_1)
        B_2(\eta_2^0,\nu^0,\sigma_2^0).
        \notag
    \end{align}
    Dividing by $A_2(\eta_2^0,\nu^0,\sigma_2^0)\neq0$ gives
    $A_1=\lambda B_1$. Fixing $(\eta_1^0,\mu^0,\sigma_1^0)$ and substituting
    $A_1=\lambda B_1$ into the product identity gives
    \begin{align}
        \lambda B_1(\eta_1^0,\mu^0,\sigma_1^0)
        A_2(\eta_2,\nu,\sigma_2)
        &=
        B_1(\eta_1^0,\mu^0,\sigma_1^0)
        B_2(\eta_2,\nu,\sigma_2).
        \notag
    \end{align}
    Dividing by $B_1(\eta_1^0,\mu^0,\sigma_1^0)\neq0$ gives
    $A_2=\lambda^{-1}B_2$.
\end{proof}

\begin{definition}[Matrix unit for coefficient extraction]%
    \label{def:peps_matrixUnit}
    \lean{TNLean.PEPS.matrixUnit}
    \leanok
    For indices $i$ and $j$, let $E_{ij}$ denote the matrix whose only
    nonzero entry is the $(i,j)$ entry, where it is equal to $1$.
\end{definition}

\begin{theorem}[Coefficient extraction for one shared bond]%
    \label{thm:peps_twoBlockInsertedCoeff_singletonBond}
    \lean{TNLean.PEPS.twoBlockInsertedCoeff_singletonBond_single}
    \leanok
    \uses{def:peps_matrixUnit}
    Suppose that the two blocks share a single virtual bond. Then inserting
    the matrix unit $E_{\mu,\nu}$ in that bond extracts the corresponding
    product of tensor coefficients:
    \begin{align}
        C_{A_1,A_2}
        (b,E_{\mu,\nu};\eta_1,\eta_2,\sigma_1,\sigma_2)
        &=
        A_1(\eta_1,\mu,\sigma_1)A_2(\eta_2,\nu,\sigma_2).
        \notag
    \end{align}
\end{theorem}
\begin{proof}\leanok
    Since there is only one shared bond, the condition that the two shared
    configurations agree away from the distinguished bond is vacuous, and
    \begin{align}
        C_{A_1,A_2}
        (b,E_{\mu,\nu};\eta_1,\eta_2,\sigma_1,\sigma_2)
        &=
        \sum_{\mu',\nu'}
        (E_{\mu,\nu})_{\mu'(b),\nu'(b)}
        A_1(\eta_1,\mu',\sigma_1)A_2(\eta_2,\nu',\sigma_2).
        \notag
    \end{align}
    The entries of the matrix unit satisfy
    $(E_{\mu,\nu})_{i,j}
        =\delta_{i,\mu(b)}\delta_{j,\nu(b)}$.
    Thus the only nonzero summand is the one with
    $(\mu',\nu')=(\mu,\nu)$.
\end{proof}

\begin{theorem}[One-bond two-injective comparison]%
    \label{thm:peps_twoInjectiveTensorInsertionComparison_singletonBond}
    \lean{TNLean.PEPS.two_injective_tensor_insertion_comparison_singletonBond}
    \leanok
    \uses{thm:peps_twoBlockInsertedCoeff_singletonBond,
        thm:peps_twoBlockPointwiseProductCancellation}
    Suppose that the two injective blocks have exactly one shared virtual
    bond and nonempty external boundary spaces. If all matrix insertions on
    that bond give equal two-block coefficients for the pairs
    $(A_1,A_2)$ and $(B_1,B_2)$, then there exists
    $\lambda\in\C^\times$ such that
    $A_1=\lambda B_1$ and $A_2=\lambda^{-1}B_2$.
\end{theorem}
\begin{proof}\leanok
    Insert the matrix unit $E_{\mu,\nu}$ and use the coefficient-extraction
    theorem to obtain the pointwise product equality
    \begin{align}
        A_1(\eta_1,\mu,\sigma_1)A_2(\eta_2,\nu,\sigma_2)
        &=
        B_1(\eta_1,\mu,\sigma_1)B_2(\eta_2,\nu,\sigma_2)
        \notag
    \end{align}
    for all indices. The pointwise product cancellation theorem gives the
    reciprocal scalar proportionality.
\end{proof}

\begin{theorem}[Generalized two-injective-tensor comparison]%
    \label{thm:peps_twoInjectiveTensorInsertionComparison}
    \lean{TNLean.PEPS.two_injective_tensor_insertion_comparison}
    \lean{TNLean.PEPS.SameAwayFromBond, TNLean.PEPS.twoBlockInsertedCoeff}
    \lean{TNLean.PEPS.IsTwoBlockInjective, TNLean.PEPS.SameTwoBlockInsertions}
    \lean{TNLean.PEPS.TwoBlockScalarProportional}
    \lean{TNLean.PEPS.TwoBlockReciprocalScalarProportional}
    \leanok
    \uses{thm:peps_twoInjectiveGaugeScalarReduction}
    Let $A_1,A_2,B_1,B_2$ be injective tensors with the same shared virtual
    bonds and nonempty external boundary spaces, and assume that there is at
    least one shared virtual bond. For the set $\mathcal B$ of shared bonds, a
    bond $b\in\mathcal B$, and $X\in\MN{D_b}$, set
    \begin{align}
        C_{A_1,A_2}(b,X;\eta_1,\eta_2,\sigma_1,\sigma_2)
        &=
        \sum_{\substack{\mu,\nu\\
        \mu|_{\mathcal B\setminus\{b\}}
        =\nu|_{\mathcal B\setminus\{b\}}}}
        X_{\mu_b,\nu_b}
        A_1(\eta_1,\mu,\sigma_1)A_2(\eta_2,\nu,\sigma_2).
        \notag
    \end{align}
    Assume, for every $b,X,\eta_1,\eta_2,\sigma_1,\sigma_2$, that
    \begin{align}
        C_{A_1,A_2}(b,X;\eta_1,\eta_2,\sigma_1,\sigma_2)
        &=
        C_{B_1,B_2}(b,X;\eta_1,\eta_2,\sigma_1,\sigma_2).
        \notag
    \end{align}
    The following diagram records the three-shared-bond proof subcase used in
    the source, with the external boundary indices restored:
    % Formula: the two inserted two-tensor coefficients agree for every shared
    % bond b and matrix X.  Ink-to-index: in the three-shared-bond proof subcase,
    % X splits the upper contraction, the two north legs are sigma_1 and sigma_2,
    % and the outer stubs expose eta_1 and eta_2.  Boundary signature: both panels
    % have (V,P)=(2,2).  Source verdict: paper_normal.tex:1093--1094 explicitly
    % reduces the arbitrary shared-bond family of eq. (9), lines 1067--1091, to
    % three shared bonds without loss of generality; the theorem's aggregate
    % external virtual boundaries are restored here.
    \begin{center}
        \tenkzkernel
        \begin{tenkz}[rows={wire}, cols=4, bonds=none]
            \tn[at=(1,1), name=twoInjAOne, label pos=s,
                ports={180:virtual:$\eta_1$, 20:virtual, 0:virtual,
                       -20:virtual, 90:physical}]{A_1}
            \tn[at=(1,4), name=twoInjATwo, label pos=s,
                ports={0:virtual:$\eta_2$, 160:virtual, 180:virtual,
                       200:virtual, 90:physical}]{A_2}
            \tn[at=0.45 n of midway twoInjAOne and twoInjATwo,
                name=twoInjAX, skin=ring, species=operator,
                ports={200:virtual, -20:virtual}]{X}
            \tn[at=0.45 s of midway twoInjAOne and twoInjATwo,
                name=twoInjALow, skin=none]{}
            \tnwire{twoInjAOne.20}{twoInjAX.200}
            \tnwire{twoInjAX.-20}{twoInjATwo.160}
            \tnwire{twoInjAOne.0}{twoInjATwo.180}
            \tnwire{twoInjAOne.-20}{twoInjALow}
            \tnwire{twoInjALow}{twoInjATwo.200}
        \end{tenkz}
        $ = $
        \begin{tenkz}[rows={wire}, cols=4, bonds=none]
            \tn[at=(1,1), name=twoInjBOne, label pos=s,
                ports={180:virtual:$\eta_1$, 20:virtual, 0:virtual,
                       -20:virtual, 90:physical}]{B_1}
            \tn[at=(1,4), name=twoInjBTwo, label pos=s,
                ports={0:virtual:$\eta_2$, 160:virtual, 180:virtual,
                       200:virtual, 90:physical}]{B_2}
            \tn[at=0.45 n of midway twoInjBOne and twoInjBTwo,
                name=twoInjBX, skin=ring, species=operator,
                ports={200:virtual, -20:virtual}]{X}
            \tn[at=0.45 s of midway twoInjBOne and twoInjBTwo,
                name=twoInjBLow, skin=none]{}
            \tnwire{twoInjBOne.20}{twoInjBX.200}
            \tnwire{twoInjBX.-20}{twoInjBTwo.160}
            \tnwire{twoInjBOne.0}{twoInjBTwo.180}
            \tnwire{twoInjBOne.-20}{twoInjBLow}
            \tnwire{twoInjBLow}{twoInjBTwo.200}
        \end{tenkz}
    \end{center}
    Then there exists $\lambda\in\C^\times$ such that
    $A_1=\lambda B_1$ and $A_2=\lambda^{-1}B_2$.
    This is the blueprint form of the two-injective tensor comparison in
    \cite[Section~3]{Molnar2018NormalPEPS}.
\end{theorem}
\begin{proof}
    \leanok
    Putting $X=\Id_{D_b}$ on every shared bond contracts all bonds diagonally and,
    by the insertion equality, equates the two fully contracted states. Writing
    each contracted state as a bipartite operator with the linearly independent
    Schmidt families supplied by injectivity, the operator-Schmidt uniqueness of
    \cite[Section~3]{Molnar2018NormalPEPS} produces an invertible gauge $g$ on the
    shared-bond configurations with
    \begin{align}
        A_1(\eta_1,\mu,\sigma_1)
        &=
        \sum_\nu g_{\mu,\nu}B_1(\eta_1,\nu,\sigma_1),
        \notag\\
        B_2(\eta_2,\nu,\sigma_2)
        &=
        \sum_\mu g_{\mu,\nu}A_2(\eta_2,\mu,\sigma_2).
        \notag
    \end{align}
    Substitute both gauge equations into the insertion equality for the
    matrix unit $X=E_{p,q}$ on a single bond $b$, which frees that bond and
    contracts the others. Separating the linearly independent families $B_1$ and
    $A_2$ gives, for every bond $b$, every pair of endpoints $p,q$, and all
    configurations $\nu,\rho$,
    \begin{align}
        [\rho_b=q]\,g_{\rho^{b\to p},\nu}
        &=
        [\nu_b=p]\,g_{\rho,\nu^{b\to q}}.
        \label{eq:peps_residual_equality}
    \end{align}
    Here $\rho^{b\to p}$ replaces the value of $\rho$ at $b$ by $p$. This is the
    bond-wise form of the residual identity of
    \cite[Section~3]{Molnar2018NormalPEPS}, where it is packaged as the equality
    $\Id_\alpha\otimes Z
        =U\otimes\Id_\gamma
        =W_{\alpha\gamma}\otimes\Id_\beta$ of the
    three residual gauges $Z$, $U$, $W$
    (Theorem~\ref{thm:peps_twoInjectiveGaugeScalarReduction}). Reading
    (\ref{eq:peps_residual_equality}) at the two endpoint choices shows that $g$
    vanishes whenever two configurations differ at a bond and that its diagonal
    value is invariant under changing any single bond. Hence
    $g=\lambda\Id$ for a scalar $\lambda$, and invertibility forces
    $\lambda\neq0$.

    Reinserting $g=\lambda\Id$ into the two gauge equations gives
    $A_1=\lambda B_1$ and $B_2=\lambda A_2$, that is
    $A_2=\lambda^{-1}B_2$.
\end{proof}


\section{Fundamental Theorem for Injective PEPS}\label{sec:peps_ft}

The edge construction is transported from single edges to arbitrary regions.
Reading a region $R$ and its complement $V\setminus R$ as two-block tensors,
when the blocked tensors of $B$ on $R$ and $V\setminus R$ are injective and the
bonds are positive, the region-insertion identities determine the boundary
conjugation uniquely: two invertible matrices $Z_1,Z_2$ realizing the same
insertion identity satisfy, for every boundary matrix $M$,
\[
    Z_1\,M\,Z_1^{-1}=Z_2\,M\,Z_2^{-1}.
\]
Comparing one vertex $v$ against its complement
then upgrades this to $A_v=\lambda_v\widetilde B_v$ at each vertex, once the edge
gauges have been absorbed into $\widetilde B$, and the per-edge gauges combine
into one gauge relating the two tensors. This proves the Fundamental Theorem for
injective PEPS, Theorem~\ref{thm:fundamentalTheorem_PEPS}, the formalization of
\cite[Theorem~2]{Molnar2018NormalPEPS}.

\begin{definition}[Region two-block tensor]%
    \label{def:peps_regionTwoBlock}
    \lean{TNLean.PEPS.regionTwoBlock}
    \leanok
    For a finite region $R\subseteq V$, the blocked tensor of $R$ may be read
    as a two-block tensor with a one-point external boundary and shared bonds
    the edges crossing the boundary of $R$:
    \[
        B_{A,R}(*,\eta,\sigma)=T_{A,R}^{\eta}(\sigma).
    \]
    This is the region block used in the one-block-against-complement step of
    \cite[Section~3]{Molnar2018NormalPEPS}.
\end{definition}

\begin{definition}[Complement-boundary identification]%
    \label{def:peps_regionComplementBoundaryConfig}
    \lean{TNLean.PEPS.regionBoundaryEdgeToCompl,
        TNLean.PEPS.regionBoundaryEdgeComplEquiv,
        TNLean.PEPS.regionComplementBoundaryConfig}
    \leanok
    The edges crossing the boundary of $R$ are the same edges crossing the
    boundary of $V\setminus R$. Thus a boundary configuration $\eta$ for $R$
    determines the boundary configuration $\eta^c$ for the complement by reading
    the same edge labels through this identification.
\end{definition}

\begin{definition}[Complement-region two-block tensor]%
    \label{def:peps_regionComplementTwoBlock}
    \lean{TNLean.PEPS.regionComplementTwoBlock}
    \leanok
    \uses{def:peps_regionComplementBoundaryConfig}
    For a finite region $R\subseteq V$, the complement block $V\setminus R$ is
    a two-block tensor over the same crossing bonds:
    \[
        B_{A,R^c}(*,\eta,\tau)
        =
        T_{A,V\setminus R}^{\eta^c}(\tau).
    \]
    This is the complementary block in the region/complement comparison of
    \cite[Section~3]{Molnar2018NormalPEPS}.
\end{definition}

\begin{theorem}[Injectivity of the complement-region two-block tensor]%
    \label{thm:peps_isTwoBlockInjective_regionComplementTwoBlock}
    \lean{TNLean.PEPS.isTwoBlockInjective_regionComplementTwoBlock}
    \leanok
    \uses{def:peps_regionComplementTwoBlock, def:peps_regionInjectivityData}
    If the blocked tensor family of $V\setminus R$ is injective, then the
    complement-region two-block tensor is injective:
    \[
        \operatorname{inj}(T_{A,V\setminus R})
        \Longrightarrow
        \operatorname{inj}(B_{A,R^c}).
    \]
\end{theorem}
\begin{proof}\leanok
    The one-point external boundary does not change the coefficient family, and
    the complement-boundary identification only reindexes the crossing bonds.
    More explicitly, if $\eta_1,\eta_2$ are boundary configurations on
    $\partial R$, then
    \[
        \eta_1^c=\eta_2^c\quad\Longrightarrow\quad \eta_1=\eta_2 .
    \]
    Thus the coefficients of $T_{A,V\setminus R}$ and $B_{A,R^c}$ differ only
    by an injective relabelling of boundary configurations.
\end{proof}

\begin{definition}[Region insertion coefficient]%
    \label{def:peps_regionInsertedCoeff}
    \lean{TNLean.PEPS.regionInsertedCoeff}
    \leanok
    \uses{def:peps_regionTwoBlock, def:peps_regionComplementTwoBlock}
    Let $f$ be an edge crossing the boundary of $R$, and let
    $M\in\MN{D_f}$.  The region insertion coefficient is
    \[
        C_A(R,f,M;\sigma,\tau)
        =
        \sum_{\substack{\mu,\nu\\
        \mu|_{\partial R\setminus\{f\}}
        =\nu|_{\partial R\setminus\{f\}}}}
        M_{\mu_f,\nu_f}\,
        T_{A,R}^{\mu}(\sigma)\,
        T_{A,V\setminus R}^{\nu^c}(\tau).
    \]
    It is the coefficient obtained by inserting $M$ on the boundary bond $f$
    between the region $R$ and its complement.
\end{definition}

\begin{theorem}[Region insertion coefficients determine the inserted matrix]%
    \label{thm:peps_regionInsertedCoeff_injective}
    \lean{TNLean.PEPS.regionInsertedCoeff_injective}
    \leanok
    \uses{def:peps_regionInsertedCoeff}
    Suppose that the blocked tensor maps of $A$ on $R$ and on
    $V\setminus R$ are injective, and that every bond dimension of $A$ is
    positive.  If $f\in\partial R$ and
    $M,M'\in\MN{D_A(f)}$ satisfy
    \[
        C_A(R,f,M;\sigma,\tau)
        =
        C_A(R,f,M';\sigma,\tau)
    \]
    for every physical configuration $\sigma$ on $R$ and $\tau$ on
    $V\setminus R$, then $M=M'$.
\end{theorem}
\begin{proof}
    \leanok
    Subtracting the two coefficients gives
    \[
        C_A(R,f,M-M';\sigma,\tau)=0
    \]
    for every $\sigma,\tau$.  Injectivity of the region and complement
    blocked tensor maps reads off every entry of $M-M'$ from configurations
    agreeing away from $f$, so $M-M'=0$ and hence $M=M'$.
\end{proof}

\begin{definition}[Bond-inserted region insert]%
    \label{def:peps_bond_inserted_region_insert}
    \lean{TNLean.PEPS.bondInsertedRegionInsert}
    \leanok
    Let $f$ be an edge crossing the boundary of $R$, and let
    $M\in\MN{D_f}$.  The bond-inserted region insert
    $I_{A,R,f,M}$ is the insert on $R$ defined by
    \[
        (I_{A,R,f,M})_{\nu,\sigma}
        =
        \sum_{\substack{\mu\in\partial R\\
        \mu|_{\partial R\setminus\{f\}}
        =\nu|_{\partial R\setminus\{f\}}}}
        M_{\mu_f,\nu_f}\,T_{A,R}^{\mu}(\sigma).
    \]
    Thus the complement boundary value $\nu$ is fixed away from $f$, while
    the matrix $M$ couples the two $f$-labels.  This is the one-bond insertion
    in the region/complement coefficient of
    \cite[Section~3]{Molnar2018NormalPEPS}.
\end{definition}

\begin{theorem}[Bond insertion as a deformed region state]
    \label{thm:peps_deformedRegionState_bondInsertedRegionInsert}
    \lean{TNLean.PEPS.deformedRegionState_bondInsertedRegionInsert}
    \leanok
    \uses{def:peps_bond_inserted_region_insert,
        def:peps_regionInsertedCoeff}
    For every physical configuration $\sigma$ on $R$ and $\tau$ on
    $V\setminus R$, the deformed state of the bond-inserted region insert is
    the region insertion coefficient:
    \[
        \Psi_{A,R,I_{A,R,f,M}}(\sigma,\tau)
        =
        C_A(R,f,M;\sigma,\tau).
    \]
\end{theorem}
\begin{proof}
    \leanok
    Expanding the deformed state gives
    \[
        \Psi_{A,R,I_{A,R,f,M}}(\sigma,\tau)
        =
        \sum_{\nu\in\partial R}
        I_{A,R,f,M}(\nu,\sigma)\,
        T_{A,V\setminus R}^{\nu^c}(\tau).
    \]
    Substituting the definition of $I_{A,R,f,M}$ and exchanging the two finite
    sums gives
    \[
    \begin{aligned}
        \Psi_{A,R,I_{A,R,f,M}}(\sigma,\tau)
        &=
        \sum_{\nu\in\partial R}
        \left(
        \sum_{\substack{\mu\in\partial R\\
        \mu|_{\partial R\setminus\{f\}}
        =\nu|_{\partial R\setminus\{f\}}}}
        M_{\mu_f,\nu_f}\,T_{A,R}^{\mu}(\sigma)
        \right)
        T_{A,V\setminus R}^{\nu^c}(\tau)\\
        &=
        C_A(R,f,M;\sigma,\tau).
    \end{aligned}
    \]
\end{proof}

\begin{theorem}[Region insertion as an edge insertion]
    \label{thm:peps_regionInsertedCoeff_eq_smul_edgeInsertedCoeff}
    \lean{TNLean.PEPS.regionInsertedCoeff_eq_smul_edgeInsertedCoeff}
    \leanok
    \uses{def:peps_regionInsertedCoeff, def:peps_edgeInsertedCoeff}
    Let $f$ be a boundary edge of $R$, let $N\in\MN{D_B(f)}$, let $\sigma$
    be a physical configuration on $R$, and let $\tau$ be a physical
    configuration on $V\setminus R$.  Put
    \[
        p_B(R)=\prod_{g\notin\partial R}D_B(g).
    \]
    Then the region insertion coefficient factors through the edge-inserted
    coefficient:
    \[
        C_B(R,f,N;\sigma,\tau)
        =
        p_B(R)\,
        E_B\bigl(f,\operatorname{Assemble}_R(\sigma,\tau),
            \mathcal O_{R,f}(N)\bigr).
    \]
    Here $\mathcal O_{R,f}$ is the boundary-edge orientation, equal to the
    identity if the left endpoint of $f$ lies in $R$ and to transpose
    otherwise.
\end{theorem}
\begin{proof}
    \leanok
    Let $\mathcal C_{R,f}(\sigma,\tau)$ denote the open-edge configurations
    compatible with the two physical configurations and with the two boundary
    readings agreeing away from $f$.  For
    $\xi\in\mathcal C_{R,f}(\sigma,\tau)$, let
    $\operatorname{Fib}_{R,f}(\xi)$ be the corresponding fibre of agreeing
    boundary readings.  Grouping the expansion of $C_B(R,f,N;\sigma,\tau)$ by
    such an open-edge configuration $\xi$ gives
    \[
    \begin{aligned}
        C_B(R,f,N;\sigma,\tau)
        &=
        \sum_{\xi\in\mathcal C_{R,f}(\sigma,\tau)}
        \#\operatorname{Fib}_{R,f}(\xi)\,
        N_{\xi_{R,f}^{(1)},\xi_{R,f}^{(2)}}
        \prod_{v\in V}
        B_v\bigl(\xi_v,\operatorname{Assemble}_R(\sigma,\tau)_v\bigr),\\
        \#\operatorname{Fib}_{R,f}(\xi)
        &=
        \prod_{g\notin\partial R}D_B(g)
        =
        p_B(R).
    \end{aligned}
    \]
    Substituting this fibre cardinality gives
    \[
    \begin{aligned}
        C_B(R,f,N;\sigma,\tau)
        &=
        p_B(R)
        \sum_{\xi\in\mathcal C_{R,f}(\sigma,\tau)}
        N_{\xi_{R,f}^{(1)},\xi_{R,f}^{(2)}}
        \prod_{v\in V}
        B_v\bigl(\xi_v,\operatorname{Assemble}_R(\sigma,\tau)_v\bigr)\\
        &=
        p_B(R)\,
        E_B\bigl(f,\operatorname{Assemble}_R(\sigma,\tau),
            \mathcal O_{R,f}(N)\bigr).
    \end{aligned}
    \]
    The second equality is the definition of the edge-inserted coefficient,
    with the inserted matrix read in the boundary-edge orientation.
\end{proof}

\begin{theorem}[Bond insertion through a corner extension]
    \label{thm:peps_regionInsertedCoeff_eq_extendInsert_bondInserted}
    \lean{TNLean.PEPS.regionInsertedCoeff_eq_extendInsert_bondInserted}
    \leanok
    \uses{thm:peps_deformedRegionState_bondInsertedRegionInsert,
        thm:peps_corner_extension_state_preserves}
    Let $R\subseteq S$, and suppose that all bond dimensions are positive.
    For every global physical configuration $\sigma$,
    \[
        C_A\bigl(R,f,M;\sigma|_R,\sigma|_{V\setminus R}\bigr)
        =
        \Psi_{A,S,
        \operatorname{Ext}_{R\subseteq S}(I_{A,R,f,M})}(\sigma).
    \]
    In particular, the one-bond insertion on $R$ may be read after extending
    the corresponding insert to any larger region $S$.
\end{theorem}
\begin{proof}
    \leanok
    The preceding theorem and state preservation for corner extension give
    \[
    \begin{aligned}
        C_A\bigl(R,f,M;\sigma|_R,\sigma|_{V\setminus R}\bigr)
        &=
        \Psi_{A,R,I_{A,R,f,M}}
            (\sigma|_R,\sigma|_{V\setminus R})\\
        &=
        \Psi_{A,S,
            \operatorname{Ext}_{R\subseteq S}(I_{A,R,f,M})}(\sigma).
    \end{aligned}
    \]
\end{proof}

\begin{theorem}[Region insertion as a two-block insertion]%
    \label{thm:peps_twoBlockInsertedCoeff_eq_regionInsertedCoeff}
    \lean{TNLean.PEPS.twoBlockInsertedCoeff_eq_regionInsertedCoeff}
    \leanok
    \uses{def:peps_regionInsertedCoeff, def:peps_regionTwoBlock,
        def:peps_regionComplementTwoBlock}
    The abstract two-block insertion coefficient of the pair
    $(B_{A,R},B_{A,R^c})$ is the region insertion coefficient:
    \[
        C_{B_{A,R},B_{A,R^c}}(f,M;*,*,\sigma,\tau)
        =
        C_A(R,f,M;\sigma,\tau).
    \]
\end{theorem}
\begin{proof}\leanok
    Expanding the abstract two-block coefficient gives exactly the two boundary
    sums in the displayed formula for $C_A(R,f,M;\sigma,\tau)$.
\end{proof}

\begin{theorem}[Region insertion under bond-dimension transport]%
    \label{thm:peps_regionInsertedCoeff_reindexTensor}
    \lean{TNLean.PEPS.regionInsertedCoeff_reindexTensor}
    \leanok
    \uses{def:peps_regionInsertedCoeff}
    If a tensor $B$ is transported along an equality of bond dimensions, then
    the region insertion coefficient is transported by the induced
    matrix-algebra identification on the inserted bond:
    \[
        C_{B^h}(R,f,N;\sigma,\tau)
        =
        C_B(R,f,\phi_f(N);\sigma,\tau).
    \]
\end{theorem}
\begin{proof}\leanok
    Reindex both boundary-configuration sums by the bond-dimension
    equivalences. The constraint away from $f$ and the two blocked-region
    weights are preserved, while the inserted matrix entry becomes the
    transported entry $\phi_f(N)$.
\end{proof}

\begin{theorem}[Same two-block insertions from region insertions]%
    \label{thm:peps_sameTwoBlockInsertions_of_regionInsertedCoeff_eq}
    \lean{TNLean.PEPS.sameTwoBlockInsertions_of_regionInsertedCoeff_eq}
    \leanok
    \uses{thm:peps_twoBlockInsertedCoeff_eq_regionInsertedCoeff,
        thm:peps_regionInsertedCoeff_reindexTensor}
    Suppose $A$ and $B$ have identified bond dimensions. Write $B^h$ for the
    transported tensor, and write $\phi_f$ for the induced transport on the
    boundary bond $f$. If
    \[
        C_A(R,f,N;\sigma,\tau)
        =
        C_B(R,f,\phi_f(N);\sigma,\tau)
    \]
    for every boundary edge $f$, matrix $N$, and physical configurations
    $\sigma,\tau$, then the two pairs
    \[
        (B_{A,R},B_{A,R^c}),\qquad
        (B_{B^h,R},B_{B^h,R^c})
    \]
    have the same one-bond insertion coefficients after the same
    bond-dimension transport.
\end{theorem}
\begin{proof}\leanok
    Rewrite the two abstract insertion coefficients as region insertion
    coefficients and use the bond-dimension transport formula:
    \[
        C_{B_{A,R},B_{A,R^c}}(f,N;*,*,\sigma,\tau)
        = C_A(R,f,N;\sigma,\tau)
        = C_B(R,f,\phi_f(N);\sigma,\tau)
        = C_{B_{B^h,R},B_{B^h,R^c}}(f,N;*,*,\sigma,\tau).
    \]
\end{proof}

\begin{theorem}[Swapping the two blocks in a one-bond insertion]
    \label{thm:peps_twoBlockInsertedCoeff_swap}
    \lean{TNLean.PEPS.twoBlockInsertedCoeff_swap}
    \leanok
    Let \(B_1,B_2\) be two blocks joined along a shared bond \(f\).  Then inserting
    a matrix \(X\) while contracting \(B_1\) against \(B_2\) is the same as
    contracting \(B_2\) against \(B_1\) with the transposed insertion:
    \[
        C_{B_1,B_2}(f,X;\eta_1,\eta_2,\sigma_1,\sigma_2)
        =
        C_{B_2,B_1}(f,X^T;\eta_2,\eta_1,\sigma_2,\sigma_1).
    \]
\end{theorem}
\begin{proof}\leanok
    Exchange the two boundary-configuration sums.  The condition that the two
    configurations agree away from \(f\) is symmetric, and
    \(X^T_{\nu_f,\mu_f}=X_{\mu_f,\nu_f}\).
\end{proof}

\begin{theorem}[Complement-side reading of a region insertion]
    \label{thm:peps_regionInsertedCoeff_complement_reading}
    \lean{TNLean.PEPS.regionInsertedCoeff_eq_complementTwoBlock}
    \leanok
    \uses{thm:peps_twoBlockInsertedCoeff_swap,
        thm:peps_twoBlockInsertedCoeff_eq_regionInsertedCoeff}
    For a boundary edge \(f\in\partial R\), the region-inserted coefficient can be
    read by contracting the complement block first:
    \[
        C_A(R,f,X;\sigma,\tau)
        =
        C_{A,R^c;R}(f,X^T;\tau,\sigma).
    \]
    Here \(C_{A,R^c;R}\) denotes the same one-bond insertion coefficient with the
    two blocks \(R^c\) and \(R\) interchanged.
\end{theorem}
\begin{proof}\leanok
    Identify \(C_A(R,f,X;\sigma,\tau)\) with the two-block insertion of
    \((B_{A,R},B_{A,R^c})\), then apply
    Theorem~\ref{thm:peps_twoBlockInsertedCoeff_swap}.
\end{proof}

\begin{theorem}[Incident form gives the transferred endpoint form]
    \label{thm:peps_hform_of_region_incident_form}
    \lean{TNLean.PEPS.hform_of_isIncidentMatrixForm}
    \leanok
    \uses{thm:peps_physical_to_virtual_insertion}
    Assume that \(A\)'s and \(B\)'s components at the in-region endpoint of \(f\)
    are linearly independent, and that \(D_B(e)>0\) for every edge \(e\).
    Let \(W\) be the virtual pullback, at the in-region endpoint of \(f\), of the
    physical operator obtained by inserting \(X^T\) on \(f\).  If
    \[
        W=\mathcal I_f(P)
    \]
    is in incident-matrix form on the boundary leg \(f\), then the transfer matrix
    read off from \(W\) reinserts to \(W\).
\end{theorem}
\begin{proof}\leanok
    The transfer matrix is defined by reading the incident matrix out of \(W\).
    Reading \(\mathcal I_f(P)\) gives \(P\), and reinserting \(P\) returns \(W\).
\end{proof}

\begin{theorem}[Virtual pullback realizes the endpoint operator]
    \label{thm:peps_region_insertion_virtual_pullback_realizes}
    \lean{TNLean.PEPS.localTensorMap_localVirtualOpOfPhysicalOpAt_regionInsertionOp}
    \leanok
    \uses{def:peps_same_state, def:IsVertexInjective, def:peps_localTensorMap,
        thm:peps_physical_to_virtual_insertion}
    Assume that \(A\) and \(B\) have the same PEPS state, are vertex-injective,
    and have positive bond dimensions.  Let \(v\) be the in-region endpoint of
    \(f\), and let \(W\) be the virtual pullback, for \(B\) at \(v\), of the
    physical operator obtained by inserting \(X^T\) on \(f\).  Then for every
    local virtual coefficient \(c\),
    \[
        \mathcal T_{B,v}(Wc)
        =
        O_{A,R,f}(X^T)\,\mathcal T_{B,v}(c).
    \]
\end{theorem}
\begin{proof}\leanok
    The physical endpoint operator preserves the image of
    \(\mathcal T_{B,v}\). The chosen left inverse defining the virtual pullback
    therefore realizes the same action after applying \(\mathcal T_{B,v}\).
\end{proof}

\begin{definition}[Out-of-region endpoint and blocked-region inverse]
    \label{def:peps_region_out_endpoint_block_inverse}
    \lean{TNLean.PEPS.regionBoundaryEdgeOutVertex,
        TNLean.PEPS.regionBoundaryEdgeOutVertex_not_mem,
        TNLean.PEPS.regionBoundaryEdgeOutVertex_mem_compl,
        TNLean.PEPS.regionBoundaryEdgeInVertex_compl_eq_outVertex,
        TNLean.PEPS.regionBlockedTensorMap,
        TNLean.PEPS.regionBlockedTensorMap_injective_of_injective,
        TNLean.PEPS.regionBlockedTensorMap_ker_eq_bot_of_injective,
        TNLean.PEPS.regionBlockedLeftInverse,
        TNLean.PEPS.regionBlockedLeftInverse_comp_regionBlockedTensorMap,
        TNLean.PEPS.regionBlockedLeftInverse_apply_regionBlockedTensorMap,
        TNLean.PEPS.regionBlockedTensorMap_apply,
        TNLean.PEPS.regionBlockedTensorMap_single,
        TNLean.PEPS.regionBlockedLeftInverse_regionBlockedWeight}
    \leanok
    \uses{def:peps_regionInjectivityData}
    The out-of-region endpoint of \(f\in\partial R\) is the endpoint outside
    \(R\).  It is the in-region endpoint of the same edge viewed as an edge of
    \(\partial(V\setminus R)\).  The blocked-region map is
    \[
        c\longmapsto \sum_{\mu\in\partial R} c_\mu\,T_{A,R}^{\mu},
    \]
    and injectivity of the blocked tensor gives a left inverse.  In particular,
    \[
        L_{A,R}\!\left(T_{A,R}^{\mu}\right)=e_\mu .
    \]
\end{definition}

\begin{theorem}[Complement-side cast and endpoint realization]
    \label{thm:peps_regionInsertedCoeff_complement_cast}
    \lean{TNLean.PEPS.mem_compl_compl_iff,
        TNLean.PEPS.regionDoubleComplVertexEquiv,
        TNLean.PEPS.regionDoubleComplPhysicalConfig,
        TNLean.PEPS.isRegionBoundaryEdge_compl_compl_iff,
        TNLean.PEPS.regionBoundaryEdgeDoubleComplEquiv,
        TNLean.PEPS.regionDoubleComplBoundaryConfig,
        TNLean.PEPS.regionBlockedWeight_doubleCompl,
        TNLean.PEPS.regionComplementBoundaryConfigEquiv,
        TNLean.PEPS.regionComplementBoundaryConfig_apply_toCompl,
        TNLean.PEPS.regionComplementBoundaryConfig_compl_compl,
        TNLean.PEPS.sameAwayFromBond_complementBoundaryConfig,
        TNLean.PEPS.regionInsertedCoeff_eq_compl,
        TNLean.PEPS.regionInsertedCoeff_eq_regionRealizationSum_compl,
        TNLean.PEPS.regionInsertedCoeff_eq_smul_op_regionStateVec_compl}
    \leanok
    \uses{def:peps_region_out_endpoint_block_inverse,
        thm:peps_regionInsertedCoeff_complement_reading}
    With \(f'\) the boundary edge \(f\) read on \(V\setminus R\), one has
    \[
        C_A(R,f,X;\sigma,\tau)
        =
        C_A(V\setminus R,f',X^T;\tau,\sigma^{cc}),
    \]
    where \(\sigma^{cc}\) is \(\sigma\) transported across
    \(V\setminus(V\setminus R)=R\).  Consequently the same coefficient is realized
    through the endpoint of \(f\) lying outside \(R\).
\end{theorem}
\begin{proof}\leanok
    Reindex the two boundary-configuration sums by the boundary-edge equivalence
    between \(R\) and its complement.  The double complement carries the blocked
    weight of \(V\setminus(V\setminus R)\) back to the blocked weight of \(R\).
    Applying the usual endpoint realization to \(V\setminus R\) gives the displayed
    out-of-region endpoint realization.
\end{proof}

\begin{theorem}[Region resonate identity and blocked read-offs]
    \label{thm:peps_region_resonate_identity_readoffs}
    \lean{TNLean.PEPS.regionInsertionOp_regionStateVec_pin_compl,
        TNLean.PEPS.region_resonate_identity,
        TNLean.PEPS.regionComplementRow,
        TNLean.PEPS.regionInsertedCoeff_eq_complement_blockedMap,
        TNLean.PEPS.regionRegionRow,
        TNLean.PEPS.regionInsertedCoeff_eq_region_blockedMap,
        TNLean.PEPS.regionBlockedLeftInverse_complement_regionInsertedCoeff,
        TNLean.PEPS.regionBlockedLeftInverse_region_regionInsertedCoeff}
    \leanok
    \uses{def:peps_region_out_endpoint_block_inverse,
        thm:peps_regionInsertedCoeff_complement_cast}
    Assume that \(A\)'s components at the in-region endpoint of \(f\) and at the
    in-complement-region endpoint of the complementary boundary edge are linearly
    independent, and that \(A\) and \(B\) define the same PEPS state.  Let
    \(P_R\) and \(P_{R^c}\) be the interior bond products of \(R\) and \(R^c\).
    Let \(O^{\mathrm{in}}_f(X)\) and \(O^{\mathrm{out}}_f(X)\) be the physical
    endpoint operators obtained from the two virtual insertions, and let
    \(\Psi_B(R;\sigma,\tau)\) and \(\Psi_B(R^c;\tau,\sigma^{cc})\) be the
    corresponding closed \(B\)-state vectors.  Then the two endpoint readings of
    the same inserted coefficient give
    \[
        P_R\,O^{\mathrm{in}}_f(X)\,\Psi_B(R;\sigma,\tau)
        =
        P_{R^c}\,O^{\mathrm{out}}_f(X)\,\Psi_B(R^c;\tau,\sigma^{cc}).
    \]
    With \(\rho^{R^c}_{X,\sigma}\) and \(\rho^{R}_{X,\tau}\) denoting the two
    row functions, one has
    \[
        C_A(R,f,X;\sigma,-)
        =
        T_{A,R^c}\!\left(\rho^{R^c}_{X,\sigma}\right),
        \qquad
        C_A(R,f,X;-,\tau)
        =
        T_{A,R}\!\left(\rho^{R}_{X,\tau}\right).
    \]
    Hence the blocked left inverses read off the rows:
    \[
        L_{A,R^c}\!\left(C_A(R,f,X;\sigma,-)\right)
        =
        \rho^{R^c}_{X,\sigma},
        \qquad
        L_{A,R}\!\left(C_A(R,f,X;-,\tau)\right)
        =
        \rho^{R}_{X,\tau}.
    \]
\end{theorem}
\begin{proof}\leanok
    The in-region and out-of-region endpoint pins both evaluate to
    \(C_A(R,f,X;\sigma,\tau)\), hence they agree.  Expanding
    \(C_A(R,f,X;\sigma,\tau)\) as a sum over boundary configurations isolates,
    respectively, the \(R^c\)-blocked and \(R\)-blocked tensor maps; the left
    inverses then read off the row functions.
\end{proof}

\begin{definition}[Incident form of the region resonance]
    \label{def:peps_region_resonate_reconcile}
    \lean{TNLean.PEPS.RegionResonateReconcile}
    \leanok
    \uses{thm:peps_region_resonate_identity_readoffs}
    The incident-form condition associated with the region resonance asserts that,
    for every matrix \(X\) inserted on \(f\), the virtual pullback \(W_X\) of the
    in-region endpoint operator is in incident-matrix form:
    \[
        \exists Y,\qquad W_X=\mathcal I_f(Y).
    \]
\end{definition}

\begin{definition}[Region insertion transfer]
    \label{def:peps_region_insertion_transfer}
    \lean{TNLean.PEPS.RegionInsertionTransfer}
    \leanok
    \uses{def:peps_regionInsertedCoeff}
    A region insertion transfer along a boundary edge $f\in\partial R$ consists
    of maps
    \[
        T_f:\MN{D_A(f)}\to\MN{D_B(f)},
        \qquad
        S_f:\MN{D_B(f)}\to\MN{D_A(f)}
    \]
    such that, for every inserted matrix and every pair of physical
    configurations,
    \[
        C_A(R,f,X;\sigma,\tau)
        =
        C_B(R,f,T_f(X);\sigma,\tau),
        \qquad
        C_B(R,f,Y;\sigma,\tau)
        =
        C_A(R,f,S_f(Y);\sigma,\tau).
    \]
    The forward map is required to satisfy
    \[
        T_f(XY)=T_f(X)T_f(Y),\qquad T_f(I)=I.
    \]
\end{definition}

\begin{definition}[Region physical-to-virtual realization]
    \label{def:peps_region_transfer_realizes}
    \lean{TNLean.PEPS.RegionTransferRealizes}
    \leanok
    \uses{thm:peps_region_insertion_virtual_pullback_realizes}
    Let $v$ be the endpoint of $f$ lying in $R$.  A matrix transfer $T_f$ is
    realized on the region if, for every matrix $X$ inserted on $A$'s boundary
    bond and every local virtual coefficient $c$ for $B$ at $v$,
    \[
        O_{A,R,f}(X^T)\,\mathcal T_{B,v}(c)
        =
        \mathcal T_{B,v}\!\left(\mathcal I_f(T_f(X)^T)c\right).
    \]
    This is the region version of the physical-to-virtual insertion relation in
    \cite[Section~3, lines~254--582]{Molnar2018NormalPEPS}.
\end{definition}

\begin{theorem}[Realized region transfers are multiplicative and unital]
    \label{thm:peps_region_transfer_realized_maps}
    \lean{TNLean.PEPS.regionTransferMatrix_mul_of_realizes,
        TNLean.PEPS.regionTransferMatrix_one_of_realizes}
    \leanok
    \uses{def:peps_region_transfer_realizes, def:peps_region_insertion_transfer}
    If a region matrix transfer is realized in the sense of
    Definition~\ref{def:peps_region_transfer_realizes}, then its forward map is
    multiplicative:
    \[
        T_f(XY)=T_f(X)T_f(Y).
    \]
    If in addition the two PEPS tensors have the same state and the same bond
    dimensions, all bond dimensions of $B$ are positive, and the blocked
    tensors of $B$ on $R$ and its complement are injective, then the identity
    insertion gives
    \[
        T_f(I)=I.
    \]
\end{theorem}
\begin{proof}\leanok
    For multiplicativity, the physical insertion for $(XY)^T$ is the composite
    of the physical insertions for $Y^T$ and $X^T$.  Realizing these insertions
    on the $B$-side gives two incident-matrix actions on the same local tensor
    image:
    \[
        \mathcal I_f(T_f(XY)^T)
        =
        \mathcal I_f((T_f(X)T_f(Y))^T).
    \]
    Reading the incident matrix on the boundary leg gives
    $T_f(XY)=T_f(X)T_f(Y)$.  For the identity, the coefficient
    $C_A(R,f,I;\sigma,\tau)$ is the interior bond product times the state
    coefficient of $A$, hence equals $C_B(R,f,I;\sigma,\tau)$ by equality of
    states and bond dimensions.  The two blocked-injectivity hypotheses and
    positivity give injectivity of the $B$-side region-inserted coefficient,
    whence $T_f(I)=I$.
\end{proof}

\begin{definition}[Region insertion transfer from realized maps]
    \label{def:peps_region_insertion_transfer_of_realizes}
    \lean{TNLean.PEPS.regionInsertionTransfer_of_realizes}
    \leanok
    \uses{def:peps_region_transfer_realizes,
        thm:peps_region_transfer_realized_maps}
    Suppose the forward and reverse boundary matrix transfers are realized on
    the two tensors, the two tensors have the same state and the same bond
    dimensions, and the region and complement-region blocked maps of $B$ are
    injective.
    Then the maps $T_f$ and $S_f$ form a region insertion transfer:
    \[
        C_A(R,f,X;\sigma,\tau)
        =
        C_B(R,f,T_f(X);\sigma,\tau),
        \qquad
        C_B(R,f,Y;\sigma,\tau)
        =
        C_A(R,f,S_f(Y);\sigma,\tau),
    \]
    together with $T_f(XY)=T_f(X)T_f(Y)$ and $T_f(I)=I$.
\end{definition}

\begin{theorem}[Incident form gives the region insertion transfer and gauge]
    \label{thm:peps_region_reconcile_to_transfer_gauge}
    \lean{TNLean.PEPS.regionTransferRealizes_of_isIncidentMatrixForm,
        TNLean.PEPS.regionTransferRealizes_of_reconcile,
        TNLean.PEPS.regionInsertionTransfer_of_reconcile,
        TNLean.PEPS.exists_regionEdgeGauge_of_reconcile}
    \leanok
    \uses{def:peps_region_resonate_reconcile,
        thm:peps_hform_of_region_incident_form, def:peps_same_state,
        def:IsVertexInjective, def:peps_regionInjectivityData}
    Assume that \(A\) and \(B\) have the same PEPS state, are vertex-injective,
    have positive bond dimensions, and have injective blocked maps for \(R\) and
    \(R^c\). If the incident-form condition holds in both directions, then the
    matrix reading gives maps
    \[
        T_f:\MN{D_A(f)}\to\MN{D_B(f)},
        \qquad
        S_f:\MN{D_B(f)}\to\MN{D_A(f)}
    \]
    realizing all region insertion coefficients. If, in addition,
    \(D_A(e)=D_B(e)\) for every edge \(e\), then there is an invertible gauge
    \(Z_f\in\GL(D_B(f))\) such that
    \[
        T_f(X)=Z_f\,\iota_f(X)\,Z_f^{-1},
    \]
    where \(\iota_f:\MN{D_A(f)}\simeq\MN{D_B(f)}\) is the identification induced
    by the equality of bond dimensions.
\end{theorem}
\begin{proof}\leanok
    The incident-form condition supplies incident-matrix form for every inserted
    matrix. Theorem~\ref{thm:peps_hform_of_region_incident_form} gives the
    realization identities. The forward and reverse realization identities
    assemble the transfer datum, and the one-bond algebra theorem reads it as
    conjugation by an invertible matrix.
\end{proof}

\begin{theorem}[Coefficient transfer gives incident form]
    \label{thm:peps_region_coeff_transfer_to_reconcile}
    \lean{TNLean.PEPS.regionInsertionOp_regionStateVec_pin_leg,
        TNLean.PEPS.regionInsertionOp_realizes_of_realizes_on_stateVec,
        TNLean.PEPS.regionInteriorBondProd_pos,
        TNLean.PEPS.regionInsertionOp_regionStateVec_eq_of_coeff_eq,
        TNLean.PEPS.isIncidentMatrixForm_of_coeff_eq,
        TNLean.PEPS.regionResonateReconcile_of_coeff_transfer}
    \leanok
    \uses{def:peps_region_resonate_reconcile}
    Assume that \(A\) and \(B\) define the same PEPS state, that \(B\) is
    vertex-injective, that \(D_B(e)>0\) for every edge \(e\), and that
    \(D_A(e)=D_B(e)\) for every edge \(e\).  Assume also that \(A\)'s and \(B\)'s
    components at the in-region endpoint of \(f\) are linearly independent.
    Suppose that for
    every matrix \(X\) on \(A\)'s boundary bond there is a matrix \(Y\) on \(B\)'s
    boundary bond such that
    \[
        C_A(R,f,X;\sigma,\tau)=C_B(R,f,Y;\sigma,\tau)
    \]
    for all \(\sigma,\tau\).  Then the virtual pullback of \(X\) is in
    incident-matrix form on \(f\).
\end{theorem}
\begin{proof}\leanok
    Equality of the coefficients gives equality on the closed state vectors, one
    physical leg at a time.  These state-vector coefficients span the local
    virtual coefficient space, so the equality extends to all coefficients.  The
    virtual pullback is therefore the incident insertion determined by \(Y\).
\end{proof}

\begin{theorem}[First coefficient through the second complement block]
    \label{thm:peps_region_first_coeff_second_complement_block}
    \lean{TNLean.PEPS.regionInteriorBondProd_smul_regionStateVec,
        TNLean.PEPS.regionBlockedWeight_update_eq_regionWeightVec,
        TNLean.PEPS.regionInteriorBondProd_smul_regionStateVec_eq_sum,
        TNLean.PEPS.regionInsertedCoeff_eq_complement_blockedMap_B,
        TNLean.PEPS.regionBlockedLeftInverse_complement_regionInsertedCoeff_B}
    \leanok
    \uses{def:peps_region_out_endpoint_block_inverse}
    Assume that \(B\)'s complement blocked tensor map is injective, that \(A\)
    and \(B\) define the same PEPS state, that \(D_A(e)=D_B(e)\) for every edge
    \(e\), and that \(A\)'s component at the in-region endpoint of \(f\) is
    linearly independent.
    The coefficient \(C_A(R,f,X;\sigma,\tau)\), as a function of the complement
    physical configuration \(\tau\), lies in the image of \(B\)'s complement
    blocked map.  Applying the complement blocked left inverse of \(B\) recovers
    the row obtained by applying \(A\)'s endpoint operator to \(B\)'s region weight
    vectors.
\end{theorem}
\begin{proof}\leanok
    The identity insertion expresses the closed state vector as a contraction of
    the region and complement blocked weights.  Substituting the endpoint operator
    into this expression gives the complement-blocked factorization, and the left
    inverse reads off the corresponding row.
\end{proof}

\begin{theorem}[Transfer coefficient and double factorization]
    \label{thm:peps_region_transfer_coeff_double_factorization}
    \lean{TNLean.PEPS.vSideRow,
        TNLean.PEPS.regionInsertedCoeff_eq_complement_blockedMap_vSideRow,
        TNLean.PEPS.regionInsertedCoeff_eq_of_complementRow_eq,
        TNLean.PEPS.regionDoubleComplBoundaryConfigEquiv,
        TNLean.PEPS.regionBlockedTensorMap_doubleCompl,
        TNLean.PEPS.complSideRow,
        TNLean.PEPS.regionInsertedCoeff_eq_region_blockedMap_B,
        TNLean.PEPS.regionRowB,
        TNLean.PEPS.regionRowB_eq_complSideRow,
        TNLean.PEPS.regionRowB_mem_range,
        TNLean.PEPS.transferCoeff,
        TNLean.PEPS.regionRowB_eq_complement_blockedMap,
        TNLean.PEPS.regionInsertedCoeff_eq_doubleSum_transferCoeff,
        TNLean.PEPS.regionInsertedCoeff_eq_of_transferCoeff_form,
        TNLean.PEPS.vSideRow_eq_region_blockedMap_transferCoeff}
    \leanok
    \uses{thm:peps_region_first_coeff_second_complement_block}
    Assume that \(B\)'s region and complement blocked tensor maps are injective,
    that \(A\)'s components at the in-region endpoint of \(f\) and at the
    in-complement-region endpoint of the complementary boundary edge are linearly
    independent, that \(A\) and \(B\) define the same PEPS state, and that
    \(D_A(e)=D_B(e)\) for every edge \(e\).
    For fixed \(X\) and \(\sigma\), define the row
    \[
        v_X^\sigma(\nu)
        =
        \left[
            O_{A,R,f}(X^T)\,
            T_{B,R}^{\tilde\nu}(\sigma)
            \right]_{\sigma(v_f)},
    \]
    where \(\tilde\nu\) is the \(R\)-boundary configuration corresponding to the
    \(R^c\)-boundary configuration \(\nu\), and \(v_f\) is the endpoint of \(f\)
    in \(R\).
    The first tensor's inserted coefficient can be written as a double contraction
    of the second tensor's region and complement blocked weights:
    \[
        C_A(R,f,X;\sigma,\tau)
        =
        \sum_{\mu,\nu}
        K_X(\mu,\nu)\,T_{B,R}^{\mu}(\sigma)\,
        T_{B,R^c}^{\nu}(\tau).
    \]
    If \(K_X\) has the incident form
    \[
        K_X(\mu,\nu)=\mathbf 1_{\mu|_{\partial R\setminus\{f\}}
        =\nu|_{\partial R\setminus\{f\}}}\,Y_{\mu_f,\nu_f},
    \]
    then \(C_A(R,f,X;\sigma,\tau)=C_B(R,f,Y;\sigma,\tau)\).
\end{theorem}
\begin{proof}\leanok
    First factor \(C_A\) through \(B\)'s complement block, then through \(B\)'s
    region block.  The two left inverses define the coefficient \(K_X\).  Substituting
    the displayed incident form of \(K_X\) turns the double contraction into the
    defining double sum for \(C_B(R,f,Y;\sigma,\tau)\).
\end{proof}

\begin{theorem}[Incident form of the transfer coefficient]
    \label{thm:peps_region_transfer_coeff_incident_form}
    \lean{TNLean.PEPS.transferCoeff_column_eq_regionBlockedLeftInverse_vSideRow,
        TNLean.PEPS.vSideRow_eq_localTensorMap_virtualPullback,
        TNLean.PEPS.vSideRow_eq_incidentSum_of_virtualPullback_incidentForm,
        TNLean.PEPS.transferCoeff_eq_incidentForm_of_virtualPullback_incidentForm}
    \leanok
    \uses{thm:peps_region_transfer_coeff_double_factorization,
        thm:peps_hform_of_region_incident_form,
        thm:peps_region_insertion_virtual_pullback_realizes}
    Under the hypotheses of
    Theorem~\ref{thm:peps_region_transfer_coeff_double_factorization}, assume in
    addition that \(A\) and \(B\) are vertex-injective, that \(D_A(e)>0\) and
    \(D_B(e)>0\) for every edge \(e\), and that \(B\)'s component at the
    in-region endpoint of \(f\) is linearly independent.
    If the virtual pullback \(W_X\) is in incident-matrix form
    \(W_X=\mathcal I_f(Y^T)\), then the transfer coefficient \(K_X\) has the
    incident form determined by \(Y\):
    \[
        K_X(\mu,\nu)=\mathbf 1_{\mu|_{\partial R\setminus\{f\}}
        =\nu|_{\partial R\setminus\{f\}}}\,Y_{\mu_f,\nu_f}.
    \]
\end{theorem}
\begin{proof}\leanok
    The column of \(K_X\) at fixed \(\nu\) is the region blocked left inverse of
    the row \(v_X^\sigma\).  If \(W_X=\mathcal I_f(Y^T)\), this row is precisely the
    incident-matrix sum of \(Y\) against the region blocked weights.  The left
    inverse sends each blocked weight to the corresponding standard boundary
    configuration, giving the displayed formula.
\end{proof}

\begin{definition}[Coefficient transfer map]
    \label{def:peps_coeff_transfer_map}
    \lean{TNLean.PEPS.coeffTransferMap}
    \leanok
    \uses{def:peps_regionInsertedCoeff}
    Suppose that every matrix $X\in\MN{D_A(f)}$ has a matrix
    $Y\in\MN{D_B(f)}$ with matching region-inserted coefficients:
    \[
        C_A(R,f,X;\sigma,\tau)
        =
        C_B(R,f,Y;\sigma,\tau).
    \]
    Choosing one such $Y$ for each $X$ defines a forward coefficient transfer
    $T_f:\MN{D_A(f)}\to\MN{D_B(f)}$.
\end{definition}

\begin{theorem}[Coefficient transfer matches coefficients and preserves the identity]
    \label{thm:peps_coeff_transfer_map_properties}
    \lean{TNLean.PEPS.coeffTransferMap_coeff,
        TNLean.PEPS.coeffTransferMap_one}
    \leanok
    \uses{def:peps_coeff_transfer_map}
    The chosen coefficient transfer satisfies, for every inserted matrix $X$ and
    every pair of physical configurations,
    \[
        C_A(R,f,X;\sigma,\tau)
        =
        C_B(R,f,T_f(X);\sigma,\tau).
    \]
    If the tensors have the same state and the same bond dimensions, all
    $B$-bond dimensions are positive, and the region and complement-region
    blocked maps of $B$ are injective, then
    \[
        T_f(I)=I.
    \]
\end{theorem}
\begin{proof}\leanok
    The coefficient equality is the defining property of the chosen transfer.
    For $X=I$, both region-inserted coefficients are the corresponding interior
    bond products times the common PEPS state coefficient:
    \[
        C_A(R,f,I;\sigma,\tau)
        =
        C_B(R,f,I;\sigma,\tau).
    \]
    The two blocked-map injectivity hypotheses and positivity of the $B$-bond
    dimensions make the $B$-side region-inserted coefficient injective.  Hence
    the chosen matrix $T_f(I)$ is equal to $I$.
\end{proof}

\begin{definition}[Region insertion transfer from coefficient transfers]
    \label{def:peps_region_insertion_transfer_of_coeff_transfer}
    \lean{TNLean.PEPS.regionInsertionTransfer_of_coeffTransfer}
    \leanok
    \uses{def:peps_region_insertion_transfer,
        def:peps_coeff_transfer_map,
        thm:peps_coeff_transfer_map_properties}
    Suppose coefficient transfers exist in both directions:
    \[
        C_A(R,f,X;\sigma,\tau)
        =
        C_B(R,f,T_f(X);\sigma,\tau),
        \qquad
        C_B(R,f,Y;\sigma,\tau)
        =
        C_A(R,f,S_f(Y);\sigma,\tau).
    \]
    Assume also that the two tensors have the same state and the same bond
    dimensions, all $B$-bond dimensions are positive, and the region and
    complement-region blocked maps of $B$ are injective.  If the forward choice
    is multiplicative,
    \[
        T_f(XY)=T_f(X)T_f(Y),
    \]
    then $T_f$ and $S_f$ form a region insertion transfer.
\end{definition}

\begin{theorem}[Per-edge gauge from coefficient transfers]
    \label{thm:peps_region_edge_gauge_of_coeff_transfer}
    \lean{TNLean.PEPS.exists_regionEdgeGauge_of_coeffTransfer}
    \leanok
    \uses{def:peps_region_insertion_transfer_of_coeff_transfer}
    Assume coefficient transfers exist in both directions, the forward transfer
    is multiplicative, the two tensors have the same state and the same bond
    dimensions, all bond dimensions are positive, and the region and
    complement-region blocked maps for both tensors are injective.  Then there
    is an invertible boundary gauge $Z_f$ such that, for every
    $X\in\MN{D_A(f)}$,
    \[
        T_f(X)=Z_f\,\iota_f(X)\,Z_f^{-1},
    \]
    where $\iota_f:\MN{D_A(f)}\simeq\MN{D_B(f)}$ is the identification induced by
    the equality of bond dimensions.
\end{theorem}
\begin{proof}\leanok
    The coefficient transfers form a region insertion transfer by
    Definition~\ref{def:peps_region_insertion_transfer_of_coeff_transfer}.  The
    insertion-transfer algebra isomorphism gives an algebra equivalence between
    the two boundary matrix algebras, and the Skolem--Noether theorem for matrix
    algebras writes this equivalence as conjugation by an invertible matrix.
\end{proof}

\begin{theorem}[Determinacy of a region insertion transfer]
    \label{thm:peps_regionInsertedCoeff_transferMap_unique}
    \lean{TNLean.PEPS.regionInsertedCoeff_transferMap_unique}
    \leanok
    \uses{def:peps_regionInsertedCoeff, def:peps_region_out_endpoint_block_inverse}
    Assume that the blocked tensors of \(B\) on \(R\) and \(R^c\) are injective
    and that all virtual bonds of \(B\) have positive dimension.  Let
    \(T_1,T_2:\MN{D_A(f)}\to\MN{D_B(f)}\) be two maps.  For \(i=1,2\), suppose
    \(T_i\) satisfies, for every inserted matrix and every pair of physical
    configurations,
    \[
        C_A(R,f,M;\sigma,\tau)
        =
        C_B(R,f,T_i(M);\sigma,\tau).
    \]
    Then \(T_1(M)=T_2(M)\) for every \(M\).
\end{theorem}
\begin{proof}\leanok
    For every pair of physical configurations, the two displayed identities give
    \[
        C_B(R,f,T_1(M);\sigma,\tau)
        =
        C_A(R,f,M;\sigma,\tau)
        =
        C_B(R,f,T_2(M);\sigma,\tau).
    \]
    Injectivity of the region and complement blocked tensors, together with
    positivity of all virtual bonds of \(B\), gives
    \[
        \left(
        \forall\sigma,\tau,\,
        C_B(R,f,T_1(M);\sigma,\tau)
        =
        C_B(R,f,T_2(M);\sigma,\tau)
        \right)
        \Longrightarrow
        T_1(M)=T_2(M).
    \]
\end{proof}

\begin{theorem}[Reindexing commutes with conjugation]
    \label{thm:peps_reindexAlgEquiv_gaugeConj}
    \lean{TNLean.PEPS.reindexAlgEquiv_gaugeConj}
    \leanok
    Assume \(D=E\), let \(\iota:\MN{D}\simeq\MN{E}\) be the canonical
    identification of matrix algebras, and let \(Z\in\GL(D)\).  Then for every
    \(N\in\MN{D}\),
    \[
        \iota(ZNZ^{-1})
        =
        \iota(Z)\,\iota(N)\,\iota(Z)^{-1}.
    \]
\end{theorem}
\begin{proof}\leanok
    The map \(\iota\) is an algebra equivalence, so it preserves products and
    inverses.  Applying this to the three factors \(Z\), \(N\), and \(Z^{-1}\)
    gives the displayed identity.
\end{proof}

\begin{theorem}[Coefficient identities determine the boundary conjugation]
    \label{thm:peps_gaugeConj_eq_of_coeffIdentities}
    \lean{TNLean.PEPS.gaugeConj_eq_of_coeffIdentities}
    \leanok
    \uses{thm:peps_regionInsertedCoeff_transferMap_unique}
    Assume that the blocked tensors of \(B\) on \(R\) and \(R^c\) are injective
    and that all virtual bonds of \(B\) have positive dimension.  Suppose that
    the equality \(D_A(f)=D_B(f)\) fixes the canonical identification
    \[
        \iota:\MN{D_A(f)}\simeq\MN{D_B(f)}.
    \]
    For \(i=1,2\), suppose that the invertible matrix \(Z_i\in\GL(D_B(f))\)
    realizes the region-insertion identity
    \[
        C_A(R,f,M;\sigma,\tau)
        =
        C_B(R,f,Z_i\,\iota(M)\,Z_i^{-1};\sigma,\tau).
    \]
    Then the two conjugations agree on every boundary matrix:
    \[
        Z_1\,\iota(M)\,Z_1^{-1}
        =
        Z_2\,\iota(M)\,Z_2^{-1}.
    \]
\end{theorem}
\begin{proof}\leanok
    For \(i=1,2\), define the transfer map
    \[
        T_i(M):=Z_i\,\iota(M)\,Z_i^{-1}.
    \]
    Its coefficient identity is
    \[
        C_A(R,f,M;\sigma,\tau)
        =
        C_B(R,f,T_i(M);\sigma,\tau).
    \]
    The preceding determinacy theorem gives
    \[
        T_1(M)=T_2(M).
    \]
    Substituting the definitions of \(T_1\) and \(T_2\) yields
    \[
        Z_1\,\iota(M)\,Z_1^{-1}
        =
        Z_2\,\iota(M)\,Z_2^{-1}.
    \]
\end{proof}

%======================================================================
% InsertSplit: the one-site quotient of the blocked-region weight
%======================================================================

\begin{theorem}[Region vertex product across an inserted site]%
    \label{thm:peps_prod_region_insert_split}
    \lean{TNLean.PEPS.prod_region_insert_split}
    \leanok
    Let $v\notin R$.  At a fixed global virtual configuration $\zeta$ and a
    physical configuration $\sigma$ on $R\cup\{v\}$, the product of the local
    tensors over the vertices of $R\cup\{v\}$ factors as the local tensor at the
    inserted site $v$ times the product over the vertices of $R$:
    \[
        \prod_{w\in R\cup\{v\}}A_w\bigl(\zeta|_{\operatorname{inc}(w)},\sigma_w\bigr)
        =
        A_v\bigl(\zeta|_{\operatorname{inc}(v)},\sigma_v\bigr)
        \prod_{w\in R}A_w\bigl(\zeta|_{\operatorname{inc}(w)},\sigma|_R{}_w\bigr).
    \]
\end{theorem}
\begin{proof}\leanok
    The inserted site $v$ is isolated from the product over $R\cup\{v\}$ by
    splitting off a single factor, and the residual product over the remaining
    vertices is reindexed through the bijection between the vertices of $R$ and
    the non-$v$ vertices of $R\cup\{v\}$.
\end{proof}

\begin{theorem}[One-site split of the blocked-region weight]%
    \label{thm:peps_regionBlockedWeight_insert_eq_sum_split}
    \lean{TNLean.PEPS.regionBlockedWeight_insert_eq_sum_split}
    \leanok
    \uses{thm:peps_prod_region_insert_split}
    Let $v\notin R$.  The blocked-region weight of $R\cup\{v\}$ is the constrained
    sum, over global virtual configurations $\zeta$ whose crossing-edge label on
    $R\cup\{v\}$ is $\mu$, of the inserted-site tensor times the vertex product
    over $R$:
    \[
        T_{A,R\cup\{v\}}^{\mu}(\sigma)
        =
        \sum_{\zeta:\,\zeta|_{\partial(R\cup\{v\})}=\mu}
        A_v\bigl(\zeta|_{\operatorname{inc}(v)},\sigma_v\bigr)
        \prod_{w\in R}A_w\bigl(\zeta|_{\operatorname{inc}(w)},\sigma|_R{}_w\bigr).
    \]
\end{theorem}
\begin{proof}\leanok
    Each summand of the blocked-region weight is the vertex product over
    $R\cup\{v\}$, which factors by
    Theorem~\ref{thm:peps_prod_region_insert_split}.
\end{proof}

\begin{theorem}[A non-incident crossing edge survives the insertion]%
    \label{thm:peps_isRegionBoundaryEdge_insert_of_not_incident}
    \lean{TNLean.PEPS.isRegionBoundaryEdge_insert_of_not_incident}
    \leanok
    Let $v\notin R$.  A crossing edge of $R$ whose two endpoints both differ from
    $v$ is also a crossing edge of $R\cup\{v\}$.
\end{theorem}
\begin{proof}\leanok
    A crossing edge of $R$ has exactly one endpoint inside $R$.  Adjoining $v$ to
    $R$ leaves both endpoints on the same side as before, since neither endpoint
    is $v$, so the edge still crosses the boundary of $R\cup\{v\}$.
\end{proof}

\begin{definition}[Boundary label of the smaller region across an insertion]%
    \label{def:peps_boundaryLabelOfInsert}
    \lean{TNLean.PEPS.boundaryLabelOfInsert}
    \leanok
    \uses{thm:peps_isRegionBoundaryEdge_insert_of_not_incident}
    Let $v\notin R$.  Given a crossing-edge label $\mu$ for $R\cup\{v\}$ and a
    local virtual configuration $\eta$ on the edges incident to $v$, the
    crossing-edge label of $R$ is read off these two data: a crossing edge of $R$
    incident to $v$ has become internal to $R\cup\{v\}$, so its label is taken
    from $\eta$; a crossing edge of $R$ not incident to $v$ is also a crossing
    edge of $R\cup\{v\}$, so its label is taken from $\mu$.
\end{definition}

\begin{theorem}[Boundary-label fiber identity across an inserted site]%
    \label{thm:peps_regionBoundaryLabel_eq_boundaryLabelOfInsert}
    \lean{TNLean.PEPS.regionBoundaryLabel_eq_boundaryLabelOfInsert}
    \leanok
    \uses{def:peps_boundaryLabelOfInsert}
    Let $v\notin R$.  If a global virtual configuration $\zeta$ restricts to $\mu$
    on the crossing edges of $R\cup\{v\}$ and to $\eta$ on the edges incident to
    $v$, then the crossing-edge label of $R$ read from $\zeta$ is exactly the
    bridge label of $\mu$ and $\eta$.
\end{theorem}
\begin{proof}\leanok
    The two labels are compared edge by edge.  On a crossing edge of $R$
    incident to $v$ both sides read $\zeta$ on an edge incident to $v$, which is
    $\eta$.  On a crossing edge of $R$ not incident to $v$ both sides read
    $\zeta$ on a crossing edge of $R\cup\{v\}$, which is $\mu$.
\end{proof}

%======================================================================
% Gauge-absorption bridge at region granularity
%======================================================================

\begin{theorem}[Region gauge absorption]%
    \label{thm:peps_regionInsertedCoeff_applyGauge}
    \lean{TNLean.PEPS.regionInsertedCoeff_applyGauge}
    \leanok
    \uses{thm:peps_regionInsertedCoeff_eq_smul_edgeInsertedCoeff,
        thm:peps_edge_gauge_absorption, def:applyGauge}
    Let $X$ be an edge-gauge family, write $B^X$ for the gauged tensor, and let
    $f$ be a crossing edge of $R$.  Inserting $M$ on $f$ between the region and
    its complement in $B^X$ equals inserting, in $B$, the matrix obtained by
    orienting $M$ to the boundary-edge convention, conjugating by the transpose
    of the endpoint gauge $X_f$, and orienting back:
    \[
        C_{B^X}(R,f,M;\sigma,\tau)
        =
        C_B\bigl(R,f,
            \mathcal O_{R,f}\bigl(X_f^{\mathsf T}\,
                \mathcal O_{R,f}(M)\,(X_f^{-1})^{\mathsf T}\bigr);
            \sigma,\tau\bigr),
    \]
    where $\mathcal O_{R,f}$ is the boundary-edge orientation.  This is the
    region-granularity analogue of
    Theorem~\ref{thm:peps_edge_gauge_absorption}.
\end{theorem}
\begin{proof}\leanok
    By Theorem~\ref{thm:peps_regionInsertedCoeff_eq_smul_edgeInsertedCoeff} the
    region/complement contraction is the single-bond cut at $f$ up to the
    interior-bond multiplicity, which is unchanged by a gauge.  At the edge
    granularity every interior edge and every crossing edge other than $f$
    cancels its gauge pairwise, while the two endpoint gauges on $f$ conjugate
    the inserted matrix.  Transporting the cancelled identity back through the
    region-to-edge identity, and using that the boundary-edge orientation is an
    involution, gives the displayed equality.
\end{proof}

%======================================================================
% Region-transport covariance and torus translation covariance
%======================================================================

\begin{theorem}[Transport of the complement-side blocked weight]%
    \label{thm:peps_regionComplementWeight_transport}
    \lean{TNLean.PEPS.regionComplementWeight_transport}
    \leanok
    \uses{def:peps_tensor_transport, def:peps_regionComplementTwoBlock}
    Let $\phi:G\simeq G'$ be a graph isomorphism.  The blocked weight of the
    transported tensor $A^\phi$ over the complement of the image region, read at
    the image complement boundary configuration and the image complement physical
    configuration, equals the blocked weight of $A$ over $V\setminus R$:
    \[
        T_{A^\phi,V'\setminus\phi(R)}^{(\phi\cdot\eta)^c}(\phi\cdot\tau)
        =
        T_{A,V\setminus R}^{\eta^c}(\tau).
    \]
\end{theorem}
\begin{proof}\leanok
    The image of the complement of $R$ is the complement of the image of $R$, and
    after this identification the blocked weight transports along $\phi$ by the
    change of summation variables on virtual configurations.
\end{proof}

\begin{theorem}[Transport preserves agreement away from a bond]%
    \label{thm:peps_sameAwayFromBond_regionBoundaryConfigMap}
    \lean{TNLean.PEPS.sameAwayFromBond_regionBoundaryConfigMap}
    \leanok
    \uses{thm:peps_regionComplementWeight_transport}
    Let $\phi:G\simeq G'$ be a graph isomorphism and let $f$ be a crossing edge
    of $R$.  Two crossing-edge labels of $R$ agree away from $f$ if and only if
    their images under the boundary reindex induced by $\phi$ agree away from the
    image of $f$.
\end{theorem}
\begin{proof}\leanok
    The boundary reindex is a bijection on crossing edges that carries $f$ to its
    image, so it carries the crossing edges other than $f$ bijectively to the
    crossing edges other than the image of $f$.  Agreement on the former family
    is therefore equivalent to agreement on the latter.
\end{proof}

\begin{theorem}[Covariance of the region insertion coefficient]%
    \label{thm:peps_regionInsertedCoeff_transport}
    \lean{TNLean.PEPS.regionInsertedCoeff_transport}
    \leanok
    \uses{def:peps_regionInsertedCoeff, def:peps_tensor_transport,
        thm:peps_regionComplementWeight_transport,
        thm:peps_sameAwayFromBond_regionBoundaryConfigMap}
    Let $\phi:G\simeq G'$ be a graph isomorphism and let $f$ be a crossing edge
    of $R$.  The region insertion coefficient of the transported tensor $A^\phi$
    over the image region, with $M$ inserted on the image of $f$ and the physical
    configurations relabelled by $\phi$, equals the region insertion coefficient
    of $A$ over $R$, with $M$ carried back to the bond of $A$ at $f$:
    \[
        C_{A^\phi}\bigl(\phi(R),\phi(f),M;\phi\cdot\sigma,\phi\cdot\tau\bigr)
        =
        C_A\bigl(R,f,\phi^{*}M;\sigma,\tau\bigr).
    \]
    No endpoint orientation enters: the matrix is inserted between the two single
    boundary values on the edge.
\end{theorem}
\begin{proof}\leanok
    The two crossing-label sums reindex by the boundary action of $\phi$, the
    agreement-away-from-$f$ constraint is preserved by
    Theorem~\ref{thm:peps_sameAwayFromBond_regionBoundaryConfigMap}, the region
    weight transports along $\phi$, and the complement weight transports by
    Theorem~\ref{thm:peps_regionComplementWeight_transport}.  The inserted matrix
    entry is read at the two boundary values, which the bond-dimension
    identification carries to the corresponding entry of $\phi^{*}M$.
\end{proof}

\begin{theorem}[Region insertion coefficient under an equality of tensors]%
    \label{thm:peps_regionInsertedCoeff_congrTensor}
    \lean{TNLean.PEPS.regionInsertedCoeff_congrTensor}
    \leanok
    \uses{def:peps_regionInsertedCoeff}
    If two tensors are equal, then their region insertion coefficients are equal,
    with the inserted matrix carried across the induced equality of bond
    dimensions:
    \[
        A_1=A_2
        \quad\Longrightarrow\quad
        C_{A_1}(R,f,M;\sigma,\tau)
        =
        C_{A_2}\bigl(R,f,\psi M;\sigma,\tau\bigr).
    \]
    This bridges the literal tensor of a translation-invariant PEPS with its
    transported copy.
\end{theorem}
\begin{proof}\leanok
    Substituting the equality of tensors identifies the two coefficients, and the
    induced equality of bond dimensions is the identity, so the carried matrix is
    the original one.
\end{proof}

\begin{theorem}[Bond dimension at a translated crossing edge]%
    \label{thm:peps_bondDim_boundaryEdgeMap_translate}
    \lean{TNLean.PEPS.bondDim_boundaryEdgeMap_translate}
    \leanok
    \uses{def:peps_torus_translation_invariant,
        def:peps_torus_geometry_translation}
    Let $A$ be translation invariant on the discrete torus, let $f$ be a crossing
    edge of $R$, and let $\tau_{a,b}$ be a translation.  The bond dimension of $A$
    at the translated crossing edge equals its bond dimension at the reference
    crossing edge:
    \[
        D_A\bigl(\tau_{a,b}(f)\bigr)=D_A(f).
    \]
\end{theorem}
\begin{proof}\leanok
    Translation carries the reference edge to its image, and translation
    invariance fixes the tensor under transport, so the bond dimension is
    unchanged.
\end{proof}

\begin{theorem}[Translation covariance of the coefficient identity]%
    \label{thm:peps_regionInsertedCoeff_translate_coeffIdentity}
    \lean{TNLean.PEPS.regionInsertedCoeff_translate_coeffIdentity}
    \leanok
    \uses{thm:peps_regionInsertedCoeff_transport,
        thm:peps_regionInsertedCoeff_congrTensor,
        thm:peps_bondDim_boundaryEdgeMap_translate,
        def:peps_torus_translation_invariant}
    Let $A$ and $B$ be translation invariant on the discrete torus, and let $f$ be
    a crossing edge of $R$.  Suppose a map $g$ realizes the coefficient identity
    at the reference edge:
    \[
        C_A(R,f,M;\sigma,\tau)
        =
        C_B\bigl(R,f,g(M);\sigma,\tau\bigr).
    \]
    Then, for every translation $\tau_{a,b}$, the same identity holds at the
    translated crossing edge of the translated region, with the inserted matrix
    carried between the translated and reference bonds:
    \[
        C_A\bigl(\tau_{a,b}(R),\tau_{a,b}(f),M';
            \tau_{a,b}\cdot\sigma,\tau_{a,b}\cdot\tau\bigr)
        =
        C_B\bigl(\tau_{a,b}(R),\tau_{a,b}(f),
            \widehat{g}(M');
            \tau_{a,b}\cdot\sigma,\tau_{a,b}\cdot\tau\bigr),
    \]
    where $\widehat{g}$ applies $g$ to the reference-bond reindex of $M'$ and
    carries the result back to the translated bond.
\end{theorem}
\begin{proof}\leanok
    Translation invariance identifies the transported tensor with the original,
    so Theorem~\ref{thm:peps_regionInsertedCoeff_transport} carries the reference
    coefficient of $A$ to the translated edge with no change of tensor;
    Theorem~\ref{thm:peps_regionInsertedCoeff_congrTensor} mediates the literal
    and transported tensors.  Applying the reference identity for $g$ and
    transporting back for $B$ in the same way gives the displayed equality, the
    inserted matrix being carried across the equal bond dimensions of
    Theorem~\ref{thm:peps_bondDim_boundaryEdgeMap_translate}.
\end{proof}

\begin{definition}[Single-vertex two-block tensor]%
    \label{def:peps_vertexTwoBlock}
    \lean{TNLean.PEPS.vertexTwoBlock}
    \leanok
    \uses{def:peps_tensor}
    Fix a vertex $v$.  The single-vertex tensor $A_v$ is regarded as a
    two-block tensor whose shared bonds are the edges incident to $v$, whose
    external boundary is a one-point space, and whose physical index is the
    physical index at $v$:
    \[
        T_v(*,\eta,\sigma)=A_v(\eta,\sigma).
    \]
\end{definition}

\begin{theorem}[Injectivity of the single-vertex two-block tensor]%
    \label{thm:peps_isTwoBlockInjective_vertexTwoBlock}
    \lean{TNLean.PEPS.isTwoBlockInjective_vertexTwoBlock}
    \leanok
    \uses{def:IsVertexInjective, def:peps_vertexTwoBlock}
    If $A$ is vertex-injective, then the single-vertex two-block tensor at every
    vertex is injective.
\end{theorem}
\begin{proof}\leanok
    Vertex injectivity is linear independence of the coefficient family
    $\eta\mapsto A_v(\eta,-)$.  The auxiliary one-point boundary only reindexes
    this same family, so the two-block injectivity condition is the same linear
    independence statement.
\end{proof}

\begin{definition}[Complement two-block tensor]%
    \label{def:peps_complementTwoBlock}
    \lean{TNLean.PEPS.complementTwoBlock}
    \leanok
    \uses{def:peps_tensor}
    For a vertex $v$, the complementary block $V\setminus\{v\}$ defines a
    two-block tensor on the same incident shared bonds. Its coefficient is the
    contraction of all tensors away from $v$, with the incident bond indices
    left open at the complement endpoints:
    \[
        T_{v^c}(*,\eta,\tau)=A_{v^c}(\eta,\tau).
    \]
\end{definition}

\begin{theorem}[Injectivity of the vertex-complement contraction]%
    \label{thm:peps_vertexComplementTensorInjective}
    \lean{TNLean.PEPS.vertexComplementTensorInjective_of_isVertexInjective}
    \leanok
    \uses{def:IsVertexInjective, def:peps_tensor}
    If $A$ is vertex-injective and every virtual bond space has positive
    dimension, then the coefficient family
    \[
        \eta\longmapsto \bigl(\tau\longmapsto A_{v^c}(\eta,\tau)\bigr)
    \]
    of the vertex-complement contraction is linearly independent.
\end{theorem}
\begin{proof}\leanok
    Let $c=(c_\eta)_\eta$ be a finite coefficient family on the open star of
    $v$. For a set $S\subseteq V\setminus\{v\}$, let $K_S(c)$ be the assertion
    that for every virtual configuration $\zeta_0$ and physical configuration
    $\tau$,
    \[
        \sum_{\zeta}
        \mathbf 1[
                \zeta_f=\zeta_{0,f}\text{ for every edge }f
                \text{ whose endpoints are outside }S]\,
        c_{\zeta|_{\operatorname{Star}(v)}}
        \prod_{w\in S} A_w(\zeta|_{\operatorname{Star}(w)},\tau_w)=0 .
    \]
    The initial vanishing of the complement contraction is exactly
    $K_{V\setminus\{v\}}(c)$. If $j\in S$ and $j\neq v$, the one-sided inverse at
    the vertex $j$ separates the local tensor $A_j$ and gives
    $K_S(c)\Rightarrow K_{S\setminus\{j\}}(c)$. Iterating this implication over
    the complement vertices gives $K_\varnothing(c)$. At $S=\varnothing$ the
    exposed virtual indices are all fixed by $\zeta_0$; positive bond dimensions
    allow $\zeta_0$ to realize any prescribed star configuration at $v$. Hence
    every coefficient $c_\eta$ is zero.
\end{proof}

\begin{theorem}[Injectivity of the complement two-block tensor]%
    \label{thm:peps_isTwoBlockInjective_complementTwoBlock}
    \lean{TNLean.PEPS.isTwoBlockInjective_complementTwoBlock}
    \leanok
    \uses{def:IsVertexInjective, def:peps_complementTwoBlock,
        thm:peps_vertexComplementTensorInjective}
    If $A$ is vertex-injective and every virtual bond space has positive
    dimension, then the complement two-block tensor at every vertex is
    injective.
\end{theorem}
\begin{proof}\leanok
    \uses{thm:peps_vertexComplementTensorInjective}
    The previous theorem gives linear independence of the coefficient family
    $\eta\mapsto A_{v^c}(\eta,-)$. The two-block tensor only adds a one-point
    external boundary, so this is the same injectivity condition after the
    harmless reindexing $(*,\eta)\leftrightarrow\eta$.
\end{proof}

\begin{theorem}[One-vertex versus complement comparison]%
    \label{thm:peps_oneVertexComplementComparison}
    \lean{TNLean.PEPS.one_vertex_complement_comparison}
    \leanok
    \uses{thm:peps_postAbsorptionEdgeInsertionEquality,
        thm:peps_twoInjectiveTensorInsertionComparison}
    Let the shared bond type and the two external boundary spaces be nonempty.
    After the edge gauges have been absorbed into the second tensor family
    $\widetilde B$, block one vertex $v$ against its complement. Assume that
    $A_v,A_{v^c},\widetilde B_v,\widetilde B_{v^c}$ are injective as two-block
    tensors. The post-absorption insertion equality gives, for every shared
    bond $b$, every matrix $X\in\MN{D_b}$, and all external and
    physical indices,
    \[
        C_{A_v,A_{v^c}}(b,X;\eta_v,\eta_{v^c},\sigma_v,\sigma_{v^c})
        =
        C_{\widetilde B_v,\widetilde B_{v^c}}
        (b,X;\eta_v,\eta_{v^c},\sigma_v,\sigma_{v^c}).
    \]
    Then
    \[
        \exists\lambda_v\in\C^\times,\qquad
        A_v=\lambda_v\widetilde B_v.
    \]
    \begin{center}
        \TNPEPSOneVertexComplementComparison
    \end{center}
    This is the final local comparison following the post-absorption insertion
    equality in \cite[Section~3]{Molnar2018NormalPEPS}.
\end{theorem}
\begin{proof}\leanok
    With $A_1=A_v$, $A_2=A_{v^c}$, $B_1=\widetilde B_v$, and
    $B_2=\widetilde B_{v^c}$, the post-absorption equality is precisely
    \[
        \forall b,X,\eta_v,\eta_{v^c},\sigma_v,\sigma_{v^c},\quad
        C_{A_1,A_2}(b,X;\eta_v,\eta_{v^c},\sigma_v,\sigma_{v^c})
        =
        C_{B_1,B_2}(b,X;\eta_v,\eta_{v^c},\sigma_v,\sigma_{v^c}).
    \]
    Theorem~\ref{thm:peps_twoInjectiveTensorInsertionComparison} gives
    \[
        \exists\lambda_v\in\mathbb C^\times,\qquad
        A_v=\lambda_v\widetilde B_v,\qquad
        A_{v^c}=\lambda_v^{-1}\widetilde B_{v^c}.
    \]
    The desired local proportionality is the first equality.
\end{proof}

\begin{remark}[Normal-region form]
    The normal-region application in \cite[Section~4]{Molnar2018NormalPEPS}
    uses the same scalar comparison after comparing two injective regions
    $R\subset S=R\cup\{v\}$. If the region comparisons give
    $A_R=\lambda_R\widetilde B_R$ and
    $A_S=\lambda_S\widetilde B_S$, then the source-paper diagram reads
    \[
        A_R A_v
        =
        A_S
        =
        \lambda_S\widetilde B_S
        =
        \lambda_S\widetilde B_R\widetilde B_v.
    \]
    Applying the left inverse supplied by injectivity of the block
    $A_R=\lambda_R\widetilde B_R$ gives
    $A_v=(\lambda_S/\lambda_R)\widetilde B_v$.
\end{remark}

\begin{theorem}[Fundamental Theorem for injective PEPS, conditional on
        bond-dimension equality]%
    \label{thm:fundamentalTheorem_PEPS_of_bondDim}
    \lean{TNLean.PEPS.fundamentalTheorem_PEPS_of_bondDim}
    \leanok
    \uses{thm:gaugeConsistency, thm:peps_oneVertexComplementComparison,
        def:GaugeEquivPEPS, def:IsVertexInjective}
    Suppose that the corresponding virtual edge spaces of $A$ and $B$ are
    already identified, every virtual bond of $A$ has positive dimension, and
    the graph is connected.
    If $A$ and $B$ are vertex-injective and give the same PEPS state, then
    there are invertible matrices $X_e$, one for each edge, such that at every
    vertex the tensor $B_v$ is obtained from $A_v$ by the oriented endpoint
    action of the incident matrices.
\end{theorem}
\begin{proof}\leanok
    Gauge consistency supplies the edge gauges directly, and the gauge
    equivalence packages them with the identified bond dimensions.
\end{proof}

\begin{theorem}[Fundamental Theorem for injective PEPS]%
    \label{thm:fundamentalTheorem_PEPS}
    \lean{TNLean.PEPS.fundamentalTheorem_PEPS}
    \leanok
    % Scope restriction: this is the positive-bond variant of Theorem 2 in
    % \cite{Molnar2018NormalPEPS}. The elimination plan for the positivity
    % hypotheses is recorded in
    % \path{docs/paper-gaps/peps_injective_ft_section3_route.tex}.
    \uses{thm:fundamentalTheorem_PEPS_of_bondDim, def:GaugeEquivPEPS,
        def:IsVertexInjective}
    Let $A$ and $B$ be vertex-injective PEPS tensors on a finite simple
    graph, and suppose every virtual bond of $A$ and $B$ has positive
    dimension and that the graph is connected. If they give the same PEPS
    state, then their corresponding
    virtual edge spaces have the same dimensions and, after identifying these
    spaces, there are invertible matrices $X_e$, one for each edge, such that
    $B_v$ is obtained from $A_v$ by the oriented endpoint action of the matrices
    incident to $v$.
\end{theorem}
\begin{proof}\leanok
    The bond dimensions agree edgewise, because blocking around an edge gives
    two injective three-site chains whose matched matrix insertions on that bond
    form an algebra equivalence between the two bond matrix algebras, forcing
    equal sizes. With the bond dimensions identified, the conditional form of
    the theorem applies.
\end{proof}

\section{Union Closure for Normal PEPS}\label{sec:peps_normal_union}

A normal PEPS is injective only after blocking certain regions, so the injective
theorem applies once those regions are available. The closure property that
makes them available is that, for positive bond dimensions, the union of two
injective regions is again injective,
$\inj(A)\wedge\inj(B)\Longrightarrow\inj(A\cup B)$.
Writing $P_0=A\setminus B$, $P_1=A\cap B$, $P_2=B\setminus A$, and
$P_3=(A\cup B)^c$, a boundary operator $X$ attached to $P_3$ that is annihilated
by the contraction over $A\cup B$ descends through the contraction chain
\begin{align}
    \mathcal C_{\{0,1,2\}}(X)=0
    \Longrightarrow \mathcal C_{\{2\}}(X)=0
    \Longrightarrow \mathcal C_{\{1,2\}}(X)=0
    \Longrightarrow X=0,
    \notag
\end{align}
so the union contraction has trivial kernel. Iterating gives injectivity of
every finite union of injective regions.

\begin{definition}[Four-region decomposition for a union of regions]%
    \label{def:peps_regionUnionDecomposition}
    \lean{TNLean.PEPS.regionOnlyLeft, TNLean.PEPS.regionOverlap}
    \lean{TNLean.PEPS.regionOnlyRight, TNLean.PEPS.regionOutsideUnion}
    \lean{TNLean.PEPS.regionUnionPart}
    \leanok
    Let $V$ be a finite vertex set. For two regions $A,B\subseteq V$,
    define the four regions $A\setminus B$, $A\cap B$,
    $B\setminus A$, and $V\setminus(A\cup B)$. These are the four
    regions used in the proof of Lemma~\ref{lem:peps_injectiveRegion_union}.
    They also form an indexed family over $\{0,1,2,3\}$ in the order
    $A\setminus B$, $A\cap B$, $B\setminus A$, $V\setminus(A\cup B)$.
\end{definition}

\begin{lemma}[Identities for the four-region decomposition]%
    \label{lem:peps_regionUnionDecomposition_identities}
    \lean{TNLean.PEPS.regionOnlyLeft_union_overlap}
    \lean{TNLean.PEPS.regionOverlap_union_onlyRight}
    \lean{TNLean.PEPS.regionOutsideUnion_eq_regionComplement_union}
    \lean{TNLean.PEPS.regionThreePart_union, TNLean.PEPS.regionFourPart_union}
    \lean{TNLean.PEPS.regionFourPart_cases}
    \lean{TNLean.PEPS.regionUnionPart_zero, TNLean.PEPS.regionUnionPart_one}
    \lean{TNLean.PEPS.regionUnionPart_two, TNLean.PEPS.regionUnionPart_three}
    \lean{TNLean.PEPS.regionUnionPart_three_eq_regionComplement_union}
    \lean{TNLean.PEPS.regionUnionPart_zero_union_one}
    \lean{TNLean.PEPS.regionUnionPart_biUnion_zero_one}
    \lean{TNLean.PEPS.regionUnionPart_one_union_two}
    \lean{TNLean.PEPS.regionUnionPart_biUnion_one_two}
    \lean{TNLean.PEPS.regionUnionPart_zero_union_one_union_two}
    \lean{TNLean.PEPS.regionUnionPart_biUnion_zero_one_two}
    \lean{TNLean.PEPS.regionUnionPart_biUnion_three}
    \lean{TNLean.PEPS.regionUnionPart_biUnion_three_eq_regionComplement_zero_one_two}
    \lean{TNLean.PEPS.regionUnionPart_three_eq_regionComplement_left_inter_right}
    \lean{TNLean.PEPS.regionUnionPart_two_union_three_eq_regionComplement_left}
    \lean{TNLean.PEPS.regionUnionPart_biUnion_two_three_eq_regionComplement_left}
    \lean{TNLean.PEPS.regionUnionPart_biUnion_two_three_eq_regionComplement_zero_one}
    \lean{TNLean.PEPS.regionUnionPart_zero_union_three_eq_regionComplement_right}
    \lean{TNLean.PEPS.regionUnionPart_biUnion_zero_three_eq_regionComplement_right}
    \lean{TNLean.PEPS.regionUnionPart_biUnion_zero_three_eq_regionComplement_one_two}
    \lean{TNLean.PEPS.regionUnionPart_two_union_three_eq_regionComplement_zero_union_one}
    \lean{TNLean.PEPS.regionUnionPart_zero_union_three_eq_regionComplement_one_union_two}
    \lean{TNLean.PEPS.regionUnionPart_zero_union_one_eq_regionComplement_two_union_three}
    \lean{TNLean.PEPS.regionUnionPart_one_union_two_eq_regionComplement_zero_union_three}
    \lean{TNLean.PEPS.regionUnionPart_biUnion}
    \lean{TNLean.PEPS.regionUnionPart_exists_unique}
    \leanok
    \uses{def:peps_regionUnionDecomposition}
    Let $V$ be a finite vertex set and let $A,B\subseteq V$. The left-only
    and overlap regions reconstruct $A$, the overlap and right-only regions
    reconstruct $B$, the three inside regions reconstruct $A\cup B$, and
    all four regions reconstruct $V$. The same identities also hold in the
    indexed four-region family: pieces $0$ and $1$ reconstruct $A$, pieces
    $1$ and $2$ reconstruct $B$, and pieces $0,1,2$ reconstruct $A\cup B$.
    These same reconstructions are available as finite unions over the
    corresponding subfamilies of the indexed four-region family.
    Piece $3$ is the complement of $A\cup B$, equivalently the complement of
    the three inside indexed pieces, and also $A^c\cap B^c$. The singleton
    subfamily $\{3\}$ gives the outside block. Pieces $2,3$
    reconstruct $A^c$, equivalently the complement of pieces $0,1$, and
    pieces $0,3$ reconstruct $B^c$, equivalently the complement of pieces
    $1,2$, again both as explicit unions and as finite unions over the
    indexed subfamilies. Conversely, pieces $0,1$ are the
    complement of pieces $2,3$,
    pieces $1,2$ are the complement of pieces $0,3$, and adding all four
    pieces reconstructs $V$.
    Equivalently, each vertex lies in one of the four regions, and the indexed
    four-region family has union $V$. In fact, every vertex belongs to exactly
    one indexed piece.
\end{lemma}
\begin{proof}\leanok
    Check membership of an arbitrary vertex in $A$ and in $B$.
\end{proof}

\begin{lemma}[Disjointness in the four-region decomposition]%
    \label{lem:peps_regionUnionDecomposition_disjoint}
    \lean{TNLean.PEPS.regionOnlyLeft_disjoint_overlap}
    \lean{TNLean.PEPS.regionOnlyRight_disjoint_overlap}
    \lean{TNLean.PEPS.regionOnlyLeft_disjoint_onlyRight}
    \lean{TNLean.PEPS.regionOnlyLeft_disjoint_outside}
    \lean{TNLean.PEPS.regionOverlap_disjoint_outside}
    \lean{TNLean.PEPS.regionOnlyRight_disjoint_outside}
    \lean{TNLean.PEPS.regionInside_disjoint_outside}
    \lean{TNLean.PEPS.regionUnionPart_zero_one_disjoint_two_three}
    \lean{TNLean.PEPS.regionUnionPart_biUnion_zero_one_disjoint_two_three}
    \lean{TNLean.PEPS.regionUnionPart_one_two_disjoint_zero_three}
    \lean{TNLean.PEPS.regionUnionPart_biUnion_one_two_disjoint_zero_three}
    \lean{TNLean.PEPS.regionUnionPart_inside_disjoint_three}
    \lean{TNLean.PEPS.regionUnionPart_pairwise_disjoint}
    \leanok
    \uses{def:peps_regionUnionDecomposition,
        lem:peps_regionUnionDecomposition_identities}
    In the four-region decomposition, the left-only and right-only parts are
    disjoint from the overlap and from each other. Each inside part, and hence
    the union of the three inside regions, is disjoint from the outside region.
    In indexed form, the pieces forming $A$ are disjoint from the pieces
    forming $A^c$, and the pieces forming $B$ are disjoint from the pieces
    forming $B^c$, both as explicit binary unions and as finite unions over
    the corresponding selected indexed subfamilies. The indexed four-region
    family is pairwise disjoint.
\end{lemma}
\begin{proof}\leanok
    Use the standard disjointness identities for set difference and
    intersection, together with the reconstruction of $A\cup B$.
\end{proof}

\begin{definition}[Contraction maps for the union injectivity proof]%
    \label{def:peps_injectiveUnionContractionMaps}
    \lean{TNLean.PEPS.InjectiveUnionContractionMaps}
    \leanok
    Let $K$, $E_{012}$, $E_2$, and $E_{12}$ be complex vector spaces. A
    contraction-map family for the union-injectivity lemma in
    \cite[Section~4]{Molnar2018NormalPEPS} consists of three linear maps
    $\mathcal C_{\{0,1,2\}}:K\to E_{012}$,
    $\mathcal C_{\{2\}}:K\to E_2$, and
    $\mathcal C_{\{1,2\}}:K\to E_{12}$.
\end{definition}

\begin{definition}[Contraction chain for the union injectivity proof]%
    \label{def:peps_injectiveUnionContractionChain}
    \lean{TNLean.PEPS.InjectiveUnionContractionChain}
    \leanok
    \uses{def:peps_injectiveUnionContractionMaps}
    A contraction chain consists of the three contraction maps of
    Definition~\ref{def:peps_injectiveUnionContractionMaps}, together with the
    three implications used in
    Lemma~\ref{lem:peps_injectiveRegion_union_contractionChain}.
\end{definition}

\begin{lemma}[Kernel criterion for the union contraction chain]%
    \label{lem:peps_injectiveRegion_union_contractionChain}
    \lean{TNLean.PEPS.InjectiveUnionContractionChain.union_contraction_ker_eq_bot}
    \lean{TNLean.PEPS.InjectiveUnionContractionChain.union_contraction_injective}
    \leanok
    \uses{def:peps_injectiveUnionContractionChain}
    Let $P_0=A\setminus B$, $P_1=A\cap B$, $P_2=B\setminus A$, and
    $P_3=(A\cup B)^c$. Let $K$ be the boundary vector space attached to
    $P_3$, and let
    $\mathcal C_{\{0,1,2\}}:K\to E_{012}$,
    $\mathcal C_{\{2\}}:K\to E_2$, and
    $\mathcal C_{\{1,2\}}:K\to E_{12}$ be the three contraction maps in the
    proof of the union-injectivity lemma in
    \cite[Section~4]{Molnar2018NormalPEPS}. If
    \begin{align}
        \mathcal C_{\{0,1,2\}}(X)=0
        &\Longrightarrow \mathcal C_{\{2\}}(X)=0,
        \notag\\
        \mathcal C_{\{2\}}(X)=0
        &\Longrightarrow \mathcal C_{\{1,2\}}(X)=0,
        \notag
    \end{align}
    and $\mathcal C_{\{1,2\}}(X)=0\Longrightarrow X=0$, then
    $\ker \mathcal C_{\{0,1,2\}}=\{0\}$.
\end{lemma}
\begin{proof}\leanok
    For $X\in K$ with $\mathcal C_{\{0,1,2\}}(X)=0$, apply the three
    implications in sequence:
    \begin{align}
        \mathcal C_{\{0,1,2\}}(X)=0
        \Longrightarrow \mathcal C_{\{2\}}(X)=0
        \Longrightarrow \mathcal C_{\{1,2\}}(X)=0
        \Longrightarrow X=0.
        \notag
    \end{align}
\end{proof}

\begin{lemma}[Union of injective regions]%
    \label{lem:peps_injectiveRegion_union}
    \notready
    \uses{def:IsVertexInjective, lem:peps_regionUnionDecomposition_identities,
        lem:peps_regionUnionDecomposition_disjoint,
        lem:peps_injectiveRegion_union_contractionChain}
    Let $V$ be a finite vertex set. Write $\inj(R)$ for the
    assertion that the PEPS tensor blocked to the region $R\subseteq V$ is
    injective. For two possibly overlapping regions $A,B\subseteq V$,
    $\inj(A)\wedge\inj(B)\Longrightarrow\inj(A\cup B)$.
    The union-injectivity proof in \cite[Section~4]{Molnar2018NormalPEPS}
    blocks the network into four parts $A\setminus B$, $A\cap B$, $B\setminus A$, and
    $(A\cup B)^c$, and applies the two available left inverses in sequence:
    % Formula: the four blocked tensors are P_0=A\setminus B, P_1=A\cap B,
    %   P_2=B\setminus A, and P_3=(A\cup B)^c.
    % Source verdict: Molnar et al., arXiv:1804.04964, page 14, lines
    %   1329--1345 of paper_normal.tex.  The source draws the complete graph
    %   on the four blocked tensors and one upper physical leg at each tensor.
    % Ink-to-index map: every one of the six joins between distinct tensor
    %   blocks contracts one virtual index; the four upper joins end at
    %   inkless boundary atoms and carry the four free physical indices.
    % Type and contraction verdict: all six internal joins are
    %   virtual--virtual, all four open-leg joins are physical--physical, and
    %   the internal virtual graph is exactly K_4.
    % Structural boundary verdict: the typed boundary atoms and their joins
    %   derive (V,P)=(0,4).  The free tier does not emit a boundary event.
    \begin{center}
        \begin{tenkzfree}[compact]
            \tnput[dot, label pos=west,
                ports={45:virtual,east:virtual,-45:virtual,north:physical}]
                {unionLeft}{(-14mm,0)}{A\setminus B}
            \tnput[dot, label pos=east,
                ports={-135:virtual,-35:virtual,-70:virtual,north:physical}]
                {unionOverlap}{(-4mm,11mm)}{A\cap B}
            \tnput[dot, label pos=east,
                ports={west:virtual,135:virtual,-135:virtual,north:physical}]
                {unionRight}{(14mm,0)}{B\setminus A}
            \tnput[dot, label pos=south,
                ports={150:virtual,110:virtual,45:virtual,north:physical}]
                {unionOutside}{(4mm,-11mm)}{(A\cup B)^c}

            \tnput[boundary, ports={south:physical}]
                {unionLeftPhysical}{(-14mm,7mm)}{}
            \tnput[boundary, ports={south:physical}]
                {unionOverlapPhysical}{(-4mm,18mm)}{}
            \tnput[boundary, ports={south:physical}]
                {unionRightPhysical}{(14mm,7mm)}{}
            \tnput[boundary, ports={south:physical}]
                {unionOutsidePhysical}{(4mm,-4mm)}{}

            \tnjoin{unionLeft.45}{unionOverlap.-135}
            \tnjoin{unionLeft.east}{unionRight.west}
            \tnjoin{unionLeft.-45}{unionOutside.150}
            \tnjoin{unionOverlap.-35}{unionRight.135}
            \tnjoin{unionOverlap.-70}{unionOutside.110}
            \tnjoin{unionRight.-135}{unionOutside.45}
            \tnjoin{unionLeft.north}{unionLeftPhysical.south}
            \tnjoin{unionOverlap.north}{unionOverlapPhysical.south}
            \tnjoin{unionRight.north}{unionRightPhysical.south}
            \tnjoin{unionOutside.north}{unionOutsidePhysical.south}
        \end{tenkzfree}
    \end{center}
    Taking a boundary tensor $X$ on $(A\cup B)^c$ with
    $\mathcal C_{\{0,1,2\}}(X)=0$, injectivity of $A$ gives
    $\mathcal C_{\{2\}}(X)=0$; reinserting $A\cap B$ and applying
    injectivity of $B$ then gives $X=0$:
    % Formula: C_{\{0,1,2\}}(X)=0 -> C_{\{2\}}(X)=0
    %   -> C_{\{1,2\}}(X)=0 -> X=0.
    % Source verdict: Molnar et al., arXiv:1804.04964, page 14, lines
    %   1346--1401 of paper_normal.tex.  The arrows are, in order, the left
    %   inverse on A, reinsertion of A\cap B, and the left inverse on B.
    % Ink-to-index map: panel 1 is K_4 on P_0,P_1,P_2,X.  Panel 2 retains the
    %   P_2--X contraction and exposes their four bonds to the removed
    %   P_0,P_1 block.  Panel 3 restores P_1 and its triangle with P_2,X,
    %   leaving the three bonds to P_0 open.  Panel 4 leaves the three virtual
    %   indices incident on X open.  Upper legs are free physical indices.
    % Type and contraction verdict: every tensor--tensor join is
    %   virtual--virtual; every open stub terminates at a typed inkless
    %   boundary atom.  The panels have respectively 6, 1, 3, and 0 internal
    %   virtual contractions.
    % Structural boundary verdict: typed boundary atoms and their joins derive,
    %   from left to right, (V,P)=(0,3), (4,1), (3,2), and (3,0).  The free
    %   tier does not emit boundary events.
    \begin{center}
        \begin{tenkzfree}[compact]
            \tnput[dot, label pos=west,
                ports={45:virtual,east:virtual,-45:virtual,north:physical}]
                {proofZeroLeft}{(-7mm,0)}{P_0}
            \tnput[dot, label pos=east,
                ports={-135:virtual,-45:virtual,-90:virtual,north:physical}]
                {proofZeroOverlap}{(0,6mm)}{P_1}
            \tnput[dot, label pos=east,
                ports={west:virtual,135:virtual,-135:virtual,north:physical}]
                {proofZeroRight}{(7mm,0)}{P_2}
            \tnput[dot, label pos=south,
                ports={135:virtual,90:virtual,45:virtual}]
                {proofZeroX}{(0,-6mm)}{X}
            \tnput[boundary, ports={south:physical}]
                {proofZeroLeftPhysical}{(-7mm,4mm)}{}
            \tnput[boundary, ports={south:physical}]
                {proofZeroOverlapPhysical}{(0,10mm)}{}
            \tnput[boundary, ports={south:physical}]
                {proofZeroRightPhysical}{(7mm,4mm)}{}

            \tnjoin{proofZeroLeft.45}{proofZeroOverlap.-135}
            \tnjoin{proofZeroLeft.east}{proofZeroRight.west}
            \tnjoin{proofZeroLeft.-45}{proofZeroX.135}
            \tnjoin{proofZeroOverlap.-45}{proofZeroRight.135}
            \tnjoin{proofZeroOverlap.-90}{proofZeroX.90}
            \tnjoin{proofZeroRight.-135}{proofZeroX.45}
            \tnjoin{proofZeroLeft.north}{proofZeroLeftPhysical.south}
            \tnjoin{proofZeroOverlap.north}{proofZeroOverlapPhysical.south}
            \tnjoin{proofZeroRight.north}{proofZeroRightPhysical.south}
        \end{tenkzfree}
        \,$=0$\quad$\xrightarrow{A^{-1}}$\quad
        \begin{tenkzfree}[compact]
            \tnput[dot, label pos=east,
                ports={-135:virtual,135:virtual,-90:virtual,north:physical}]
                {proofOneRight}{(4mm,1mm)}{P_2}
            \tnput[dot, label pos=south,
                ports={45:virtual,90:virtual,150:virtual}]
                {proofOneX}{(-2mm,-6mm)}{X}
            \tnput[boundary, ports={east:virtual}]
                {proofOneRightLower}{(-3mm,-1mm)}{}
            \tnput[boundary, ports={-45:virtual}]
                {proofOneRightUpper}{(-2mm,7mm)}{}
            \tnput[boundary, ports={east:virtual}]
                {proofOneXLower}{(-8mm,-3mm)}{}
            \tnput[boundary, ports={south:virtual}]
                {proofOneXUpper}{(-2mm,1mm)}{}
            \tnput[boundary, ports={south:physical}]
                {proofOneRightPhysical}{(4mm,7mm)}{}

            \tnjoin{proofOneRight.-90}{proofOneX.45}
            \tnjoin{proofOneRight.-135}{proofOneRightLower.east}
            \tnjoin{proofOneRight.135}{proofOneRightUpper.-45}
            \tnjoin{proofOneX.150}{proofOneXLower.east}
            \tnjoin{proofOneX.90}{proofOneXUpper.south}
            \tnjoin{proofOneRight.north}{proofOneRightPhysical.south}
        \end{tenkzfree}
        \,$=0$\quad$\xrightarrow{A\cap B}$\quad
        \begin{tenkzfree}[compact]
            \tnput[dot, label pos=east,
                ports={-135:virtual,-90:virtual,-45:virtual,north:physical}]
                {proofTwoOverlap}{(0,6mm)}{P_1}
            \tnput[dot, label pos=east,
                ports={135:virtual,-135:virtual,west:virtual,north:physical}]
                {proofTwoRight}{(7mm,0)}{P_2}
            \tnput[dot, label pos=south,
                ports={90:virtual,45:virtual,135:virtual}]
                {proofTwoX}{(0,-6mm)}{X}
            \tnput[boundary, ports={45:virtual}]
                {proofTwoOverlapFree}{(-5mm,2mm)}{}
            \tnput[boundary, ports={east:virtual}]
                {proofTwoRightFree}{(-5mm,0)}{}
            \tnput[boundary, ports={-45:virtual}]
                {proofTwoXFree}{(-5mm,-2mm)}{}
            \tnput[boundary, ports={south:physical}]
                {proofTwoOverlapPhysical}{(0,10mm)}{}
            \tnput[boundary, ports={south:physical}]
                {proofTwoRightPhysical}{(7mm,4mm)}{}

            \tnjoin{proofTwoOverlap.-135}{proofTwoOverlapFree.45}
            \tnjoin{proofTwoRight.west}{proofTwoRightFree.east}
            \tnjoin{proofTwoX.135}{proofTwoXFree.-45}
            \tnjoin{proofTwoOverlap.-90}{proofTwoX.90}
            \tnjoin{proofTwoOverlap.-45}{proofTwoRight.135}
            \tnjoin{proofTwoRight.-135}{proofTwoX.45}
            \tnjoin{proofTwoOverlap.north}{proofTwoOverlapPhysical.south}
            \tnjoin{proofTwoRight.north}{proofTwoRightPhysical.south}
        \end{tenkzfree}
        \,$=0$\quad$\xrightarrow{B^{-1}}$\quad
        \begin{tenkzfree}[compact]
            \tnput[dot, label pos=south, ports={135:virtual,90:virtual,45:virtual}]
                {proofThreeX}{(0,0)}{X}
            \tnput[boundary, ports={-45:virtual}]
                {proofThreeLeft}{(-4mm,5mm)}{}
            \tnput[boundary, ports={south:virtual}]
                {proofThreeMiddle}{(0,6mm)}{}
            \tnput[boundary, ports={-135:virtual}]
                {proofThreeRight}{(4mm,5mm)}{}

            \tnjoin{proofThreeX.135}{proofThreeLeft.-45}
            \tnjoin{proofThreeX.90}{proofThreeMiddle.south}
            \tnjoin{proofThreeX.45}{proofThreeRight.-135}
        \end{tenkzfree}
        \,$=0$.
    \end{center}
\end{lemma}
\begin{proof}
    Put $P_0=A\setminus B$, $P_1=A\cap B$, $P_2=B\setminus A$, and
    $P_3=(A\cup B)^c$. Then $P_0\cup P_1=A$, $P_1\cup P_2=B$, and
    $P_0\cup P_1\cup P_2=A\cup B$. Let $\mathcal C_I$ denote contraction of
    the blocked tensor over the region $\bigcup_{i\in I}P_i$. To prove
    injectivity of $A\cup B$, take a boundary tensor $X$ on $P_3$ such that
    $\mathcal C_{\{0,1,2\}}(X)=0$. Since $P_0\cup P_1=A$ is injective, its
    left inverse gives
    $\mathcal C_{\{0,1,2\}}(X)=0\Longrightarrow
    \mathcal C_{\{2\}}(X)=0$.
    The remaining two steps are
    $\mathcal C_{\{2\}}(X)=0\Longrightarrow
    \mathcal C_{\{1,2\}}(X)=0$ and
    $\mathcal C_{\{1,2\}}(X)=0\Longrightarrow X=0$.
    The first implication is obtained by contracting with the tensor on
    $P_1$, and the second is the left inverse for the injective region
    $P_1\cup P_2=B$.
    Thus $\ker \mathcal C_{\{0,1,2\}}=\{0\}$, and therefore
    $P_0\cup P_1\cup P_2=A\cup B$ is injective.
\end{proof}

\begin{lemma}[Union of the edge-centred blocks]%
    \label{lem:peps_injectiveRegion_union_edgeBlocked}
    \lean{TNLean.PEPS.regionBlockedTensorInjective_union}
    \leanok
    \uses{def:peps_regionInjectivityDataOf}
    Fix an edge $e$ and a partition of the vertex set into the three
    edge-centred blocks $R$, $B$, $C$ (red, blue, complementary) of the normal
    PEPS proof, with the red block carrying one endpoint of $e$ and the blue
    block the other. The injectivity of the blue and complementary blocks is
    part of the edge-centred blocking data; if every virtual bond dimension
    is positive, then their union $B\cup C=V\setminus R$ is injective.
\end{lemma}
\begin{proof}\leanok
    This is Lemma~\ref{lem:peps_injectiveRegion_union} for the
    edge-centred triple, where the two regions are the blue and complementary
    blocks and their union is the set complement of the red block. A coefficient
    family annihilating the blocked weights of $V\setminus R$ is stripped of the
    blue block by its left inverse, leaving a complement-coupling combination
    that vanishes for every blue boundary configuration. The complementary block
    then carries each coupling coefficient, read as a function of the
    complementary physical leg, into a combination of its blocked weights whose
    coefficient at a boundary configuration is the host residual coefficient
    reconstructed from the blue and complementary boundary data; injectivity of
    the complementary block forces that coefficient to vanish wherever a global
    configuration realizes the two boundary labels. Since every host residual is
    realized once the bond dimensions are positive, the coefficient family is
    zero.
\end{proof}

\begin{lemma}[Union of two disjoint injective regions]%
    \label{lem:peps_injectiveRegion_union_disjoint}
    \lean{TNLean.PEPS.regionBlockedTensorInjective_union_disjoint}
    \leanok
    \uses{lem:peps_injectiveRegion_union_edgeBlocked}
    Let $S,T\subseteq V$ be two disjoint regions. If the blocked tensors of $S$
    and of $T$ are each injective and every virtual bond dimension is positive,
    then the blocked tensor of their union $S\cup T$ is injective.
\end{lemma}
\begin{proof}\leanok
    This is Lemma~\ref{lem:peps_injectiveRegion_union} in the case
    where the two regions are disjoint, so that the overlap $S\cap T$ is empty
    and the four-region picture collapses to the three blocks $S$, $T$, and the
    set complement $(S\cup T)^c$. The argument places $S$ and $T$ as the blue and
    complementary blocks of a three-block decomposition whose third block is the
    set complement of their union, and applies the two-step inverse argument of
    the edge-centred Lemma~\ref{lem:peps_injectiveRegion_union_edgeBlocked}
    to that decomposition. No injectivity hypothesis on the set complement and no
    distinguished edge are needed.
\end{proof}

\begin{lemma}[Union of two injective regions with positive bonds]%
    \label{lem:peps_injectiveRegion_union_pos}
    \lean{TNLean.PEPS.regionBlockedTensorInjective_union_overlap}
    \leanok
    \uses{lem:peps_injectiveRegion_union_edgeBlocked}
    Let $A,B\subseteq V$ be two possibly overlapping regions. If the blocked
    tensors of $A$ and of $B$ are each injective and every virtual bond dimension
    is positive, then the blocked tensor of their union $A\cup B$ is injective.
    This is Lemma~\ref{lem:peps_injectiveRegion_union} with the explicit
    positive-bond hypothesis carried by the formalization; the overlap is no longer
    required to be empty.
\end{lemma}
\begin{proof}\leanok
    Write the four parts $P_0=A\setminus B$, $P_1=A\cap B$, $P_2=B\setminus A$.
    A coefficient family annihilating the blocked weights of $A\cup B$ is stripped
    of $A$ by its left inverse, leaving the coupling through $P_2$ vanishing for
    every $A$-boundary configuration. The source then reinserts the overlap $P_1$
    and inverts $B$. The legs running from $P_0$ to the outside are the open legs
    of the final tensor; they stay uncontracted through the reinsertion. Carrying
    those legs as a separate label, fix it to a reference value and restrict the
    coefficient family to the matching fiber. The fiber-restricted first strip
    still vanishes, so reinserting $P_1$ and inverting $B$ forces the fiber row to
    vanish. The union boundary edges partition into the $B$-boundary edges and the
    $P_0$-outer edges, so fixing both the reference label and the $B$-boundary
    label of a configuration realizing a host label selects that single host term,
    forcing the coefficient to vanish at every realizable host label. Surjectivity
    of the union host boundary, available once the bond dimensions are positive,
    realizes every label, so the coefficient family is zero.
\end{proof}

\begin{lemma}[Finite disjoint union of injective regions]%
    \label{lem:peps_injectiveRegion_biUnion_disjoint}
    \lean{TNLean.PEPS.regionBlockedTensorInjective_biUnion_disjoint}
    \leanok
    \uses{lem:peps_injectiveRegion_union_disjoint}
    Let $(R_i)_{i\in s}$ be a nonempty finite family of pairwise-disjoint regions.
    If the blocked tensor of each $R_i$ is injective and every virtual bond
    dimension is positive, then the blocked tensor of their union
    $\bigcup_{i\in s} R_i$ is injective.
\end{lemma}
\begin{proof}\leanok
    Induct on the nonempty family. The singleton case is the injectivity of the
    one block. For the inductive step, the next block is disjoint from the union
    of the remaining blocks, so
    Lemma~\ref{lem:peps_injectiveRegion_union_disjoint} combines the two
    injective pieces.
\end{proof}

\begin{lemma}[Finite union of injective regions]%
    \label{lem:peps_injectiveRegion_biUnion}
    \lean{TNLean.PEPS.regionBlockedTensorInjective_biUnion_overlap}
    \leanok
    \uses{lem:peps_injectiveRegion_union_pos}
    Let $(R_i)_{i\in s}$ be a nonempty finite family of possibly overlapping
    regions. If the blocked tensor of each $R_i$ is injective and every virtual
    bond dimension is positive, then the blocked tensor of their union
    $\bigcup_{i\in s} R_i$ is injective. This is the finite iteration of the
    overlapping union lemma, with no disjointness assumption.
\end{lemma}
\begin{proof}\leanok
    Induct on the nonempty family. The singleton case is the injectivity of the
    one block. For the inductive step, combine the next block with the union of
    the remaining blocks by Lemma~\ref{lem:peps_injectiveRegion_union_pos}; no
    disjointness is needed.
\end{proof}

\begin{definition}[Injectivity predicate for blocked regions]%
    \label{def:peps_regionInjectivityData}
    \lean{TNLean.PEPS.RegionInjectivityData}
    \lean{TNLean.PEPS.RegionInjectivityUnionClosure}
    \leanok
    A region-injectivity hypothesis assigns to each finite vertex region the
    assertion that the tensor obtained by blocking that region is injective.
    This predicate is used for regions such as $R$, $S$, $T$, and their
    complements in the normal PEPS proof. The union-closure hypothesis records
    the conclusion supplied by the source union-of-injective-regions lemma:
    the union of two injective regions is injective.
\end{definition}

\begin{definition}[Concrete region-injectivity predicate of a tensor]%
    \label{def:peps_regionInjectivityDataOf}
    \lean{TNLean.PEPS.regionInjectivityDataOf}
    \leanok
    \uses{def:peps_regionInjectivityData}
    The region-injectivity predicate of a fixed PEPS tensor $A$ sends a finite
    region $R$ to the assertion that the blocked-region tensor of $R$ is
    injective.
\end{definition}

\begin{theorem}[Gauged blocked-region weight as a global sum]
    \label{thm:peps_regionBlockedWeight_applyGauge_eq_globalSum}
    \lean{TNLean.PEPS.regionBlockedWeight_applyGauge_eq_globalSum}
    \leanok
    \uses{def:applyGauge}
    Let $\widetilde B=\operatorname{Gauge}_X(B)$ be obtained from $B$ by a gauge
    family $X$. For a region $R$, a boundary configuration $\mu$, and a
    physical configuration $\tau$ on $R$, the blocked weight of $\widetilde B$
    is
    \begin{align}
        T_{\widetilde B,R}^{\mu}(\tau)
        &=
        \sum_{\zeta}
        \left(\prod_{f\in\partial R}G_f^X(\mu_f,\zeta_f)\right)
        \prod_{v\in R}B_v(\zeta|_{\inc(v)},\tau_v),
        \label{eq:peps_gauged_global_sum}
    \end{align}
    where the sum ranges over global virtual configurations $\zeta$, and
    $G_f^X$ is the surviving gauge on the boundary edge $f$.
\end{theorem}
\begin{proof}\leanok
    In the contraction defining $T_{\widetilde B,R}^{\mu}(\tau)$, the two
    half-edge gauge factors on every edge internal to $R$ combine as
    \begin{align}
        \sum_a (X_e)_{\alpha,a}(X_e^{-1})_{a,\beta}
        &=
        (X_eX_e^{-1})_{\alpha,\beta}
        =
        \delta_{\alpha,\beta}.
        \notag
    \end{align}
    Thus every interior edge contributes the identity and its gauge cancels.
    The only remaining factors are the boundary gauges
    $G_f^X(\mu_f,\zeta_f)$, giving (\ref{eq:peps_gauged_global_sum}).
\end{proof}

\begin{theorem}[Boundary-coupling form of the gauged blocked weight]
    \label{thm:peps_regionBlockedWeight_applyGauge}
    \lean{TNLean.PEPS.regionBlockedWeight_applyGauge}
    \leanok
    \uses{thm:peps_regionBlockedWeight_applyGauge_eq_globalSum}
    With the same notation, set
    $K_R^X(\mu,\nu):=\prod_{f\in\partial R}G_f^X(\mu_f,\nu_f)$.
    Then the gauged blocked weight is the original blocked weight multiplied by
    the boundary-coupling matrix:
    \begin{align}
        T_{\widetilde B,R}^{\mu}(\tau)
        &=
        \sum_{\nu}K_R^X(\mu,\nu)\,T_{B,R}^{\nu}(\tau).
        \label{eq:peps_boundary_coupling}
    \end{align}
    This is the gauge absorption of the blocked tensor in
    \cite[Section~4, proof of Theorem~3]{Molnar2018NormalPEPS}.
\end{theorem}
\begin{proof}\leanok
    Start from Theorem~\ref{thm:peps_regionBlockedWeight_applyGauge_eq_globalSum}
    and group the global virtual configurations by their boundary label $\nu$:
    \begin{align}
        \sum_{\zeta}
        K_R^X(\mu,\zeta|_{\partial R})
        \prod_{v\in R}B_v(\zeta|_{\inc(v)},\tau_v)
        &=
        \sum_{\nu}K_R^X(\mu,\nu)\,T_{B,R}^{\nu}(\tau).
        \notag
    \end{align}
\end{proof}

\begin{theorem}[The boundary coupling has a right inverse]
    \label{thm:peps_sum_regionBoundaryCoupling_mul_inv}
    \lean{TNLean.PEPS.sum_regionBoundaryCoupling_mul_inv}
    \leanok
    Let
    $K_R^X(\mu,\nu)=\prod_{f\in\partial R}G_f^X(\mu_f,\nu_f)$ and
    $(K_R^X)^{-1}(\nu,\rho)
    =\prod_{f\in\partial R}(G_f^X)^{-1}(\nu_f,\rho_f)$.
    Then, for boundary configurations $\mu$ and $\rho$,
    \begin{align}
        \sum_{\nu}K_R^X(\mu,\nu)\,(K_R^X)^{-1}(\nu,\rho)
        &=
        \begin{cases}
            1, & \mu=\rho,\\
            0, & \mu\neq\rho.
        \end{cases}
        \label{eq:peps_boundary_coupling_inverse}
    \end{align}
\end{theorem}
\begin{proof}\leanok
    Exchange the finite sum over $\nu$ with the product over boundary edges:
    \begin{align}
        \sum_{\nu}K_R^X(\mu,\nu)\,(K_R^X)^{-1}(\nu,\rho)
        &=
        \prod_{f\in\partial R}
        \sum_a G_f^X(\mu_f,a)\,(G_f^X)^{-1}(a,\rho_f).
        \notag
    \end{align}
    Since
    \begin{align}
        \sum_a G_f^X(\mu_f,a)\,(G_f^X)^{-1}(a,\rho_f)
        &=
        (G_f^X(G_f^X)^{-1})_{\mu_f,\rho_f}
        =
        \delta_{\mu_f,\rho_f},
        \notag
    \end{align}
    each factor reduces to the Kronecker symbol at $(\mu_f,\rho_f)$. The
    product of these one-edge Kronecker symbols is $\delta_{\mu,\rho}$.
\end{proof}

\begin{theorem}[Gauge preservation of blocked-region injectivity]
    \label{thm:peps_regionBlockedTensorInjective_applyGauge}
    \lean{TNLean.PEPS.regionBlockedTensorInjective_applyGauge}
    \leanok
    \uses{def:applyGauge, thm:peps_regionBlockedWeight_applyGauge,
        thm:peps_sum_regionBoundaryCoupling_mul_inv}
    Let $\widetilde B=\operatorname{Gauge}_X(B)$. If the blocked tensor family
    $\{T_{B,R}^{\nu}\}_{\nu}$ is linearly independent, then the blocked tensor
    family $\{T_{\widetilde B,R}^{\mu}\}_{\mu}$ is linearly independent.
    Equivalently,
    $\inj(T_{B,R})\Longrightarrow\inj(T_{\widetilde B,R})$.
\end{theorem}
\begin{proof}\leanok
    By (\ref{eq:peps_boundary_coupling}),
    $T_{\widetilde B,R}^{\mu}
    =\sum_{\nu}K_R^X(\mu,\nu)\,T_{B,R}^{\nu}$.
    Suppose $\sum_\mu c_\mu T_{\widetilde B,R}^{\mu}=0$. Substituting
    (\ref{eq:peps_boundary_coupling}) and using the linear independence of
    $\{T_{B,R}^{\nu}\}_{\nu}$ gives
    $\sum_{\mu}c_\mu K_R^X(\mu,\nu)=0$ for every $\nu$.
    Contracting this equality with the right inverse of $K_R^X$ yields, for
    every $\rho$,
    \begin{align}
        c_\rho
        &=
        \sum_{\nu}
        \left(\sum_{\mu}c_\mu K_R^X(\mu,\nu)\right)
        (K_R^X)^{-1}(\nu,\rho)
        =
        0.
        \notag
    \end{align}
    Thus all coefficients $c_\rho$ vanish, proving the linear independence of
    $\{T_{\widetilde B,R}^{\mu}\}_{\mu}$.
\end{proof}

\begin{theorem}[Kernel descent for blocked regions]%
    \label{thm:peps_regionBlockedTensorInjective_of_isVertexInjective}
    \lean{TNLean.PEPS.regionBlockedTensorInjective_of_isVertexInjective}
    \leanok
    \uses{def:IsVertexInjective, def:peps_regionInjectivityDataOf}
    Let $A$ be vertex-injective, and assume that every virtual bond dimension is
    positive. Then every finite region $R$ has an injective blocked tensor:
    $\inj(R)$.
% Paper-gap note: docs/paper-gaps/peps_injective_ft_section3_route.tex.
\end{theorem}
\begin{proof}\leanok
    Let $(c_\rho)_\rho$ be a finitely supported coefficient family on the
    boundary configurations of $R$, and assume that, for every physical
    configuration $\tau$,
    $\sum_{\rho} c_\rho\,T_{A,R}^{\rho}(\tau)=0$.
    The kernel descent removes the vertices of $R$ one at a time and preserves
    the vanishing contraction. After all vertices have been removed, the
    terminal contraction gives $c_\rho=0$ for every boundary configuration
    $\rho$. Hence the family $\{T_{A,R}^{\rho}\}_\rho$ is linearly independent,
    which is $\inj(R)$.
\end{proof}

\begin{theorem}[Union closure for a vertex-injective tensor]%
    \label{thm:peps_regionInjectivityUnionClosure_of_isVertexInjective}
    \lean{TNLean.PEPS.regionInjectivityUnionClosure_of_isVertexInjective}
    \leanok
    \uses{def:IsVertexInjective, def:peps_regionInjectivityDataOf,
        thm:peps_regionBlockedTensorInjective_of_isVertexInjective}
    Let $A$ be vertex-injective with every bond dimension positive. Then the
    concrete region-injectivity predicate of $A$ satisfies union closure: for
    any two finite regions $R$ and $S$,
    $\inj(R)\wedge\inj(S)\Longrightarrow\inj(R\cup S)$.
% Paper-gap note: docs/paper-gaps/peps_injective_ft_section3_route.tex.
\end{theorem}
\begin{proof}\leanok
    The kernel descent theorem gives $\inj(Q)$ for every finite region $Q$.
    Taking $Q=R\cup S$ gives $\inj(R\cup S)$, so the union-closure implication
    holds independently of the injectivity of $R$ and $S$.
\end{proof}

\begin{theorem}[Union closure from the overlapping union lemma]%
    \label{thm:peps_regionInjectivityUnionClosure_of_overlap}
    \lean{TNLean.PEPS.regionInjectivityUnionClosure_of_overlap}
    \leanok
    \uses{def:peps_regionInjectivityDataOf,
        lem:peps_injectiveRegion_union_pos}
    Let $A$ be a PEPS tensor with every bond dimension positive. Then the
    concrete region-injectivity predicate of $A$ satisfies union closure: for
    any two finite regions $R$ and $S$,
    $\inj(R)\wedge\inj(S)\Longrightarrow\inj(R\cup S)$.
    Thus injectivity of the two regions implies injectivity of their union;
    no single-vertex injectivity hypothesis is needed.
    % Source-faithful scope note: this is the union-of-injective-regions
    % lemma, rather than the stronger vertex-injective specialization.
\end{theorem}
\begin{proof}\leanok
    The union-closure implication is the overlapping union lemma
    (Lemma~\ref{lem:peps_injectiveRegion_union_pos}): injectivity of $R$ and of
    $S$, together with the positive bond dimensions, gives injectivity of
    $R\cup S$.
\end{proof}

\begin{lemma}[Finite unions of injective regions]%
    \label{lem:peps_regionInjectivityUnionClosure_finite}
    \lean{TNLean.PEPS.RegionInjectivityUnionClosure.biUnion_injective}
    \leanok
    \uses{def:peps_regionInjectivityData}
    Under the union-closure hypothesis for region injectivity, any nonempty
    finite union of injective finite regions is injective.
\end{lemma}
\begin{proof}\leanok
    Induct on the nonempty finite indexing set. The singleton case is the
    given injectivity hypothesis, and the induction step is one application of
    binary union closure.
\end{proof}

\begin{lemma}[Indexed inside union under closure]%
    \label{lem:peps_regionInjectivityUnionClosure_inside}
    \lean{TNLean.PEPS.RegionInjectivityUnionClosure.regionUnionPart_inside_injective}
    \lean{TNLean.PEPS.RegionInjectivityUnionClosure.regionUnionPart_inside_binary_injective}
    \leanok
    \uses{def:peps_regionInjectivityData,
        lem:peps_regionUnionDecomposition_identities}
    Under the union-closure hypothesis for region injectivity, if the indexed
    subfamilies $\{0,1\}$ and $\{1,2\}$ in the four-region decomposition
    are injective, then the inside subfamily $\{0,1,2\}$ is injective.
    Equivalently, injectivity of the two source-paper regions
    $(A\setminus B)\cup(A\cap B)$ and
    $(A\cap B)\cup(B\setminus A)$ gives injectivity of
    $(A\setminus B)\cup(A\cap B)\cup(B\setminus A)=A\cup B$.
\end{lemma}
\begin{proof}\leanok
    Rewrite the three selected indexed unions as $A$, $B$, and $A\cup B$,
    then apply binary union closure.
\end{proof}

\section{Square-Lattice Normal PEPS}\label{sec:peps_normal_square}
On the square lattice the injective regions are built from $2\times3$ and
$3\times2$ coordinate rectangles. This section fixes the geometry: the nested
regions $R\subseteq S$ with $S\setminus R$ a single site, which realize the
one-site-different injective regions the normal theorem requires at each vertex,
together with the local window $T$ used around an edge. The union closure of
Section~\ref{sec:peps_normal_union} builds these regions from injective
rectangles, so a rectangular injectivity hypothesis yields their injectivity.

\begin{definition}[Complement of a finite region]%
    \label{def:peps_regionComplement}
    \lean{TNLean.PEPS.regionComplement}
    \lean{TNLean.PEPS.mem_regionComplement}
    \leanok
    For a finite vertex region $R\subseteq V$, its complement is
    $V\setminus R$. Equivalently, a vertex lies in the complement exactly
    when it does not lie in $R$.
\end{definition}

\begin{definition}[Square-lattice coordinate regions]%
    \label{def:peps_squareLatticeCoordinateRegions}
    \lean{TNLean.PEPS.SquareLatticeVertex}
    \lean{TNLean.PEPS.squareLatticeHorizontalNeighbor}
    \lean{TNLean.PEPS.squareLatticeVerticalNeighbor}
    \lean{TNLean.PEPS.squareLatticeGraph}
    \lean{TNLean.PEPS.IsHorizontalSquareLatticeEdge}
    \lean{TNLean.PEPS.IsVerticalSquareLatticeEdge}
    \lean{TNLean.PEPS.squareLatticeRectangle}
    \lean{TNLean.PEPS.squareLatticeFull}
    \lean{TNLean.PEPS.squareLatticeCoordinateInterval}
    \lean{TNLean.PEPS.squareLatticeContiguousRectangle}
    \lean{TNLean.PEPS.IsTwoByThreeSquareLatticeProduct}
    \lean{TNLean.PEPS.IsThreeByTwoSquareLatticeProduct}
    \lean{TNLean.PEPS.IsTwoByThreeContiguousSquareLatticeRectangle}
    \lean{TNLean.PEPS.IsThreeByTwoContiguousSquareLatticeRectangle}
    \lean{TNLean.PEPS.normalSquareRegionR}
    \lean{TNLean.PEPS.normalSquareRegionS}
    \lean{TNLean.PEPS.normalSquareRegionTVerticalBlock}
    \lean{TNLean.PEPS.normalSquareRegionTHorizontalBlock}
    \lean{TNLean.PEPS.normalSquareRegionTHole}
    \lean{TNLean.PEPS.normalSquareRegionT}
    \lean{TNLean.PEPS.normalSquareEdgeComplementRegion}
    \lean{TNLean.PEPS.squareLatticeGraph_adj}
    \lean{TNLean.PEPS.squareLatticeGraph_adj_right}
    \lean{TNLean.PEPS.squareLatticeGraph_adj_up}
    \lean{TNLean.PEPS.squareLatticeRightEdge}
    \lean{TNLean.PEPS.squareLatticeUpEdge}
    \lean{TNLean.PEPS.squareLatticeRightEdge_isHorizontal}
    \lean{TNLean.PEPS.squareLatticeUpEdge_isVertical}
    \lean{TNLean.PEPS.squareLatticeEdge_horizontal_or_vertical}
    \lean{TNLean.PEPS.squareLatticeEdge_not_horizontal_and_vertical}
    \lean{TNLean.PEPS.horizontalSquareLatticeEdge_coords}
    \lean{TNLean.PEPS.horizontalSquareLatticeEdge_eq_rightEdge}
    \lean{TNLean.PEPS.verticalSquareLatticeEdge_coords}
    \lean{TNLean.PEPS.verticalSquareLatticeEdge_eq_upEdge}
    \leanok
    The finite $n\times m$ square lattice is represented by pairs of
    coordinates. The corresponding square-lattice graph has horizontal and
    vertical nearest-neighbor edges; the two elementary adjacency lemmas record
    the right-neighbor and upper-neighbor cases used for translated edge
    blockings, and the corresponding coordinate edges name those two
    nearest-neighbor edges as elements of the edge set. These named right and
    upward edges have the expected horizontal and vertical orientations.
    Every edge of this
    graph is horizontal or vertical, but not both. Since edges are represented
    by ordered endpoint pairs, a horizontal edge is oriented from left to right
    and is exactly the coordinate right edge from its left endpoint; similarly,
    a vertical edge is oriented from bottom to top and is exactly the
    coordinate upward edge from its lower endpoint.
    A coordinate rectangle is a product of a horizontal coordinate set and a
    vertical coordinate set. The $2\times3$ and
    $3\times2$ coordinate-product predicates assert only that these two
    coordinate sets have cardinalities $2,3$ and $3,2$, respectively.
    Contiguous coordinate rectangles are products of contiguous coordinate
    intervals, and the contiguous $2\times3$ and $3\times2$ predicates name the
    rectangular blocks used in the normal square-lattice theorem. The
    displayed region $R$ is the union of a $2\times3$ contiguous rectangle
    with the $3\times2$ contiguous rectangle occupying its upper two rows and
    extending one column to the right. The displayed region $S$ is the union
    of two $2\times3$ contiguous rectangles shifted by one horizontal
    coordinate, hence the full $3\times3$ square spanned by them. The
    displayed local-window model for $T$ is the complement, inside a local
    $5\times6$ coordinate window, of the vertical $2\times3$ edge block
    and the horizontal $3\times2$ edge block shown in the source picture.
    The finite-lattice complementary block is the whole finite square lattice
    with those two edge blocks removed; this is the block denoted $A_3$ in
    the proof of \cite[Theorem~3]{Molnar2018NormalPEPS}.
    The source shapes are:
    \begin{tenkzequation}
        % Source/formula verdict: Molnar et al., arXiv:1804.04964,
        % paper_normal.tex:1405--1429 and source PDF p. 15 depict
        % R=(2-4,2-4)\{(4,4)} and S=(2-4,2-4).
        % Ink-to-index verdict: every dot is an ordinary tensor; every lattice
        % edge is a virtual index, and the selected fill records membership only.
        % Structural count verdict (not mechanically emitted as a lattice
        % signature): each 5-by-5 panel has 25 sites, 40 internal virtual
        % contractions, and 20 cut virtual bonds; (4,4) remains ordinary.
        % Structurally derived boundary/type verdict: no physical index is
        % exposed, so each panel has (V,P)=(20,0).
        \tenkzkernel
        \begin{tenkz}[rows={ket,ket,ket,ket,ket}, cols=5,
            west=open, east=open, north=open, south=open]
            \tnmark[form=enclosure, slot=selected, tint, label pos=nw]
                {(2,2) .. (4,4) - (4,4)}{$R$}
        \end{tenkz}
        $\qquad\text{and}\qquad$
        \begin{tenkz}[rows={ket,ket,ket,ket,ket}, cols=5,
            west=open, east=open, north=open, south=open]
            \tnmark[form=enclosure, slot=selected, tint, label pos=nw]
                {(2,2) .. (4,4)}{$S$}
        \end{tenkz}
    \end{tenkzequation}
    \begin{tenkzequation}
        % Formula verdict: in row-column coordinates the exact local formal
        % region is T=(1-6,1-5) minus (3-5,1-2) and (2-3,3-5); the two
        % removed blocks are respectively 3-by-2 and 2-by-3 cell rectangles.
        % Ink-to-index verdict: every dot is an ordinary tensor, every lattice
        % edge is a virtual contraction, and the selected fill records only
        % membership in T; it adds no contraction or tensor operation.
        % Contraction/boundary verdict: the 6-by-5 panel has 49 internal
        % virtual contractions and 22 cut virtual bonds, with no exposed
        % physical indices, hence boundary type (V,P)=(22,0).
        % Source verdict: arXiv:1804.04964, paper_normal.tex:1430--1445
        % depicts the same 5-by-6 window complement; this panel uses the exact
        % cell set of normalSquareRegionT rather than the source's rounded ink.
        \tenkzkernel
        \begin{tenkz}[rows={ket,ket,ket,ket,ket,ket}, cols=5,
            west=open, east=open, north=open, south=open]
            \tnmark[form=enclosure, slot=selected, tint, label pos=n]
                {(1,1) .. (6,5) - (3,1) .. (5,2) - (2,3) .. (3,5)}{$T$}
        \end{tenkz}
    \end{tenkzequation}
    % Maintainer note: injectivity of this finite-lattice complementary block
    % is not yet proved; this definition only names the region.
\end{definition}

\begin{lemma}[Basic identities for square-lattice coordinate regions]%
    \label{lem:peps_squareLatticeCoordinateRegion_identities}
    \lean{TNLean.PEPS.mem_squareLatticeRectangle}
    \lean{TNLean.PEPS.mem_squareLatticeCoordinateInterval}
    \lean{TNLean.PEPS.mem_squareLatticeContiguousRectangle}
    \lean{TNLean.PEPS.mem_normalSquareRegionR}
    \lean{TNLean.PEPS.mem_normalSquareRegionS}
    \lean{TNLean.PEPS.mem_normalSquareRegionTVerticalBlock}
    \lean{TNLean.PEPS.mem_normalSquareRegionTHorizontalBlock}
    \lean{TNLean.PEPS.mem_normalSquareRegionTHole}
    \lean{TNLean.PEPS.mem_normalSquareRegionT}
    \lean{TNLean.PEPS.mem_normalSquareEdgeComplementRegion}
    \lean{TNLean.PEPS.isTwoByThreeContiguousSquareLatticeRectangle_of_bounds}
    \lean{TNLean.PEPS.isThreeByTwoContiguousSquareLatticeRectangle_of_bounds}
    \lean{TNLean.PEPS.card_squareLatticeRectangle}
    \lean{TNLean.PEPS.squareLatticeFull_eq_univ}
    \leanok
    \uses{def:peps_squareLatticeCoordinateRegions}
    Membership in a coordinate rectangle is coordinatewise membership, the
    contiguous coordinate interval has the expected coordinate inequalities,
    membership in a contiguous coordinate rectangle is the conjunction of the
    four bounding inequalities, and membership in $R$ is the disjunction of
    membership in the $2\times3$ lower-left piece or in the shifted
    $3\times2$ upper piece. Membership in $S$ is the disjunction of
    membership in the left $2\times3$ piece or in the horizontally shifted
    $2\times3$ piece. Membership in the local-window $T$ region is
    membership in the local $5\times6$ window together with nonmembership in
    the union of the two displayed edge blocks. Bounded contiguous coordinate
    rectangles of sizes $2\times3$ and $3\times2$ satisfy, respectively,
    the predicates for
    contiguous $2\times3$ and contiguous $3\times2$ square-lattice
    rectangles. The cardinality of a coordinate rectangle is the
    product of the two coordinate cardinalities, and the full coordinate
    rectangle is the whole finite vertex set.
\end{lemma}
\begin{proof}\leanok
    The coordinate-rectangle statements follow from membership and cardinality
    identities for Cartesian products of finite coordinate sets. The
    contiguous-rectangle statement then unfolds the two coordinate intervals.
    The membership statements for $R$, $S$, and the two edge blocks removed
    from $T$ additionally unfold membership in a union of two such
    rectangles, while membership in $T$ also unfolds the set difference from
    the local $5\times6$ window. The shape predicates follow by using the
    displayed lower-left corners as witnesses.
\end{proof}

\begin{lemma}[The normal regions $R$ and $S$ are unions of contiguous
        rectangles]%
    \label{lem:peps_normalSquareRegionRS_rectangular_decomposition}
    \lean{TNLean.PEPS.normalSquareRegionR_rectangular_decomposition}
    \lean{TNLean.PEPS.normalSquareRegionS_rectangular_decomposition}
    \leanok
    \uses{def:peps_squareLatticeCoordinateRegions,
        lem:peps_squareLatticeCoordinateRegion_identities}
    If the displayed $3\times3$ window containing $R$ and $S$ lies in
    the finite square lattice, then $R$ is the union of a contiguous
    $2\times3$ rectangle and a contiguous $3\times2$ rectangle, while
    $S$ is the union of two contiguous $2\times3$ rectangles.
\end{lemma}
\begin{proof}\leanok
    This is the defining description of $R$ and $S$, together with the
    coordinate bounds which make the displayed rectangles contiguous
    $2\times3$ and $3\times2$ rectangles.
\end{proof}

\begin{lemma}[The edge blocks removed from $T$ are contiguous rectangles]%
    \label{lem:peps_normalSquareRegionTHole_rectangular_decomposition}
    \lean{TNLean.PEPS.normalSquareRegionTVerticalBlock_rectangular}
    \lean{TNLean.PEPS.normalSquareRegionTHorizontalBlock_rectangular}
    \lean{TNLean.PEPS.normalSquareRegionTHole_rectangular_decomposition}
    \leanok
    \uses{def:peps_squareLatticeCoordinateRegions,
        lem:peps_squareLatticeCoordinateRegion_identities}
    If the two edge blocks cut out of the displayed $T$-region lie in the
    finite square lattice, then they are respectively a contiguous
    $2\times3$ rectangle and a contiguous $3\times2$ rectangle.
\end{lemma}
\begin{proof}\leanok
    This is the defining description of the two removed edge blocks, together
    with the coordinate bounds which make them valid contiguous rectangles.
\end{proof}

\begin{lemma}[The local $T$-region and its removed edge blocks fill the window]%
    \label{lem:peps_normalSquareRegionT_window_complement}
    \lean{TNLean.PEPS.normalSquareRegionTHole_subset_window}
    \lean{TNLean.PEPS.normalSquareRegionTVerticalBlock_disjoint_horizontalBlock}
    \lean{TNLean.PEPS.normalSquareRegionT_disjoint_THole}
    \lean{TNLean.PEPS.normalSquareRegionT_disjoint_verticalBlock}
    \lean{TNLean.PEPS.normalSquareRegionT_disjoint_horizontalBlock}
    \lean{TNLean.PEPS.normalSquareRegionT_union_THole}
    \lean{TNLean.PEPS.normalSquareRegionT_union_verticalBlock_union_horizontalBlock}
    \leanok
    \uses{def:peps_squareLatticeCoordinateRegions,
        lem:peps_squareLatticeCoordinateRegion_identities}
    The two edge blocks removed from the displayed local $T$-region lie inside
    the local $5\times6$ coordinate window. They are disjoint from each other
    and from this local $T$ region, and their union with it is exactly that
    local window.
\end{lemma}
\begin{proof}\leanok
    This follows by unfolding the definition of $T$ as the set-theoretic
    complement of the two displayed edge blocks inside the local window.
\end{proof}

\begin{lemma}[The edge-complementary block fills the finite lattice]%
    \label{lem:peps_normalSquareEdgeComplementRegion}
    \lean{TNLean.PEPS.normalSquareRegionT_subset_edgeComplementRegion}
    \lean{TNLean.PEPS.normalSquareEdgeComplementRegion_disjoint_THole}
    \lean{TNLean.PEPS.normalSquareEdgeComplementRegion_union_THole}
    \lean{TNLean.PEPS.normalSquareEdgeComplementRegion_union_verticalBlock_union_horizontalBlock}
    \leanok
    \uses{def:peps_squareLatticeCoordinateRegions,
        lem:peps_squareLatticeCoordinateRegion_identities}
    The finite-lattice complementary block around the edge is disjoint from
    the two displayed edge blocks, and its union with them is the full finite
    square lattice. The local-window $T$ region is contained in this
    finite-lattice complementary block.
\end{lemma}
\begin{proof}\leanok
    This is the set-theoretic complement of the two displayed edge blocks in
    the full finite square lattice. The containment of the local-window
    $T$ region follows because its points also avoid the two removed edge
    blocks.
\end{proof}

\begin{definition}[Rectangular covers of square-lattice regions]%
    \label{def:peps_normalSquareRectangleCover}
    \lean{TNLean.PEPS.SquareLatticeRectangleCover}
    \lean{TNLean.PEPS.NormalSquareRegionTRectangleCover}
    \lean{TNLean.PEPS.NormalSquareEdgeComplementRectangleCover}
    \lean{TNLean.PEPS.normalSquareVerticalEdgeComplementRegion}
    \lean{TNLean.PEPS.NormalSquareVerticalEdgeComplementRectangleCover}
    \leanok
    \uses{def:peps_squareLatticeCoordinateRegions,
        lem:peps_normalSquareEdgeComplementRegion}
    A rectangular cover of a square-lattice region is a nonempty finite family
    of regions whose union is exactly the given region, and each member is a
    contiguous $2\times3$ or $3\times2$ rectangle. The two cases used in the
    proof would be a cover of the local displayed $T$-region and a cover of the
    finite-lattice complementary block $A_3$. The panel below shows only the
    two permitted shapes of a cover member; it does not assert that these
    particular rectangles cover either region.
    \begin{tenkzequation}
        % Formula verdict: the selected cells form one 2-by-3 rectangle and
        % the secondary cells form one 3-by-2 rectangle.  This is a schematic
        % of the two allowed member shapes, not a union equation or a cover.
        % Ink-to-index verdict: every dot is an ordinary tensor and every
        % lattice edge is a virtual contraction; the outlined fills record
        % rectangle membership only and remain independently customizable.
        % Contraction/boundary verdict: the 4-by-4 panel has 24 internal
        % virtual contractions and 16 cut virtual bonds, with no exposed
        % physical indices, hence boundary type (V,P)=(16,0).
        % Source verdict: arXiv:1804.04964, paper_normal.tex:1405--1445 uses
        % 2-by-3 and 3-by-2 injective pieces.  The following formal lemma proves
        % the local no-cover obstruction for the present origin-based T and A3,
        % so no stronger geometric claim is encoded by this schematic.
        \tenkzkernel
        \begin{tenkz}[rows={ket,ket,ket,ket}, cols=4,
            west=open, east=open, north=open, south=open]
            \tnmark[form=enclosure, slot=selected]{(1,1) .. (2,3)}{$2\times3$}
            \tnmark[form=enclosure, slot=secondary]{(2,3) .. (4,4)}{$3\times2$}
        \end{tenkz}
    \end{tenkzequation}
\end{definition}
% Maintainer note: this definition records sufficient coordinate criteria.
% The concrete local T cover remains to be constructed.

\begin{lemma}[The present local $T$-window and $A_3$ cover obstructions]%
    \label{lem:peps_normalSquareRegionT_no_five_by_six_cover}
    \lean{TNLean.PEPS.not_squareLatticeRectangleCover_of_forced_points}
    \lean{TNLean.PEPS.not_normalSquareRegionT_rectangleCover_at_origin}
    \lean{TNLean.PEPS.not_normalSquareRegionT_rectangleCover_five_by_six}
    \lean{TNLean.PEPS.not_normalSquareVerticalRegionT_rectangleCover_at_origin}
    \lean{TNLean.PEPS.not_normalSquareVerticalRegionT_rectangleCover_six_by_five}
    \lean{TNLean.PEPS.not_normalSquareEdgeComplementRectangleCover_at_origin}
    \lean{TNLean.PEPS.not_normalSquareEdgeComplementRectangleCover_five_by_seven}
    \lean{TNLean.PEPS.not_normalSquareVerticalEdgeComplementRectangleCover_at_origin}
    \lean{TNLean.PEPS.not_normalSquareVerticalEdgeComplementRectangleCover_seven_by_five}
    \leanok
    \uses{def:peps_normalSquareRectangleCover,
        lem:peps_normalSquareRegionT_window_complement}
    The present origin-based local-window model for the displayed $T$-region
    admits no nonempty finite cover by contained contiguous $2\times3$ and
    $3\times2$ rectangles whenever the ambient lattice contains the
    $5\times6$ window. The origin-based rotated local-window model used for
    vertical edges has the analogous obstruction whenever the ambient lattice
    contains the $6\times5$ window. The normalized $5\times6$ and
    $6\times5$ cases are recorded as corollaries. The present origin-based
    horizontal edge-complement model has the analogous obstruction whenever
    the ambient lattice contains the $5\times7$ frame; the normalized
    $5\times7$ case is again recorded as a corollary. The present
    origin-based vertical edge-complement model has the corresponding
    obstruction whenever the ambient lattice contains the $7\times5$ frame;
    the normalized $7\times5$ case is recorded as a corollary.
\end{lemma}
% Maintainer note: this is a local-model obstruction, not a statement of the
% source theorem. It records that the current coordinate model is not yet the
% source cover needed for the $T$-injectivity sentence.
\begin{proof}\leanok
    The lower-left point of the local window belongs to the present $T$
    model. Any contiguous $2\times3$ rectangle containing it also contains
    a point of the removed vertical edge block, while any contiguous
    $3\times2$ rectangle containing it also contains a point of the removed
    horizontal edge block. The rotated statement is the same coordinate
    argument with the two shapes interchanged. The horizontal
    edge-complement statement uses the same lower-left point, and so does the
    normalized vertical edge-complement statement. These cases are all
    instances of the same point-forcing obstruction to an exact rectangular
    cover.
\end{proof}

\begin{lemma}[Injectivity from a rectangular cover]%
    \label{lem:peps_normalSquareRectangleCover_injective}
    \lean{TNLean.PEPS.NormalSquareLatticeRectangleInjectivityHypotheses.injective_of_rectangleCover}
    \leanok
    \uses{def:peps_normalSquareRectangleCover,
        def:peps_normalRectangleInjectivityHypotheses,
        lem:peps_regionInjectivityUnionClosure_finite}
    If a square-lattice region has a nonempty finite cover by contiguous
    $2\times3$ and $3\times2$ rectangles, then rectangular injectivity and
    finite union closure imply that the region is injective.
\end{lemma}
\begin{proof}\leanok
    Apply rectangular injectivity to every rectangle in the cover, then apply
    finite union closure to the nonempty family.
\end{proof}

\begin{lemma}[The normal regions $R$ and $S$ differ in one site]%
    \label{lem:peps_normalSquareRegionR_regionS_difference}
    \lean{TNLean.PEPS.normalSquareRegionR_subset_regionS}
    \lean{TNLean.PEPS.mem_normalSquareRegionS_sdiff_regionR}
    \leanok
    \uses{def:peps_squareLatticeCoordinateRegions,
        lem:peps_squareLatticeCoordinateRegion_identities}
    For the displayed regions $R$ and $S$ with the same lower-left corner,
    $R$ is contained in $S$. The difference $S\setminus R$ is exactly
    the lower-right site of the $3\times3$ square: in coordinates, the site
    one obtains by shifting two steps horizontally and zero steps vertically
    from the lower-left corner.
\end{lemma}
\begin{proof}\leanok
    Expand the two displayed unions. The only point of the second
    $2\times3$ piece of $S$ which is not already in the union defining
    $R$ is its lower-right point.
\end{proof}

\begin{definition}[Normal rectangular injectivity hypotheses]%
    \label{def:peps_normalRectangleInjectivityHypotheses}
    \lean{TNLean.PEPS.NormalRectangleInjectivityHypotheses}
    \lean{TNLean.PEPS.NormalSquareLatticeRectangleInjectivityHypotheses}
    \leanok
    \uses{def:peps_regionInjectivityData,
        def:peps_squareLatticeCoordinateRegions}
    In the square-lattice normal theorem, every $2\times 3$ and
    $3\times 2$ rectangular region is assumed injective. The corresponding
    abstract hypotheses specify which finite regions have these two shapes and
    assert injectivity for all of them. In the coordinate square-lattice
    specialization, these two families are the contiguous coordinate
    $2\times 3$ and $3\times 2$ rectangles.
\end{definition}

\begin{lemma}[Coordinate rectangular injectivity gives abstract rectangular
        injectivity]%
    \label{lem:peps_squareLatticeRectangleInjectivity_toRectangular}
    \lean{TNLean.PEPS.NormalSquareLatticeRectangleInjectivityHypotheses.toRectangular}
    \leanok
    \uses{def:peps_normalRectangleInjectivityHypotheses}
    The coordinate square-lattice rectangular injectivity hypotheses determine
    abstract rectangular injectivity hypotheses, with the two abstract
    rectangular families taken to be the contiguous coordinate $2\times 3$
    and $3\times 2$ rectangles.
\end{lemma}
\begin{proof}\leanok
    This is the direct specialization of the abstract rectangular families to
    the two contiguous coordinate-rectangle predicates.
\end{proof}

\begin{lemma}[Coordinate rectangular injectivity gives injective rectangles]%
    \label{lem:peps_squareLatticeRectangleInjectivity_rectangles}
    \lean{TNLean.PEPS.NormalSquareLatticeRectangleInjectivityHypotheses.rect23_injective}
    \lean{TNLean.PEPS.NormalSquareLatticeRectangleInjectivityHypotheses.rect32_injective}
    \leanok
    \uses{def:peps_normalRectangleInjectivityHypotheses,
        lem:peps_squareLatticeCoordinateRegion_identities}
    Under the coordinate square-lattice rectangular injectivity hypotheses,
    every bounded contiguous $2\times3$ rectangle and every bounded
    contiguous $3\times2$ rectangle is injective.
\end{lemma}
\begin{proof}\leanok
    This is the defining rectangular injectivity hypothesis applied to the
    coordinate witnesses for the two rectangle shapes.
\end{proof}

\begin{lemma}[Rectangular injectivity gives injective $T$-edge blocks]%
    \label{lem:peps_normalSquareRegionT_edgeBlocks_injective}
    \lean{TNLean.PEPS.NormalSquareLatticeRectangleInjectivityHypotheses.tVerticalBlock_injective}
    \lean{TNLean.PEPS.NormalSquareLatticeRectangleInjectivityHypotheses.tHorizontalBlock_injective}
    \lean{TNLean.PEPS.NormalSquareLatticeRectangleInjectivityHypotheses.tHole_injective_of_union}
    \leanok
    \uses{def:peps_regionInjectivityData,
        def:peps_normalRectangleInjectivityHypotheses,
        lem:peps_normalSquareRegionTHole_rectangular_decomposition}
    For the displayed local $T$-region inside its $5\times6$ coordinate window,
    if the smaller edge-block bounds
    $x_0 + 5 \leq n$ and $y_0 + 5 \leq m$ hold in the finite $n\times m$
    square lattice, then the coordinate square-lattice rectangular injectivity
    hypotheses imply that the vertical $2\times3$ block and horizontal
    $3\times2$ block removed from the window are injective. If, in addition,
    region injectivity is closed under unions, then the union of these two
    removed edge blocks is injective.
\end{lemma}
% Maintainer note: this does not prove injectivity of the finite-lattice
% complementary block used as the third region in
% \cite[Theorem~3]{Molnar2018NormalPEPS}.
\begin{proof}\leanok
    Apply the rectangular injectivity hypotheses to the already identified
    contiguous rectangle shapes of the two edge blocks. The final assertion is
    one application of union closure to these two injective blocks.
\end{proof}

\begin{lemma}[Injectivity from a rectangular cover of the local $T$-region]%
    \label{lem:peps_normalSquareRegionT_injective_of_cover}
    \lean{TNLean.PEPS.NormalSquareLatticeRectangleInjectivityHypotheses.regionT_injective_of_cover}
    \lean{TNLean.PEPS.NormalSquareVerticalRegionTRectangleCover}
    \lean{TNLean.PEPS.NormalSquareLatticeRectangleInjectivityHypotheses.verticalT_inj_of_cover}
    \leanok
    \uses{def:peps_normalSquareRectangleCover,
        def:peps_normalRectangleInjectivityHypotheses,
        lem:peps_normalSquareRectangleCover_injective,
        lem:peps_regionInjectivityUnionClosure_finite}
    If the displayed local $T$-region has a nonempty finite cover by
    contiguous $2\times3$ and $3\times2$ rectangles, then rectangular
    injectivity and finite union closure imply that $T$ is injective. The
    rotated vertical local $T$-region has the same conditional injectivity
    criterion.
\end{lemma}
% Current source-comparison status: docs/paper-gaps/peps_normal_ft_section3_route.tex.
\begin{proof}\leanok
    Apply rectangular injectivity to every rectangle in the cover, then apply
    finite union closure to the nonempty family.
\end{proof}

\begin{lemma}[Injectivity from a rectangular cover of the edge complement]%
    \label{lem:peps_normalSquareEdgeComplementRegion_injective_of_cover}
    \lean{TNLean.PEPS.NormalSquareLatticeRectangleInjectivityHypotheses.edgeComplement_injective}
    \lean{TNLean.PEPS.NormalSquareLatticeRectangleInjectivityHypotheses.%
        verticalEdgeComplement_injective}
    \lean{TNLean.PEPS.NormalSquareLatticeRectangleInjectivityHypotheses.%
        edgeComplement_injective_of_T_cover}
    \leanok
    \uses{def:peps_normalSquareRectangleCover,
        def:peps_normalRectangleInjectivityHypotheses,
        lem:peps_normalSquareRectangleCover_injective,
        lem:peps_regionInjectivityUnionClosure_finite,
        lem:peps_normalSquareEdgeComplementRegion_coverOfT}
    If the finite-lattice complementary block $A_3$ has a nonempty finite
    cover by contiguous $2\times3$ and $3\times2$ rectangles, then
    rectangular injectivity and finite union closure imply that $A_3$ is
    injective. The same criterion applies to every vertical edge-complementary
    block: a nonempty finite cover of that block by contiguous $2\times3$ and
    $3\times2$ rectangles implies injectivity of the vertical
    edge-complementary block. In the normalized horizontal-edge $5\times7$
    frame, it is enough to give such a cover of the local $T$-region, since
    adjoining the two top-collar $3\times2$ rectangles gives a cover of $A_3$.
\end{lemma}
% Current source-comparison status: docs/paper-gaps/peps_normal_ft_section3_route.tex.
\begin{proof}\leanok
    Apply rectangular injectivity to every rectangle in the cover, then apply
    finite union closure to the nonempty family. For the normalized
    $5\times7$ edge complement, first extend the local $T$-cover by the
    two top-collar rectangles.
\end{proof}

\begin{lemma}[A cover of $T$ extends to a cover of the edge complement]%
    \label{lem:peps_normalSquareEdgeComplementRegion_coverOfT}
    \lean{TNLean.PEPS.normalSquareEdgeComplementRectangleCoverOfT}
    \lean{TNLean.PEPS.normalSquareVerticalEdgeComplementRectangleCoverOfVerticalT}
    \leanok
    \uses{def:peps_normalSquareRectangleCover,
        lem:peps_normalSquareEdgeComplementRegion_topCollar}
    In the normalized horizontal-edge $5\times7$ frame with lower-left corner
    $(0,0)$, any rectangular cover of the local $T$-region extends to a
    rectangular cover of the edge-complementary block $A_3$ by adjoining the
    two top-collar contiguous $3\times2$ rectangles. In the normalized
    vertical-edge $7\times5$ frame with lower-left corner $(0,0)$, any
    rectangular cover of the rotated local $T$-region extends to a rectangular
    cover of the vertical edge-complementary block by adjoining the two
    right-collar contiguous $2\times3$ rectangles.
\end{lemma}
\begin{proof}\leanok
    Index the new cover by the old cover indices together with two extra
    indices for the collar rectangles. For the horizontal block, the
    decomposition is
    \[
        A_3 =
        T \cup (\{0,1,2\}\times\{5,6\})
        \cup (\{2,3,4\}\times\{5,6\}).
    \]
    For the vertical block, the decomposition is
    \[
        A_3^{\mathrm v} =
        T^{\mathrm v} \cup (\{5,6\}\times\{0,1,2\})
        \cup (\{5,6\}\times\{2,3,4\}).
    \]
\end{proof}

\begin{lemma}[The normalized edge complements have collars]%
    \label{lem:peps_normalSquareEdgeComplementRegion_topCollar}
    \lean{TNLean.PEPS.normalSquareEdgeComplementRegion_eq_T_union_topCollar}
    \lean{TNLean.PEPS.normalSquareEdgeComplementRegion_eq_T_union_rightCollar}
    \lean{TNLean.PEPS.normalSquareVerticalRegionTHole}
    \lean{TNLean.PEPS.normalSquareVerticalRegionT}
    \lean{TNLean.PEPS.mem_normalSquareVerticalRegionTHole}
    \lean{TNLean.PEPS.mem_normalSquareVerticalRegionT}
    \lean{TNLean.PEPS.normalSquareVerticalEdgeComplement_eq_verticalT_union_rightCollar}
    \lean{TNLean.PEPS.NormalSquareLatticeRectangleInjectivityHypotheses.topCollar_injective}
    \lean{TNLean.PEPS.NormalSquareLatticeRectangleInjectivityHypotheses.rightCollar_injective}
    \lean{TNLean.PEPS.NormalSquareLatticeRectangleInjectivityHypotheses.verticalComp_injective}
    \lean{TNLean.PEPS.NormalSquareLatticeRectangleInjectivityHypotheses.verticalComp_inj_of_cover}
    \leanok
    \uses{lem:peps_normalSquareEdgeComplementRegion,
        lem:peps_normalSquareRegionT_window_complement,
        def:peps_regionInjectivityData,
        lem:peps_squareLatticeRectangleInjectivity_rectangles}
    In the normalized horizontal-edge $5\times7$ frame, the finite-lattice
    complementary block $A_3$ is the local-window $T$-region together with
    two top-collar contiguous $3\times2$ rectangles. Consequently, if the
    local $T$-region is injective, then rectangular injectivity and union
    closure imply that $A_3$ is injective. There is also a transposed
    $7\times5$ horizontal-frame collar identity with two contiguous
    $2\times3$ right-collar rectangles. For the rotated vertical
    edge-blocking regions, the corresponding complement is the rotated local
    $T$-region together with the same two right-collar rectangles; hence rotated
    $T$-injectivity, and hence a rectangular cover of rotated $T$, implies
    injectivity of the vertical complementary block.
    \begin{tenkzequation}
        % Formula verdict: in the normalized 7-by-5 row-column frame,
        % A3=((1-7,1-5)-(3-5,1-2)-(2-3,3-5)) is exactly the local T
        % region together with collars C1=(6-7,1-3) and C2=(6-7,3-5).
        % Ink-to-index verdict: every dot is an ordinary tensor and every
        % lattice edge is a virtual contraction; region fills encode only the
        % stated cell-set equality and do not alter the contraction graph.
        % Contraction/boundary verdict: each 7-by-5 panel has 58 internal
        % virtual contractions and 24 cut virtual bonds, with no exposed
        % physical indices, hence boundary type (V,P)=(24,0).
        % Source verdict: arXiv:1804.04964, paper_normal.tex:1430--1445 and
        % 1475--1499 motivate the complementary block.  The collar equality is
        % the exact local formal decomposition in NormalEdgeComplementCover.lean.
        \begin{tenkzlattice}[rows=7, cols=5, boundary legs]
            \tnregion[slot=selected, label=T, label at=north]
                {(1-6,1-5) - (3-5,1-2) - (2-3,3-5)}
            \tnregion[slot=collar, outline, label={C_1},
                label at=center]{(6-7,1-3)}
            \tnregion[slot=collar, outline, inset=1, label={C_2},
                label at=center]{(6-7,3-5)}
        \end{tenkzlattice}
        $\;=\;$
        \begin{tenkzlattice}[rows=7, cols=5, boundary legs]
            \tnregion[slot=complement, label={A_3}, label at=north]
                {(1-7,1-5) - (3-5,1-2) - (2-3,3-5)}
        \end{tenkzlattice}
    \end{tenkzequation}
\end{lemma}
% Current source-comparison status: docs/paper-gaps/peps_normal_ft_section3_route.tex.
\begin{proof}\leanok
    The two set identities follow by expanding the coordinate inequalities.
    The injectivity statements apply union closure first to the local
    $T$-region and the first collar rectangle, and then to the second collar
    rectangle.
\end{proof}

\begin{lemma}[A large rectangle is injective]%
    \label{lem:peps_normalSquareWideRectangle_injective}
    \lean{TNLean.PEPS.NormalSquareLatticeRectangleInjectivityHypotheses.wideRectangle_injective}
    \lean{TNLean.PEPS.NormalSquareLatticeRectangleInjectivityHypotheses.shortRectangle_injective}
    \leanok
    \uses{lem:peps_injectiveRegion_union,
        lem:peps_squareLatticeRectangleInjectivity_rectangles}
    A contiguous rectangle whose width is at least two and whose height is at
    least three is the union of the source $2\times3$ rectangles sliding over
    it, so rectangular injectivity and union closure prove it injective. The
    transposed statement holds for a contiguous rectangle of width at least
    three and height at least two.
\end{lemma}
% Current source-comparison status: docs/paper-gaps/peps_normal_ft_section3_route.tex.
\begin{proof}\leanok
    Each cell of the rectangle lies in one of the contained source windows, so
    the rectangle equals the finite union of those windows; union closure then
    gives injectivity.
\end{proof}

\begin{lemma}[Interior edge-complement injectivity]%
    \label{lem:peps_normalSquareInteriorEdgeComplement_injective}
    \lean{TNLean.PEPS.NormalSquareLatticeRectangleInjectivityHypotheses.interiorEdgeComplement_injective}
    \lean{TNLean.PEPS.NormalSquareLatticeRectangleInjectivityHypotheses.interiorVerticalEdgeComplement_injective}
    \leanok
    \uses{lem:peps_normalSquareEdgeComplementRegion,
        lem:peps_normalSquareWideRectangle_injective,
        lem:peps_squareLatticeRectangleInjectivity_rectangles}
    Let the two edge blocks be removed at an interior position of a sufficiently
    large lattice, with one row and column of room below and to the left and
    further room above and to the right. Then the complementary block $A_3$ is
    the union of four full-length bands around the removed blocks --- below,
    above, left, and right --- together with two filler rectangles inside the
    bounding box of the removed blocks. Each band is a large contiguous
    rectangle and each filler is a source rectangle, all avoiding the removed
    blocks, so rectangular injectivity and union closure prove $A_3$ injective
    with no rectangular-cover hypothesis. The transposed statement holds for the
    rotated vertical complementary block.

    \emph{Scope restriction (interior edge).} The removed blocks must sit in the
    interior. The fixed origin window has its corner forced and so has no
    rectangular cover; the source object is the full-lattice complement $A_3$,
    which has room only at an interior edge. The source comparison is recorded
    in \path{docs/paper-gaps/peps_normal_ft_section3_route.tex}.
\end{lemma}
% Current source-comparison status: docs/paper-gaps/peps_normal_ft_section3_route.tex.
\begin{proof}\leanok
    The six pieces cover the complement of the removed blocks, by expanding the
    coordinate inequalities. Each band is injective by the large-rectangle
    lemma and each filler by rectangular injectivity, so union closure gives
    injectivity of the union.
\end{proof}

\begin{lemma}[Cover-free interior every-edge blocking]%
    \label{lem:peps_normalSquareInteriorEdgeBlocking}
    \lean{TNLean.PEPS.horizontalSquareLatticeEdge_blockingDatum_interior}
    \lean{TNLean.PEPS.verticalSquareLatticeEdge_blockingDatum_interior}
    \lean{TNLean.PEPS.NormalSquareInteriorEdgeDatum}
    \lean{TNLean.PEPS.normalSquareEdgeBlockingHypotheses_of_interiorData}
    \leanok
    \uses{lem:peps_normalSquareInteriorEdgeComplement_injective,
        lem:peps_normalSquareHorizontalTranslatedEdge_blockingData,
        def:peps_normalEdgeBlockingHypotheses}
    For an interior horizontal or vertical edge with the corresponding interior
    margins, the translated edge-blocking datum is built with no rectangular
    cover: its complementary-block injectivity is the interior $A_3$
    injectivity. A family of these cover-free interior data over all edges
    assembles the normal edge-blocking hypotheses, with the injective chain,
    endpoint membership, pairwise disjointness, and coverage at every edge.

    \emph{Scope restriction (interior edge).} The construction supplies data
    only at edges with interior margins. Edges near the lattice boundary of the
    open rectangular model lack such margins; on a translation-invariant lattice
    every edge is interior. The source comparison is recorded in
    \path{docs/paper-gaps/peps_normal_ft_section3_route.tex}.
\end{lemma}
% Current source-comparison status: docs/paper-gaps/peps_normal_ft_section3_route.tex.
\begin{proof}\leanok
    Each interior edge is horizontal or vertical; the corresponding interior
    blocking datum discharges its complementary-block injectivity through the
    interior $A_3$ injectivity. Assembling these per-edge data over all edges
    gives the edge-blocking hypotheses.
\end{proof}

\begin{lemma}[Normalized edge-blocking data]%
    \label{lem:peps_normalSquareHorizontalEdge_blockingData}
    \lean{TNLean.PEPS.normalSquareHorizontalEdge}
    \lean{TNLean.PEPS.normalSquareHorizontalEdge_isHorizontal}
    \lean{TNLean.PEPS.normalSquareHorizontalEdgeRed}
    \lean{TNLean.PEPS.normalSquareHorizontalEdgeBlue}
    \lean{TNLean.PEPS.normalSquareHorizontalEdgeComplement}
    \lean{TNLean.PEPS.normalSquareHorizontalEdge_blockingDatum}
    \lean{TNLean.PEPS.normalSquareHorizontalEdge_blockingData}
    \lean{TNLean.PEPS.normalSquareHorizontalEdge_blockingData_of_TCover}
    \lean{TNLean.PEPS.normalSquareHorizontalEdge_injective_chain}
    \lean{TNLean.PEPS.normalSquareHorizontalEdge_blockingDatum_of_TCover}
    \lean{TNLean.PEPS.normalSquareHorizontalEdge_injective_chain_of_TCover}
    \lean{TNLean.PEPS.normalSquareVerticalEdge}
    \lean{TNLean.PEPS.normalSquareVerticalEdge_isVertical}
    \lean{TNLean.PEPS.normalSquareVerticalEdgeRed}
    \lean{TNLean.PEPS.normalSquareVerticalEdgeBlue}
    \lean{TNLean.PEPS.normalSquareVerticalEdgeComplement}
    \lean{TNLean.PEPS.normalSquareVerticalEdge_blockingDatum}
    \lean{TNLean.PEPS.normalSquareVerticalEdge_blockingData}
    \lean{TNLean.PEPS.normalSquareVerticalEdge_blockingData_of_verticalT}
    \lean{TNLean.PEPS.normalSquareVerticalEdge_blockingData_of_verticalTCover}
    \lean{TNLean.PEPS.normalSquareVerticalEdge_blockingDatum_of_verticalT}
    \lean{TNLean.PEPS.normalSquareVerticalEdge_blockingDatum_of_verticalTCover}
    \lean{TNLean.PEPS.normalSquareVerticalEdge_injective_chain}
    \lean{TNLean.PEPS.normalSquareVerticalEdge_injective_chain_of_verticalT}
    \lean{TNLean.PEPS.normalSquareVerticalEdge_injective_chain_of_verticalTCover}
    \leanok
    \uses{def:peps_squareLatticeCoordinateRegions,
        lem:peps_normalSquareEdgeComplementRegion,
        lem:peps_normalSquareRegionT_edgeBlocks_injective,
        lem:peps_normalSquareEdgeComplementRegion_topCollar}
    In the normalized $5\times7$ horizontal-edge frame, the red block is the
    vertical $2\times3$ block, the blue block is the horizontal $3\times2$
    block, and the complementary block is the finite-lattice edge complement
    $A_3$. The distinguished edge is a horizontal edge of the square-lattice
    graph. The left endpoint of the distinguished edge lies in the red block,
    the right endpoint lies in the blue block, the three regions are pairwise
    disjoint, and their union is the whole finite square lattice.
    If the local $T$-region is injective, rectangular injectivity and union
    closure give injectivity of all three blocks; in particular a rectangular
    cover of $T$ suffices. The normalized $7\times5$ vertical-edge frame
    has the analogous red and blue rectangular blocks, a distinguished vertical
    edge, and the same set-theoretic blocking data; at this stage its
    complementary-block injectivity is either kept as an explicit hypothesis or
    obtained from
    rotated $T$-injectivity, and therefore from a rectangular cover of
    rotated $T$, by the vertical collar lemma. In both cases the named
    blocking datum supplies the corresponding three injective regions for the
    edge-blocked chain.
\end{lemma}
% Maintainer note: horizontal and vertical complement injectivity are reduced
% to local and rotated T-covers, respectively. The translated every-edge
% construction remains a separate step.
\begin{proof}\leanok
    The endpoint memberships are coordinate checks. Rectangular injectivity
    gives the rectangular red and blue blocks. For the horizontal edge, the
    collar lemma gives the complement, and the rectangular cover criterion
    removes the direct $T$-injectivity input when such a cover is supplied.
    For the vertical edge, complement injectivity is an explicit input or is
    obtained from the rotated $T$-cover. The disjointness and cover
    statements are finite-set complement identities.
\end{proof}

\begin{lemma}[Translated horizontal edge-blocking data]%
    \label{lem:peps_normalSquareHorizontalTranslatedEdge_blockingData}
    \lean{TNLean.PEPS.normalSquareHorizontalTranslatedEdge}
    \lean{TNLean.PEPS.normalSquareHorizontalTranslatedEdgeRed}
    \lean{TNLean.PEPS.normalSquareHorizontalTranslatedEdgeBlue}
    \lean{TNLean.PEPS.normalSquareHorizontalTranslatedEdgeComplement}
    \lean{TNLean.PEPS.normalSquareHorizontalTranslatedEdge_isHorizontal}
    \lean{TNLean.PEPS.normalSquareHorizontalTranslatedEdge_eq_rightEdge}
    \lean{TNLean.PEPS.normalSquareHorizontalTranslatedEdge_blockingDatum}
    \lean{TNLean.PEPS.normalSquareHorizontalTranslatedEdge_injective_chain}
    \lean{TNLean.PEPS.normalSquareHorizontalTranslatedEdge_blockingDatum_of_complementCover}
    \lean{TNLean.PEPS.normalSquareHorizontalTranslatedEdge_injective_chain_of_complementCover}
    \lean{TNLean.PEPS.horizontalTranslatedEdge_zero}
    \lean{TNLean.PEPS.horizontalTranslatedRed_zero}
    \lean{TNLean.PEPS.horizontalTranslatedBlue_zero}
    \lean{TNLean.PEPS.horizontalTranslatedComp_zero}
    \leanok
    \uses{def:peps_squareLatticeCoordinateRegions,
        lem:peps_normalSquareEdgeComplementRegion,
        lem:peps_normalSquareEdgeComplementRegion_injective_of_cover}
    For any coordinate origin whose horizontal blocking window lies in the
    finite square lattice, there is a distinguished horizontal edge with the
    same relative position as in the normalized picture. The translated red
    region is the $2\times3$ block, the translated blue region is the
    $3\times2$ block, and the complementary region is the finite-lattice
    complement of their union. If this complementary region is injective,
    these three regions form the one-edge blocking datum; if it has a
    rectangular cover, rectangular injectivity and union closure supply that
    complementary injectivity. In both forms the datum supplies the three
    injective regions used in the edge-blocked chain. The distinguished edge
    is the corresponding coordinate right edge. At coordinate origin $0,0$,
    the translated edge and the three translated regions specialize to the
    normalized horizontal picture.
\end{lemma}
\begin{proof}\leanok
    The endpoint memberships and horizontal orientation are coordinate
    checks. Rectangular injectivity gives the red and blue blocks. The
    complementary block is supplied either as an input or from a rectangular
    cover, and the disjointness and covering identities follow from the
    finite-set complement formulas.
\end{proof}

\begin{lemma}[Translated vertical edge-blocking data]%
    \label{lem:peps_normalSquareVerticalTranslatedEdge_blockingData}
    \lean{TNLean.PEPS.normalSquareVerticalTranslatedEdge}
    \lean{TNLean.PEPS.normalSquareVerticalTranslatedEdgeRed}
    \lean{TNLean.PEPS.normalSquareVerticalTranslatedEdgeBlue}
    \lean{TNLean.PEPS.normalSquareVerticalTranslatedEdgeComplement}
    \lean{TNLean.PEPS.normalSquareVerticalTranslatedEdge_isVertical}
    \lean{TNLean.PEPS.normalSquareVerticalTranslatedEdge_eq_upEdge}
    \lean{TNLean.PEPS.normalSquareVerticalTranslatedEdge_blockingDatum}
    \lean{TNLean.PEPS.normalSquareVerticalTranslatedEdge_injective_chain}
    \lean{TNLean.PEPS.normalSquareVerticalTranslatedEdge_blockingDatum_of_complementCover}
    \lean{TNLean.PEPS.normalSquareVerticalTranslatedEdge_injective_chain_of_complementCover}
    \lean{TNLean.PEPS.verticalTranslatedEdge_zero}
    \lean{TNLean.PEPS.verticalTranslatedRed_zero}
    \lean{TNLean.PEPS.verticalTranslatedBlue_zero}
    \lean{TNLean.PEPS.verticalTranslatedComp_zero}
    \leanok
    \uses{def:peps_squareLatticeCoordinateRegions,
        lem:peps_normalSquareEdgeComplementRegion_injective_of_cover}
    For any coordinate origin whose vertical blocking window lies in the
    finite square lattice, there is a distinguished vertical edge with the same
    relative position as in the normalized vertical picture. The translated
    red region is the $3\times2$ block, the translated blue region is the
    $2\times3$ block, and the complementary region is the finite-lattice
    complement of their union. If this complementary region is injective,
    these three regions form the one-edge blocking datum; if it has a
    rectangular cover, rectangular injectivity and union closure supply that
    complementary injectivity. In both forms the datum supplies the three
    injective regions used in the edge-blocked chain. The distinguished edge
    is the corresponding coordinate upward edge. At coordinate origin $0,0$,
    the translated edge and the three translated regions specialize to the
    normalized vertical picture.
\end{lemma}
\begin{proof}\leanok
    The endpoint memberships and vertical orientation are coordinate checks.
    Rectangular injectivity gives the red and blue blocks. The complementary
    block is supplied either as an input or from a rectangular cover, and the
    disjointness and covering identities follow from the finite-set complement
    formulas.
\end{proof}

\begin{definition}[Translated edge windows]%
    \label{def:peps_normalSquareTranslatedEdgeWindow}
    \lean{TNLean.PEPS.NormalSquareTranslatedEdgeWindow}
    \lean{TNLean.PEPS.NormalSquareTranslatedEdgeWindow.blockingDatum}
    \lean{TNLean.PEPS.NormalSquareTranslatedEdgeWindow.injective_chain}
    \lean{TNLean.PEPS.NormalSquareTranslatedEdgeWindow.endpoint_disjoint_cover}
    \lean{TNLean.PEPS.normalSquareTranslatedEdgeWindow_rightEdge_margins}
    \lean{TNLean.PEPS.normalSquareTranslatedEdgeWindow_upEdge_margins}
    \lean{TNLean.PEPS.not_leftBoundaryRightEdge_window}
    \lean{TNLean.PEPS.not_leftBoundaryRightEdge_window_seven}
    \lean{TNLean.PEPS.not_rightMarginRightEdge_window}
    \lean{TNLean.PEPS.not_rightMarginRightEdge_window_seven}
    \lean{TNLean.PEPS.not_bottomBoundaryUpEdge_window}
    \lean{TNLean.PEPS.not_bottomBoundaryUpEdge_window_seven}
    \lean{TNLean.PEPS.not_topMarginUpEdge_window}
    \lean{TNLean.PEPS.not_topMarginUpEdge_window_seven}
    \lean{TNLean.PEPS.not_fourSidedBoundaryWindow_seven}
    \lean{TNLean.PEPS.not_forall_normalSquareTranslatedEdgeWindow_seven}
    \lean{TNLean.PEPS.normalSquareHorizontalTranslatedEdgeWindow}
    \lean{TNLean.PEPS.normalSquareVerticalTranslatedEdgeWindow}
    \lean{TNLean.PEPS.normalSquareHorizontalTranslatedEdge_sub_eq_rightEdge}
    \lean{TNLean.PEPS.normalSquareVerticalTranslatedEdge_sub_eq_upEdge}
    \lean{TNLean.PEPS.squareLatticeRightEdgeWindow}
    \lean{TNLean.PEPS.squareLatticeUpEdgeWindow}
    \lean{TNLean.PEPS.squareLatticeRightEdgeWindow_of_margins}
    \lean{TNLean.PEPS.squareLatticeUpEdgeWindow_of_margins}
    \lean{TNLean.PEPS.horizontalSquareLatticeEdgeWindow}
    \lean{TNLean.PEPS.horizontalSquareLatticeEdgeWindow_of_margins}
    \lean{TNLean.PEPS.verticalSquareLatticeEdgeWindow}
    \lean{TNLean.PEPS.verticalSquareLatticeEdgeWindow_of_margins}
    \lean{TNLean.PEPS.IsNormalSquareHorizontalEdgeMargins}
    \lean{TNLean.PEPS.IsNormalSquareVerticalEdgeMargins}
    \lean{TNLean.PEPS.NormalSquareEdgeMarginCover}
    \lean{TNLean.PEPS.NormalSquareEdgeMarginCover.window}
    \lean{TNLean.PEPS.NormalSquareEdgeMarginCover.blockingDatum}
    \lean{TNLean.PEPS.NormalSquareEdgeMarginCover.injective_chain}
    \lean{TNLean.PEPS.NormalSquareEdgeMarginCover.endpoint_disjoint_cover}
    \lean{TNLean.PEPS.normalSquareEdgeMarginCover_rightEdge_bounds}
    \lean{TNLean.PEPS.normalSquareEdgeMarginCover_upEdge_bounds}
    \lean{TNLean.PEPS.not_leftBoundaryRightEdge_marginCover}
    \lean{TNLean.PEPS.not_leftBoundaryRightEdge_marginCover_seven}
    \lean{TNLean.PEPS.not_rightMarginRightEdge_marginCover}
    \lean{TNLean.PEPS.not_rightMarginRightEdge_marginCover_seven}
    \lean{TNLean.PEPS.not_bottomBoundaryUpEdge_marginCover}
    \lean{TNLean.PEPS.not_bottomBoundaryUpEdge_marginCover_seven}
    \lean{TNLean.PEPS.not_topMarginUpEdge_marginCover}
    \lean{TNLean.PEPS.not_topMarginUpEdge_marginCover_seven}
    \lean{TNLean.PEPS.not_fourSidedBoundaryMarginCover_seven}
    \leanok
    \uses{lem:peps_normalSquareHorizontalTranslatedEdge_blockingData,
        lem:peps_normalSquareVerticalTranslatedEdge_blockingData}
    A translated edge window records that a chosen edge is one of the
    translated horizontal or vertical edge pictures, together with a
    rectangular cover of the complementary block. Such a window gives the
    one-edge blocking datum for that edge, hence the three injective regions
    used in the edge-blocked chain. The coordinate right and upward edges
    obtain such windows directly from the corresponding translated picture.
    Equivalently, an actual coordinate right or upward edge obtains such a
    window when it has the necessary room around it for the translated
    $5\times7$ or $7\times5$ frame and the complementary cover is supplied.
    The same criterion applies to an arbitrary horizontal or vertical
    square-lattice edge after identifying it with its coordinate right or
    upward representative. The oriented margin-and-cover datum is the same
    criterion written as per-edge data in the rectangular coordinate model;
    it directly supplies the one-edge blocking datum, its three injective
    regions, and the identities
    $u_e\in R_e$, $v_e\in B_e$,
    $R_e\cap B_e=R_e\cap C_e=B_e\cap C_e=\varnothing$, and
    $R_e\cup B_e\cup C_e=V$. The boundary obstruction lemmas show that these
    translated-window and margin-cover criteria are interior criteria, not the
    final every-edge construction in the open square-lattice model.
    % Detailed source-comparison status:
    % docs/paper-gaps/peps_normal_ft_section3_route.tex.
\end{definition}

\begin{theorem}[Construction from translated edge windows]%
    \label{thm:peps_normalSquareTranslatedEdgeBlockingHypotheses_of_windows}
    \lean{TNLean.PEPS.normalSquareTranslatedEdgeBlockingHypotheses_of_windows}
    \lean{TNLean.PEPS.normalSquareTranslatedEdgeBlockingHypotheses_blockingData_of_windows}
    \lean{TNLean.PEPS.normalSquareTranslatedEdgeBlockingHypotheses_injective_chain_of_windows}
    \lean{TNLean.PEPS.%
        normalSquareTranslatedEdgeBlockingHypotheses_endpoint_disjoint_cover_of_windows}
    \lean{TNLean.PEPS.normalSquareEdgeBlockingHypotheses_of_marginCovers}
    \lean{TNLean.PEPS.normalSquareEdgeBlockingHypotheses_blockingData_of_marginCovers}
    \lean{TNLean.PEPS.normalSquareEdgeBlockingHypotheses_injective_chain_of_marginCovers}
    \lean{TNLean.PEPS.normalSquareEdgeBlockingHypotheses_endpoint_disjoint_cover_of_marginCovers}
    \lean{TNLean.PEPS.not_forall_normalSquareEdgeMarginCover_seven}
    \leanok
    \uses{def:peps_normalSquareTranslatedEdgeWindow,
        def:peps_normalEdgeBlockingHypotheses}
    Given rectangular injectivity hypotheses, finite union closure, and a
    translated edge window at every edge of the finite square lattice, the
    corresponding one-edge blockings determine the normal edge-blocking
    hypotheses. This is the conditional construction;
    constructing such windows for all edges of the $7\times7$ lattice remains
    separate; the current open rectangular translated-window criterion does not
    by itself supply such a family. The one-edge datum recovered from the
    resulting hypotheses is the datum supplied by the corresponding translated
    window, and the resulting hypotheses give the three injective regions at
    every edge, together with
    $u_e\in R_e$, $v_e\in B_e$,
    $R_e\cap B_e=R_e\cap C_e=B_e\cap C_e=\varnothing$, and
    $R_e\cup B_e\cup C_e=V$. Equivalently, it is enough to supply the oriented
    margin-and-cover data at every edge; the resulting hypotheses recover the
    corresponding datum, give the same three injectivity conclusions, and give
    the same endpoint, disjointness, and covering conclusions. The obstruction
    theorem records that this interior margin-cover criterion does not by
    itself supply data at every edge of the open $7\times7$ lattice.
    % Detailed source-comparison status:
    % docs/paper-gaps/peps_normal_ft_section3_route.tex.
\end{theorem}
\begin{proof}\leanok
    Apply the general construction for one-edge blocking data to the datum
    supplied by the chosen translated window at each edge.
\end{proof}

\begin{lemma}[Conditional injectivity of the normal $R$ and $S$ regions]%
    \label{lem:peps_normalSquareRegionRS_injective_of_unionClosure}
    \lean{TNLean.PEPS.NormalSquareLatticeRectangleInjectivityHypotheses.regionR_injective_of_union}
    \lean{TNLean.PEPS.NormalSquareLatticeRectangleInjectivityHypotheses.regionS_injective_of_union}
    \leanok
    \uses{def:peps_regionInjectivityData,
        def:peps_normalRectangleInjectivityHypotheses,
        lem:peps_normalSquareRegionRS_rectangular_decomposition}
    If region injectivity is closed under unions, then the coordinate
    square-lattice rectangular injectivity hypotheses imply that the displayed
    regions $R$ and $S$ are injective.
\end{lemma}
\begin{proof}\leanok
    Region $R$ is the union of one injective $2\times3$ rectangle and one
    injective $3\times2$ rectangle. Region $S$ is the union of two
    injective $2\times3$ rectangles. Apply union closure in each case.
\end{proof}

\begin{lemma}[Injectivity of the normal $R$, $S$, and $T$ regions]%
    \label{lem:peps_normalSquareRegionRST_injective}
    \lean{TNLean.PEPS.regionBlockedTensorInjective_normalSquareRegionR}
    \lean{TNLean.PEPS.regionBlockedTensorInjective_normalSquareRegionS}
    \lean{TNLean.PEPS.regionBlockedTensorInjective_normalSquareRegionT}
    \leanok
    \uses{thm:peps_regionInjectivityUnionClosure_of_overlap,
        lem:peps_normalSquareRegionRS_injective_of_unionClosure,
        lem:peps_normalSquareRegionT_injective_of_cover}
    Fix a square-lattice PEPS tensor with every bond dimension positive whose
    blocked tensor is injective on every contiguous $2\times3$ and $3\times2$
    rectangle. Then the displayed regions $R$ and $S$ are injective. If
    moreover the displayed region $T$ has a cover by such rectangles, then $T$
    is injective too. The union closure used here is supplied by the
    overlapping union lemma, so no single-vertex injectivity is assumed.
\end{lemma}
\begin{proof}\leanok
    Apply the conditional injectivity of $R$ and $S$
    (Lemma~\ref{lem:peps_normalSquareRegionRS_injective_of_unionClosure}) and the
    conditional cover criterion for $T$
    (Lemma~\ref{lem:peps_normalSquareRegionT_injective_of_cover}), with the union
    closure discharged by
    Theorem~\ref{thm:peps_regionInjectivityUnionClosure_of_overlap}.
\end{proof}

% The comparison square and comparison region at a parametric offset, used to
% place the one-site comparison pair $R\subseteq S=R\cup\{v\}$ afresh at every
% vertex of the open square lattice.

\begin{definition}[Comparison square, comparison region, and interior window]%
    \label{def:peps_normalSquareComparisonRegions}
    \lean{TNLean.PEPS.normalSquareComparisonSquare}
    \lean{TNLean.PEPS.normalSquareComparisonRegion}
    \lean{TNLean.PEPS.IsNormalSquareComparisonWindow}
    \leanok
    \uses{def:peps_squareLatticeCoordinateRegions}
    At a coordinate offset $(a,b)$ the comparison square is the $3\times3$
    contiguous coordinate rectangle with lower-left corner $(a,b)$. The
    comparison region is this square with its lower-left corner $(a,b)$
    removed, presented as the union of the $2\times3$ contiguous rectangle at
    $(a+1,b)$ and the $3\times2$ contiguous rectangle at $(a,b+1)$. This is an
    auxiliary one-site comparison pair carrying its omitted site at the
    lower-left corner; it is the reflected counterpart of the displayed regions
    $R$ and $S$, whose omitted site $S\setminus R$ sits at the lower-right
    corner, placed at an arbitrary offset rather than at a fixed origin. The
    interior window at $(a,b)$ for a finite $n\times m$ lattice is the placement
    margin
    $4\le a$, $a+9\le n$, $4\le b$, $b+9\le m$, under which every boundary edge
    of the comparison square and of the comparison region carries the interior
    margins of the open edge-blocking geometry.
\end{definition}

\begin{lemma}[Membership in the comparison region]%
    \label{lem:peps_normalSquareComparisonRegion_membership}
    \lean{TNLean.PEPS.mem_normalSquareComparisonRegion}
    \lean{TNLean.PEPS.notMem_normalSquareComparisonRegion}
    \lean{TNLean.PEPS.insert_normalSquareComparisonRegion}
    \leanok
    \uses{def:peps_normalSquareComparisonRegions,
        lem:peps_squareLatticeCoordinateRegion_identities}
    For any offset $(a,b)$ a vertex with coordinates $(x,y)$ lies in the
    comparison region at $(a,b)$ if and only if
    \begin{align}
        &(a+1\le x<a+3\wedge b\le y<b+3)
        \lor
        (a\le x<a+3\wedge b+1\le y<b+3).
        \label{eq:peps_comp_region_mem}
    \end{align}
    Let $v$ be a lattice vertex and place the comparison region at $v$'s own
    coordinates. Then $v$ lies outside that region, and inserting $v$ back in
    recovers the full $3\times3$ comparison square placed at $v$'s coordinates.
\end{lemma}
\begin{proof}\leanok
    Membership unfolds the union of the two contiguous rectangles into
    (\ref{eq:peps_comp_region_mem}). Place the region at $v$'s own
    coordinates, so $(a,b)=(x,y)=(v_1,v_2)$. The corner offset $(v_1,v_2)$ fails
    the first disjunct because $v_1+1\le v_1$ is false, and fails the second
    because $v_2+1\le v_2$ is false:
    \begin{align}
        &\neg\bigl((v_1+1\le v_1<v_1+3)
            \wedge(v_2\le v_2<v_2+3)\bigr)
        \notag\\
        &\qquad\wedge
        \neg\bigl((v_1\le v_1<v_1+3)
            \wedge(v_2+1\le v_2<v_2+3)\bigr).
        \notag
    \end{align}
    Thus $v$ lies outside the region. Inserting $v$ therefore yields the region
    together with $\{(v_1,v_2)\}$, and a vertex lies in this set exactly when it
    equals $v$ or already lies in the region. That is the membership condition
    of the $3\times3$ square, the corner being the one cell of the square that
    both pieces omit.
\end{proof}

\begin{lemma}[Band decomposition of the comparison square complement]%
    \label{lem:peps_normalSquareComparisonSquare_complementBands}
    \lean{TNLean.PEPS.compl_normalSquareComparisonSquare_eq_bands}
    \leanok
    \uses{def:peps_normalSquareComparisonRegions}
    On a finite $n\times m$ lattice, if $a+3\le n$ and $b+3\le m$ then the
    complement of the comparison square at $(a,b)$ is the union of the four
    coordinate bands
    \begin{align}
        &\{(x,y):0\le x<a,\ 0\le y<m\}
        \cup
        \{(x,y):a+3\le x<n,\ 0\le y<m\}
        \notag\\
        &\qquad\cup
        \{(x,y):a\le x<a+3,\ 0\le y<b\}
        \cup
        \{(x,y):a\le x<a+3,\ b+3\le y<m\}.
        \label{eq:peps_comp_square_bands}
    \end{align}
\end{lemma}
\begin{proof}\leanok
    A vertex $(x,y)$ lies in the comparison square exactly when
    $a\le x<a+3$ and $b\le y<b+3$. With $0\le x<n$ and $0\le y<m$ on the
    lattice, the negation of that conjunction is membership in
    (\ref{eq:peps_comp_square_bands}):
    \begin{align}
        &\neg(a\le x<a+3\wedge b\le y<b+3)
        \notag\\
        &\quad\Longleftrightarrow
        (0\le x<a\wedge 0\le y<m)
        \lor
        (a+3\le x<n\wedge 0\le y<m)
        \notag\\
        &\qquad\lor
        (a\le x<a+3\wedge 0\le y<b)
        \lor
        (a\le x<a+3\wedge b+3\le y<m).
        \notag
    \end{align}
\end{proof}

\begin{lemma}[Band decomposition of the comparison region complement]%
    \label{lem:peps_normalSquareComparisonRegion_complementBands}
    \lean{TNLean.PEPS.compl_normalSquareComparisonRegion_eq_bands}
    \leanok
    \uses{def:peps_normalSquareComparisonRegions,
        lem:peps_normalSquareComparisonRegion_membership}
    On a finite $n\times m$ lattice, if $2\le a$, $1\le b$, $a+3\le n$, and
    $b+3\le m$ then the complement of the comparison region at $(a,b)$ is the
    union of the same four bands of
    Lemma~\ref{lem:peps_normalSquareComparisonSquare_complementBands} together
    with the $3\times2$ corner rectangle
    $\{(x,y):a-2\le x<a+1,\ b-1\le y<b+1\}$,
    which contains the omitted corner $(a,b)$.
\end{lemma}
\begin{proof}\leanok
    By (\ref{eq:peps_comp_region_mem}), a vertex $(x,y)$ lies in the complement
    exactly when both disjuncts fail:
    \begin{align}
        &\neg(a+1\le x<a+3\wedge b\le y<b+3)
        \notag\\
        &\qquad\wedge
        \neg(a\le x<a+3\wedge b+1\le y<b+3).
        \notag
    \end{align}
    Together with $0\le x<n$ and $0\le y<m$, this places the vertex either in
    one of the four bands of
    (\ref{eq:peps_comp_square_bands}) or in the $3\times2$ corner rectangle.
    Writing $R$ for the comparison region and $S$ for the comparison square,
    \begin{align}
        V\setminus R
        &=
        \{(x,y):a-2\le x<a+1,\ b-1\le y<b+1\}
        \cup(V\setminus S),
        \notag
    \end{align}
    since the comparison region is the comparison square with the corner $(a,b)$
    removed, and that corner lies in the displayed $3\times2$ rectangle. The
    hypotheses $2\le a$ and $1\le b$ keep the corner rectangle's lower-left
    coordinates $(a-2,b-1)$ within the lattice.
\end{proof}

\begin{lemma}[Injectivity of the comparison region and the comparison square]%
    \label{lem:peps_normalSquareComparison_injective}
    \lean{TNLean.PEPS.NormalSquareLatticeRectangleInjectivityHypotheses.comparisonRegion_injective}
    \lean{TNLean.PEPS.NormalSquareLatticeRectangleInjectivityHypotheses.comparisonSquare_injective}
    \leanok
    \uses{def:peps_normalRectangleInjectivityHypotheses,
        def:peps_regionInjectivityData,
        lem:peps_squareLatticeRectangleInjectivity_rectangles,
        lem:peps_normalSquareWideRectangle_injective}
    Assume rectangular injectivity and union closure. On a finite $n\times m$
    lattice, if $a+3\le n$ and $b+3\le m$ then the comparison region at $(a,b)$
    is injective, being the union of a $2\times3$ and a $3\times2$ rectangle,
    and the comparison square at $(a,b)$ is injective, being a $3\times3$
    rectangle.
\end{lemma}
\begin{proof}\leanok
    Apply union closure to the rectangular injectivity of the contained
    $2\times3$ and $3\times2$ rectangles for the region, the square being a
    single large rectangle.
\end{proof}

\begin{lemma}[Injectivity of the comparison complements]%
    \label{lem:peps_normalSquareComparison_complement_injective}
    \lean{TNLean.PEPS.NormalSquareLatticeRectangleInjectivityHypotheses.compl_comparisonSquare_injective}
    \lean{TNLean.PEPS.NormalSquareLatticeRectangleInjectivityHypotheses.compl_comparisonRegion_injective}
    \leanok
    \uses{def:peps_normalRectangleInjectivityHypotheses,
        def:peps_regionInjectivityData,
        lem:peps_normalSquareWideRectangle_injective,
        lem:peps_normalSquareComparisonSquare_complementBands,
        lem:peps_normalSquareComparisonRegion_complementBands}
    Assume rectangular injectivity and union closure. On a finite $n\times m$
    lattice, if $2\le a$, $2\le b$, $a+5\le n$, and $b+5\le m$ then the
    complement of the comparison square at $(a,b)$ is injective, being the union
    of its four surrounding bands, and the complement of the comparison region
    at $(a,b)$ is injective, being that union together with one further corner
    rectangle.
\end{lemma}
\begin{proof}\leanok
    Rewrite each complement by its band decomposition and apply union closure to
    the bands, each a large contiguous rectangle, and to the corner rectangle.
\end{proof}

\begin{definition}[Normal edge-blocking hypotheses]%
    \label{def:peps_normalEdgeBlockingHypotheses}
    \lean{TNLean.PEPS.NormalEdgeBlockingData}
    \lean{TNLean.PEPS.NormalEdgeBlockingHypotheses}
    \leanok
    \uses{def:peps_edge, def:peps_regionInjectivityData}
    Around one edge, the normal proof uses a blocking into three injective
    regions: a red region containing one endpoint, a blue region containing the
    other endpoint, and an injective complementary region. The three regions
    are pairwise disjoint and cover the whole vertex set. The edge-blocking
    hypotheses are precisely this one-edge datum at every edge; the red,
    blue, and complementary regions are read off from it. The local blocking
    picture used in the proof of
    \cite[Theorem~3]{Molnar2018NormalPEPS}:
    \begin{tenkzequation}
        % Formula verdict: in the source 7-by-7 frame, A1=(4-6,2-3),
        % A2=(3-4,4-6), and A3 is their cell-set complement; the marked
        % horizontal edge is (4,3)-(4,4), joining one endpoint in each block.
        % Ink-to-index verdict: every dot is an ordinary tensor, every lattice
        % edge is a virtual contraction, the three fills encode the disjoint
        % cover, and the distinguished edge is the edge being blocked.
        % Contraction/boundary verdict: the panel has 84 internal virtual
        % contractions and 28 cut virtual bonds, with no exposed physical
        % indices, hence boundary type (V,P)=(28,0).
        % Source verdict: exact discrete translation of arXiv:1804.04964,
        % paper_normal.tex:1475--1499; A1 is red, A2 blue, and A3 the rest.
        \begin{tenkzlattice}[rows=7, cols=7, boundary legs]
            \tnregion[slot=selected, label={$A_1$},
                label at=center]{(4-6,2-3)}
            \tnregion[slot=secondary, label={$A_2$},
                label at=center]{(3-4,4-6)}
            \tnregion[slot=complement, outline, label={$A_3$},
                label at=north]
                {(1-7,1-7) - (4-6,2-3) - (3-4,4-6)}
            \tnedge[distinguished]{(4,3)-(4,4)}
        \end{tenkzlattice}
    \end{tenkzequation}
\end{definition}

\begin{theorem}[Injectivity supplied by normal edge blocking]%
    \label{thm:peps_normalEdgeBlockingHypotheses_injectiveChain}
    \lean{TNLean.PEPS.NormalEdgeBlockingData.injective_chain}
    \lean{TNLean.PEPS.NormalEdgeBlockingData.endpoint_disjoint_cover}
    \lean{TNLean.PEPS.NormalEdgeBlockingData.endpoint_injective_disjoint_cover}
    \lean{TNLean.PEPS.NormalEdgeBlockingHypotheses.ofBlockingData}
    \lean{TNLean.PEPS.NormalEdgeBlockingHypotheses.ofBlockingData_blockingData}
    \lean{TNLean.PEPS.NormalEdgeBlockingHypotheses.injective_chain_at_edge}
    \lean{TNLean.PEPS.NormalEdgeBlockingHypotheses.endpoint_disjoint_cover_at_edge}
    \lean{TNLean.PEPS.%
        NormalEdgeBlockingHypotheses.endpoint_injective_disjoint_cover_at_edge}
    \leanok
    \uses{def:peps_normalEdgeBlockingHypotheses}
    A family of one-edge blocking data is precisely the normal edge-blocking
    hypotheses. If $H$ is assembled from the family $d_e$, then the one-edge
    datum that $H$ associates with $e$ is $d_e$. Write the endpoints of $e$
    as $u_e<v_e$ in the canonical vertex ordering. If the three regions are
    $R_e$, $B_e$, and $C_e$, then
    \begin{align}
        u_e&\in R_e,
        \notag\\
        v_e&\in B_e,
        \notag\\
        R_e\cap B_e&=R_e\cap C_e=B_e\cap C_e=\varnothing,
        \notag\\
        R_e\cup B_e\cup C_e&=V,
        \label{eq:peps_edge_blocking_cover}
    \end{align}
    and the underlying PEPS tensor is injective on $R_e$, $B_e$, and $C_e$.
\end{theorem}
\begin{proof}\leanok
    By construction, the one-edge datum at each edge $e$ is $d_e$.
    The endpoint, disjointness, and coverage statements in
    (\ref{eq:peps_edge_blocking_cover}), together with the injectivity
    statements, are exactly the corresponding assumptions of that one-edge
    datum.
\end{proof}

\begin{definition}[Normal one-site separation hypotheses]%
    \label{def:peps_normalOneSiteSeparationHypotheses}
    \lean{TNLean.PEPS.NormalOneSiteSeparationHypotheses}
    \leanok
    \uses{def:peps_regionInjectivityData, def:peps_regionComplement}
    For each site, the normal theorem assumes two injective regions with
    injective complements that differ exactly at that site. This is the
    one-site comparison used after the edge gauges have been obtained.
    Diagrammatically, in the square-lattice proof this is the comparison
    between the source regions $R$ and $S=R\cup\{v\}$:
    \begin{tenkzequation}
        % Source verdict: paper_normal.tex:1519--1544 compares the injective
        % regions R and S=R\cup\{v\}; their complements are injective as well.
        % Type verdict: this is a lattice/set schematic, not a tensor
        % contraction. The marked site v remains present with all its bonds.
        % Cell-set audit: R=(2-4,2-4)-(4,4) has eight selected cells, while
        % S=(2-4,2-4) has nine; v=(4,4) is marked in both pictures.
        % Contraction/boundary audit: there are no tensor contractions and no
        % tensor boundary signature; the boundary legs are cut lattice edges.
        \tenkzkernel
        \begin{tenkz}[rows={ket,ket,ket,ket,ket}, cols=5,
            west=open, east=open, north=open, south=open]
            \tn[at=(4,4), species=marked, name=vSite, label pos=sw]{v}
            \tnmark[form=enclosure, slot=selected, tint, label pos=nw]
                {(2,2) .. (4,4) - (4,4)}{$R$}
        \end{tenkz}
        $\qquad$
        \begin{tenkz}[rows={ket,ket,ket,ket,ket}, cols=5,
            west=open, east=open, north=open, south=open]
            \tn[at=(4,4), species=marked, name=vSite, label pos=sw]{v}
            \tnmark[form=enclosure, slot=selected, tint, label pos=nw]
                {(2,2) .. (4,4)}{$S$}
        \end{tenkz}
        $\qquad R\subset S=R\cup\{v\}$
    \end{tenkzequation}
\end{definition}

\begin{theorem}[Membership in a one-site separated region]%
    \label{thm:peps_normalOneSiteSeparation_mem_withSite_iff}
    \lean{TNLean.PEPS.NormalOneSiteSeparationHypotheses.mem_withSite_iff}
    \leanok
    \uses{def:peps_normalOneSiteSeparationHypotheses}
    For the two regions associated to a site $v$, a vertex $w$ belongs to
    the region containing $v$ if and only if either $w=v$, or $w$
    belongs to the comparison region omitting $v$.
\end{theorem}
\begin{proof}\leanok
    If $w=v$, use the distinguished-site membership and nonmembership
    assumptions. If $w\neq v$, use the equality of the two regions away
    from the distinguished site.
\end{proof}

\begin{definition}[Normal square-lattice blocking regions]%
    \label{def:peps_normalSquareBlockingRegions}
    \lean{TNLean.PEPS.NormalSquareBlockingRegions}
    \lean{TNLean.PEPS.normalSquareBlockingRegions_of_TCover}
    \leanok
    \uses{def:peps_normalRectangleInjectivityHypotheses,
        def:peps_regionInjectivityData,
        lem:peps_normalSquareRegionRS_injective_of_unionClosure,
        lem:peps_normalSquareRegionT_injective_of_cover}
    In the translationally invariant square-lattice proof, the source paper
    assumes that every $2\times 3$ and $3\times 2$ region is injective.
    The formal statement here is the corresponding abstract, scope-restricted
    hypothesis bundle: it records the rectangular injectivity assumptions, the
    size bound $n,m\geq 7$ with $|V|=nm$, and three finite regions
    $R$, $S$, and $T$ whose injectivity is assumed directly. It does
    not by itself assert that these regions are the displayed square-lattice
    unions, nor that their translated variants give the edge-centred blocking
    around every edge. Conversely, for the coordinate square lattice with
    $n,m\geq7$, the rectangular hypotheses, union closure, and a rectangular
    cover of the displayed $T$-region yield this abstract $R,S,T$
    structure, with $R$ and $S$ obtained from their rectangular
    decompositions and $T$ obtained from the supplied cover.
    % Current source-comparison status: docs/paper-gaps/peps_normal_ft_section3_route.tex.
    The source regions are:
    \begin{tenkzequation}
        % Source/formula verdict: Molnar et al., arXiv:1804.04964,
        % paper_normal.tex:1405--1429 and source PDF p. 15 depict
        % R=(2-4,2-4)\{(4,4)} and S=(2-4,2-4).
        % Ink-to-index verdict: every dot is an ordinary tensor; every lattice
        % edge is a virtual index, and the selected fill records membership only.
        % Structural count verdict (not mechanically emitted as a lattice
        % signature): each 5-by-5 panel has 25 sites, 40 internal virtual
        % contractions, and 20 cut virtual bonds; (4,4) remains ordinary.
        % Structurally derived boundary/type verdict: no physical index is
        % exposed, so each panel has (V,P)=(20,0).
        \tenkzkernel
        \begin{tenkz}[rows={ket,ket,ket,ket,ket}, cols=5,
            west=open, east=open, north=open, south=open]
            \tnmark[form=enclosure, slot=selected, tint, label pos=nw]
                {(2,2) .. (4,4) - (4,4)}{$R$}
        \end{tenkz}
        $\qquad\text{and}\qquad$
        \begin{tenkz}[rows={ket,ket,ket,ket,ket}, cols=5,
            west=open, east=open, north=open, south=open]
            \tnmark[form=enclosure, slot=selected, tint, label pos=nw]
                {(2,2) .. (4,4)}{$S$}
        \end{tenkz}
    \end{tenkzequation}
    \begin{tenkzequation}
        % Formula verdict: in row-column coordinates the exact local formal
        % region is T=(1-6,1-5) minus the vertical block (3-5,1-2) and
        % the shifted horizontal block (2-3,3-5).
        % Ink-to-index verdict: every dot is an ordinary tensor, every lattice
        % edge is a virtual contraction, and the selected fill records only
        % membership in T; it adds no contraction or tensor operation.
        % Contraction/boundary verdict: the 6-by-5 panel has 49 internal
        % virtual contractions and 22 cut virtual bonds, with no exposed
        % physical indices, hence boundary type (V,P)=(22,0).
        % Source verdict: arXiv:1804.04964, paper_normal.tex:1430--1445
        % depicts the same 5-by-6 window complement; this panel uses the exact
        % cell set of normalSquareRegionT rather than the source's rounded ink.
        \tenkzkernel
        \begin{tenkz}[rows={ket,ket,ket,ket,ket,ket}, cols=5,
            west=open, east=open, north=open, south=open]
            \tnmark[form=enclosure, slot=selected, tint, label pos=n]
                {(1,1) .. (6,5) - (3,1) .. (5,2) - (2,3) .. (3,5)}{$T$}
        \end{tenkz}
    \end{tenkzequation}
\end{definition}

\begin{theorem}[Normal edge blocking to a three-site injective chain]%
    \label{thm:peps_normalEdgeBlockingToInjectiveChain}
    \notready
    \uses{lem:peps_injectiveRegion_union,
        def:peps_normalSquareBlockingRegions,
        def:peps_normalEdgeBlockingHypotheses,
        thm:peps_normalEdgeBlockingHypotheses_injectiveChain,
        thm:peps_edgeBlockedInsertionAlgebraIsomorphism}
    Under the square-lattice normality hypotheses of
    \cite[Theorem~3]{Molnar2018NormalPEPS}, the regions $R$, $S$, and
    $T$, together with their translated variants, give an edge-centred
    blocking into red, blue, and complementary injective regions. These are
    the regions on which the three-site injective-chain insertion argument is
    applied:
    \begin{tenkzequation}
        % Formula verdict: A1=(4-6,2-3), A2=(3-4,4-6), and A3 is their
        % complement in the 7-by-7 frame; blocking contracts each region to
        % one injective chain site and carries the marked A1--A2 edge across.
        % Ink-to-index verdict: left-hand dots are ordinary tensors and lattice
        % edges are virtual contractions; right-hand boxes are the three
        % blocked tensors, joined by virtual bonds with physical legs exposed.
        % Contraction/boundary verdict: the lattice has 84 internal virtual
        % contractions and boundary type (V,P)=(28,0); the cyclic blocked
        % network has no open virtual bonds and three physical legs, so
        % (V,P)=(0,3).
        % Source verdict: exact discrete translation and contraction summary of
        % arXiv:1804.04964, paper_normal.tex:1475--1499.
        \tenkzkernel
        \begin{tenkz}[rows={ket,ket,ket,ket,ket,ket,ket}, cols=7,
            west=open, east=open, north=open, south=open]
            % The authored marked wire restates the (4,3)--(4,4) grid bond.
            \tnwire[species=marked]{(4,3)}{(4,4)}
            \tnmark[form=enclosure, slot=selected, tint]
                {(4,2) .. (6,3)}{$A_1$}
            \tnmark[form=enclosure, slot=secondary, tint]
                {(3,4) .. (4,6)}{$A_2$}
            \tnmark[form=enclosure, slot=complement, label pos=n]
                {(1,1) .. (7,7) - (4,2) .. (6,3) - (3,4) .. (4,6)}{$A_3$}
        \end{tenkz}
        $\quad\Rightarrow\quad$
        \begin{tenkz}[rows={ket}, cols=3, physical=up, boundary=periodic]
            \tn[skin=box]{A_1} & \tn[skin=box]{A_2} & \tn[skin=box]{A_3}
            \tnwire[species=marked]{(1,1)}{(1,2)}
            \tnmark[form=enclosure, slot=selected]{(1,1)}{}
            \tnmark[form=enclosure, slot=secondary]{(1,2)}{}
            \tnmark[form=enclosure, slot=complement]{(1,3)}{}
        \end{tenkz}
    \end{tenkzequation}
\end{theorem}

\begin{theorem}[Translationally invariant normal gauge absorption]%
    \label{thm:peps_tiNormalGaugeAbsorption}
    \notready
    \uses{thm:peps_normalEdgeBlockingToInjectiveChain,
        thm:peps_edge_gauge_absorption}
    In the translationally invariant normal square-lattice case, the edge
    gauges obtained after blocking are represented by one horizontal matrix
    $X$ and one vertical matrix $Y$, unique up to scalar. Absorbing these
    matrices into $B$ gives a tensor $\widetilde B$, equivalently the final
    local gauge formula of \cite[Theorem~3]{Molnar2018NormalPEPS}:
    % Formula (paper_normal.tex:1451--1471):
    % B = lambda (Y on north)(Y^{-1} on south)(X on west)(X^{-1} on east) A.
    % Ink-to-index: the four gauge capsules act on the four virtual indices of A;
    % the diagonal north-east stub is the unchanged physical index.
    % Boundary signature: both sides have (V,P)=(4,1).  Source verdict: exact
    % translation-invariant square-lattice gauge absorption in house style.
    \begin{tenkzequation}
        \tenkzkernel
        % Each unbonded typed port renders its own stub, so the four virtual
        % faces and the diagonal physical index need no boundary atoms.
        \begin{tenkz}[rows={wire}, cols=1, bonds=none]
            \tn[label pos=sw,
                ports={180:virtual, 0:virtual, 90:virtual, 270:virtual,
                       45:physical}]{B}
        \end{tenkz}
        $ = \lambda $
        \begin{tenkz}[rows={wire}, cols=1, bonds=none]
            \tn[name=tiAbsA, label pos=sw,
                ports={180:virtual, 0:virtual, 90:virtual, 270:virtual,
                       45:physical}]{A}
            \tn[at=1 w of tiAbsA, name=tiAbsX, skin=ring, species=operator,
                ports={0:virtual, 180:virtual}]{X}
            \tn[at=1 e of tiAbsA, name=tiAbsXinv, skin=ring,
                species=operator,
                ports={180:virtual, 0:virtual}]{X^{-1}}
            \tn[at=1 n of tiAbsA, name=tiAbsY, skin=ring, species=operator,
                ports={270:virtual, 90:virtual}]{Y}
            \tn[at=1 s of tiAbsA, name=tiAbsYinv, skin=ring,
                species=operator,
                ports={90:virtual, 270:virtual}]{Y^{-1}}
            \tnwire{tiAbsA.180}{tiAbsX.0}
            \tnwire{tiAbsA.0}{tiAbsXinv.180}
            \tnwire{tiAbsA.90}{tiAbsY.270}
            \tnwire{tiAbsA.270}{tiAbsYinv.90}
        \end{tenkz}
    \end{tenkzequation}
    The torus form of this absorption, with the gauge family carried onto
    every edge by the translations from one matrix per orientation class, is
    part of Theorem~\ref{thm:fundamentalTheorem_normalTorusPEPS_unconditional}.
    The present open-lattice entry awaits a translation-invariance predicate
    on the open lattice, which would make the per-edge gauges of the edge
    blocking one horizontal and one vertical matrix.
\end{theorem}

\begin{theorem}[Fundamental Theorem for normal square-lattice PEPS]%
    \label{thm:fundamentalTheorem_normalSquarePEPS}
    \notready
    \uses{thm:peps_tiNormalGaugeAbsorption,
        thm:peps_oneVertexComplementComparison,
        def:GaugeEquivPEPS}
    Let $A$ and $B$ be translationally invariant normal square-lattice PEPS
    tensors such that every $2\times 3$ and $3\times 2$ region is
    injective. If they generate the same state on an $n\times m$ region with
    $n,m\geq 7$, then they are related by horizontal and vertical gauges
    $X,Y$, up to a scalar $\lambda$ satisfying
    $\lambda^{nm}=1$. This is
    \cite[Theorem~3]{Molnar2018NormalPEPS}.
    Theorem~\ref{thm:fundamentalTheorem_normalSquarePEPS_unconditional} proves
    an interior-window form of this statement from the source hypotheses alone,
    without translation invariance and with per-vertex scalars in place of the
    single $\lambda$; the gauge-equivalence conclusion on a general connected
    graph, from region injectivity alone, is
    Theorem~\ref{thm:fundamentalTheorem_normalPEPS_connected}. The present
    source-exact entry awaits a translation-invariance predicate on the open
    lattice, which would make the per-edge gauges one horizontal and one
    vertical matrix and the per-vertex scalars a single $\lambda$ with
    $\lambda^{nm}=1$.
\end{theorem}

\begin{theorem}[Fundamental Theorem for normal square-lattice PEPS,
        interior-window form]%
    \label{thm:fundamentalTheorem_normalSquarePEPS_unconditional}
    \lean{TNLean.PEPS.fundamentalTheorem_normalSquarePEPS_unconditional}
    \leanok
    \uses{thm:fundamentalTheorem_normalSquarePEPS,
        def:peps_normalRectangleInjectivityHypotheses,
        def:applyGauge,
        def:peps_edgeInsertedCoeff}
    Let $A$ and $B$ be PEPS tensors on the open $n\times m$ square lattice
    with the same positive physical dimension, equal and positive bond
    dimensions, generating the same state, and such that every contiguous
    $2\times 3$ and $3\times 2$ coordinate rectangle of either tensor is
    injective. Then there is a family of invertible matrices, one on each
    edge, such that:
    \begin{enumerate}
        \item at every edge with interior margins (for a horizontal edge with
              left endpoint $(x,y)$: $3\le x$, $x+7\le n$, $4\le y$,
              $y+6\le m$; for a vertical edge the mirrored margins), inserting
              any matrix on the edge of $A$ gives the same coefficient as
              inserting it on the corresponding edge of the gauge-absorbed $B$,
              for every physical configuration; and
        \item every vertex $v=(x,y)$ of the interior comparison window
              $4\le x\le n-9$, $4\le y\le m-9$ carries a nonzero scalar
              $\lambda_v$ such that $A_v$ equals $\lambda_v$ times the gauge
              action of the family on $B$ at $v$, on every virtual configuration
              of the incident edges and every physical index.
    \end{enumerate}
    This is the open-lattice form of \cite[Theorem~3]{Molnar2018NormalPEPS},
    restricted in two ways. First, the conclusion reaches only the interior:
    the absorbed equality holds at the edges with interior margins and the
    per-vertex relation at the window vertices, while vertices near the
    lattice boundary are not reached, the blocking frames of the open lattice
    requiring interior margins. Second, the scalars of distinct vertices are
    not asserted equal, and no condition $\lambda^{nm}=1$ is asserted: the
    single $\lambda$ of the source comes from translation invariance, which is
    not assumed here, and without it a single scalar is false --- rescaling
    the components of $A$ at two interior vertices by $\mu$ and $\mu^{-1}$
    preserves the state and the rectangular injectivity but separates the two
    scalars. No lower bound on $n$ or $m$ is assumed; both conclusions are
    empty when the lattice is too small to contain interior edges or window
    vertices.
    \begin{proof}
        \leanok
        \uses{lem:peps_injectiveRegion_union_pos}
        The union closure of injective regions follows from the positive bond
        dimensions and Lemma~\ref{lem:peps_injectiveRegion_union_pos}. The
        translated blocking frames at the edges with interior margins produce
        per-edge gauges, whose absorption into $B$ gives the absorbed
        inserted-coefficient equality at every such edge. A gauge preserves
        blocked-region linear independence, so at every window vertex the
        region comparison applies to the gauge-absorbed tensor: comparing the
        $3\times 3$ square at the vertex minus its corner with its one-site
        completion yields two proportionality scalars $c_R$ and $c_S$, every
        boundary edge of either region carrying interior margins, and the
        inserted-site extraction gives the per-vertex relation with the ratio
        $\lambda_v = c_S/c_R$.
    \end{proof}
\end{theorem}


\section{Torus Normal PEPS}\label{sec:peps_torus}
A translationally invariant PEPS lives on a torus: a finite square lattice whose
opposite boundaries are identified, so that the two coordinates are read modulo
the periods $n$ and $m$.
\begin{center}
    % Formula/source: Papers/1804.04964/paper_normal.tex, lines 104--112,
    % defines a translationally invariant PEPS by repeating one tensor on a
    % regular lattice and contracting according to the lattice geometry;
    % lines 1451--1471 state the normal-PEPS theorem on an n-by-m region.
    % Ink-to-index map: the nine upper legs are the free physical indices of
    % the displayed 3-by-3 window.  Every horizontal and vertical line is a
    % virtual bond; the marked A is a representative of the tensor repeated
    % at every site.
    % Contraction set: twelve nearest-neighbour bonds inside the window are
    % summed, and six further bonds identify the three west/east endpoint
    % pairs and the three north/south endpoint pairs.  No virtual index stays
    % open, while all nine physical indices remain free.
    % Boundary signature: virtual=0 | physical=9.
    % Source verdict: the source has no separate torus-geometry picture.  This
    % is the house-style realization of the chapter's cyclic-coordinate
    % definition; unlike the former open-boundary schematic, it displays the
    % periodic contraction in both coordinate directions.
    \begin{tenkzlattice}[
        rows=3, cols=3, sheets={ket}, outer legs=top, frame=oblique,
        west=trace, east=trace, north=trace, south=trace]
        \tnsite[role=marked, label=$A$]{(1,2)}
    \end{tenkzlattice}
\end{center}
Translations act on the vertices and edges. Transporting a tensor along a graph
isomorphism $\phi$ relabels its coefficients by
$\Coeff_{A^\phi}(\sigma)=\Coeff_A(\sigma\circ\phi)$,
so two tensors generating the same state still do after transport. This is the
setting in which the normal Fundamental Theorem specializes to a single
site-independent tensor and produces the global phase of
\cite[Theorem~3]{Molnar2018NormalPEPS}.

\begin{definition}[Torus geometry and translations]%
    \label{def:peps_torus_geometry_translation}
    \lean{TNLean.PEPS.TorusVertex}
    \lean{TNLean.PEPS.torusGraph}
    \lean{TNLean.PEPS.torusGraph_adj_right}
    \lean{TNLean.PEPS.torusGraph_adj_up}
    \lean{TNLean.PEPS.Edge.map}
    \lean{TNLean.PEPS.Edge.equiv}
    \lean{TNLean.PEPS.IncidentEdge.equiv}
    \lean{TNLean.PEPS.translateEquiv}
    \lean{TNLean.PEPS.translate}
    \lean{TNLean.PEPS.translateEdge}
    \lean{TNLean.PEPS.translateEdge_isHorizontal}
    \lean{TNLean.PEPS.translateEdge_isVertical}
    \lean{TNLean.PEPS.translateIncidentEdge}
    \lean{TNLean.PEPS.torusEdge_horizontal_or_vertical}
    \leanok
    The discrete $n\times m$ torus has vertex set $\Z/n\Z\times\Z/m\Z$.
    Its graph joins $(x,y)$ to $(x+1,y)$ and to $(x,y+1)$, with the
    coordinates read cyclically; every edge is of one of these two forms,
    so every edge is horizontal or vertical. Translation by $(a,b)$ is the
    graph automorphism $\tau_{a,b}(x,y)=(x+a,y+b)$.
    The induced action on edges sends an edge with unordered endpoint pair
    $\{u,v\}$ to the edge with unordered endpoint pair
    $\{\tau_{a,b}(u),\tau_{a,b}(v)\}$; in particular it preserves the
    horizontal and vertical orientation classes. For a general graph
    isomorphism $\phi:G\simeq G'$, the same endpoint action gives a
    bijection on edges and a bijection from the edges incident to $v$ to
    the edges incident to $\phi(v)$.
\end{definition}

\begin{definition}[Transport of a PEPS tensor]%
    \label{def:peps_tensor_transport}
    \lean{TNLean.PEPS.Tensor.transport}
    \lean{TNLean.PEPS.Tensor.transport_bondDim}
    \leanok
    \uses{def:peps_tensor, def:peps_torus_geometry_translation}
    Let $\phi:G\simeq G'$ be a graph isomorphism. The transported tensor
    $A^\phi$ on $G'$ has bond dimension
    $D^\phi_{e'}=D_{\phi^{-1}(e')}$ at every edge $e'$ of $G'$. Its local
    tensor at $w\in G'$ is the local tensor of $A$ at $\phi^{-1}(w)$, with
    virtual indices moved through the incident-edge bijection induced by
    $\phi$.
\end{definition}

\begin{theorem}[State coefficient under transport]%
    \label{thm:peps_stateCoeff_transport}
    \lean{TNLean.PEPS.stateCoeff_transport}
    \leanok
    \uses{def:peps_stateCoeff, def:peps_tensor_transport}
    For every graph isomorphism $\phi:G\simeq G'$ and every physical
    configuration $\sigma$ on the vertices of $G'$,
    $\Coeff_{A^\phi}(\sigma)=\Coeff_A(\sigma\circ\phi)$.
\end{theorem}
\begin{proof}\leanok
    The edge bijection induced by $\phi$ identifies virtual configurations
    of $A^\phi$ with virtual configurations of $A$. After this change of
    summation variables, the local factor at $w\in G'$ becomes the local
    factor of $A$ at $\phi^{-1}(w)$. Reindexing the finite product over
    vertices by $\phi$ gives the coefficient equality.
\end{proof}

\begin{theorem}[Transport preserves equality of PEPS states]%
    \label{thm:peps_sameState_transport}
    \lean{TNLean.PEPS.SameState.transport}
    \leanok
    \uses{def:peps_same_state, thm:peps_stateCoeff_transport}
    If two PEPS tensors on $G$ have equal state coefficients for every
    physical configuration, then their transports along any graph
    isomorphism $\phi:G\simeq G'$ have equal state coefficients for every
    physical configuration on $G'$.
\end{theorem}
\begin{proof}\leanok
    By Theorem~\ref{thm:peps_stateCoeff_transport},
    \begin{align}
        \Coeff_{A^\phi}(\sigma)
        &=\Coeff_A(\sigma\circ\phi)
          =\Coeff_B(\sigma\circ\phi)
          =\Coeff_{B^\phi}(\sigma),
        \notag
    \end{align}
    where the middle equality is the assumed equality of state coefficients.
\end{proof}

\begin{definition}[Translation-invariant torus tensor]%
    \label{def:peps_torus_translation_invariant}
    \lean{TNLean.PEPS.IsTorusTranslationInvariant}
    \leanok
    \uses{def:peps_tensor_transport, def:peps_torus_geometry_translation}
    A PEPS tensor $A$ on the discrete torus is translation invariant if
    $A^{\tau_{a,b}}=A$ for every translation $\tau_{a,b}$. This formalizes
    the setting of \cite[Theorem~3]{Molnar2018NormalPEPS}, where one tensor
    is repeated at every site of a periodic lattice.
\end{definition}

\begin{theorem}[State coefficient of a translation-invariant torus tensor]%
    \label{thm:peps_stateCoeff_translationInvariant}
    \lean{TNLean.PEPS.stateCoeff_translationInvariant}
    \leanok
    \uses{def:peps_torus_translation_invariant,
        thm:peps_stateCoeff_transport}
    If $A$ is translation invariant on the discrete torus, then
    $\Coeff_A(\sigma\circ\tau_{a,b})=\Coeff_A(\sigma)$ for every translation
    $\tau_{a,b}$ and every physical configuration $\sigma$.
\end{theorem}
\begin{proof}\leanok
    Theorem~\ref{thm:peps_stateCoeff_transport} gives
    \begin{align}
        \Coeff_A(\sigma\circ\tau_{a,b})
        &=\Coeff_{A^{\tau_{a,b}}}(\sigma)
          =\Coeff_A(\sigma),
        \notag
    \end{align}
    where the second equality substitutes $A^{\tau_{a,b}}=A$ from translation
    invariance.
\end{proof}

\begin{definition}[Orientation-uniform bond dimensions on the torus]%
    \label{def:peps_torus_uniform_bond_dim}
    \lean{TNLean.PEPS.torusRightEdge}
    \lean{TNLean.PEPS.torusUpEdge}
    \lean{TNLean.PEPS.TorusUniformBondDim}
    \leanok
    \uses{def:peps_torus_geometry_translation}
    A bond-dimension function on the torus is orientation-uniform with
    horizontal dimension $D_h$ and vertical dimension $D_v$ if
    $D_{e_h}=D_h$ for every horizontal edge $e_h$ and $D_{e_v}=D_v$ for
    every vertical edge $e_v$. The reference horizontal edge is the right
    edge at the origin, and the reference vertical edge is the up edge at
    the origin.
\end{definition}

\begin{theorem}[Translation invariance gives orientation-uniform bond dimensions]%
    \label{thm:peps_torusUniformBondDim_of_translationInvariant}
    \lean{TNLean.PEPS.translateEdge_eq_map}
    \lean{TNLean.PEPS.bondDim_translateEdge_of_translationInvariant}
    \lean{TNLean.PEPS.isHorizontalTorusEdge_eq_rightEdge}
    \lean{TNLean.PEPS.isVerticalTorusEdge_eq_upEdge}
    \lean{TNLean.PEPS.translateEdge_torusRightEdge}
    \lean{TNLean.PEPS.translateEdge_torusUpEdge}
    \lean{TNLean.PEPS.bondDim_torusRightEdge_const}
    \lean{TNLean.PEPS.bondDim_torusUpEdge_const}
    \lean{TNLean.PEPS.torusUniformBondDim_of_translationInvariant}
    \leanok
    \uses{def:peps_torus_translation_invariant,
        def:peps_torus_uniform_bond_dim}
    If $A$ is translation invariant on the discrete torus, then its bond
    dimensions are orientation-uniform. Explicitly, with
    $D_h=D_{e_h}$ at the reference right edge $e_h$ and
    $D_v=D_{e_v}$ at the reference up edge $e_v$, one has
    $D_{f_h}=D_h$ for every horizontal edge $f_h$ and $D_{f_v}=D_v$ for
    every vertical edge $f_v$.
\end{theorem}
\begin{proof}\leanok
    Translation invariance gives $A^{\tau_{a,b}}=A$, hence
    $D_A(\tau_{a,b}(e))=D_A(e)$ for every edge $e$. Every horizontal edge
    $f_h$ has the form $f_h=\tau_{a,b}(e_h)$ for a translate of the reference
    right edge, so
    \begin{align}
        D_{f_h}
        &=D_A(\tau_{a,b}(e_h))
          =D_A(e_h)
          =D_h.
        \notag
    \end{align}
    The same argument with the reference up edge gives $D_{f_v}=D_v$.
\end{proof}

\paragraph{Failure of the row-and-column reduction.}

\begin{definition}[Per-edge product matrix on two vertical bonds]%
    \label{def:peps_row_cut_per_edge_product}
    \lean{TNLean.PEPS.RowColumnReductionObstruction.IsPerEdgeProduct}
    \leanok
    A matrix $M\in \MN{4}$ is a per-edge product on two collected vertical
    bonds if, under the identification
    $(r,s)\longmapsto 2r+s$ between
    $\{0,1\}\times\{0,1\}$ and $\{0,1,2,3\}$, there are matrices
    $Y_0,Y_1\in \MN{2}$ such that, for all
    $r,s,r',s'\in\{0,1\}$,
    \begin{align}
        M_{(r,s),(r',s')}
        &=(Y_0)_{r,r'}(Y_1)_{s,s'}.
        \notag
    \end{align}
\end{definition}

\begin{definition}[Per-edge factorization of the row-cut gauge]%
    \label{def:peps_row_cut_gauge_factorizes}
    \lean{TNLean.PEPS.RowColumnReductionObstruction.RowCutGaugeFactorizes}
    \leanok
    \uses{def:mps_tensor, def:l_blk_injective, def:same_mpv,
        def:peps_row_cut_per_edge_product}
    The row-cut factorization assertion is the following statement. Let $A$
    and $B$ be MPS tensors with physical alphabet $\{0,\ldots,15\}$ and
    bond space $\C^4$. If $A$ and $B$ are $1$-block injective and generate
    the same MPV family, then there exist $Z\in \GL_4(\C)$ and
    $\lambda\in\C$ such that $Z$ is a per-edge product matrix and
    $B^i=\lambda Z^{-1}A^iZ$ for all physical indices
    $i\in\{0,\ldots,15\}$. This is the additional product conclusion needed
    by a row-wise one-dimensional reduction before it could imply the per-edge
    conclusion of \cite[Theorem~3]{Molnar2018NormalPEPS}.
\end{definition}

\begin{theorem}[A row-cut gauge need not factor per edge]%
    \label{thm:peps_row_cut_gauge_factorizes_false}
    \lean{TNLean.PEPS.RowColumnReductionObstruction.rowCutGaugeFactorizes_false}
    \leanok
    \uses{def:peps_row_cut_gauge_factorizes}
    The row-cut factorization assertion is false. Explicitly, let
    $E_{ab}\in\MN{4}$ be the matrix units and set
    \begin{align}
        A^{(a,b)}
        &=E_{ab}, \notag\\
        G
        &=
        \begin{pmatrix}
            1&0&1&0\\
            0&1&0&0\\
            0&0&1&0\\
            0&0&0&1
        \end{pmatrix}, \notag\\
        B^{(a,b)}
        &=G^{-1}E_{ab}G.
        \notag
    \end{align}
    Then $A$ and $B$ are $1$-block injective and generate the same MPV
    family, but there are no $Z\in\GL_4(\C)$ and $\lambda\in\C$ with
    $Z$ a per-edge product matrix and
    $B^{(a,b)}=\lambda Z^{-1}E_{ab}Z$ for all
    $a,b\in\{0,\ldots,3\}$.
\end{theorem}

\begin{proof}
    \leanok
    The matrix units span $\MN{4}$, hence $A$ is $1$-block injective; conjugation
    by $G$ preserves this property for $B$. For every word
    $w=((a_1,b_1),\ldots,(a_N,b_N))$, one has
    $B^w=G^{-1}A^wG$ and $\tr(B^w)=\tr(A^w)$, so the two MPV families are
    equal.

    Suppose that $Z$ and $\lambda$ satisfy
    $B^{(a,b)}=\lambda Z^{-1}E_{ab}Z$ for every $a,b\in\{0,\ldots,3\}$.
    With $W=ZG^{-1}$, multiplication by $Z$ on the left and by $G^{-1}$
    on the right gives $WE_{ab}=\lambda E_{ab}W$ for every matrix unit
    $E_{ab}$. Taking the $(x,y)$ entry gives
    $W_{x,a}\delta_{b,y}=\lambda\delta_{x,a}W_{b,y}$ for all
    $x,y,a,b\in\{0,\ldots,3\}$. With $y=b$, this gives
    $x\ne a\Rightarrow W_{x,a}=0$ and
    $W_{a,a}=\lambda W_{b,b}$.

    Since $Z$ and $G$ are invertible, so is $W$. Thus the diagonal entries
    cannot all vanish; the case $a=b$ gives $\lambda=1$, and then
    $W_{a,a}=W_{b,b}$ for all $a,b$. Hence $W=c\Id_4$ and $\lambda=1$ for
    some $c\ne0$, so $Z=cG$.

    For a matrix $M\in\MN{4}$, write $M_{rr'}\in\MN{2}$ for the block
    whose $(s,s')$ entry is $M_{(r,s),(r',s')}$. If $Z$ were a per-edge
    product, then $Z_{rr'}=(Y_0)_{r,r'}Y_1$ for some
    $Y_0,Y_1\in\MN{2}$. In particular,
    $(Z_{00})_{s,s'}(Z_{01})_{1,1}
    =(Z_{01})_{s,s'}(Z_{00})_{1,1}$.
    Since $Z=cG$ and $c\ne0$, scaling cancels from this block
    proportionality condition. The corresponding blocks of $G$ are
    \begin{align}
        G_{00}
        &=
        \begin{pmatrix}
            1&0\\
            0&1
        \end{pmatrix}, \notag\\
        G_{01}
        &=
        \begin{pmatrix}
            1&0\\
            0&0
        \end{pmatrix}.
        \notag
    \end{align}
    Substitution into
    $(G_{00})_{00}(G_{01})_{11}
    =(G_{01})_{00}(G_{00})_{11}$ gives $0=1$. Hence no per-edge product
    conjugator can exist.
\end{proof}


\paragraph{Staircase-window geometry for a reference edge.}

\begin{lemma}[Left endpoint of the staircase reference edge]%
    \label{lem:peps_reference_edge_left_endpoint_staircase_window}
    \lean{TNLean.PEPS.referenceEdge_leftEndpoint_mem_staircaseWindow}
    \leanok
    Let $L,K>0$, let $s=(a,b)$ be a staircase corner on the torus, and assume
    $a+2L\le n$ and $2K\le m$. Let $e$ be the horizontal reference edge
    from $(a+L-1,b+K-1)$ to $(a+L,b+K-1)$, and let $W_j$ be the $j$th
    staircase window based at $s$. Then the left endpoint
    $u=(a+L-1,b+K-1)$ lies in $W_j$ if and only if $1\le j$.
\end{lemma}
\begin{proof}\leanok
    The endpoint formula gives the two coordinate distances from the corner:
    the column distance is $L-1$ and the row distance is $K-1$. Substituting
    these distances into the staircase-window membership criterion gives exactly
    $1\le j$.
\end{proof}

\begin{lemma}[Right endpoint of the staircase reference edge]%
    \label{lem:peps_reference_edge_right_endpoint_staircase_window}
    \lean{TNLean.PEPS.referenceEdge_rightEndpoint_mem_staircaseWindow}
    \leanok
    In the same notation, the right endpoint $w=(a+L,b+K-1)$ lies in $W_j$
    if and only if $j<L$.
\end{lemma}
\begin{proof}\leanok
    The right endpoint has column distance $L$ and row distance $K-1$ from
    $s$. The row condition is automatic, while the column band contains this
    distance precisely for $j<L$.
\end{proof}

\begin{lemma}[Boundary windows for the staircase reference edge]%
    \label{lem:peps_reference_edge_boundary_windows}
    \lean{TNLean.PEPS.isRegionBoundaryEdge_staircaseWindow_referenceEdge}
    \leanok
    \uses{lem:peps_reference_edge_left_endpoint_staircase_window,
        lem:peps_reference_edge_right_endpoint_staircase_window}
    If $j=0$ or $L\le j$, then the reference edge $e$ is a boundary edge
    of $W_j$: exactly one of its two endpoints lies in the window.
\end{lemma}
\begin{proof}\leanok
    For $j=0$, Lemmas~\ref{lem:peps_reference_edge_left_endpoint_staircase_window}
    and~\ref{lem:peps_reference_edge_right_endpoint_staircase_window} show that
    only $w$ lies in $W_j$. For $L\le j$, the same two endpoint tests show
    that only $u$ lies in $W_j$.
\end{proof}

\begin{lemma}[Interior windows for the staircase reference edge]%
    \label{lem:peps_reference_edge_interior_windows}
    \lean{TNLean.PEPS.referenceEdge_endpoints_mem_staircaseWindow_of_interior}
    \leanok
    \uses{lem:peps_reference_edge_left_endpoint_staircase_window,
        lem:peps_reference_edge_right_endpoint_staircase_window}
    If $1\le j<L$, then both endpoints of $e$ lie in $W_j$.
\end{lemma}
\begin{proof}\leanok
    The left-endpoint criterion gives $u\in W_j$ from $1\le j$, and the
    right-endpoint criterion gives $w\in W_j$ from $j<L$.
\end{proof}

\begin{lemma}[Interior windows are not boundary windows]%
    \label{lem:peps_reference_edge_not_boundary_interior_windows}
    \lean{TNLean.PEPS.not_isRegionBoundaryEdge_staircaseWindow_referenceEdge_of_interior}
    \leanok
    \uses{lem:peps_reference_edge_interior_windows}
    If $1\le j<L$, then $e$ is not a boundary edge of $W_j$.
\end{lemma}
\begin{proof}\leanok
    By Lemma~\ref{lem:peps_reference_edge_interior_windows}, both endpoints of
    $e$ lie in $W_j$. The definition of a boundary edge requires exactly one
    endpoint in the region, so $e$ is not a boundary edge of $W_j$.
\end{proof}

\subsection{Staircase patch geometry}
\label{subsec:peps_torus_staircase_patch_geometry}

\begin{definition}[Horizontal staircase patch and its corner]
    \label{def:peps_horizontal_staircase_patch}
    \lean{TNLean.PEPS.horizontalStaircasePatch}
    \lean{TNLean.PEPS.horizontalStaircaseLeftWindow}
    \lean{TNLean.PEPS.horizontalStaircaseRightWindow}
    \lean{TNLean.PEPS.horizontalStaircaseEndPair}
    \lean{TNLean.PEPS.horizontalStaircaseCorner}
    \lean{TNLean.PEPS.horizontalStaircaseCompletedCorner}
    \leanok
    \uses{def:peps_torus_geometry_translation}
    Fix a horizontal reference edge and write the staircase corner as
    $s=(a,b)$.  For coordinate intervals write
    \[
        I(u,v)=\{t\mid u\le t<v\}.
    \]
    The two end windows of the overlapping-window chain are
    \[
        W_{\mathrm{left}}=I(a,a+L)\times I(b,b+K),
        \qquad
        W_{\mathrm{right}}
        =I(a+L,a+2L)\times I(b+K-1,b+2K-1).
    \]
    Their end pair is $S=W_{\mathrm{left}}\cup W_{\mathrm{right}}$.  The
    staircase patch is the union
    \[
        P=
        \bigl(I(a,a+2L)\times I(b+K-1,b+2K-1)\bigr)
        \cup
        \bigl(I(a,a+L)\times I(b,b+2K-1)\bigr),
    \]
    and the corner cut out by the two end windows is
    \[
        C=I(a,a+L)\times I(b+K,b+2K-1).
    \]
    Completing $C$ by one row gives the $L\times K$ rectangle
    \[
        \widehat C=I(a,a+L)\times I(b+K,b+2K).
    \]
    These cyclic rectangles are the regions used in the staircase chaining and
    corner-stripping step of
    \cite[proof sketch, lines~2320--2445]{Molnar2018NormalPEPS}.
\end{definition}

\begin{lemma}[The staircase end windows are disjoint]
    \label{lem:peps_horizontal_staircase_end_pair_disjoint}
    \lean{TNLean.PEPS.horizontalStaircaseEndPair_disjoint}
    \leanok
    \uses{def:peps_horizontal_staircase_patch}
    If $2L\le W$, then the two end windows of the horizontal staircase are
    disjoint:
    \[
        W_{\mathrm{left}}\cap W_{\mathrm{right}}=\varnothing.
    \]
\end{lemma}
\begin{proof}
    \leanok
    Their horizontal coordinate intervals are $I(a,a+L)$ and $I(a+L,a+2L)$.
    The bound $2L\le W$ prevents wraparound overlap.
\end{proof}

\begin{lemma}[Injectivity of the staircase end windows and completed corner]
    \label{lem:peps_horizontal_staircase_window_injectivity}
    \lean{TNLean.PEPS.NormalTorusArcWindowInjectivityHypotheses.horizontalStaircaseLeftWindow_injective}
    \lean{TNLean.PEPS.NormalTorusArcWindowInjectivityHypotheses.horizontalStaircaseRightWindow_injective}
    \lean{TNLean.PEPS.NormalTorusArcWindowInjectivityHypotheses.horizontalStaircaseCompletedCorner_injective}
    \leanok
    \uses{def:peps_horizontal_staircase_patch}
    Under the one-orientation $L\times K$ torus-window injectivity hypotheses,
    the two end windows and the completed corner are injective regions:
    \[
        W_{\mathrm{left}},\quad W_{\mathrm{right}},\quad \widehat C
        \quad\text{are injective.}
    \]
\end{lemma}
\begin{proof}
    \leanok
    Each of the three regions is an $L\times K$ cyclic rectangle, so the
    window-injectivity hypothesis applies directly.
\end{proof}

\begin{lemma}[Patch decomposition by the end pair and the corner]
    \label{lem:peps_horizontal_staircase_patch_decomposition}
    \lean{TNLean.PEPS.horizontalStaircaseCorner_disjoint_endPair}
    \lean{TNLean.PEPS.horizontalStaircasePatch_eq_endPair_union_corner}
    \lean{TNLean.PEPS.horizontalStaircasePatch_sdiff_endPair}
    \leanok
    \uses{def:peps_horizontal_staircase_patch,
        lem:peps_horizontal_staircase_end_pair_disjoint}
    If $L,K>0$, $2L\le W$, and $2K\le H$, then the corner is disjoint from
    the end pair and the patch decomposes as
    \[
        P=S\cup C,
        \qquad
        P\setminus S=C.
    \]
\end{lemma}
\begin{proof}
    \leanok
    In the no-wraparound rectangle, the rows and columns of the two end
    windows and the corner are the displayed intervals above.  The interval
    calculation gives the disjointness and identifies every point of $P$ with
    either an end-window point or a corner point.
\end{proof}

\begin{theorem}[Fundamental Theorem for normal PEPS on the torus]%
    \label{thm:fundamentalTheorem_normalTorusPEPS}
    \lean{TNLean.PEPS.fundamentalTheorem_normalTorusPEPS_unconditional}
    \lean{TNLean.PEPS.torusAbsorbedGauge_unique_scalar}
    \lean{TNLean.PEPS.torusCovariantAbsorbedGauge_unique_classScalar}
    \lean{TNLean.PEPS.torusGauge_unique_scalar_of_perVertex}
    \leanok
    \uses{thm:fundamentalTheorem_normalSquarePEPS,
        thm:peps_oneVertexComplementComparison,
        def:peps_torus_translation_invariant,
        thm:peps_torusUniformBondDim_of_translationInvariant,
        def:GaugeEquivPEPS}
    Let $A$ and $B$ be translationally invariant normal PEPS on the discrete
    torus such that every $2\times 3$ and $3\times 2$ region is injective. If
    they generate the same state on an $n\times m$ region with $n,m\geq 7$, then
    they are related by horizontal and vertical gauges $X,Y$, up to a scalar
    $\lambda$ satisfying $\lambda^{nm}=1$, and $X$ and $Y$ are unique up to a
    multiplicative constant. This is
    \cite[Theorem~3]{Molnar2018NormalPEPS} on the torus geometry, where every
    edge is interior and translation acts freely.
    Theorem~\ref{thm:fundamentalTheorem_normalTorusPEPS_unconditional} proves
    this at the source's sizes $n,m\geq 7$ with the gauge family described
    by its translation covariance, and the uniqueness of $X$ and $Y$ up to a
    multiplicative constant is proved at the same sizes in
    Theorem~\ref{thm:torusAbsorbedGauge_unique_scalar} per edge,
    Theorem~\ref{thm:torusCovariantAbsorbedGauge_unique_classScalar} on the
    two orientation classes, and
    Theorem~\ref{thm:torusGauge_unique_scalar_of_perVertex} against the
    per-vertex relation. In the torus
    edge convention the two class matrices $X$ and $Y$ are read through this
    translation-covariant family: seam edges carry the transposed inverse of
    the corresponding class matrix, so this is the orientation-faithful form
    of the source's horizontal and vertical gauges.
\end{theorem}

\begin{theorem}[Fundamental Theorem for normal PEPS on the torus,
        translation-covariant form]%
    \label{thm:fundamentalTheorem_normalTorusPEPS_unconditional}
    \lean{TNLean.PEPS.fundamentalTheorem_normalTorusPEPS_unconditional}
    \leanok
    \uses{thm:fundamentalTheorem_normalSquarePEPS,
        thm:peps_oneVertexComplementComparison,
        thm:peps_sameState_transport,
        thm:peps_stateCoeff_translationInvariant,
        thm:peps_torusUniformBondDim_of_translationInvariant,
        def:GaugeEquivPEPS}
    Let $A$ and $B$ be translationally invariant PEPS tensors on the discrete
    $n\times m$ torus with $n,m\geq 7$, with positive physical dimension, equal
    and positive bond dimensions,
    generating the same state, and such that every contiguous $2\times 3$ and
    $3\times 2$ coordinate rectangle of either tensor is injective. Then there
    are a family of invertible matrices, one on each edge, carried onto each
    edge by the torus translations from a single matrix on a reference
    horizontal edge and a single matrix on a reference vertical edge, and a
    single scalar $\lambda$ with $\lambda^{nm}=1$, such that at every site the
    tensor $A$ equals $\lambda$ times the gauge action of the family on $B$.
    The family also realizes the absorbed inserted-coefficient equality at
    every edge: inserting any matrix on an edge of $A$ gives the same
    coefficient as inserting it on the corresponding edge of the gauge-absorbed
    $B$.
    This is \cite[Theorem~3]{Molnar2018NormalPEPS} on the torus at the source's
    sizes $n,m\geq 7$, with the horizontal and vertical gauges $X$, $Y$
    described by translation covariance: the stored orientation of a
    wraparound edge is reversed, so on the seam the family carries the
    transposed inverse of the class matrix, and a family that is literally one
    matrix per orientation class cannot realize the absorbed equality there.
    The translation covariance of the family, with the scalar $\lambda$ kept
    separate rather than absorbed, realizes the remark of
    \cite{Molnar2018NormalPEPS} following the general normal theorem: for a
    translationally invariant system the gauges are also translationally
    invariant, provided the proportionality constants are not absorbed into
    the gauges.
    The uniqueness of $X$ and $Y$ up to a multiplicative constant is not part
    of this statement; the uniqueness clause is
    Theorem~\ref{thm:torusAbsorbedGauge_unique_scalar} (per edge, against the
    absorbed equality), Theorem~\ref{thm:torusCovariantAbsorbedGauge_unique_classScalar}
    (one constant per orientation class), and
    Theorem~\ref{thm:torusGauge_unique_scalar_of_perVertex} (against the
    per-vertex relation itself).
    \begin{proof}
        \leanok
        \uses{lem:peps_injectiveRegion_union_pos}
        The union closure of injective regions follows from the positive bond
        dimensions and Lemma~\ref{lem:peps_injectiveRegion_union_pos}. The
        reference blockings are anchored against the far seam: the removed
        blocks keep a margin of at least two rows and columns below and to
        the left and end exactly at the seam above and to the right, so the
        complementary region decomposes into wraparound-free rectangular
        bands on every torus with $n,m\geq 7$. The
        coherent blocking frames at every edge produce per-edge gauges, whose
        absorption into $B$ gives the absorbed inserted-coefficient equality at
        every edge; the absorbing family is built by transporting one reference
        gauge per orientation class, so it commutes with the torus
        translations. A gauge preserves blocked-region linear independence, so
        the region comparison applies to the gauge-absorbed tensor: comparing
        the corner region with its one-site completion yields proportionality
        scalars $c_R$ and $c_S$ such that, with $\lambda=c_S/c_R$,
        \[
            A_v(\eta,\sigma)
            =
            \lambda\,
            \bigl(\text{gauge action of }X\text{ on }B\bigr)_v(\eta,\sigma).
        \]
        Translation covariance carries the same scalars to every site. Taking
        the product over all vertices and using the invariance of the closed
        state coefficient gives
        \[
            \lambda^{nm}=1.
        \]
    \end{proof}
\end{theorem}

\begin{corollary}[Torus scalar condition from the per-vertex two-block proportionalities]%
    \label{cor:peps_lambda_pow_card_torus_eq_one_of_twoBlockProportional}
    \lean{TNLean.PEPS.lambda_pow_card_torus_eq_one_of_twoBlockProportional}
    \leanok
    \uses{cor:peps_component_eq_gaugeVertex_of_twoBlockProportional}
    Let $A$ and $B$ be PEPS tensors on the discrete $n\times m$ torus with $n,m>1$,
    with equal bond dimensions and positive bond dimensions, generating the same
    state, and let $X$ be a family of invertible matrices, one on each edge. Suppose
    that at every site $v$ a region $R_v$ not containing $v$ is supplied together with
    the two two-block scalar proportionalities of the blocked weights of $A$ and of
    the gauge-absorbed second tensor over $R_v$ and over $R_v\cup\{v\}$, with ratios
    $c_R(v)$ and $c_S(v)$ satisfying $c_R(v)\ne 0$ and whose quotient
    $c_S(v)/c_R(v)$ is a single scalar $\lambda$ independent of $v$, and with the
    gauge-absorbed region block linearly independent. If, in addition, a single
    comparison region whose block and complement block are both blocked-tensor
    injective is supplied, then $\lambda$ raised to the number of torus sites is one:
    \[
        \lambda^{nm}=1.
    \]
\end{corollary}
\begin{proof}\leanok
    At each of the $nm$ sites $v$ the inserted-site scalar extraction of
    Corollary~\ref{cor:peps_component_eq_gaugeVertex_of_twoBlockProportional} turns
    the two two-block proportionalities into the per-vertex gauge relation
    \[
        A_v(\eta,\sigma)
        =
        \lambda\,\widetilde B_v(\eta,\sigma),
    \]
    where $\widetilde B_v$ is the gauge action of $X$ on $B$ at $v$. Contracting
    both sides over the whole torus pulls one factor of $\lambda$ out of each of
    the $nm$ sites, so the closed torus contractions of $A$ and of $\widetilde B$
    satisfy
    \[
        \langle A\rangle_{\mathrm{torus}}
        =
        \lambda^{nm}\,\langle \widetilde B\rangle_{\mathrm{torus}}.
    \]
    The single comparison region and its complement being injective makes this
    common torus contraction nonzero; by hypothesis the quotient
    $c_S(v)/c_R(v)$ is the same scalar $\lambda$ at every site, so
    $\lambda^{nm}=1$.
\end{proof}

\begin{theorem}[Per-edge uniqueness of the torus gauge up to a constant]%
    \label{thm:torusAbsorbedGauge_unique_scalar}
    \lean{TNLean.PEPS.torusAbsorbedGauge_unique_scalar}
    \leanok
    \uses{def:applyGauge, def:peps_edgeInsertedCoeff}
    Let $A$ and $B$ be translationally invariant PEPS tensors on the discrete
    $n\times m$ torus with $n,m\geq 7$, with positive physical dimension,
    equal and positive bond dimensions, generating the same state, and such
    that every contiguous
    $2\times 3$ and $3\times 2$ coordinate rectangle of either tensor is
    injective.  Let $X$ and $X'$ be two families of invertible matrices, one
    on each edge of the torus, and let $e$ be an edge at which both families
    realize the absorbed inserted-coefficient equality: inserting any matrix
    on the edge $e$ of $A$ gives the same coefficient as inserting it on the
    edge $e$ of the gauge-absorbed $B$, for every physical configuration.
    Then the two families are proportional at $e$: $X'_e = c\cdot X_e$ for a
    nonzero scalar $c$.
    This is the uniqueness clause of \cite[Theorem~3]{Molnar2018NormalPEPS}
    read at one edge, at the source's sizes $n,m\geq 7$ as in
    Theorem~\ref{thm:fundamentalTheorem_normalTorusPEPS_unconditional}; the
    source proves it through the determinacy of the conjugator in the
    isomorphism lemma, whose realizing matrix is unique up to a multiplicative
    constant.
    \begin{proof}
        \leanok
        \uses{lem:peps_injectiveRegion_union_pos}
        The rectangle hypotheses produce, along the orientation class of $e$,
        a comparison region having $e$ as its single boundary edge on which
        the second tensor and its complement are injective; the union closure
        of injective regions follows from the positive bond dimensions and
        Lemma~\ref{lem:peps_injectiveRegion_union_pos}.  The bare-edge
        absorbed equality of each family lifts to the region
        inserted-coefficient identity over this region, with the inserted
        matrix conjugated by the orientation-adapted absorption of the
        family's matrix at $e$.  The region-insertion transfer map is
        determined by this identity, so both orientation-adapted absorbing
        matrices induce the same conjugation on every bond matrix:
        \[
            X_{\mathrm{abs}}\,N\,X_{\mathrm{abs}}^{-1}
            =
            X'_{\mathrm{abs}}\,N\,\bigl(X'_{\mathrm{abs}}\bigr)^{-1}
            \qquad\text{for every bond matrix }N.
        \]
        A matrix commuting with the full bond matrix algebra is a scalar
        multiple of the identity, so $X'_{\mathrm{abs}}=c\,X_{\mathrm{abs}}$
        for some $c\ne0$. Unwinding the orientation adaptation gives
        $X'_e=c\cdot X_e$.
    \end{proof}
\end{theorem}

\begin{theorem}[Class constancy of the gauge comparison scalars]%
    \label{thm:torusCovariantAbsorbedGauge_unique_classScalar}
    \lean{TNLean.PEPS.torusCovariantAbsorbedGauge_unique_classScalar}
    \leanok
    \uses{def:applyGauge, def:peps_edgeInsertedCoeff}
    Let $A$, $B$, $X$, $X'$ be as in
    Theorem~\ref{thm:torusAbsorbedGauge_unique_scalar}, with both families
    realizing the absorbed inserted-coefficient equality at every edge, and
    suppose moreover that both families are translation covariant: the matrix
    at any translate of an edge is the matrix at the edge, transposed-inverted
    exactly when the translation swaps the stored endpoint order.  Then there
    are nonzero scalars $c_h$ and $c_v$ such that on every horizontal edge
    stored in the natural cyclic order (the stored second endpoint is the
    cyclic right neighbour of the first) $X'_e = c_h\cdot X_e$, on every
    horizontal edge wrapping the seam (stored with swapped endpoints)
    $X'_e = c_h^{-1}\cdot X_e$, and likewise on the vertical class with $c_v$.
    The inversion on the seam is forced by the ordered-edge convention: the
    covariance carries the transposed inverse of the class matrix onto a seam
    edge, and the transposed inverse of a proportionality inverts its
    constant.  A single constant on the whole class is therefore not the true
    statement in this convention; this is the class form of the uniqueness
    clause of \cite[Theorem~3]{Molnar2018NormalPEPS}.
    \begin{proof}
        \leanok
        \uses{thm:torusAbsorbedGauge_unique_scalar}
        The constants are the per-edge scalars of
        Theorem~\ref{thm:torusAbsorbedGauge_unique_scalar} at the two
        reference edges.  Every edge of a class is the translate of its
        reference edge by a unique translation $\tau$. At an interior edge,
        covariance gives $X_e=\tau\cdot X_{\mathrm{ref}}$, hence
        \[
            X'_e
            =
            \tau\cdot X'_{\mathrm{ref}}
            =
            \tau\cdot(c_h\,X_{\mathrm{ref}})
            =
            c_h\,(\tau\cdot X_{\mathrm{ref}})
            =
            c_h\,X_e.
        \]
        At a seam edge, covariance gives $X_e=(X_{\mathrm{ref}})^{-\top}$,
        hence
        \[
            X'_e
            =
            (X'_{\mathrm{ref}})^{-\top}
            =
            (c_h\,X_{\mathrm{ref}})^{-\top}
            =
            c_h^{-1}(X_{\mathrm{ref}})^{-\top}
            =
            c_h^{-1}X_e.
        \]
        The marker $e_{1,x}+1=e_{2,x}$ distinguishes the two branches, and the
        vertical class is identical with $c_v$ in place of $c_h$.
    \end{proof}
\end{theorem}

\begin{theorem}[Uniqueness of the gauges against the per-vertex relation]%
    \label{thm:torusGauge_unique_scalar_of_perVertex}
    \lean{TNLean.PEPS.torusGauge_unique_scalar_of_perVertex}
    \leanok
    \uses{def:gaugeVertex, def:applyGauge}
    Let $A$ and $B$ be translationally invariant PEPS tensors on the discrete
    $n\times m$ torus with $n,m\geq 7$, with positive physical dimension,
    equal and positive bond dimensions, generating the same state, and such
    that every contiguous
    $2\times 3$ and $3\times 2$ coordinate rectangle of either tensor is
    injective.  Suppose the families $X$ and $X'$ of invertible matrices, one
    on each edge, each realize the per-vertex relation of
    Theorem~\ref{thm:fundamentalTheorem_normalTorusPEPS_unconditional}: at
    every site the tensor $A$ equals $\lambda$ (respectively $\lambda'$)
    times the gauge action of the family on $B$, with $\lambda^{nm} = 1$
    (respectively $\lambda'^{nm} = 1$).  Then the two families are
    proportional at every edge: $X'_e = c\cdot X_e$ for a nonzero scalar $c$
    depending on the edge.
    This is the uniqueness clause of \cite[Theorem~3]{Molnar2018NormalPEPS}
    read against the displayed relation $B = \lambda\cdot(X,Y)A$ itself, at
    the source's sizes $n,m\geq 7$.
    \begin{proof}
        \leanok
        \uses{thm:torusAbsorbedGauge_unique_scalar,
            def:peps_edgeInsertedCoeff}
        The inserted-edge coefficient is multilinear in the site tensors, one
        factor per site, so the per-vertex relation scales every inserted
        coefficient by $\lambda^{nm} = 1$: each family realizes the absorbed
        inserted-coefficient equality at every edge.
        Theorem~\ref{thm:torusAbsorbedGauge_unique_scalar} then gives the
        proportionality at each edge.
    \end{proof}
\end{theorem}

\subsection{Overlapping-window corner cancellation}
\label{subsec:peps_torus_overlap_corner_cancellation}

\begin{definition}[Corner extension of a region insert]
    \label{def:peps_corner_extension}
    \lean{TNLean.PEPS.nestedThreeBlockGeometry}
    \lean{TNLean.PEPS.bareExtendInsert}
    \lean{TNLean.PEPS.extendInsert}
    \leanok
    Let $R\subseteq S$ be regions in a finite graph, and let $C$ be a region
    insert on $R$.  Write $\Phi_{R,S}$ for the coefficient obtained by
    contracting the added block $S\setminus R$ between the boundary of $R$ and
    the boundary of $S$.  The bare corner extension is
    \[
        \overline{\Ext}_{R\subseteq S}(C)_{\nu,\sigma}
        =
        \sum_{\mu\in \partial R}
        C_{\mu,\sigma|_R}\,
        \Phi_{R,S}\bigl(\mu,\sigma|_{S\setminus R},\nu\bigr).
    \]
    The normalized corner extension of $C$ from $R$ to $S$ is the region insert
    \[
        \Ext_{R\subseteq S}(C)_{\nu,\sigma}
        =
        P_{V\setminus S}^{-1}
        \sum_{\mu\in \partial R}
        C_{\mu,\sigma|_R}\,
        \Phi_{R,S}\bigl(\mu,\sigma|_{S\setminus R},\nu\bigr),
    \]
    where $P_{V\setminus S}$ is the product of bond dimensions over the edges
    not crossing the boundary of $V\setminus S$.  This is the open-boundary
    form of the extension used in the overlapping-window proof sketch of
    \cite[Section~4, proof of Theorem~3]{Molnar2018NormalPEPS}.
\end{definition}

\begin{theorem}[State preservation under corner extension]
    \label{thm:peps_corner_extension_state_preserves}
    \lean{TNLean.PEPS.deformedRegionState_extend}
    \leanok
    \uses{def:peps_corner_extension}
    If all bond dimensions are positive, then corner extension from
    $R\subseteq S$ preserves the deformed region state:
    \[
        \Psi_{A,R,C}
        =
        \Psi_{A,S,\Ext_{R\subseteq S}(C)}.
    \]
\end{theorem}
\begin{proof}
    \leanok
    Expanding the deformed state on $S$ separates the added block
    $S\setminus R$ from the complement $V\setminus S$.  The corner-extension
    sum contracts the added block first, giving the pointwise identity
    \[
        \Psi_{A,S,\Ext_{R\subseteq S}(C)}(\tau)
        =
        P_{V\setminus S}^{-1}\,P_{V\setminus S}\,
        \Psi_{A,R,C}(\tau)
        =
        \Psi_{A,R,C}(\tau).
    \]
\end{proof}

\begin{theorem}[Consecutive horizontal windows have equal corner extensions]
    \label{thm:peps_horizontal_consecutive_window_extend_eq}
    \lean{TNLean.PEPS.horizontalConsecutiveWindow_extend_eq}
    \leanok
    \uses{def:peps_corner_extension,
        thm:peps_corner_extension_state_preserves}
    Assume the one-orientation $L\times K$ window-injectivity hypotheses,
    union closure for injective regions, positive bond dimensions, and the
    bounds
    \[
        2\le L,\qquad 2\le K,\qquad 2L+1\le W,\qquad 2K+1\le H.
    \]
    Suppose two horizontally consecutive $L\times K$ windows carry inserts
    $C_1$ and $C_2$ whose deformed states agree.  Let $U$ be their
    $(L+1)\times K$ union.  Then their extensions to $U$ agree:
    \[
        \Ext_{W_1\subseteq U}(C_1)
        =
        \Ext_{W_2\subseteq U}(C_2).
    \]
\end{theorem}
\begin{proof}
    \leanok
    Theorem~\ref{thm:peps_corner_extension_state_preserves} rewrites each
    single-window deformed state as the deformed state on $U$ with the
    corresponding extended insert:
    \[
        \Psi_{A,U,\Ext_{W_1\subseteq U}(C_1)}
        =
        \Psi_{A,W_1,C_1}
        =
        \Psi_{A,W_2,C_2}
        =
        \Psi_{A,U,\Ext_{W_2\subseteq U}(C_2)}.
    \]
    Since $U$ is injective, equality of the two $U$-states implies equality
    of the two inserts:
    \[
        \Psi_{A,U,\Ext_{W_1\subseteq U}(C_1)}
        =
        \Psi_{A,U,\Ext_{W_2\subseteq U}(C_2)}
        \quad\Longrightarrow\quad
        \Ext_{W_1\subseteq U}(C_1)
        =
        \Ext_{W_2\subseteq U}(C_2).
    \]
\end{proof}

\paragraph{Uniform interior-bond products for staircase windows.}

For a PEPS tensor $A$ and a region $R$, write
\[
    P_A(R)=\prod_{e\notin\partial R} D_A(e)
\]
for the product of bond dimensions over the non-boundary edges of $R$.

\begin{theorem}[Interior-bond product under a graph isomorphism]%
    \label{thm:peps_regionInteriorBondProd_map}
    \lean{TNLean.PEPS.regionInteriorBondProd_map}
    \leanok
    Let $\phi:G\simeq G'$ be an isomorphism of finite graphs, and let $R$ be
    a region of $G$.  If $A'$ is a PEPS tensor on $G'$, then
    \[
        P_{A'}(\phi(R))
        =
        \prod_{e\notin\partial R} D_{A'}(\phi(e)).
    \]
\end{theorem}
\begin{proof}
    \leanok
    An edge of $\phi(R)$ is non-boundary exactly when its preimage is
    non-boundary for $R$:
    \[
        e'\notin\partial\phi(R)
        \Longleftrightarrow
        \phi^{-1}(e')\notin\partial R.
    \]
    Reindexing the product over non-boundary edges by the edge bijection
    $e\mapsto\phi(e)$ gives the displayed equality.
\end{proof}

\begin{theorem}[Translation invariance of the interior-bond product]%
    \label{thm:peps_regionInteriorBondProd_translate}
    \lean{TNLean.PEPS.regionInteriorBondProd_translate_of_translationInvariant}
    \leanok
    \uses{def:peps_torus_translation_invariant,
        thm:peps_regionInteriorBondProd_map}
    Let $A$ be translation invariant on a torus whose periods $W,H$ satisfy
    $1<W$ and $1<H$.  For every translation $\tau_{a,b}$ and every region $R$,
    \[
        P_A(\tau_{a,b}(R))=P_A(R).
    \]
\end{theorem}
\begin{proof}
    \leanok
    Theorem~\ref{thm:peps_regionInteriorBondProd_map} gives
    \[
        P_A(\tau_{a,b}(R))
        =
        \prod_{e\notin\partial R}D_A(\tau_{a,b}(e)).
    \]
    Translation invariance gives $D_A(\tau_{a,b}(e))=D_A(e)$ for every
    edge $e$, so the product is $P_A(R)$.
\end{proof}

\begin{theorem}[Interior-bond product of a cyclic rectangle is start-independent]%
    \label{thm:peps_regionInteriorBondProd_torusArcRectangle_eq}
    \lean{TNLean.PEPS.regionInteriorBondProd_torusArcRectangle_eq}
    \leanok
    \uses{thm:peps_regionInteriorBondProd_translate}
    Let $A$ be translation invariant on a torus whose periods $W,H$ satisfy
    $1<W$ and $1<H$.  If $Q(s;x,y)$ denotes the cyclic $x\times y$ rectangle
    starting at $s$, then for any starts $s,s'$,
    \[
        P_A(Q(s;x,y))=P_A(Q(s';x,y)).
    \]
\end{theorem}
\begin{proof}
    \leanok
    If $s=(s_x,s_y)$ and $s'=(s'_x,s'_y)$, then the translation by
    $(s_x-s'_x,s_y-s'_y)$ carries one rectangle to the other:
    \[
        \tau_{s_x-s'_x,s_y-s'_y}(Q(s';x,y))=Q(s;x,y).
    \]
    Apply Theorem~\ref{thm:peps_regionInteriorBondProd_translate}.
\end{proof}

\begin{theorem}[Uniform interior-bond product of the staircase windows]%
    \label{thm:peps_regionInteriorBondProd_staircaseWindow_eq}
    \lean{TNLean.PEPS.regionInteriorBondProd_staircaseWindow_eq}
    \leanok
    \uses{thm:peps_regionInteriorBondProd_torusArcRectangle_eq}
    Let $A$ be translation invariant on a torus whose periods $W,H$ satisfy
    $1<W$ and $1<H$.  If $W_j$ and $W_{j'}$ are two windows of the staircase
    based at the same corner $s$ with window dimensions $L\times K$, then
    \[
        P_A(W_j)=P_A(W_{j'}).
    \]
\end{theorem}
\begin{proof}
    \leanok
    Each staircase window is an $L\times K$ cyclic rectangle:
    \[
        W_j=Q(s_j;L,K),
        \qquad
        W_{j'}=Q(s_{j'};L,K).
    \]
    Theorem~\ref{thm:peps_regionInteriorBondProd_torusArcRectangle_eq}
    gives $P_A(Q(s_j;L,K))=P_A(Q(s_{j'};L,K))$.
\end{proof}

\begin{theorem}[Uniform horizontal per-step bond transport]%
    \label{thm:peps_horizontal_staircase_consecutive_window_bond_transport_uniform}
    \lean{TNLean.PEPS.horizontalStaircaseConsecutiveWindow_bondTransport_extend_eq_uniform}
    \leanok
    \uses{def:peps_corner_extension,
        thm:peps_regionInteriorBondProd_staircaseWindow_eq}
    Let $B$ be translation invariant on a torus whose periods $W,H$ satisfy
    $1<W$ and $1<H$.  Assume the one-orientation $L\times K$
    window-injectivity hypotheses, union closure for injective regions,
    positive bond dimensions, and the bounds
    \[
        2\le L,\qquad 2\le K,\qquad 2L+1\le W,\qquad 2K+1\le H.
    \]
    Let $j<L$, let $U_j=W_j\cup W_{j+1}$ be the horizontal consecutive-window
    union, and let $g$ be a boundary edge of both $W_j$ and $W_{j+1}$.  If the
    matrices obtained from $M_1$ and $M_2$ by the orientation convention at
    $g$ agree, then, with $P_j=P_B(W_j)$,
    \[
        \operatorname{Ext}_{W_j\subseteq U_j}
            \bigl(P_j\,I_{B,W_j,g,M_1}\bigr)
        =
        \operatorname{Ext}_{W_{j+1}\subseteq U_j}
            \bigl(P_j\,I_{B,W_{j+1},g,M_2}\bigr).
    \]
\end{theorem}
\begin{proof}
    \leanok
    The consecutive-window transport gives the cross-rescaled equality
    \[
        \operatorname{Ext}_{W_j\subseteq U_j}
            \bigl(P_B(W_{j+1})\,I_{B,W_j,g,M_1}\bigr)
        =
        \operatorname{Ext}_{W_{j+1}\subseteq U_j}
            \bigl(P_B(W_j)\,I_{B,W_{j+1},g,M_2}\bigr).
    \]
    By Theorem~\ref{thm:peps_regionInteriorBondProd_staircaseWindow_eq},
    \[
        P_B(W_{j+1})=P_B(W_j)=P_j.
    \]
    Substituting this equality gives the displayed common-scalar form.
\end{proof}

\begin{theorem}[Uniform vertical per-step bond transport]%
    \label{thm:peps_vertical_staircase_consecutive_window_bond_transport_uniform}
    \lean{TNLean.PEPS.verticalStaircaseConsecutiveWindow_bondTransport_extend_eq_uniform}
    \leanok
    \uses{def:peps_corner_extension,
        thm:peps_regionInteriorBondProd_staircaseWindow_eq}
    Let $B$ be translation invariant on a torus whose periods $W,H$ satisfy
    $1<W$ and $1<H$.  Assume the one-orientation $L\times K$
    window-injectivity hypotheses, union closure for injective regions,
    positive bond dimensions, and the bounds
    \[
        2\le L,\qquad 2\le K,\qquad 2L+1\le W,\qquad 2K+1\le H.
    \]
    Let $L\le j$ and $j+1<L+K$, let $U_j=W_j\cup W_{j+1}$ be the vertical
    consecutive-window union, and let $g$ be a boundary edge of both $W_j$ and
    $W_{j+1}$.  If the matrices obtained from $M_1$ and $M_2$ by the
    orientation convention at $g$ agree, then, with $P_j=P_B(W_j)$,
    \[
        \operatorname{Ext}_{W_j\subseteq U_j}
            \bigl(P_j\,I_{B,W_j,g,M_1}\bigr)
        =
        \operatorname{Ext}_{W_{j+1}\subseteq U_j}
            \bigl(P_j\,I_{B,W_{j+1},g,M_2}\bigr).
    \]
\end{theorem}
\begin{proof}
    \leanok
    The descending-arm transport gives the same cross-rescaled equality as in
    the horizontal case, with $U_j$ now the $L\times(K+1)$ union:
    \[
        \operatorname{Ext}_{W_j\subseteq U_j}
            \bigl(P_B(W_{j+1})\,I_{B,W_j,g,M_1}\bigr)
        =
        \operatorname{Ext}_{W_{j+1}\subseteq U_j}
            \bigl(P_B(W_j)\,I_{B,W_{j+1},g,M_2}\bigr).
    \]
    The uniformity theorem gives $P_B(W_{j+1})=P_B(W_j)=P_j$.  Substituting
    into the displayed cross-rescaled form gives
    \[
        \operatorname{Ext}_{W_j\subseteq U_j}
            \bigl(P_j\,I_{B,W_j,g,M_1}\bigr)
        =
        \operatorname{Ext}_{W_{j+1}\subseteq U_j}
            \bigl(P_j\,I_{B,W_{j+1},g,M_2}\bigr).
    \]
\end{proof}


\begin{theorem}[Patch-extended equality of the staircase end windows]
    \label{thm:peps_horizontal_staircase_patch_extend_eq}
    \lean{TNLean.PEPS.horizontalStaircase_patch_extend_eq}
    \leanok
    \uses{def:peps_horizontal_staircase_patch,
        thm:peps_corner_extension_state_preserves}
    Assume positive bond dimensions and the bounds
    \[
        0<L,\qquad 0<K,\qquad 2L\le W,\qquad 2K\le H.
    \]
    If the deformed states of the two end-window inserts agree, then their
    extensions to the staircase patch $P$ have equal deformed states:
    \[
        \Psi_{A,P,\Ext_{W_{\mathrm{right}}\subseteq P}(C_0)}
        =
        \Psi_{A,P,\Ext_{W_{\mathrm{left}}\subseteq P}(C_{\mathrm{end}})}.
    \]
\end{theorem}
\begin{proof}
    \leanok
    Applying Theorem~\ref{thm:peps_corner_extension_state_preserves} to each
    end window gives
    \[
        \Psi_{A,P,\Ext_{W_{\mathrm{right}}\subseteq P}(C_0)}
        =
        \Psi_{A,W_{\mathrm{right}},C_0}
        =
        \Psi_{A,W_{\mathrm{left}},C_{\mathrm{end}}}
        =
        \Psi_{A,P,\Ext_{W_{\mathrm{left}}\subseteq P}(C_{\mathrm{end}})},
    \]
    where the middle equality is the hypothesis.
\end{proof}

% The end-pair stripping under full complement injectivity has been retired: at the
% corollary's minimal size $V\setminus S$ is not blocked-tensor injective, so its hypothesis
% is not derivable.  The faithful Step~3 is the open-boundary shared-corner cancellation
% \ref{thm:peps_staircase_pair_insert_eq_open}, which inverts only the injective completed
% corner $Q$ and never $V\setminus S$.

\begin{lemma}[Two-block merge collapse]
    \label{lem:peps_two_block_merge_collapse}
    \lean{TNLean.PEPS.mergeCollapse2}
    \leanok
    Let $B_1,K_1\subseteq V\setminus T$ and $B_2,H_2\subseteq T$.  Fix
    boundary data $c_1$ on $V\setminus K_1$ and $c_2$ on $H_2$, and physical
    labels $\sigma_1$ on $B_1$ and $\sigma_2$ on $B_2$.  Set
    \[
        F_1(\xi)=\prod_{x\in B_1}A_x(\xi|_{\partial x},\sigma_1(x)),
        \qquad
        F_2(\xi)=\prod_{x\in B_2}A_x(\xi|_{\partial x},\sigma_2(x)).
    \]
    Then the
    boundary-compatible double sum collapses to the multiplicity of the
    interior bonds of $T$ times the single merged sum:
    \[
        \sum_{\substack{
            \zeta_2|_{\partial(V\setminus K_1)}=c_1,\;
            \zeta_1|_{\partial H_2}=c_2\\
            \zeta_1|_{\partial T}=\zeta_2|_{\partial T}}}
        F_1(\zeta_2)F_2(\zeta_1)
        =
        P_T
        \sum_{\substack{
            \eta|_{\partial(V\setminus K_1)}=c_1,\;
            \eta|_{\partial H_2}=c_2}}
        F_1(\eta)F_2(\eta).
    \]
\end{lemma}
\begin{proof}
    \leanok
    Each merged configuration has exactly $P_T$ preimages, indexed by the
    virtual bonds strictly inside $T$.  Since $B_1$ and $K_1$ lie outside
    $T$, while $B_2$ and $H_2$ lie inside $T$, the two products and the two
    boundary conditions are unchanged under the merge.
\end{proof}

\begin{theorem}[Composition of corner-extension coefficient kernels]
    \label{thm:peps_corner_extension_coefficient_composition}
    \lean{TNLean.PEPS.threeBlockBlueCoeff_comp}
    \leanok
    \uses{def:peps_corner_extension, lem:peps_two_block_merge_collapse}
    For nested regions $R\subseteq S\subseteq P$, the coefficient kernels
    satisfy
    \[
        \sum_{\nu\in\partial S}
        \Phi_{R,S}\bigl(\mu,\sigma|_{S\setminus R},\nu\bigr)
        \Phi_{S,P}\bigl(\nu,\sigma|_{P\setminus S},\omega\bigr)
        =
        P_{V\setminus S}\,
        \Phi_{R,P}\bigl(\mu,\sigma|_{P\setminus R},\omega\bigr).
    \]
    This is the two-block merge collapse in the proof of
    \cite[Section~4, proof of Theorem~3]{Molnar2018NormalPEPS}.
\end{theorem}

\begin{theorem}[Linearity and transitivity of corner extension]
    \label{thm:peps_corner_extension_linear_transitive}
    \lean{TNLean.PEPS.bareExtendInsert_const_smul}
    \lean{TNLean.PEPS.extendInsert_const_smul}
    \lean{TNLean.PEPS.extendInsert_add}
    \lean{TNLean.PEPS.extendInsert_sub}
    \lean{TNLean.PEPS.bareExtendInsert_trans}
    \lean{TNLean.PEPS.extendInsert_trans}
    \leanok
    \uses{def:peps_corner_extension,
        thm:peps_corner_extension_coefficient_composition}
    The corner extension is linear in the inserted boundary coefficients.  For
    nested regions $R\subseteq S\subseteq P$, the bare extension satisfies
    \[
        \overline{\Ext}_{S\subseteq P}
        \bigl(\overline{\Ext}_{R\subseteq S}(C)\bigr)
        =
        P_{V\setminus S}\,
        \overline{\Ext}_{R\subseteq P}(C)
    \]
    without a positivity assumption.  If all bond dimensions are positive, the
    normalized corner extension is transitive:
    \[
        \Ext_{S\subseteq P}
        \bigl(\Ext_{R\subseteq S}(C)\bigr)
        =
        \Ext_{R\subseteq P}(C).
    \]
\end{theorem}
\begin{proof}
    \leanok
    The scalar and additive laws follow from the defining sum.  For
    transitivity, substitute
    Theorem~\ref{thm:peps_corner_extension_coefficient_composition} into the
    defining sum.  This gives the bare formula, and the factors
    $P_{V\setminus S}$ cancel in the normalized extension.
\end{proof}

\begin{definition}[Assembly of subregion physical configurations]
    \label{def:peps_assemble_subregion_physical_config}
    \lean{TNLean.PEPS.assembleSubRegionσ}
    \leanok
    If $R\subseteq S$, a physical configuration on $R$ and a physical
    configuration on $S\setminus R$ assemble to a physical configuration on
    $S$:
    \[
        \operatorname{As}_{R,S}(\sigma_R,\sigma_Q)(v)
        =
        \begin{cases}
        \sigma_R(v), & v\in R,\\
        \sigma_Q(v), & v\in S\setminus R.
        \end{cases}
    \]
\end{definition}

\begin{theorem}[Restrictions of assembled subregion configurations]
    \label{thm:peps_assemble_subregion_physical_config_restricts}
    \lean{TNLean.PEPS.restrictSubRegionσ_assembleSubRegionσ}
    \lean{TNLean.PEPS.restrictSubRegionσ_sdiff_assembleSubRegionσ}
    \leanok
    \uses{def:peps_assemble_subregion_physical_config}
    The two restrictions recover the original configurations:
    \[
        \operatorname{As}_{R,S}(\sigma_R,\sigma_Q)|_R=\sigma_R,
        \qquad
        \operatorname{As}_{R,S}(\sigma_R,\sigma_Q)|_{S\setminus R}=\sigma_Q.
    \]
\end{theorem}

\begin{theorem}[Kernel triviality of corner extension from the added block]
    \label{thm:peps_extend_insert_kernel_trivial_added_injective}
    \lean{TNLean.PEPS.extendInsert_kernel_trivial_of_addedInjective}
    \leanok
    \uses{def:peps_corner_extension,
        def:peps_assemble_subregion_physical_config,
        thm:peps_assemble_subregion_physical_config_restricts}
    Let $R\subseteq S$ be finite regions.  Suppose that all bond dimensions are
    positive and that the added block $Q=S\setminus R$ is blocked-tensor
    injective.  If a region insert $C$ on $R$ has zero corner extension,
    \[
        \Ext_{R\subseteq S}(C)=0,
    \]
    then $C=0$.
\end{theorem}
\begin{proof}
    \leanok
    Since all bond dimensions are positive, $P_{V\setminus S}\ne 0$.
    Thus $\Ext_{R\subseteq S}(C)=0$ is equivalent to
    $\overline{\Ext}_{R\subseteq S}(C)=0$.  Fix the physical
    configuration $\sigma_R$ on $R$.  Let
    \[
        \rho:\partial R \simeq \partial(V\setminus R)
        \qquad\text{and}\qquad
        \tau:\partial S \simeq \partial(V\setminus S)
    \]
    be the complement-boundary bijections, and define
    \[
        c_\eta=C_{\rho^{-1}(\eta),\sigma_R}
        \qquad(\eta\in \partial(V\setminus R)).
    \]
    Evaluating the bare zero extension at
    $\operatorname{As}_{R,S}(\sigma_R,\sigma_Q)$ gives, for every added-block
    configuration $\sigma_Q$ and every boundary value
    $\omega\in\partial(V\setminus S)$,
    \[
        \sum_{\eta\in\partial(V\setminus R)}
        c_\eta\,\Phi_{R,S}(\rho^{-1}(\eta),\sigma_Q,\tau^{-1}(\omega))=0.
    \]
    For each outer boundary value $\omega$ and each boundary value $\beta$ of $Q$,
    define
    \[
        f_\omega(\beta)
        =
        \sum_{\eta\in\partial(V\setminus R)}
        c_\eta\,
        \mathbf 1_{\{\exists q:\,
            \partial_{V\setminus R}q=\eta,\,
            \partial_{V\setminus S}q=\omega,\,
            \partial_Qq=\beta\}}.
    \]
    Let $\Gamma$ be the product of bond dimensions on the crossing edges
    separated when the $Q$-block is read first.  Unfolding $\Phi_{R,S}$ as a sum
    over the blocked weights of $Q$ gives, for every $\sigma_Q$,
    \[
    \begin{aligned}
        (T_Q f_\omega)(\sigma_Q)
        &=
        \sum_{\beta\in\partial Q} f_\omega(\beta)\,W_Q(\beta,\sigma_Q) \\
        &=
        \sum_{\beta\in\partial Q}
        \left(
        \sum_{\eta\in\partial(V\setminus R)}
        c_\eta\,
        \mathbf 1_{\{\exists q:\,
            \partial_{V\setminus R}q=\eta,\,
            \partial_{V\setminus S}q=\omega,\,
            \partial_Qq=\beta\}}
        \right) W_Q(\beta,\sigma_Q) \\
        &=
        \Gamma
        \sum_{\eta\in\partial(V\setminus R)}
        c_\eta\,\Phi_{R,S}(\rho^{-1}(\eta),\sigma_Q,\tau^{-1}(\omega))
        =
        0.
    \end{aligned}
    \]
    Hence $T_Q(f_\omega)=0$.  Injectivity of the blocked tensor of $Q$ gives
    \[
        T_Q(f_\omega)=0 \Longrightarrow f_\omega=0.
    \]
    Now fix $\mu\in\partial R$.  Positivity of bond dimensions gives a global
    virtual configuration $q$ whose residual boundary on $V\setminus R$ is
    $\rho(\mu)$.  Set
    \[
        \omega=\partial_{V\setminus S}q,\qquad
        \beta=\partial_Q q.
    \]
    Evaluating $f_\omega=0$ at $\beta$, the residual-boundary uniqueness identity
    leaves only the term $\eta=\rho(\mu)$:
    \[
        0
        =
        f_\omega(\beta)
        =
        c_{\rho(\mu)}\cdot 1
        =
        C_{\mu,\sigma_R}.
    \]
    Since $\sigma_R$ and $\mu$ were arbitrary, $C=0$.
\end{proof}

\begin{theorem}[Cancellation of a shared injective corner]
    \label{thm:peps_extend_insert_cancel_added_injective}
    \lean{TNLean.PEPS.extendInsert_cancel_addedInjective}
    \leanok
    \uses{thm:peps_corner_extension_linear_transitive,
        thm:peps_extend_insert_kernel_trivial_added_injective}
    Let $R\subseteq S$, and suppose that all bond dimensions are positive and
    that $S\setminus R$ is blocked-tensor injective.  If two inserts on $R$
    have the same corner extension to $S$,
    \[
        \Ext_{R\subseteq S}(C_1)
        =
        \Ext_{R\subseteq S}(C_2),
    \]
    then $C_1=C_2$.
\end{theorem}
\begin{proof}
    \leanok
    By linearity,
    \[
        \Ext_{R\subseteq S}(C_1-C_2)=0.
    \]
    Theorem~\ref{thm:peps_extend_insert_kernel_trivial_added_injective}
    applied to the added block $S\setminus R$ gives $C_1-C_2=0$, hence
    $C_1=C_2$.
\end{proof}

\begin{theorem}[Staircase-pair cancellation on the torus]
    \label{thm:peps_staircase_pair_cancel}
    \lean{TNLean.PEPS.staircasePair_cancel}
    \leanok
    \uses{thm:peps_extend_insert_cancel_added_injective}
    Let the torus have width $W$ and height $H$, with $1<W$ and $1<H$.
    Let $S$ be the staircase end pair around a horizontal edge, with seam
    dimensions $L$ and $K$, and let $Q$ be the completed corner rectangle.
    Assume the size bounds
    \[
        2L\le W,\qquad 2K\le H,
    \]
    that $Q$ is blocked-tensor injective, and that all bond dimensions are
    positive.
    If two inserts $X$ and $Y$ on $S$ have equal extensions to $S\cup Q$,
    \[
        \Ext_{S\subseteq S\cup Q}(X)
        =
        \Ext_{S\subseteq S\cup Q}(Y),
    \]
    then $X=Y$.  Thus the injective completed corner can be cancelled from an
    open-boundary equality, without assuming injectivity of the torus
    complement of $S$.
\end{theorem}
\begin{proof}
    \leanok
    The size bounds imply that the staircase end pair is disjoint from the
    completed corner, and therefore
    \[
        (S\cup Q)\setminus S=Q.
    \]
    The claimed equality is exactly
    Theorem~\ref{thm:peps_extend_insert_cancel_added_injective} applied to the
    inclusion $S\subseteq S\cup Q$.
\end{proof}

\subsection{Single-bond peeling and the window witness}
\label{subsec:peps_torus_single_bond_peeling_window_witness}

\begin{definition}[Reference edge of the horizontal staircase]
    \label{def:peps_horizontal_staircase_reference_edge}
    \lean{TNLean.PEPS.horizontalStaircaseReferenceEdge}
    \leanok
    \uses{def:peps_horizontal_staircase_patch}
    In the staircase coordinates $s=(a,b)$, the reference edge is the
    horizontal edge joining $(a+L-1,b+K-1)$ to $(a+L,b+K-1)$.
    It is the edge between the left and right end windows.
\end{definition}

\begin{theorem}[The reference edge is the unique crossing of the two end windows]
    \label{thm:peps_horizontal_staircase_reference_edge_crossing}
    \lean{TNLean.PEPS.isCrossingEdge_horizontalStaircaseEndWindows}
    \lean{TNLean.PEPS.isRegionBoundaryEdge_horizontalStaircaseLeftWindow_referenceEdge}
    \lean{TNLean.PEPS.isRegionBoundaryEdge_horizontalStaircaseRightWindow_referenceEdge}
    \lean{TNLean.PEPS.not_isRegionBoundaryEdge_horizontalStaircaseEndPair_referenceEdge}
    \leanok
    \uses{def:peps_horizontal_staircase_reference_edge, def:peps_horizontal_staircase_patch}
    Assume
    $1<W$, $1<H$, $0<L$, $0<K$, $1\le a$, $a+2L\le W$, and
    $b+2K-1\le H$.
    Then an edge crosses from the left end window to the right end window if
    and only if it is the reference edge $e$:
    \begin{align}
        f\in\partial W_{\mathrm{left}}\cap\partial W_{\mathrm{right}}
        \quad\Longleftrightarrow\quad
        f=e.
        \notag
    \end{align}
    Thus $e$ is a boundary edge of each end window, but both endpoints of $e$
    lie in the end pair $S$, so $e\notin\partial S$.
\end{theorem}
\begin{proof}
    \leanok
    The endpoint calculation identifies $e$ with the unique horizontal edge
    joining the two displayed rectangles. Since one endpoint lies in
    $W_{\mathrm{left}}$ and the other lies in $W_{\mathrm{right}}$, the same
    edge is internal to their union $S$.
\end{proof}

\begin{lemma}[Boundary edges of the end pair come from the end windows]
    \label{lem:peps_horizontal_staircase_end_pair_boundary_from_windows}
    \lean{TNLean.PEPS.isRegionBoundaryEdge_endPair_of_leftWindow}
    \lean{TNLean.PEPS.isRegionBoundaryEdge_endPair_of_rightWindow}
    \lean{TNLean.PEPS.isRegionBoundaryEdge_window_of_endPair}
    \lean{TNLean.PEPS.ne_referenceEdge_of_isRegionBoundaryEdge_endPair}
    \leanok
    \uses{thm:peps_horizontal_staircase_reference_edge_crossing}
    Under the bounds
    $1<W$, $1<H$, $0<L$, $0<K$, $1\le a$, $a+2L\le W$, and
    $b+2K-1\le H$,
    a boundary edge of the end pair is a boundary edge of one of the two end
    windows, and it is not the reference edge. Conversely, any boundary edge
    of one end window other than $e$ is a boundary edge of $S$.
\end{lemma}
\begin{proof}
    \leanok
    A boundary edge of $S$ has one endpoint in $S$; that endpoint lies in one
    of the two end windows. The other endpoint is outside $S$, hence outside
    that window. Conversely, a boundary edge of one window which is not $e$
    cannot enter the other window, because $e$ is the unique crossing edge.
\end{proof}

\begin{theorem}[Boundary of the staircase end pair]
    \label{thm:peps_horizontal_staircase_end_pair_boundary}
    \lean{TNLean.PEPS.eq_referenceEdge_of_isRegionBoundaryEdge_both_endWindows}
    \lean{TNLean.PEPS.isRegionBoundaryEdge_endPair_iff}
    \lean{TNLean.PEPS.isRegionBoundaryEdge_endPair_exactlyOne_window}
    \leanok
    \uses{lem:peps_horizontal_staircase_end_pair_boundary_from_windows}
    Under the bounds
    $1<W$, $1<H$, $0<L$, $0<K$, $1\le a$, $a+2L\le W$, and
    $b+2K-1\le H$,
    an edge is a boundary edge of the staircase end pair $S$ if and only if it
    is a boundary edge of one of the two end windows and is not the reference
    edge:
    \begin{align}
        f\in\partial S
        \quad\Longleftrightarrow\quad
        \bigl(f\in\partial W_{\mathrm{left}}
        \quad\text{or}\quad
        f\in\partial W_{\mathrm{right}}\bigr)
        \ \text{and}\ f\ne e.
        \notag
    \end{align}
    Moreover each edge in $\partial S$ lies on exactly one of the two window
    boundaries.
\end{theorem}
\begin{proof}
    \leanok
    Lemma~\ref{lem:peps_horizontal_staircase_end_pair_boundary_from_windows}
    gives both implications. If an edge lay on both end-window boundaries,
    the unique-crossing theorem would identify it with $e$, which is not a
    boundary edge of $S$.
\end{proof}

\begin{theorem}[Boundary configurations of the staircase end windows]
    \label{thm:peps_horizontal_staircase_end_window_boundary_config}
    \lean{TNLean.PEPS.horizontalStaircaseLeftWindowBoundaryConfigEquiv}
    \lean{TNLean.PEPS.horizontalStaircaseRightWindowBoundaryConfigEquiv}
    \lean{TNLean.PEPS.horizontalStaircaseEndPairBoundaryConfigEquiv}
    \leanok
    \uses{thm:peps_horizontal_staircase_end_pair_boundary, thm:peps_horizontal_staircase_reference_edge_crossing}
    Under the same geometric bounds, write
    $\partial_{\mathrm{left}}S$ and $\partial_{\mathrm{right}}S$
    for the boundary edges of $S$ which lie on the corresponding end-window
    boundary.  Then the boundary-configuration spaces have canonical
    factorizations
    \begin{align}
        \operatorname{Conf}(\partial W_{\mathrm{left}})
          &\cong \operatorname{Conf}(e)
            \times\operatorname{Conf}(\partial_{\mathrm{left}}S),
        \notag \\
        \operatorname{Conf}(\partial W_{\mathrm{right}})
          &\cong \operatorname{Conf}(e)
            \times\operatorname{Conf}(\partial_{\mathrm{right}}S),
        \notag \\
        \operatorname{Conf}(\partial S)
          &\cong
            \operatorname{Conf}(\partial_{\mathrm{left}}S)
              \times\operatorname{Conf}(\partial_{\mathrm{right}}S).
        \notag
    \end{align}
    These equivalences preserve every underlying lattice-edge label.  Thus
    the reference edge is the only contracted virtual index shared by the
    two highlighted end windows; every other virtual index remains on the
    corresponding open boundary of $S$.
\end{theorem}
\begin{proof}
    \leanok
    Split each end-window configuration at the reference edge.  A remaining
    boundary edge of either window is a boundary edge of $S$, and conversely
    every boundary edge of $S$ belongs to exactly one end window.  Reindexing
    the dependent products along these edge identifications gives the three
    displayed equivalences.
\end{proof}

\begin{lemma}[Complement split for an end window]
    \label{lem:peps_horizontal_staircase_end_window_complement_split}
    \lean{TNLean.PEPS.horizontalStaircaseEndPair_sdiff_leftWindow}
    \lean{TNLean.PEPS.horizontalStaircaseEndPair_sdiff_rightWindow}
    \lean{TNLean.PEPS.compl_leftWindow_eq_rightWindow_union_complEndPair}
    \lean{TNLean.PEPS.compl_rightWindow_eq_leftWindow_union_complEndPair}
    \leanok
    \uses{def:peps_horizontal_staircase_patch, lem:peps_horizontal_staircase_end_pair_disjoint}
    If $2L\le W$, the complement of either end window splits as the opposite
    end window and the complement of the end pair:
    \begin{align}
        V\setminus W_{\mathrm{left}}
         & =W_{\mathrm{right}}\cup (V\setminus S),
        \notag                                     \\
        V\setminus W_{\mathrm{right}}
         & =W_{\mathrm{left}}\cup (V\setminus S).
        \notag
    \end{align}
\end{lemma}
\begin{proof}
    \leanok
    Since the two end windows are disjoint and
    $S=W_{\mathrm{left}}\cup W_{\mathrm{right}}$, removing one window from
    $S$ leaves the other. Splitting $V\setminus W$ through the inclusion
    $W\subseteq S$ gives the displayed decompositions.
\end{proof}

\begin{theorem}[The reference edge bounds the two end windows]
    \label{thm:peps_staircase_end_reference_edge_boundary}
    \lean{TNLean.PEPS.isRegionBoundaryEdge_staircaseWindow_zero_referenceEdge}
    \lean{TNLean.PEPS.isRegionBoundaryEdge_staircaseWindow_last_referenceEdge}
    \leanok
    Let the torus have width $W$ and height $H$, with $1<H$.
    Let $e$ be the horizontal reference edge of a staircase window family with
    dimensions $L,K$ and lower-left coordinate $(a,b)$. Suppose
    $0<L$, $0<K$, $1\le a$, $a+2L\le W$, and $b+2K-1\le H$.
    Then $e$ is a boundary edge of the first staircase window $W_0$ and of the
    last staircase window $W_{L+K-1}$:
    \begin{align}
        e & \in\partial W_0,
        \notag                     \\
        e & \in\partial W_{L+K-1}.
        \notag
    \end{align}
\end{theorem}
\begin{proof}
    \leanok
    The first staircase window is the right end window, and the last
    staircase window is the left end window:
    \begin{align}
        W_0       & =W_{\mathrm{right}},
        \notag                           \\
        W_{L+K-1} & =W_{\mathrm{left}}.
        \notag
    \end{align}
    The reference edge is a boundary edge of these two end windows by their
    defining geometry.
\end{proof}

\begin{theorem}[Open-boundary equality on the staircase end pair]
    \label{thm:peps_staircase_pair_insert_eq_open}
    \lean{TNLean.PEPS.staircasePair_insert_eq_open}
    \leanok
    \uses{thm:peps_staircase_pair_cancel, thm:peps_corner_extension_linear_transitive}
    Let $S=W_0\cup W_{L+K-1}$ be the staircase end pair. Suppose that the
    one-orientation $L\times K$ window injectivity hypotheses hold for a
    tensor $B$, unions of injective regions are injective, all bond dimensions
    of $B$ are positive, and
    $2\le L$, $2\le K$, $2L+1\le W$, and $2K+1\le H$.
    If $C_j$ are inserts on the staircase windows whose deformed states all
    equal one common state, then the two end-window inserts have equal corner
    extensions to $S$:
    \begin{align}
        \Ext_{W_0\subseteq S}(C_0)
        =
        \Ext_{W_{L+K-1}\subseteq S}(C_{L+K-1}).
        \notag
    \end{align}
\end{theorem}
\begin{proof}
    \leanok
    The consecutive-window equalities are composed across the staircase patch.
    Both sides are then extended through the completed corner, and
    transitivity of corner extension rewrites the two extension chains. The
    completed corner is an injective $L\times K$ window, so
    Theorem~\ref{thm:peps_staircase_pair_cancel} cancels it and leaves the
    equality on the end pair $S$.
\end{proof}

\begin{theorem}[Lemma 5 comparison of the two end blocks]
    \label{thm:peps_existsUnique_virtualOperation_of_endPair}
    \lean{TNLean.PEPS.existsUnique_virtualOperation_of_endPair}
    \leanok
    Let $K$ be the basis of the virtual space on the bond between two
    injective blocks, and let
    $\{A_{1,\mu}\}_{\mu\in K}$ and
    $\{A_{2,\mu}\}_{\mu\in K}$ be linearly independent families.  Suppose
    that the two end physical operations $O_1$ and $O_3$ produce families
    $C_1$ and $C_2$ satisfying the open contraction identity
    \begin{align}
        \sum_{\mu\in K} C_{1,\mu}\otimes A_{2,\mu}
        =
        \sum_{\nu\in K} A_{1,\nu}\otimes C_{2,\nu}.
        \notag
    \end{align}
    Then there is a unique matrix $X\in\MN{|K|}$ such that
    \begin{align}
        C_{1,\nu}
        &=\sum_{\mu\in K} A_{1,\mu}X_{\mu,\nu},
        &
        C_{2,\mu}
        &=\sum_{\nu\in K} X_{\mu,\nu}A_{2,\nu}.
        \notag
    \end{align}
    Thus the two end operations are obtained by inserting the same virtual
    operation $X$ on their common bond.  No injectivity is required of the
    modified families $C_1$ and $C_2$.  This is the ``compare the two ends''
    step in Lemma~5 of \cite{Molnar2018NormalPEPS}.
\end{theorem}
\begin{proof}
    \leanok
    Apply a dual functional to the linearly independent family
    $\{A_{2,\mu}\}$ in the open contraction identity.  This expresses every
    $C_{1,\nu}$ as a linear combination of the $A_{1,\mu}$ with a coefficient
    matrix $X$.  Substitution into the same identity and linear independence
    of $\{A_{1,\mu}\}$ gives the second displayed relation with the same
    matrix.  Linear independence of the first genuine family also makes the
    coefficient matrix unique.
\end{proof}

\begin{theorem}[Single-bond peeling of the staircase end equality]
    \label{thm:peps_regionInsertedCoeff_endWindows_eq_of_staircase}
    \lean{TNLean.PEPS.regionInsertedCoeff_endWindows_eq_of_staircase}
    \leanok
    \uses{def:peps_bond_inserted_region_insert, thm:peps_regionInsertedCoeff_eq_extendInsert_bondInserted, thm:peps_staircase_end_reference_edge_boundary, thm:peps_staircase_pair_insert_eq_open}
    Let the torus have width $W$ and height $H$.
    Suppose that the one-orientation $L\times K$ window injectivity
    hypotheses hold for a tensor $B$, unions of injective regions are
    injective, and all bond dimensions of $B$ are positive. Assume
    \begin{align}
        2\le L,\qquad 2\le K,\qquad 1\le a,\qquad
        a+2L\le W,\qquad b+2K-1\le H,
        \notag \\
        2L+1\le W,\qquad 2K+1\le H.
        \notag
    \end{align}
    Let $C_j$ be inserts on the staircase windows $W_j$ whose deformed states
    all equal one common state. Suppose that
    $C_0=I_{B,W_0,e,M}$ and
    $C_{L+K-1}=I_{B,W_{L+K-1},e,M'}$.
    Then for every global physical configuration $\sigma$,
    \begin{align}
        C_B(W_0,e,M;\sigma|_{W_0},\sigma|_{V\setminus W_0})
         & =
        C_B\bigl(W_{L+K-1},e,M';
        \sigma|_{W_{L+K-1}},\sigma|_{V\setminus W_{L+K-1}}\bigr).
        \notag
    \end{align}
\end{theorem}
\begin{proof}
    \leanok
    Theorem~\ref{thm:peps_regionInsertedCoeff_eq_extendInsert_bondInserted}
    rewrites the two coefficients as the deformed states on the staircase end
    pair of the two corner-extended inserts:
    \begin{align}
        C_B(W_0,e,M;\sigma|_{W_0},\sigma|_{V\setminus W_0})
         & =
        \Psi_{B,S,\Ext_{W_0\subseteq S}(C_0)}(\sigma),
        \notag \\
        C_B(W_{L+K-1},e,M';
        \sigma|_{W_{L+K-1}},\sigma|_{V\setminus W_{L+K-1}})
         & =
        \Psi_{B,S,
            \Ext_{W_{L+K-1}\subseteq S}(C_{L+K-1})}(\sigma).
        \notag
    \end{align}
    The open-boundary equality on the staircase end pair gives
    $\Ext_{W_0\subseteq S}(C_0)
        =\Ext_{W_{L+K-1}\subseteq S}(C_{L+K-1})$, and hence the two coefficients
    are equal.
\end{proof}

\paragraph{Window-independence of the bond insertion.}

\begin{lemma}[Restriction and reassembly of a physical configuration]
    \label{lem:peps_assembleRegionSigma_restrict}
    \lean{TNLean.PEPS.assembleRegionσ_restrict}
    \leanok
    Let $\omega$ be a physical configuration on the full vertex set.  If
    $\omega_R$ and $\omega_{V\setminus R}$ are its restrictions to $R$ and
    to the complement, then
    \[
        \operatorname{Assemble}_R(\omega_R,\omega_{V\setminus R})=\omega.
    \]
\end{lemma}
\begin{proof}
    \leanok
    At a vertex in $R$ the assembled configuration reads $\omega_R$, and at a
    vertex outside $R$ it reads $\omega_{V\setminus R}$.  These restrictions
    are both inherited from $\omega$, so the two configurations agree at every
    vertex.
\end{proof}

\begin{theorem}[Bond-inserted state as an edge insertion]
    \label{thm:peps_deformedRegionStateAssembled_bondInserted_eq_smul_edgeInsertedCoeff}
    \lean{TNLean.PEPS.deformedRegionStateAssembled_bondInserted_eq_smul_edgeInsertedCoeff}
    \leanok
    \uses{thm:peps_deformedRegionState_bondInsertedRegionInsert,
        thm:peps_regionInsertedCoeff_eq_smul_edgeInsertedCoeff,
        lem:peps_assembleRegionSigma_restrict}
    Let $R$ be a region, let $f$ be a boundary edge of $R$, let
    $M\in\MN{D_B(f)}$, and let $\omega$ be a physical configuration on the
    full vertex set.  Then the assembled deformed state of the bond-inserted
    region insert is the non-boundary bond product times the corresponding
    edge-inserted coefficient:
    \[
        \Psi^{\mathrm{as}}_{B,R,I_{B,R,f,M}}(\omega)
        =
        p_B(R)\,
        E_B\bigl(f,\omega,\mathcal O_{R,f}(M)\bigr).
    \]
    Here $p_B(R)=\prod_{g\notin\partial R}D_B(g)$, and $\mathcal O_{R,f}$ is
    the boundary-edge orientation.
\end{theorem}
\begin{proof}
    \leanok
    Theorem~\ref{thm:peps_deformedRegionState_bondInsertedRegionInsert}
    rewrites the assembled deformed state as
    \[
        C_B(R,f,M;\omega_R,\omega_{V\setminus R}).
    \]
    Theorem~\ref{thm:peps_regionInsertedCoeff_eq_smul_edgeInsertedCoeff}
    turns this coefficient into
    \[
        p_B(R)\,
        E_B\bigl(f,\operatorname{Assemble}_R
            (\omega_R,\omega_{V\setminus R}),\mathcal O_{R,f}(M)\bigr).
    \]
    Lemma~\ref{lem:peps_assembleRegionSigma_restrict} identifies the assembled
    configuration with $\omega$.
\end{proof}

\begin{theorem}[Window-independence of a bond insertion]
    \label{thm:peps_bondInserted_windowIndependent}
    \lean{TNLean.PEPS.bondInserted_windowIndependent}
    \leanok
    \uses{thm:peps_deformedRegionStateAssembled_bondInserted_eq_smul_edgeInsertedCoeff}
    Let $R_1$ and $R_2$ be two regions for which the same edge $e$ is a
    boundary edge.  Let $M_1,M_2\in\MN{D_B(e)}$, and suppose that the two
    oriented insertions agree,
    \[
        \mathcal O_{R_1,e}(M_1)=\mathcal O_{R_2,e}(M_2).
    \]
    Then for every physical configuration $\omega$ on the full vertex set,
    \[
    \begin{aligned}
        p_B(R_2)\,
        \Psi^{\mathrm{as}}_{B,R_1,I_{B,R_1,e,M_1}}(\omega)
        &=
        p_B(R_1)\,
        \Psi^{\mathrm{as}}_{B,R_2,I_{B,R_2,e,M_2}}(\omega).
    \end{aligned}
    \]
\end{theorem}
\begin{proof}
    \leanok
    Applying
    Theorem~\ref{thm:peps_deformedRegionStateAssembled_bondInserted_eq_smul_edgeInsertedCoeff}
    to each of the two regions gives, for $i=1,2$,
    \[
        \Psi^{\mathrm{as}}_{B,R_i,I_{B,R_i,e,M_i}}(\omega)
        =
        p_B(R_i)\,E_B(e,\omega,\mathcal O_{R_i,e}(M_i)).
    \]
    The oriented insertions are equal by hypothesis, so both sides of the
    claimed identity are
    \[
        p_B(R_1)p_B(R_2)\,E_B(e,\omega,\mathcal O_{R_1,e}(M_1)).
    \]
\end{proof}

\paragraph{One virtual operation on every staircase window.}

The construction below formalizes the forward virtual/physical-operation
correspondence in the proof sketch of the two-dimensional corollary
\cite[corollary after Lemma~5, lines~2297--2444]{Molnar2018NormalPEPS}.

\begin{definition}[Partial state across a region cut]%
    \label{def:peps_region_partial_state}
    \lean{TNLean.PEPS.regionPartialState}
    \leanok
    For a region $R$ and a physical configuration $\tau$ on its complement,
    define the partial state $\psi^\tau_{A,R}$ on the physical space of $R$
    by
    \begin{align}
        \psi^\tau_{A,R}(\sigma)
         & =c_A(\omega_{R,\sigma,\tau}).
        \notag
    \end{align}
\end{definition}

\begin{theorem}[Partial states of equal closed states]%
    \label{thm:peps_region_partial_state_same_state}
    \lean{TNLean.PEPS.regionPartialState_sameState}
    \leanok
    \uses{def:peps_region_partial_state}
    If $A$ and $B$ generate the same closed state, then, for every region
    $R$ and every complement configuration $\tau$,
    $\psi^\tau_{A,R}=\psi^\tau_{B,R}$.
\end{theorem}
\begin{proof}
    \leanok
    Evaluate both functions at a physical configuration $\sigma$ on $R$ and
    apply the equality of the two closed-state coefficients at
    $\omega_{R,\sigma,\tau}$.
\end{proof}

\begin{theorem}[Blocked-tensor expansion of a partial state]%
    \label{thm:peps_region_partial_state_blocked_expansion}
    \lean{TNLean.PEPS.regionInteriorBondProd_smul_regionPartialState}
    \leanok
    \uses{def:peps_region_partial_state,
        def:peps_regionComplementBoundaryConfig}
    Let $P_A(R)$ be the product of the dimensions of the bonds internal to
    $R$. For every complement configuration $\tau$,
    \begin{align}
        P_A(R)\psi^\tau_{A,R}
         & =\sum_{\mu\in\partial R}
        T_{A,V\setminus R}(\bar\mu)(\tau)T_{A,R}(\mu).
        \label{eq:peps_partial_state_blocked_expansion}
    \end{align}
\end{theorem}
\begin{proof}
    \leanok
    Split the closed contraction across $R$. If $\mu$ and $\nu$ are boundary
    labels on $R$, the complement-boundary identification satisfies
    \begin{align}
        \bar\mu=\bar\nu
         & \quad\Longleftrightarrow\quad \mu=\nu.
        \notag
    \end{align}
    Hence the double boundary sum collapses to its diagonal, which gives the
    sum in~(\ref{eq:peps_partial_state_blocked_expansion}); each bond internal
    to $R$ contributes its dimension to the factor $P_A(R)$.
\end{proof}

\begin{definition}[Physical operator of a region insert]%
    \label{def:peps_physical_operator_of_region_insert}
    \lean{TNLean.PEPS.physicalOpOfRegionInsert}
    \leanok
    \uses{def:peps_region_out_endpoint_block_inverse}
    Let $R$ be an injective region and let
    $C=\{C(\mu)\}_{\mu\in\partial R}$ be any family in its physical
    space, indexed by the virtual boundary configurations.  If $L_{A,R}$
    is a left inverse of the blocked tensor map, define
    \begin{align}
        O_C
         & =\left(c\longmapsto\sum_{\mu}c_\mu C(\mu)\right)L_{A,R}.
        \notag
    \end{align}
\end{definition}

\begin{theorem}[Open-boundary realization of a region insert]%
    \label{thm:peps_physical_operator_of_region_insert_realizes}
    \lean{TNLean.PEPS.physicalOpOfRegionInsert_realizes}
    \leanok
    \uses{def:peps_physical_operator_of_region_insert}
    For every virtual boundary configuration $\mu$,
    \begin{align}
        O_C\bigl(T_{A,R}(\mu)\bigr) & =C(\mu).
        \notag
    \end{align}
    Thus the insert is obtained from the genuine region block by one
    physical operator before the region is contracted with its complement.
\end{theorem}
\begin{proof}
    \leanok
    The left inverse sends $T_{A,R}(\mu)$ to the standard basis vector
    indexed by $\mu$, and the displayed linear combination sends that basis
    vector to $C(\mu)$.
\end{proof}

\begin{definition}[Region insert of a physical operator]%
    \label{def:peps_region_insert_of_physical_operator}
    \lean{TNLean.PEPS.regionInsertOfPhysicalOp}
    \leanok
    Let $O$ be a physical operator on the physical space of a region $R$.
    Its action on the genuine open-boundary blocks defines the insert
    \begin{align}
        C^O_{A,R}(\mu) & =O\bigl(T_{A,R}(\mu)\bigr).
        \notag
    \end{align}
\end{definition}

\begin{theorem}[Closed state of a physical operator]%
    \label{thm:peps_closed_state_of_physical_operator}
    \lean{TNLean.PEPS.deformedRegionState_regionInsertOfPhysicalOp}
    \leanok
    \uses{def:peps_region_insert_of_physical_operator,
        thm:peps_region_partial_state_blocked_expansion}
    For a complement physical configuration $\tau$, let
    $\psi^\tau_{A,R}$ be the partial state across the cut. Then
    \begin{align}
        \Psi_{A,R,C^O_{A,R}}(\sigma,\tau)
         & =\bigl[O\bigl(P_A(R)\psi^\tau_{A,R}\bigr)\bigr](\sigma).
        \label{eq:peps_physical_op_partial_state}
    \end{align}
\end{theorem}
\begin{proof}
    \leanok
    Apply $O$ to~(\ref{eq:peps_partial_state_blocked_expansion}) and evaluate
    at $\sigma$. This proves
    (\ref{eq:peps_physical_op_partial_state}).
\end{proof}

\begin{theorem}[Physical operator transferred between equal closed states]%
    \label{thm:peps_physical_operator_same_state_transfer}
    \lean{TNLean.PEPS.deformedRegionState_regionInsertOfPhysicalOp_sameState}
    \leanok
    \uses{thm:peps_closed_state_of_physical_operator,
        thm:peps_region_partial_state_same_state}
    If $A$ and $B$ generate the same closed state, then
    \begin{align}
        P_B(R)\Psi_{A,R,C^O_{A,R}}(\sigma,\tau)
         & =P_A(R)\Psi_{B,R,C^O_{B,R}}(\sigma,\tau).
        \label{eq:peps_physical_op_same_state}
    \end{align}
    The two virtual boundary spaces need not be identified.
\end{theorem}
\begin{proof}
    \leanok
    Equality of the closed states gives
    $\psi^\tau_{A,R}=\psi^\tau_{B,R}$. Apply
    (\ref{eq:peps_physical_op_partial_state}) to both tensors and
    cross-multiply the two interior-bond products.
\end{proof}

\begin{definition}[Interior-bond region insert]%
    \label{def:peps_interior_bond_inserted_region_insert}
    \lean{TNLean.PEPS.interiorBondInsertedRegionInsert}
    \leanok
    \uses{def:peps_bond_inserted_region_insert, def:peps_corner_extension}
    Let $e=\{u,w\}$ be an ordered edge whose two endpoints lie in a region
    $R$, with $u$ its ordered left endpoint, and let
    $X\in\MN{D_A(e)}$.  Write $U=\{u\}$ and let $I_{A,U,e,X}$ be the
    boundary-bond insert on $U$.  The interior-bond insert on $R$ is
    \begin{align}
        I^{\mathrm{int}}_{A,R,e,X}
         & =
        \frac{P_A(R)}{P_A(U)}
        \Ext_{U\subseteq R}(I_{A,U,e,X}).
        \notag
    \end{align}
    Thus the genuine tensors at the vertices of $R\setminus U$ are adjoined
    to a boundary insertion of $X$ on $e$, with the normalization changed
    from the non-boundary bond product of $U$ to that of $R$.
\end{definition}

\begin{theorem}[Closed state of an interior-bond insert]%
    \label{thm:peps_interior_bond_inserted_state}
    \lean{TNLean.PEPS.deformedRegionStateAssembled_interiorBondInserted_eq_smul_edgeInsertedCoeff}
    \leanok
    \uses{def:peps_interior_bond_inserted_region_insert, thm:peps_corner_extension_state_preserves, thm:peps_deformedRegionStateAssembled_bondInserted_eq_smul_edgeInsertedCoeff}
    Suppose all bond dimensions are positive.  For every global physical
    configuration $\omega$,
    \begin{align}
        \Psi^{\mathrm{as}}_{A,R,I^{\mathrm{int}}_{A,R,e,X}}(\omega)
         & =P_A(R)E_A(e,\omega,X).
        \notag
    \end{align}
\end{theorem}
\begin{proof}
    \leanok
    Corner extension preserves the closed state.  On the singleton $U$, the
    boundary orientation is the identity because $U$ contains the ordered
    left endpoint of $e$.  Hence its boundary insert has assembled state
    $P_A(U)E_A(e,\omega,X)$.  The factor $P_A(R)/P_A(U)$ cancels $P_A(U)$
    and gives the displayed identity.
\end{proof}

\begin{definition}[Staircase family for one virtual operation]%
    \label{def:peps_staircase_virtual_operation_insert}
    \lean{TNLean.PEPS.staircaseVirtualOperationInsert}
    \leanok
    \uses{def:peps_bond_inserted_region_insert, def:peps_interior_bond_inserted_region_insert, lem:peps_reference_edge_boundary_windows, lem:peps_reference_edge_interior_windows}
    Let $e$ be the horizontal reference edge and let
    $X\in\MN{D_B(e)}$.  Define an insert $C_j(X)$ on each staircase window
    $W_j$, for $0\le j<L+K$, by
    \begin{align}
        C_j(X)
         & =
        \begin{cases}
            I_{B,W_0,e,X^{\mathsf T}},    & j=0,        \\
            I^{\mathrm{int}}_{B,W_j,e,X}, & 1\le j<L,   \\
            I_{B,W_j,e,X},                & L\le j<L+K.
        \end{cases}
        \notag
    \end{align}
    The first window contains only the ordered right endpoint of $e$, the
    middle windows contain both endpoints, and the remaining windows contain
    only the ordered left endpoint.  The transpose in the first case makes
    all three boundary orientations represent the same matrix $X$.
\end{definition}

\begin{theorem}[Physical operators on the staircase windows]%
    \label{thm:peps_staircase_virtual_operation_physical_realization}
    \lean{TNLean.PEPS.staircaseVirtualOperationPhysicalOp}
    \lean{TNLean.PEPS.staircaseVirtualOperationPhysicalOp_realizes}
    \leanok
    \uses{def:peps_staircase_virtual_operation_insert, def:peps_physical_operator_of_region_insert, thm:peps_physical_operator_of_region_insert_realizes}
    Suppose every $L\times K$ window is injective.  For every
    $0\leq j<L+K$, there is a physical operator $O_j(X)$ on the physical
    space of $W_j$ such that
    \begin{align}
        O_j(X)\bigl(T_{B,W_j}(\mu)\bigr) & =C_j(X)(\mu)
        \qquad (\mu\in\partial W_j).
        \notag
    \end{align}
    Hence the virtual operation is realized locally on each window, before
    any contraction with the complementary tensor network.
\end{theorem}
\begin{proof}
    \leanok
    Each staircase window is an $L\times K$ translate and is therefore
    injective.  Apply the open-boundary realization theorem to $C_j(X)$.
\end{proof}

\begin{lemma}[The two end inserts of the staircase family]%
    \label{lem:peps_staircase_virtual_operation_end_inserts}
    \lean{TNLean.PEPS.staircaseVirtualOperationInsert_zero}
    \lean{TNLean.PEPS.staircaseVirtualOperationInsert_last}
    \leanok
    \uses{def:peps_staircase_virtual_operation_insert}
    The end members of the family are
    \begin{align}
        C_0(X)       & =I_{B,W_0,e,X^{\mathsf T}},
                     &
        C_{L+K-1}(X) & =I_{B,W_{L+K-1},e,X}.
        \notag
    \end{align}
\end{lemma}
\begin{proof}
    \leanok
    The first identity is the $j=0$ case of the definition.  Since $K>0$,
    the last index satisfies $L\le L+K-1$, which gives the second identity.
\end{proof}

\begin{theorem}[Common state of the staircase virtual-operation family]%
    \label{thm:peps_staircase_virtual_operation_common_state}
    \lean{TNLean.PEPS.deformedRegionStateAssembled_staircaseVirtualOperationInsert}
    \leanok
    \uses{def:peps_staircase_virtual_operation_insert, lem:peps_staircase_virtual_operation_end_inserts, lem:peps_reference_edge_left_endpoint_staircase_window, thm:peps_interior_bond_inserted_state, thm:peps_deformedRegionStateAssembled_bondInserted_eq_smul_edgeInsertedCoeff, thm:peps_regionInteriorBondProd_staircaseWindow_eq}
    Let $B$ be translation invariant and suppose all its bond dimensions are
    positive.  Then for $0\le j<L+K$ and every global physical configuration
    $\omega$,
    \begin{align}
        \Psi^{\mathrm{as}}_{B,W_j,C_j(X)}(\omega)
         & =P_B(W_0)E_B(e,\omega,X).
        \notag
    \end{align}
    In particular, all $L+K$ inserts have one common deformed state.
\end{theorem}
\begin{proof}
    \leanok
    On $W_0$, the boundary orientation sends $X^{\mathsf T}$ to $X$.
    On $W_1,\ldots,W_{L-1}$, apply the interior-bond identity.  On the
    remaining windows the ordered left endpoint lies in the window, so the
    boundary orientation fixes $X$.  Each case gives
    $P_B(W_j)E_B(e,\omega,X)$.  Translation invariance identifies
    $P_B(W_j)$ with $P_B(W_0)$.
\end{proof}

\begin{theorem}[Common staircase state transferred between tensors]%
    \label{thm:peps_staircase_physical_operation_same_state}
    \lean{TNLean.PEPS.deformedRegionStateAssembled_staircasePhysicalOp_sameState}
    \leanok
    \uses{thm:peps_physical_operator_same_state_transfer,
        thm:peps_staircase_virtual_operation_physical_realization,
        thm:peps_staircase_virtual_operation_common_state,
        thm:peps_regionInteriorBondProd_staircaseWindow_eq}
    Let $L,K>0$, and choose staircase coordinates $(a,b)$ with
    $a+2L\leq\mathrm{width}$ and $2K\leq\mathrm{height}$. Let $A$ and $B$
    be translation invariant tensors with the same physical dimension and the
    same closed state. Suppose every cyclic $L\times K$ window of $A$ is
    injective and $D_A(g)>0$ for every edge $g$. For
    $X\in\MN{D_A(e)}$, let $O_j^A(X)$ be the physical operators obtained from
    the staircase insert family of $A$, and define inserts on the windows of
    $B$ by
    \begin{align}
        C_j^{A\to B}(X)(\mu)
         & =O_j^A(X)\bigl(T_{B,W_j}(\mu)\bigr).
        \notag
    \end{align}
    Then, for all $j,j'\in\mathbb{N}$,
    \begin{align}
        \Psi^{\mathrm{as}}_{B,W_j,C_j^{A\to B}(X)}
         & =\Psi^{\mathrm{as}}_{B,W_{j'},C_{j'}^{A\to B}(X)}.
        \notag
    \end{align}
    The source family has $0\leq j<L+K$. In the definition used here the
    truncated offsets stabilize at $W_{L+K-1}$ for $j\geq L+K-1$; the indices
    are not read modulo $L+K$. No equality between the edge-dependent bond
    dimensions of $A$ and $B$, and no positivity assumption on those of $B$,
    is required.
\end{theorem}
\begin{proof}
    \leanok
    For each $j$, (\ref{eq:peps_physical_op_same_state}) and the
    open-boundary realization of $O_j^A(X)$ give
    \begin{align}
        P_B(W_j)\Psi^{\mathrm{as}}_{A,W_j,C_j^A(X)}
         & =P_A(W_j)\Psi^{\mathrm{as}}_{B,W_j,C_j^{A\to B}(X)}.
        \notag
    \end{align}
    The staircase family of $A$ has one common deformed state. Translation
    invariance gives $P_A(W_j)=P_A(W_0)$ and
    $P_B(W_j)=P_B(W_0)$. Hence the right-hand side is independent of $j$.
    Positivity of the bond dimensions makes $P_A(W_0)$ nonzero, so it may be
    cancelled.
\end{proof}

\begin{definition}[The left staircase end operation $O_1^A(X)$]%
    \label{def:peps_staircase_O1_physical_operation}
    \lean{TNLean.PEPS.staircaseO1PhysicalOp}
    \leanok
    \uses{thm:peps_staircase_virtual_operation_physical_realization}
    For $X\in\MN{D_A(e)}$, define
    $O_1^A(X)=O_{L+K-1}^A(X)$ after identifying the last staircase
    window with the left end window.
\end{definition}

\begin{definition}[The underlying right staircase end operation $O_3^A(X)$]%
    \label{def:peps_staircase_O3_physical_operation}
    \lean{TNLean.PEPS.staircaseO3PhysicalOp}
    \leanok
    \uses{thm:peps_staircase_virtual_operation_physical_realization}
    For $X\in\MN{D_A(e)}$, define $O_3^A(X)=O_0^A(X)$ after identifying
    the first staircase window with the right end window.
    The relation (\ref{eq:peps_staircase_o3_cross_tensor_relation}) has the
    paper's transposed $O_3^{\mathsf T}$ orientation; no transpose is chosen
    on the abstract physical endomorphism space.
\end{definition}

\begin{theorem}[Cross-tensor virtual operation on the staircase edge]%
    \label{thm:peps_cross_tensor_virtual_operation_of_staircase}
    \lean{TNLean.PEPS.existsUnique_crossTensorVirtualOperation_of_staircasePhysicalOp_sameState}
    \leanok
    \uses{def:peps_staircase_O1_physical_operation,
        def:peps_staircase_O3_physical_operation}
    Let $L,K\geq2$, and choose the displayed non-wrapping horizontal
    staircase coordinates $(a,b)$ with
    \begin{align}
        1    & \leq a,
             & a+2L                 & \leq \mathrm{width},
             & b+2K-1               & \leq \mathrm{height},
        \notag                                              \\
        2L+1 & \leq \mathrm{width},
             & 2K+1                 & \leq \mathrm{height}.
        \notag
    \end{align}
    Suppose that $A$ and $B$ are translation invariant, generate the same
    closed state, every cyclic $L\times K$ window of either tensor is
    injective, and every bond dimension of either tensor is positive. For
    $X\in\MN{D_A(e)}$, there is a unique $Y\in\MN{D_B(e)}$ such that, for
    every external boundary configuration $\eta_{\mathrm L}$ and
    $\eta_{\mathrm R}$, and every reference-edge index $k$,
    \begin{align}
        O_1^A(X)\bigl(T_{B,W_{\mathrm L}}(k,\eta_{\mathrm L})\bigr)
         & =\sum_jY_{jk}T_{B,W_{\mathrm L}}(j,\eta_{\mathrm L}),
        \label{eq:peps_staircase_o1_cross_tensor_relation}       \\
        O_3^A(X)\bigl(T_{B,W_{\mathrm R}}(k,\eta_{\mathrm R})\bigr)
         & =\sum_jY_{kj}T_{B,W_{\mathrm R}}(j,\eta_{\mathrm R}).
        \label{eq:peps_staircase_o3_cross_tensor_relation}
    \end{align}
    Thus the relation (\ref{eq:peps_staircase_o3_cross_tensor_relation}) has the
    paper's $O_3^{\mathsf T}$ orientation.
    No equality between $D_A(e)$ and $D_B(e)$ is assumed or concluded, and
    this statement contains no multiplication or gauge assertion.
\end{theorem}
\begin{proof}
    \leanok
    \uses{thm:peps_staircase_physical_operation_same_state,
        thm:peps_staircase_pair_insert_eq_open,
        thm:peps_existsUnique_virtualOperation_of_horizontal_staircase_end_pair,
        thm:peps_regionInjectivityUnionClosure_of_overlap}
    The transferred physical inserts form one common $B$-state. The
    overlapping-union lemma supplies the union closure needed to compare
    consecutive staircase windows, and corner completion gives equality of
    the two end inserts with open boundary. At the two ends these inserts are
    precisely the actions of $O_3^A(X)$ and $O_1^A(X)$ on the genuine
    $B$-window tensors. Injectivity of those two windows and the two-end
    comparison give the displayed relations and the uniqueness of $Y$.
\end{proof}

\begin{theorem}[Composition of the canonical staircase end operations]%
    \label{thm:peps_staircase_end_physical_operation_multiplication}
    \lean{TNLean.PEPS.staircaseO1PhysicalOp_mul}
    \lean{TNLean.PEPS.staircaseO3PhysicalOp_mul}
    \leanok
    \uses{def:peps_staircase_O1_physical_operation,
        def:peps_staircase_O3_physical_operation}
    Let $L,K>0$, and choose staircase coordinates $(a,b)$ with
    $a+2L\leq\mathrm{width}$ and $2K\leq\mathrm{height}$. Suppose every
    cyclic $L\times K$ window of $A$ is injective. For matrices
    $X,Y\in\MN{D_A(e)}$, the canonical end operations satisfy
    \begin{align}
        O_1^A(XY) & =O_1^A(X)O_1^A(Y),
        \notag                         \\
        O_3^A(XY) & =O_3^A(Y)O_3^A(X).
        \notag
    \end{align}
    Thus the right relation preserves multiplication in ordinary order after
    passing to the orientation $(O_3^A)^{\mathsf T}$ used in Lemma~5.
\end{theorem}
\begin{proof}
    \leanok
    \uses{lem:peps_staircase_virtual_operation_end_inserts,
        thm:peps_lemma5_end_operation_multiplication}
    Split the virtual boundary of either end window into the reference-edge
    coordinate and the remaining external coordinates.  Let $E$ be this
    splitting map, and write $C_M$ for insertion of $M$ on the reference edge.
    On every fixed external fibre,
    \begin{align}
        C_M\bigl(E^{-1}(j,\eta)\bigr)
         & =\sum_i M_{ij}\,T_{A,W}\bigl(E^{-1}(i,\eta)\bigr).
        \notag
    \end{align}
    The chosen left inverse of the injective blocked tensor is a left inverse
    on the whole boundary-coordinate space, so this identity gives
    $O_W(M)O_W(N)=O_W(MN)$.  The first window carries $X^{\mathsf T}$; hence
    $(XY)^{\mathsf T}=Y^{\mathsf T}X^{\mathsf T}$ reverses the order of the
    underlying right operations.
\end{proof}

\begin{definition}[Canonical cross-tensor matrix assignment]%
    \label{def:peps_staircase_cross_tensor_virtual_operation}
    \lean{TNLean.PEPS.staircaseCrossTensorVirtualOperation}
    \lean{TNLean.PEPS.staircaseCrossTensorVirtualOperation_spec}
    \leanok
    \uses{thm:peps_cross_tensor_virtual_operation_of_staircase}
    Under the hypotheses of
    Theorem~\ref{thm:peps_cross_tensor_virtual_operation_of_staircase}, in the
    paper's notation the input and uniquely recovered matrices are
    called $X$ and $Y$.  Writing $Y_{A\to B}(X)$ only to distinguish the two
    tensor families, define
    \begin{align}
        Y_{A\to B}(X) & \in\MN{D_B(e)}
        \notag
    \end{align}
    to be the unique matrix satisfying
    (\ref{eq:peps_staircase_o1_cross_tensor_relation}) and
    (\ref{eq:peps_staircase_o3_cross_tensor_relation}) for every left and
    right external boundary configuration.  This definition makes no choice
    of an external boundary configuration.
\end{definition}

\begin{theorem}[Multiplicativity of the cross-tensor assignment]%
    \label{thm:peps_staircase_cross_tensor_virtual_operation_multiplication}
    \lean{TNLean.PEPS.staircaseCrossTensorVirtualOperation_mul}
    \leanok
    \uses{def:peps_staircase_cross_tensor_virtual_operation}
    Under the hypotheses of
    Theorem~\ref{thm:peps_cross_tensor_virtual_operation_of_staircase}, for all
    $X_1,X_2\in\MN{D_A(e)}$,
    \begin{align}
        Y_{A\to B}(X_1X_2)
         & =Y_{A\to B}(X_1)Y_{A\to B}(X_2).
        \notag
    \end{align}
\end{theorem}
\begin{proof}
    \leanok
    \uses{thm:peps_staircase_end_physical_operation_multiplication,
        thm:peps_lemma5_end_operation_multiplication}
    Set $Y_r=Y_{A\to B}(X_r)$ for $r=1,2$.  For every left external
    boundary configuration, the two $O_1$ relations compose in ordinary order:
    \begin{align}
        O_1^A(X_1)\Bigl(O_1^A(X_2)
        \bigl(T_{B,W_{\mathrm L}}(k,\eta_{\mathrm L})\bigr)\Bigr)
         & =\sum_i (Y_1Y_2)_{ik}\,
        T_{B,W_{\mathrm L}}(i,\eta_{\mathrm L}).
        \notag
    \end{align}
    For every right external boundary configuration, the underlying $O_3$
    operations compose in reverse order:
    \begin{align}
        O_3^A(X_2)\Bigl(O_3^A(X_1)
        \bigl(T_{B,W_{\mathrm R}}(k,\eta_{\mathrm R})\bigr)\Bigr)
         & =\sum_j (Y_1Y_2)_{kj}\,
        T_{B,W_{\mathrm R}}(j,\eta_{\mathrm R}).
        \notag
    \end{align}
    This is ordinary multiplication in the $O_3^{\mathsf T}$ orientation.  The
    preceding theorem identifies the left sides with the canonical operations
    for $X_1X_2$.  Both end relations therefore realize $Y_1Y_2$, and global
    uniqueness identifies this product with $Y_{A\to B}(X_1X_2)$.
\end{proof}

\begin{theorem}[End-window identity for one virtual operation]%
    \label{thm:peps_end_window_identity_staircase_virtual_operation}
    \lean{TNLean.PEPS.regionInsertedCoeff_endWindows_eq_staircaseVirtualOperation}
    \leanok
    \uses{lem:peps_assembleRegionSigma_restrict, lem:peps_reference_edge_left_endpoint_staircase_window, thm:peps_regionInsertedCoeff_eq_smul_edgeInsertedCoeff, thm:peps_regionInteriorBondProd_staircaseWindow_eq}
    Suppose the tensor is translation invariant and the staircase satisfies
    the stated size bounds.  Then for every $X\in\MN{D_B(e)}$ and every
    global physical configuration $\omega$,
    \begin{align}
        C_B(W_0,e,X^{\mathsf T};\omega|_{W_0},
        \omega|_{V\setminus W_0})
         & =
        C_B(W_{L+K-1},e,X;\omega|_{W_{L+K-1}},
        \omega|_{V\setminus W_{L+K-1}}).
        \notag
    \end{align}
\end{theorem}
\begin{proof}
    \leanok
    The region-to-edge identity sends both sides to the edge coefficient
    with $X$ on the same reference edge: the first boundary orientation
    transposes $X^{\mathsf T}$, while the last fixes $X$.  Translation
    invariance identifies the two non-boundary bond products.
\end{proof}


\begin{definition}[Window-region coefficient-identity witness]
    \label{def:peps_window_edge_coeff_identity_witness}
    \lean{TNLean.PEPS.windowEdgeCoeffIdentityWitness}
    \leanok
    \uses{def:peps_regionInsertedCoeff, thm:peps_staircase_end_reference_edge_boundary}
    Let $W$ be the left end window of a horizontal staircase with parameters
    $(L,K)$ and lower-left coordinate $(a,b)$, and let $e$ be its reference
    edge. Suppose that
    \begin{align}
        0<L,\qquad 0<K,\qquad
        1\le a,\qquad
        a+2L\le W_{\mathrm{torus}},\qquad
        b+2K-1\le H_{\mathrm{torus}}.
        \notag
    \end{align}
    Then $e\in\partial W$. Suppose also that the blocked tensors of $B$ on
    $W$ and $V\setminus W$ are injective, that all bond dimensions of $B$ are
    positive, and that the bond dimensions of $A$ and $B$ at $e$ are
    identified by $\iota_e:\MN{D_A(e)}\simeq\MN{D_B(e)}$. If
    $Z\in\GL(D_B(e))$ satisfies, for every matrix $M$ and physical
    configurations $\sigma,\tau$,
    \begin{align}
        C_A(W,e,M;\sigma,\tau)
        =
        C_B(W,e,Z\,\iota_e(M)\,Z^{-1};\sigma,\tau),
        \notag
    \end{align}
    then these objects form the coefficient-identity witness at $e$ with
    witness region $W$ and with both gauge entries equal to $Z$.
\end{definition}

\begin{definition}[Window-region witness from window injectivity]
    \label{def:peps_window_edge_coeff_identity_witness_of_hypotheses}
    \lean{TNLean.PEPS.windowEdgeCoeffIdentityWitness_of_hypotheses}
    \leanok
    \uses{def:peps_window_edge_coeff_identity_witness}
    Let $W$ be the left end window around the reference edge $e$, placed with
    lower-left coordinate $(a,b)$. Suppose that every cyclic $L\times K$
    window of $B$ is blocked-tensor injective, unions of injective regions are
    injective, all bond dimensions of $B$ are positive, and
    \begin{align}
        2\le L,\qquad 2\le K,\qquad
        2L+1\le W_{\mathrm{torus}},\qquad
        2K+1\le H_{\mathrm{torus}},
        \notag \\
        1\le a,\qquad
        a+2L\le W_{\mathrm{torus}},\qquad
        b+2K-1\le H_{\mathrm{torus}}.
        \notag
    \end{align}
    Suppose also that the bond dimensions of $A$ and $B$ at $e$ are
    identified by $\iota_e:\MN{D_A(e)}\simeq\MN{D_B(e)}$, and let
    $Z\in\GL(D_B(e))$. If the conjugation identity
    \begin{align}
        C_A(W,e,M;\sigma,\tau)
        =
        C_B(W,e,Z\,\iota_e(M)\,Z^{-1};\sigma,\tau)
        \notag
    \end{align}
    holds for every $M,\sigma,\tau$, then the coefficient-identity witness of
    Definition~\ref{def:peps_window_edge_coeff_identity_witness} is obtained
    for the region $W$. The hypotheses supply the injectivity of both $W$ and
    its complement at the minimal torus size.
\end{definition}

\paragraph{Window realization and witness read-off.}

\begin{lemma}[Window-level edge-insertion injectivity]%
    \label{lem:peps_windowRegionInsertedCoeff_injective}
    \lean{TNLean.PEPS.windowRegionInsertedCoeff_injective}
    \leanok
    \uses{thm:peps_regionInsertedCoeff_injective,
        lem:peps_horizontal_staircase_window_injectivity,
        thm:peps_horizontal_staircase_reference_edge_crossing}
    Let $B$ be a PEPS tensor on a torus whose periods $W,H$ satisfy
    $1<W$ and $1<H$.  Assume the one-orientation $L\times K$
    window-injectivity hypotheses for $B$, union closure for injective
    regions, positive bond dimensions, and
    \[
        2\le L,\qquad 2\le K,\qquad 1\le a,\qquad
        a+2L\le W,\qquad b+2K-1\le H,
    \]
    together with
    \[
        2L+1\le W,\qquad 2K+1\le H.
    \]
    Let $W_{\mathrm L}$ be the left end window and let $e$ be the reference
    edge.  If $M,M'\in\MN{D_B(e)}$ satisfy
    \[
        C_B(W_{\mathrm L},e,M;\sigma,\tau)
        =
        C_B(W_{\mathrm L},e,M';\sigma,\tau)
    \]
    for every physical configuration $\sigma$ on $W_{\mathrm L}$ and $\tau$
    on $V\setminus W_{\mathrm L}$, then $M=M'$.
\end{lemma}
\begin{proof}
    \leanok
    The window-injectivity hypotheses give injectivity of the blocked tensor
    map on $W_{\mathrm L}$.  The single-window complement is also injective at
    the displayed size.  Since $e\in\partial W_{\mathrm L}$,
    Theorem~\ref{thm:peps_regionInsertedCoeff_injective} applied to the
    region $W_{\mathrm L}$ gives $M=M'$.
\end{proof}

\begin{theorem}[Window witness from a region insertion transfer]%
    \label{thm:peps_windowEdgeCoeffIdentityWitness_of_transfer}
    \lean{TNLean.PEPS.exists_windowEdgeCoeffIdentityWitness_of_transfer}
    \leanok
    \uses{def:peps_region_insertion_transfer,
        def:peps_window_edge_coeff_identity_witness,
        def:peps_window_edge_coeff_identity_witness_of_hypotheses}
    Let $A$ and $B$ be PEPS tensors on a torus whose periods $W,H$ satisfy
    $1<W$ and $1<H$.  Assume the one-orientation $L\times K$
    window-injectivity hypotheses for both tensors, union closure for
    injective regions for both tensors, positive bond dimensions, and
    \[
        2\le L,\qquad 2\le K,\qquad 1\le a,\qquad
        a+2L\le W,\qquad b+2K-1\le H,
    \]
    together with
    \[
        2L+1\le W,\qquad 2K+1\le H.
    \]
    Let $W_{\mathrm L}$ be the left end window and let $e$ be the reference
    edge.  If $T_e:\MN{D_A(e)}\to\MN{D_B(e)}$ is a region insertion transfer
    on $W_{\mathrm L}$ at $e$, then there exist
    $Z,Z_{\mathrm{ref}}\in\GL(D_B(e))$ and a bond-dimension identification
    $\iota_e:\MN{D_A(e)}\simeq\MN{D_B(e)}$ such that the coefficient-identity
    witness at $e$ is nonempty.  The coefficient identity has the form
    \[
        C_A(W_{\mathrm L},e,M;\sigma,\tau)
        =
        C_B(W_{\mathrm L},e,Z\,\iota_e(M)\,Z^{-1};\sigma,\tau).
    \]
\end{theorem}
\begin{proof}
    \leanok
    The transfer determines a gauge $Z$ by the region insertion algebra:
    \[
        T_e(M)=Z\,\iota_e(M)\,Z^{-1}.
    \]
    Its coefficient-matching property gives the displayed identity.  The
    window and complement injectivity hypotheses are those required by
    Definition~\ref{def:peps_window_edge_coeff_identity_witness_of_hypotheses},
    which yields the witness with $Z_{\mathrm{ref}}=Z$.
\end{proof}

\begin{theorem}[Window witness from single-vertex realization]%
    \label{thm:peps_windowEdgeCoeffIdentityWitness_of_realizes}
    \lean{TNLean.PEPS.exists_windowEdgeCoeffIdentityWitness_of_realizes}
    \leanok
    \uses{def:peps_region_transfer_realizes,
        def:peps_region_insertion_transfer_of_realizes,
        thm:peps_windowEdgeCoeffIdentityWitness_of_transfer}
    Let $A$ and $B$ be PEPS tensors on a torus whose periods $W,H$ satisfy
    $1<W$ and $1<H$.  Assume the one-orientation $L\times K$
    window-injectivity hypotheses for both tensors, union closure for
    injective regions for both tensors, matched bond dimensions, equality of
    the PEPS state, positive bond dimensions, and
    \[
        2\le L,\qquad 2\le K,\qquad 1\le a,\qquad
        a+2L\le W,\qquad b+2K-1\le H,
    \]
    together with
    \[
        2L+1\le W,\qquad 2K+1\le H.
    \]
    Let $W_{\mathrm L}$ be the left end window and let $e$ be the reference
    edge.  Suppose the endpoint tensor of each family at the in-region
    endpoint of $e$ has linearly independent columns, and suppose the boundary
    matrix transfer is realized on $W_{\mathrm L}$ in both directions.  Then
    there exist $Z,Z_{\mathrm{ref}}\in\GL(D_B(e))$ and a bond-dimension
    identification $\iota_e:\MN{D_A(e)}\simeq\MN{D_B(e)}$ such that the
    coefficient-identity witness at $e$ is nonempty.
\end{theorem}
\begin{proof}
    \leanok
    The two realized boundary transfers determine a region insertion transfer:
    \[
        C_A(W_{\mathrm L},e,M;\sigma,\tau)
        =
        C_B(W_{\mathrm L},e,T_e(M);\sigma,\tau),
        \qquad
        T_e(MM')=T_e(M)T_e(M').
    \]
    Theorem~\ref{thm:peps_windowEdgeCoeffIdentityWitness_of_transfer} then
    gives the two gauges and the nonempty coefficient-identity witness.
\end{proof}

\paragraph{Window witness from bond locality.}

\begin{theorem}[Forward window coefficient transfer from bond locality]
    \label{thm:peps_windowCoeffTransferAB_of_bondLocal}
    \lean{TNLean.PEPS.windowCoeffTransferAB_of_bondLocal}
    \leanok
    \uses{def:peps_regionInsertedCoeff,
        lem:peps_horizontal_staircase_window_injectivity,
        thm:peps_horizontal_staircase_reference_edge_crossing,
        thm:peps_region_transfer_coeff_double_factorization,
        thm:peps_region_transfer_coeff_incident_form}
    Let $A$ and $B$ be PEPS tensors on a torus whose periods
    $W_{\mathrm{torus}},H_{\mathrm{torus}}$ satisfy $1<W_{\mathrm{torus}}$
    and $1<H_{\mathrm{torus}}$.  Assume the one-orientation $L\times K$
    window-injectivity hypotheses and union closure for both tensors, equality
    of the PEPS state, matched bond dimensions, positive bond dimensions, and
    \[
        2\le L,\qquad 2\le K,\qquad 1\le a,\qquad
        a+2L\le W_{\mathrm{torus}},\qquad
        b+2K-1\le H_{\mathrm{torus}},
    \]
    together with
    \[
        2L+1\le W_{\mathrm{torus}},\qquad
        2K+1\le H_{\mathrm{torus}}.
    \]
    Let $W_{\mathrm L}$ be the left end window and let $e$ be the reference
    edge.  Suppose the $A$-to-$B$ transfer kernel on $W_{\mathrm L}$ at $e$ is
    incident along the boundary leg $e$.  Then for every
    $M\in\MN{D_A(e)}$ there is a matrix $N\in\MN{D_B(e)}$ such that
    \[
        C_A(W_{\mathrm L},e,M;\sigma,\tau)
        =
        C_B(W_{\mathrm L},e,N;\sigma,\tau)
    \]
    for every physical configuration $\sigma$ on $W_{\mathrm L}$ and $\tau$
    on $V\setminus W_{\mathrm L}$.
\end{theorem}
\begin{proof}
    \leanok
    The window hypotheses give injectivity of the blocked tensor maps on
    $W_{\mathrm L}$ for $A$ and $B$, while union closure gives injectivity on
    the complements.  The coefficient-transfer theorem for the blocked region
    and its complement rewrites the first tensor's inserted coefficient as a
    double contraction through the second tensor.  The incident form of the
    transfer kernel along $e$ turns this double contraction into the single
    region insertion $C_B(W_{\mathrm L},e,N;\sigma,\tau)$.
\end{proof}

\begin{theorem}[Backward window coefficient transfer from bond locality]
    \label{thm:peps_windowCoeffTransferBA_of_bondLocal}
    \lean{TNLean.PEPS.windowCoeffTransferBA_of_bondLocal}
    \leanok
    \uses{def:peps_regionInsertedCoeff,
        lem:peps_horizontal_staircase_window_injectivity,
        thm:peps_horizontal_staircase_reference_edge_crossing,
        thm:peps_region_transfer_coeff_double_factorization,
        thm:peps_region_transfer_coeff_incident_form}
    Under the hypotheses of
    Theorem~\ref{thm:peps_windowCoeffTransferAB_of_bondLocal}, suppose instead
    that the $B$-to-$A$ transfer kernel on $W_{\mathrm L}$ at $e$ is incident
    along the boundary leg $e$.  Then for every $N\in\MN{D_B(e)}$ there is a
    matrix $M\in\MN{D_A(e)}$ such that
    \[
        C_B(W_{\mathrm L},e,N;\sigma,\tau)
        =
        C_A(W_{\mathrm L},e,M;\sigma,\tau)
    \]
    for every physical configuration $\sigma$ on $W_{\mathrm L}$ and $\tau$
    on $V\setminus W_{\mathrm L}$.
\end{theorem}
\begin{proof}
    \leanok
    Apply the forward theorem with the two tensors interchanged.  The same
    window and complement injectivity hypotheses are available in the reversed
    direction, and the incident transfer kernel now supplies the matrix $M$ on
    the $A$ boundary bond.
\end{proof}

\begin{theorem}[Window witness from bond locality and multiplicativity]
    \label{thm:peps_windowEdgeCoeffIdentityWitness_of_bondLocal}
    \lean{TNLean.PEPS.exists_windowEdgeCoeffIdentityWitness_of_bondLocal}
    \leanok
    \uses{thm:peps_windowCoeffTransferAB_of_bondLocal,
        thm:peps_windowCoeffTransferBA_of_bondLocal,
        def:peps_coeff_transfer_map,
        thm:peps_region_edge_gauge_of_coeff_transfer,
        def:peps_window_edge_coeff_identity_witness,
        def:peps_window_edge_coeff_identity_witness_of_hypotheses}
    Assume the hypotheses of
    Theorem~\ref{thm:peps_windowCoeffTransferAB_of_bondLocal} and its
    backward analogue
    Theorem~\ref{thm:peps_windowCoeffTransferBA_of_bondLocal}.  Let
    $T_e:\MN{D_A(e)}\to\MN{D_B(e)}$ be the coefficient transfer chosen from
    the forward theorem.  If
    \[
        T_e(MM')=T_e(M)T_e(M')
    \]
    for all $M,M'\in\MN{D_A(e)}$, then there exist
    $Z,Z_{\mathrm{ref}}\in\GL(D_B(e))$ and a bond-dimension identification
    $\iota_e:\MN{D_A(e)}\simeq\MN{D_B(e)}$ such that the window-region
    coefficient-identity witness at $e$ is nonempty.  In particular the
    coefficient identity has the form
    \[
        C_A(W_{\mathrm L},e,M;\sigma,\tau)
        =
        C_B(W_{\mathrm L},e,Z\,\iota_e(M)\,Z^{-1};\sigma,\tau).
    \]
\end{theorem}
\begin{proof}
    \leanok
    The two bond-locality hypotheses give the forward and backward coefficient
    transfers by
    Theorems~\ref{thm:peps_windowCoeffTransferAB_of_bondLocal}
    and~\ref{thm:peps_windowCoeffTransferBA_of_bondLocal}.  With the displayed
    multiplicativity, the coefficient-transfer construction gives a conjugating
    gauge $Z$ and the coefficient identity on $W_{\mathrm L}$.  The window and
    complement injectivity hypotheses then supply the window-region
    coefficient-identity witness with $Z_{\mathrm{ref}}=Z$.
\end{proof}



\section{Normal PEPS Fundamental Theorem}\label{sec:peps_normal_capstone}
The pieces now combine. The blocking hypotheses of a normal PEPS, that the
tensors are injective after blocking around every edge and that each site
separates two injective regions whose complements are also injective, are what
the injective theorem needs once union closure is in hand. The Fundamental
Theorem for normal PEPS, Theorem~\ref{thm:fundamentalTheorem_normalPEPS_connected},
follows on a connected graph; this is the general normal-PEPS theorem in
\cite[Section~4]{Molnar2018NormalPEPS}, following its Theorem~3, and its
gauges are unique modulo balanced edge scalars.
The corollaries specialize it to closed chains, translation-invariant matrix
product states, and the square lattice.

\begin{definition}[Normal PEPS blocking hypotheses]%
    \label{def:peps_normalPEPSBlockingHypotheses}
    \lean{TNLean.PEPS.NormalPEPSBlockingHypotheses}
    \lean{TNLean.PEPS.NormalPEPSBlockingHypotheses.injective_chain_at_edge}
    \lean{TNLean.PEPS.%
        NormalPEPSBlockingHypotheses.endpoint_disjoint_cover_at_edge}
    \lean{TNLean.PEPS.%
        NormalPEPSBlockingHypotheses.endpoint_injective_disjoint_cover_at_edge}
    \lean{TNLean.PEPS.NormalPEPSBlockingHypotheses.mem_withSite_iff}
    \leanok
    \uses{def:peps_normalEdgeBlockingHypotheses,
        def:peps_normalOneSiteSeparationHypotheses}
    The general normal theorem assumes both the edge-centred blockings into
    three injective regions and the one-site comparison regions with injective
    complements. These two assumptions are collected as the normal PEPS
    blocking hypotheses; the same hypotheses give the corresponding edge and
    one-site consequences without restating the separate edge and one-site
    assumptions. The following two panels show their square-lattice source
    realization; they do not restrict the abstract definition to this lattice.
    \begin{tenkzequation}
        % Formula verdict: A1=(4-6,2-3), A2=(3-4,4-6), and A3 is their
        % complement in the 7-by-7 source frame; the marked edge
        % (4,3)-(4,4) has one endpoint in A1 and the other in A2.
        % Ink-to-index verdict: every dot is an ordinary tensor, every lattice
        % edge is a virtual contraction, and the three fills encode the
        % edge-centred disjoint cover used by the blocking hypothesis.
        % Contraction/boundary verdict: the panel has 84 internal virtual
        % contractions and 28 cut virtual bonds, with no exposed physical
        % indices, hence boundary type (V,P)=(28,0).
        % Source verdict: exact square-lattice specialization of the first
        % hypothesis in arXiv:1804.04964, paper_normal.tex:1475--1499.
        \begin{tenkzlattice}[rows=7, cols=7, boundary legs]
            \tnregion[slot=selected, label={$A_1$},
                label at=center]{(4-6,2-3)}
            \tnregion[slot=secondary, label={$A_2$},
                label at=center]{(3-4,4-6)}
            \tnregion[slot=complement, outline, label={$A_3$},
                label at=north]
                {(1-7,1-7) - (4-6,2-3) - (3-4,4-6)}
            \tnedge[distinguished]{(4,3)-(4,4)}
        \end{tenkzlattice}
        $\qquad\text{edge blocking}$
    \end{tenkzequation}
    \begin{tenkzequation}
        % Formula verdict: R=(2-4,2-4)-(4,4), S=(2-4,2-4), and
        % S=R union {v} for v=(4,4); the marked site is present in both panels.
        % Ink-to-index verdict: every dot is an ordinary tensor and every
        % lattice edge is a virtual contraction; region fills encode only the
        % one-site set comparison, while the marked site identifies v.
        % Contraction/boundary verdict: each 5-by-5 panel has 40 internal
        % virtual contractions and 20 cut virtual bonds, with no exposed
        % physical indices, hence boundary type (V,P)=(20,0).
        % Source verdict: exact discrete form of the one-site comparison in
        % arXiv:1804.04964, paper_normal.tex:1405--1429 and 1519--1544.
        $\text{and}\qquad$
        \begin{tenkzlattice}[rows=5, cols=5, boundary legs]
            \tnregion[slot=selected, label=$R$, label at=north west]
                {(2-4,2-4) - (4,4)}
            \tnsite[role=marked]{(4,4)}
        \end{tenkzlattice}
        $\qquad\subset\qquad$
        \begin{tenkzlattice}[rows=5, cols=5, boundary legs]
            \tnregion[slot=selected, label=$S$, label at=north west]{(2-4,2-4)}
            \tnsite[role=marked]{(4,4)}
        \end{tenkzlattice}
        $\qquad S=R\cup\{v\}$
    \end{tenkzequation}
\end{definition}

\begin{theorem}[Square-lattice normal blocking hypotheses from margin covers]%
    \label{thm:peps_normalSquarePEPSBlockingHypotheses_of_marginCovers}
    \lean{TNLean.PEPS.normalSquarePEPSBlockingHypotheses_of_marginCovers}
    \lean{TNLean.PEPS.normalSquarePEPSBlockingHypotheses_edgeBlocking_of_marginCovers}
    \lean{TNLean.PEPS.normalSquarePEPSBlockingHypotheses_blockingData_of_marginCovers}
    \lean{TNLean.PEPS.%
        normalSquarePEPSBlockingHypotheses_oneSiteSeparation_of_marginCovers}
    \lean{TNLean.PEPS.%
        normalSquarePEPSBlockingHypotheses_injective_chain_of_marginCovers}
    \lean{TNLean.PEPS.%
        normalSquarePEPSBlockingHypotheses_endpoint_disjoint_cover_of_marginCovers}
    \lean{TNLean.PEPS.%
        normalSquarePEPSBlockingHypotheses_endpoint_injective_disjoint_cover_of_marginCovers}
    \lean{TNLean.PEPS.%
        normalSquarePEPSBlockingHypotheses_mem_withSite_iff_of_marginCovers}
    \leanok
    \uses{thm:peps_normalSquareTranslatedEdgeBlockingHypotheses_of_windows,
        def:peps_normalPEPSBlockingHypotheses,
        def:peps_normalOneSiteSeparationHypotheses}
    Given rectangular injectivity hypotheses, finite union closure, oriented
    margin-and-cover hypotheses at every edge, and one-site separation, the finite
    square lattice satisfies the general normal blocking hypotheses. The
    edge-blocking hypotheses are exactly the margin-cover construction, the
    one-edge datum at an edge is the datum supplied there, and the one-site
    separation is exactly the supplied one-site separation. Hence these hypotheses
    give the three injective regions, the endpoint-disjoint-cover facts at
    each edge, and the one-site membership formula as consequences of the same
    normal blocking hypotheses.
\end{theorem}
\begin{proof}\leanok
    Combine the edge-blocking hypotheses constructed from the margin-cover
    hypotheses with the given one-site separation hypotheses.
\end{proof}

\begin{theorem}[Fundamental Theorem for normal PEPS]%
    \label{thm:fundamentalTheorem_normalPEPS}
    \notready
    \uses{thm:peps_normalEdgeBlockingHypotheses_injectiveChain,
        def:peps_normalPEPSBlockingHypotheses,
        thm:peps_oneVertexComplementComparison,
        def:GaugeEquivPEPS}
    Suppose two normal PEPS generate the same state, can be blocked into
    three-site injective chains around every edge, and for every site admit
    two injective regions with injective complements that differ exactly at
    that site. Then their defining tensors are related by a local gauge, and
    the gauges are unique up to the scalar freedom allowed by the graph. This
    is the general normal-PEPS theorem in
    \cite[Section~4]{Molnar2018NormalPEPS}, following its Theorem~3.
    A connected-graph form, with the blocking hypothesis stated as the
    source's proof uses it (one blocking shared by the two states, each
    distinguished edge the entire bond of its chain), is
    Theorem~\ref{thm:fundamentalTheorem_normalPEPS_connected}; its per-edge
    uniqueness clause is
    Theorem~\ref{thm:fundamentalTheorem_normalPEPS_gauge_unique}.
    The source's remark following the theorem, that for a translationally
    invariant system the gauges are also translationally invariant when the
    proportionality constants are not absorbed into them, is realized on the
    torus by
    Theorem~\ref{thm:fundamentalTheorem_normalTorusPEPS_unconditional}: its
    gauge family is carried onto every edge by the translations from one
    reference matrix per orientation class, and its scalar $\lambda$ is kept
    separate rather than absorbed.
\end{theorem}

\begin{definition}[Doubly injective regions]%
    \label{def:peps_regionInjectivityDataPair}
    \lean{TNLean.PEPS.regionInjectivityDataPair}
    \leanok
    \uses{def:peps_regionInjectivityData}
    For two region-injectivity assignments on the same vertex set, a region is
    doubly injective when it is injective for each of the two. The general
    normal theorem blocks two states by one shared geometry, and its final
    comparison contracts both tensors over the same regions, so its blocking
    hypotheses are instantiated at the doubly injective assignment of the two
    tensors.
\end{definition}

\begin{theorem}[Equal bond dimensions from the normal blocking hypotheses]%
    \label{thm:peps_bondDim_eq_of_normalBlocking}
    \lean{TNLean.PEPS.bondDim_eq_of_normalBlocking}
    \lean{TNLean.PEPS.bondDim_apply_eq_of_blockingData}
    \leanok
    \uses{def:peps_normalPEPSBlockingHypotheses,
        def:peps_regionInjectivityDataPair}
    Let $A$ and $B$ be PEPS tensors on a finite simple graph with the same
    positive physical dimension and positive bond dimensions, generating the
    same state and satisfying the normal blocking hypotheses at the doubly
    injective assignment of the two tensors, with each distinguished edge the
    only edge of the graph joining its two adjacent blocks. Then the bond
    dimensions of $A$ and $B$ agree on every edge.
    In the source the equality is part of the isomorphism lemma of
    \cite[Section~3]{Molnar2018NormalPEPS}: the insertion correspondence
    $X\mapsto Y$ on the chosen bond is an algebra isomorphism between the two
    full bond matrix algebras, so the bond dimensions on the two sides agree;
    the normal case inherits the argument through the blocked three-site
    chains. The single-crossing condition is the source's blocking around an
    edge made explicit: the chosen edge is the entire bond between the first
    two parties of the blocked chain.
    \begin{proof}
        \leanok
        \uses{thm:peps_edgeBlockedThreeSiteInjective,
            thm:peps_edgeBlockedInsertionAlgebraIsomorphism}
        Fix an edge. Blocking both networks into the shared three regions
        around it produces two three-site chains on the three-vertex complete
        graph whose super-sites carry the blocked-region weight families,
        injective by hypothesis. Each coarse state coefficient is a positive
        merge-collapse constant, a product of bond dimensions over the edges
        not incident to each region, times the corresponding original state
        coefficient, so the two coarse states are proportional with ratio the
        two constants. Multiplying one super-site of each chain by the other
        chain's constant yields chains with equal states, unchanged bond
        dimensions, and the super-site injectivities preserved. The matched
        matrix insertions on the distinguished coarse bond of the two scaled
        chains form an algebra isomorphism between the two full bond matrix
        algebras
        (Theorem~\ref{thm:peps_edgeBlockedInsertionAlgebraIsomorphism}), and
        an isomorphism of full matrix algebras forces equal matrix sizes.
        Under the single-crossing hypothesis the coarse bond carries the
        original bond dimension of the distinguished edge for each tensor, so
        the two original bond dimensions agree.
    \end{proof}
\end{theorem}

\begin{theorem}[Fundamental Theorem for normal PEPS on a connected graph]%
    \label{thm:fundamentalTheorem_normalPEPS_connected}
    \lean{TNLean.PEPS.fundamentalTheorem_normalPEPS}
    \leanok
    % Reading note: this is the general normal theorem stated after Theorem 3
    % of \cite{Molnar2018NormalPEPS} on a connected graph.  The shared
    % blockings and the single-crossing hypothesis are the source's blocking
    % hypothesis made explicit, not restrictions; both readings are recorded
    % in \path{docs/paper-gaps/peps_normal_ft_section3_route.tex}.
    \uses{def:peps_normalPEPSBlockingHypotheses,
        def:peps_regionInjectivityDataPair, def:GaugeEquivPEPS}
    Let $A$ and $B$ be PEPS tensors on a connected finite simple graph with
    the same positive physical dimension and positive bond dimensions,
    generating the same state. Suppose the normal blocking
    hypotheses hold at the doubly injective assignment of the two tensors:
    around every edge the network blocks into three doubly injective regions
    with the edge endpoints in the first two, and every site admits two doubly
    injective regions, with doubly injective complements, differing exactly at
    that site. Suppose moreover that at every edge the distinguished edge is
    the only edge of the graph joining its two adjacent blocks. Then the
    defining tensors are related by a local gauge: there are invertible
    matrices, one on each edge, whose endpoint actions transform $A_v$ into
    $B_v$ at every vertex, with no leftover scalar.
    The shared blocking gives a common tripartition for the two tensors, with
    all three regions injective for both. For an edge $e$ between the first
    two regions, the coefficient transfer $T_e$ supplied by the three-site
    isomorphism argument satisfies
    \begin{align}
        C_A(R_1,e,X;\sigma,\tau)
        &=C_B(R_1,e,T_e(X);\sigma,\tau),
        \notag\\
        T_e(X)
        &=Z_e\,\iota_e(X)\,Z_e^{-1},
        \notag
    \end{align}
    for an invertible $Z_e$ and every virtual insertion $X$ on $e$.
    The single-crossing hypothesis identifies $e$ with the entire bond between
    $R_1$ and $R_2$, so $Z_e$ acts on the original edge.
    The equality of the bond dimensions is not assumed; it follows from the
    blocking hypotheses by
    Theorem~\ref{thm:peps_bondDim_eq_of_normalBlocking}. Positive bond
    dimensions are required in the injectivity argument, while connectivity
    permits the per-vertex constants to be absorbed into the edge gauges; see
    Theorem~\ref{thm:fundamentalTheorem_PEPS}. This is the blocking argument
    of \cite[Section~4]{Molnar2018NormalPEPS}.
    % Source-faithful reading: Section 4 of \cite{Molnar2018NormalPEPS} uses
    % one common blocking of the two states and enlarges the endpoint regions
    % so that they meet only along the distinguished edge.
    \begin{proof}
        \leanok
        \uses{thm:peps_bondDim_eq_of_normalBlocking,
            lem:peps_injectiveRegion_union_pos,
            thm:peps_region_edge_gauge_of_coeff_transfer,
            thm:gaugeConsistency}
        The bond dimensions of the two tensors agree on every edge by
        Theorem~\ref{thm:peps_bondDim_eq_of_normalBlocking}.
        The union closure of injective regions for either tensor follows from
        the positive bond dimensions and
        Lemma~\ref{lem:peps_injectiveRegion_union_pos}. At every edge the
        shared blocking datum restricts to a datum for each tensor over the
        same three regions, so the per-edge gauge of the region insertion
        transfer (Theorem~\ref{thm:peps_region_edge_gauge_of_coeff_transfer})
        and its orientation-adapted absorption give a family of invertible
        matrices realizing the absorbed inserted-coefficient equality at every
        edge. A gauge preserves blocked-region linear independence, so at
        every site the one-site comparison applies to the gauge-absorbed
        tensor: comparing the smaller region with its one-site completion
        yields two proportionality scalars whose ratio is a nonzero scalar
        $\lambda_v$ with $A_v$ equal to $\lambda_v$ times the gauge action on
        $B$ at $v$. When the smaller region is empty its block is the count
        of global virtual configurations, and when the completion is the full
        vertex set its block is the closed state coefficient; in these
        degenerate cases the proportionality holds with scalar one, and every
        other comparison region has a boundary edge by connectivity. The
        per-vertex relations against the nonzero closed state force
        $\prod_v \lambda_v = 1$, so on the connected graph a spanning-tree
        construction realizes the reciprocals of the scalars as oriented
        incidence products of edge scalars, which fold into the gauges as in
        the absorption step of the injective Fundamental Theorem
        (Theorem~\ref{thm:gaugeConsistency}).
    \end{proof}
\end{theorem}

\begin{theorem}[Uniqueness of the normal-PEPS gauges up to a constant]%
    \label{thm:fundamentalTheorem_normalPEPS_gauge_unique}
    \lean{TNLean.PEPS.fundamentalTheorem_normalPEPS_gauge_unique}
    \leanok
    \uses{def:peps_normalPEPSBlockingHypotheses,
        def:peps_regionInjectivityDataPair, def:gaugeVertex}
    Let $A$ and $B$ be PEPS tensors on a finite simple graph with equal bond
    dimensions, the bond dimensions of the first tensor positive, satisfying
    the normal blocking hypotheses at the doubly injective assignment of the
    two tensors. Suppose two families of invertible matrices, one on each
    edge, each realize the scalar-free gauge relation of
    Theorem~\ref{thm:fundamentalTheorem_normalPEPS_connected}: at every vertex
    the endpoint actions of the family transform $A_v$ into $B_v$. Then the
    two families are proportional at every edge: $X'_e = c\cdot X_e$ for a
    nonzero scalar $c$ depending on the edge. This is the uniqueness clause
    of the general normal-PEPS theorem in
    \cite[Section~4]{Molnar2018NormalPEPS}, following its Theorem~3, read
    against the scalar-free relation; neither equality of the two states nor
    connectivity is needed.
    \begin{proof}
        \leanok
        \uses{thm:peps_gaugeConj_eq_of_coeffIdentities,
            lem:peps_injectiveRegion_union_pos}
        The inserted-edge coefficient is multilinear in the site tensors, one
        factor per site, so each scalar-free per-vertex relation makes every
        inserted coefficient of $B$ equal the corresponding inserted
        coefficient of the gauge-absorbed $A$: each family realizes the
        absorbed inserted-coefficient equality at every edge, with the roles
        of the two tensors exchanged. On the first block of the edge's frame,
        whose blocked map and complement blocked map are injective for $A$ by
        the blocking hypotheses and union closure
        (Lemma~\ref{lem:peps_injectiveRegion_union_pos}), the region insertion
        transfer is determined by the absorbed equality
        (Theorem~\ref{thm:peps_gaugeConj_eq_of_coeffIdentities}), so the two
        orientation-adapted absorbing gauges induce the same conjugation and
        differ by a nonzero scalar; unwinding the orientation adaptation gives
        the proportionality of the gauges themselves.
    \end{proof}
\end{theorem}

\section{Fundamental Theorem for normal matrix product states}%
\label{sec:peps_normal_mps}

The closed chain is the one-dimensional PEPS: the cycle graph on $n$ vertices
carrying one tensor of $D \times D$ matrices per site, so that a physical
configuration $\sigma = (i_1, \ldots, i_n)$ receives the amplitude
$\tr(A^{i_1} \cdots A^{i_n})$. The Fundamental Theorem for normal PEPS
therefore specializes to matrix product states once its blocking hypotheses are
read off the chain, and this section records the resulting statements,
following the matrix-product-state corollaries of
\cite[Section~4]{Molnar2018NormalPEPS} and the sharper system-size bound in its
closing part \emph{Normal MPS: alternative proof}.

The development proceeds in four steps. The closed chain is first identified
with a cycle of tensors, and the normal blocking hypotheses are obtained from
the injectivity of every arc of $L$ consecutive sites, giving the Fundamental
Theorem for normal matrix product states on a ring of $n \ge 3L$ sites with one
invertible gauge per bond. The translation-invariant case then collapses the
per-bond gauges to one gauge $Z$ and one scalar $\lambda$, with the length-$n$
trace identity forcing $\lambda^n = 1$.

For the sharper bound, the matrix-product-state form of
\cite[Lemma~5]{Molnar2018NormalPEPS} compares the $L + 1$ overlapping
windows around a fixed bond. If $C_j$ is the deformation of the $j$-th
window, then for length-$L$ words $a$ and $b$ the family of
consecutive-window trace identities strips to
$C_0(a)\,B^b = B^a\,C_L(b)$, and hence to
$C_0(a) = B^a X$ and $C_L(b) = X B^b$.
This is the step that lowers the system size from $n \ge 3L$ to
$n \ge 2L + 1$. Finally the same argument is run for site-dependent
chains, one tensor per site, which is the setting of
\cite[Lemma~5]{Molnar2018NormalPEPS}.

\subsection{Matrix product states on the closed chain}

\begin{theorem}[Cycle blocking hypotheses from arc injectivity]%
    \label{thm:peps_cycleNormalPEPSBlockingHypotheses}
    \lean{TNLean.PEPS.cycleArcFrom}
    \lean{TNLean.PEPS.NormalCycleArcInjectivityHypotheses}
    \lean{TNLean.PEPS.cycleEdgeBlockingData}
    \lean{TNLean.PEPS.cycleOneSiteSeparation}
    \lean{TNLean.PEPS.cycleNormalPEPSBlockingHypotheses}
    \lean{TNLean.PEPS.isCrossingEdge_cycleEdge_iff}
    \leanok
    \uses{def:peps_normalPEPSBlockingHypotheses,
        def:peps_regionInjectivityData}
    A closed chain of $n$ sites is the cycle graph on $n$ vertices, and a
    block of consecutive sites is an arc: the vertices reached clockwise
    from a starting site. Suppose $0 < L$, $3L \le n$, every arc
    of $L$ consecutive sites is injective for a region-injectivity
    assignment, and injectivity is closed under unions. Then the chain
    satisfies the normal blocking hypotheses: around every edge it blocks
    into the $L$ sites on one side, the $L$ sites on the other, and the
    remaining $n - 2L \ge L$ sites, all injective and covering the chain;
    and at every site the $L + 1$ sites starting there and the $L$ sites
    starting at its successor are injective regions with injective
    complements differing exactly at that site. Moreover the only
    nearest-neighbour pair joining the two blocks around an edge is the edge
    itself, since the far ends of the two blocks are separated by the
    remaining $n - 2L \ge L$ sites.
\end{theorem}
\begin{proof}\leanok
    Membership in an arc is a cyclic-distance inequality, so the
    disjointness, covering, and unique-crossing claims are modular
    arithmetic on the distances. Arcs of length between $L$ and $n$ are
    unions of arcs of length $L$, hence injective by union closure; this
    covers the two complements at a site and the third block around an
    edge.
\end{proof}

\begin{corollary}[Fundamental Theorem for normal MPS on a closed chain]%
    \label{cor:fundamentalTheorem_normalMPS_cycle}
    \lean{TNLean.PEPS.fundamentalTheorem_normalMPS_cycle}
    \leanok
    \uses{def:peps_regionInjectivityDataOf, def:GaugeEquivPEPS}
    Let $A$ and $B$ be tensors on the closed chain of $n$ sites (the cycle
    graph on $n$ vertices, one tensor per site with two virtual legs), with
    the same positive physical dimension and positive bond dimensions, and
    let $0 < L$ with $n \ge 3L$. Suppose blocking any $L$ consecutive sites
    gives an injective tensor, for $A$ and for $B$, and the two MPS generate
    the same state. Then the defining tensors are related by a local gauge:
    there are invertible matrices, one on each bond, whose endpoint actions
    transform $A_i$ into $B_i$ at every site with no leftover scalar. This
    is the source's family $Z_i$ ($i = 1, \ldots, n$, $n + 1 \equiv 1$) with
    $B_i = Z_i^{-1} A_i Z_{i+1}$, the first corollary following the general
    normal-PEPS theorem in \cite[Section~4]{Molnar2018NormalPEPS}, the
    statement for MPS without translation invariance at a fixed system size.
    The positive physical and bond dimensions are the counterexample-backed
    corrections carried over from the injective Fundamental Theorem
    (Theorem~\ref{thm:fundamentalTheorem_PEPS}), and $0 < L$ is implicit in
    the source's blocking of $L$ consecutive sites; the equality of the bond
    dimensions is not assumed but derived, and the connectivity of the cycle
    is proved, not assumed.
\end{corollary}
\begin{proof}\leanok
    \uses{thm:peps_cycleNormalPEPSBlockingHypotheses,
        thm:fundamentalTheorem_normalPEPS_connected,
        lem:peps_injectiveRegion_union_pos}
    Positive bond dimensions give union closure of injective regions for
    each tensor (Lemma~\ref{lem:peps_injectiveRegion_union_pos}), hence for
    the doubly injective assignment. The arc-injectivity hypothesis then
    yields the normal blocking hypotheses on the cycle with the unique
    crossing property
    (Theorem~\ref{thm:peps_cycleNormalPEPSBlockingHypotheses}), the cycle is
    connected, and the Fundamental Theorem for normal PEPS on a connected
    graph (Theorem~\ref{thm:fundamentalTheorem_normalPEPS_connected})
    supplies the scalar-free local gauge.
\end{proof}

\begin{corollary}[Uniqueness of the closed-chain gauges up to a constant]%
    \label{cor:fundamentalTheorem_normalMPS_cycle_gauge_unique}
    \lean{TNLean.PEPS.fundamentalTheorem_normalMPS_cycle_gauge_unique}
    \leanok
    \uses{def:peps_regionInjectivityDataOf, def:gaugeVertex}
    In the setting of
    Corollary~\ref{cor:fundamentalTheorem_normalMPS_cycle}, with the bond
    dimensions of the two tensors equal and positive and the arcs of $L$
    consecutive sites injective for each tensor, any two families of
    invertible bond matrices realizing the scalar-free gauge relation are
    proportional at every bond: $Z'_i = c \cdot Z_i$ for a nonzero scalar
    $c$ depending on the bond. This is the uniqueness clause of the first
    corollary following the general normal-PEPS theorem of
    \cite[Section~4]{Molnar2018NormalPEPS}: the gauges $Z_i$ are unique up
    to a multiplicative constant. Neither equality of the two states nor
    connectivity is needed, as in the general uniqueness clause.
\end{corollary}
\begin{proof}\leanok
    \uses{thm:peps_cycleNormalPEPSBlockingHypotheses,
        thm:fundamentalTheorem_normalPEPS_gauge_unique,
        lem:peps_injectiveRegion_union_pos}
    The arc-injectivity hypothesis and union closure from the positive bond
    dimensions (Lemma~\ref{lem:peps_injectiveRegion_union_pos}) assemble the
    normal blocking hypotheses on the cycle
    (Theorem~\ref{thm:peps_cycleNormalPEPSBlockingHypotheses}), and the
    per-edge uniqueness of the general theorem
    (Theorem~\ref{thm:fundamentalTheorem_normalPEPS_gauge_unique}) gives the
    proportionality at every bond.
\end{proof}

\begin{definition}[Cycle tensor of a matrix tensor]%
    \label{def:cycleTensorOfMPS}
    \lean{TNLean.PEPS.cycleSuccEdge}
    \lean{TNLean.PEPS.cycleLeftIncident}
    \lean{TNLean.PEPS.cycleRightIncident}
    \lean{TNLean.PEPS.cycleTensorOfMPS}
    \leanok
    \uses{def:peps_tensor, def:mps_tensor}
    Every edge of the cycle graph on $n \ge 3$ vertices joins a site to its
    cyclic successor, so sites and bonds are in bijection, and every site has
    exactly two incident bonds: the bond shared with its cyclic predecessor
    and the bond shared with its cyclic successor. The \emph{cycle tensor}
    of a tensor $A = (A^i)_{i=0}^{d-1}$ of $D \times D$ matrices places one
    copy of $A$ at every site of the closed chain of $n$ sites, every bond of
    dimension $D$, with the row index of $A^i$ read from the bond shared with
    the predecessor and the column index from the bond shared with the
    successor.
\end{definition}

\begin{lemma}[State coefficients of the cycle tensor]%
    \label{lem:stateCoeff_cycleTensorOfMPS}
    \lean{MPSTensor.evalWord_ofFn_apply}
    \lean{MPSTensor.trace_evalWord_eq_sum_cyclic}
    \lean{TNLean.PEPS.stateCoeff_cycleTensorOfMPS}
    \leanok
    \uses{def:cycleTensorOfMPS, def:peps_stateCoeff, def:mpv}
    For $n \ge 3$ and every configuration
    $\sigma = (i_1, \ldots, i_n)$, the state coefficient of the cycle
    tensor of $A$ at $\sigma$ is the matrix product vector component
    $V^{(n)}(A)_\sigma = \tr\!(A^{i_1} \cdots A^{i_n})$.
\end{lemma}
\begin{proof}\leanok
    An entry of a word product is a sum over paths of bond indices pinned at
    the two ends of products of matrix entries read along the path, by
    induction on the word; closing the two pinned ends into a loop, the trace
    of a closed word product is the sum over all cyclic bond configurations
    of the products of entries around the loop. Contracting the cycle tensor
    sums the per-site entries over all assignments of bond indices, which is
    that cyclic sum after matching each bond with the site on its
    predecessor side.
\end{proof}

\begin{lemma}[Arc injectivity of the cycle tensor from block injectivity]%
    \label{lem:regionBlockedTensorInjective_cycleTensorOfMPS}
    \lean{TNLean.PEPS.isRegionBoundaryEdge_cycleArcFrom_iff}
    \lean{TNLean.PEPS.regionBlockedWeight_cycleTensorOfMPS}
    \lean{TNLean.PEPS.matrix_eq_zero_of_isNBlkInjective}
    \lean{TNLean.PEPS.regionBlockedTensorInjective_cycleTensorOfMPS}
    \leanok
    \uses{def:cycleTensorOfMPS, def:l_blk_injective,
        def:peps_regionInjectivityDataOf}
    Let $0 < L < n$ with $n \ge 3$, let the bond dimension be positive, and
    suppose $A$ is $L$-block injective. Then for every arc of $L$
    consecutive sites of the closed chain, the blocked tensor of the arc of
    the cycle tensor of $A$ is injective.
\end{lemma}
\begin{proof}\leanok
    An arc of $L < n$ consecutive sites has exactly two boundary-crossing
    bonds, the bond entering its first site and the bond leaving its last
    site, by cyclic-distance arithmetic. The blocked weight of the arc at a
    boundary assignment $(a, b)$ and a physical word $x$ is
    $D^{\,n - (L+1)}\,(A^{x_1} \cdots A^{x_L})_{a b}$: the bonds inside
    the arc contract to the matrix product, the two boundary bonds stay
    pinned, and each bond away from the arc contributes a free factor $D$.
    A linear relation among the blocked components over the boundary pairs
    is a matrix $C$ with $\tr\!(C^{\mathsf T} A^{x_1} \cdots A^{x_L}) = 0$
    for every word $x$; the word products span the full matrix algebra by
    $L$-block injectivity, and the trace pairing is nondegenerate, so
    $C = 0$.
\end{proof}

\subsection{Translation-invariant matrix product states}

When the same tensor sits at every site, the closed chain ties the per-bond
gauges together. If one gauge $Z$ and one scalar $\lambda$ satisfy
$B^i = \lambda\,Z^{-1}A^iZ$ for every physical index $i$, then for a
configuration $\sigma=(i_1,\ldots,i_n)$,
\begin{align}
    \tr(B^{i_1}\cdots B^{i_n})
    &= \lambda^n\,\tr(Z^{-1}A^{i_1}\cdots A^{i_n}Z)
     = \lambda^n\,\tr(A^{i_1}\cdots A^{i_n}).
    \notag
\end{align}
Equality of the matrix product states therefore forces $\lambda^n = 1$.
Thus $Z$ and $\lambda$ capture precisely the freedom in the
translation-invariant corollary of \cite[Section~4]{Molnar2018NormalPEPS}.

\begin{corollary}[Fundamental Theorem for translation-invariant normal MPS
        on a closed chain, matrix form]%
    \label{cor:fundamentalTheorem_normalMPS}
    \lean{TNLean.PEPS.fundamentalTheorem_normalMPS}
    \leanok
    \uses{def:l_blk_injective, def:mpv}
    Let $n$ be positive, let $0 < L$ with $2L + 1 \le n$, and let $A$ and $B$
    be tensors of $D \times D$ matrices with $D>0$, over physical dimension
    $d$, each
    $L$-block injective, whose matrix product vectors agree at system size
    $n$: $V^{(n)}(A)_\sigma = V^{(n)}(B)_\sigma$ for every configuration
    $\sigma$. Then there are invertible matrices $Z_v$, one per bond of the
    closed chain ($v = 1, \ldots, n$, $n + 1 \equiv 1$), with
    $B^i = Z_v^{-1}\, A^i\, Z_{v+1}$ for every site $v$ and every $i$.
    This is the first corollary following the general normal-PEPS theorem of
    \cite[Section~4]{Molnar2018NormalPEPS} specialized to one
    site-independent tensor: the source states the corollary for
    site-dependent families, and its conclusion is the per-bond gauge family
    delivered here. The system size is the optimal $n \ge 2L + 1$ of the
    source's alternative proof (its Section ``Normal MPS: alternative
    proof''), rather than the $n \ge 3L$ of the blocking route. The source's
    translation-invariant corollary (a single gauge $Z$ and a constant
    $\lambda$ with $\lambda^n = 1$) refines this conclusion and is not part
    of this statement. A vanishing physical dimension is already excluded
    by block injectivity.
\end{corollary}
\begin{proof}\leanok
    \uses{cor:fundamentalTheorem_normalMPS_of_overlap}
    The per-bond form of the strengthened closed-chain corollary
    (Corollary~\ref{cor:fundamentalTheorem_normalMPS_of_overlap}), proved
    through the overlapping windows at $n \ge 2L + 1$, gives the per-bond
    gauge family directly.
\end{proof}

\begin{corollary}[Uniqueness of the per-bond gauges up to one constant,
        matrix form]%
    \label{cor:fundamentalTheorem_normalMPS_gauge_unique}
    \lean{TNLean.PEPS.fundamentalTheorem_normalMPS_gauge_unique}
    \leanok
    \uses{def:l_blk_injective}
    Let $n$ be positive, let $0 < L$ with $3L \le n$, and let $A$ and $B$ be
    tensors of $D \times D$ matrices with $D>0$, each $L$-block injective.
    If two
    families $Z$ and $Z'$ of invertible matrices on the bonds of the closed
    chain both satisfy $B^i = Z_v^{-1} A^i Z_{v+1}$ at every site, then they
    are proportional by a single constant: there is a nonzero scalar $c$
    with $Z'_v = c\,Z_v$ at every bond. This is the uniqueness clause of
    the first corollary following the general normal-PEPS theorem of
    \cite[Section~4]{Molnar2018NormalPEPS}, for one site-independent
    tensor. Equality of the two states is not needed.
\end{corollary}
\begin{proof}\leanok
    \uses{cor:fundamentalTheorem_normalMPS_cycle_gauge_unique,
        lem:regionBlockedTensorInjective_cycleTensorOfMPS,
        def:cycleTensorOfMPS, def:gaugeVertex}
    Each per-bond family converts back into a per-edge family - the
    transpose of the bond gauge away from the seam, the inverse on the seam
    edge - realizing the per-site gauge identity between the cycle tensors,
    so the closed-chain uniqueness clause
    (Corollary~\ref{cor:fundamentalTheorem_normalMPS_cycle_gauge_unique})
    gives a constant on every edge; undoing the conversion gives a constant
    per bond. The relation $B^i = Z_v^{-1} A^i Z_{v+1}$ for both families
    forces the scalar $\mu_v^{-1}\mu_{v+1}$ to fix the nonzero tensor $B$,
    so consecutive constants agree, and going around the cycle merges them
    into one.
\end{proof}

\begin{lemma}[Consecutive per-bond gauges are proportional]%
    \label{lem:gaugeFamily_succ_proportional}
    \lean{TNLean.PEPS.gaugeFamily_succ_proportional}
    \leanok
    \uses{def:l_blk_injective, def:eval_word}
    Let $n$ be positive, let $A$ and $B$ be tensors of $D \times D$ matrices
    over physical dimension $d$ with $A$ being $L$-block injective, and let
    $Z_v$ ($v = 1, \ldots, n$, $n + 1 \equiv 1$) be invertible matrices with
    $B^i = Z_v^{-1} A^i Z_{v+1}$ at every site $v$ and every $i$. Then for
    every site $v$ there is a nonzero scalar $c$ with $Z_{v+1} = c\,Z_v$.
\end{lemma}
\begin{proof}\leanok
    \uses{def:eval_word}
    Iterating the per-bond relation along a word $x$ of length $L$ from the
    starting site $v$ gives
    $B^{x_1} \cdots B^{x_L} = Z_v^{-1} A^{x_1} \cdots A^{x_L} Z_{v+L}$.
    Comparing the iterations from $v$ and from $v + 1$ over the word
    products $A^{x_1} \cdots A^{x_L}$, which span the full matrix algebra by
    $L$-block injectivity, the comparison extends linearly to every matrix
    $W$: $Z_v^{-1} W Z_{v+L} = Z_{v+1}^{-1} W Z_{v+1+L}$. The empty word
    ($W$ the identity) pins the two bond transports
    $Z_v^{-1} Z_{v+L} = Z_{v+1}^{-1} Z_{v+1+L}$ to the same invertible
    matrix; cancelling it leaves $Z_v^{-1} W Z_v = Z_{v+1}^{-1} W Z_{v+1}$
    for every $W$, so $Z_{v+1} Z_v^{-1}$ commutes with the full matrix
    algebra, hence is a scalar, nonzero by invertibility.
\end{proof}

\begin{corollary}[Fundamental Theorem for translation-invariant normal MPS
        on a closed chain, single-gauge form]%
    \label{cor:fundamentalTheorem_normalMPS_translationInvariant}
    \lean{TNLean.PEPS.fundamentalTheorem_normalMPS_translationInvariant}
    \leanok
    \uses{def:l_blk_injective, def:mpv}
    Let $n$ be positive, let $0 < L$ with $2L + 1 \le n$, and let $A$ and $B$
    be tensors of $D \times D$ matrices with $D>0$, over physical dimension
    $d$, each
    $L$-block injective, whose matrix product vectors agree at system size
    $n$: $V^{(n)}(A)_\sigma = V^{(n)}(B)_\sigma$ for every configuration
    $\sigma$. Then there are an invertible matrix $Z$ and a constant
    $\lambda$ with $\lambda^n = 1$ such that
    $B^i = \lambda\, Z^{-1} A^i Z$ for every $i$.
    This is the translation-invariant corollary of
    \cite[Section~4]{Molnar2018NormalPEPS}, the statement for TI MPS
    following the first corollary after the general normal-PEPS theorem. The
    system size is the optimal $n \ge 2L + 1$ of the source's alternative
    proof.
\end{corollary}
\begin{proof}\leanok
    \uses{cor:fundamentalTheorem_normalMPS_translationInvariant_of_overlap}
    This is the single-gauge form of the strengthened closed-chain corollary
    (Corollary~\ref{cor:fundamentalTheorem_normalMPS_translationInvariant_of_overlap}),
    proved through the overlapping windows at $n \ge 2L + 1$.
\end{proof}

\begin{corollary}[Uniqueness of the translation-invariant gauge up to a
        constant]%
    \label{cor:fundamentalTheorem_normalMPS_translationInvariant_gauge_unique}
    \lean{TNLean.PEPS.fundamentalTheorem_normalMPS_translationInvariant_gauge_unique}
    \leanok
    \uses{def:l_blk_injective}
    Let $0 < L$ and $0<D$, and let $A$ and $B$ be tensors of $D \times D$ matrices,
    each $L$-block injective. If invertible matrices $Z$, $Z'$ and
    constants $\lambda$, $\lambda'$ satisfy
    $B^i = \lambda\, Z^{-1} A^i Z$ and $B^i = \lambda'\, Z'^{-1} A^i Z'$ for
    every $i$, then there is a nonzero scalar $c$ with $Z' = c\,Z$. This is
    the uniqueness clause of the translation-invariant corollary of
    \cite[Section~4]{Molnar2018NormalPEPS} for TI MPS. No system size
    enters the statement: equality of the two states is not needed, nor are
    the root-of-unity conditions on $\lambda$ and $\lambda'$.
\end{corollary}
\begin{proof}\leanok
    \uses{def:eval_word}
    The scalar $\lambda$ is nonzero because $B$ is $L$-block injective with
    positive bond dimension, so some $B^{i_0} \neq 0$. Iterating each relation
    along a word $x$ of length $L$ gives
    $B^{x_1} \cdots B^{x_L} = \lambda^L\, Z^{-1} A^{x_1} \cdots A^{x_L} Z$
    and the primed analog; the word products span the full matrix algebra
    by $L$-block injectivity, so
    $\lambda^L\, Z^{-1} W Z = \lambda'^L\, Z'^{-1} W Z'$ for every matrix
    $W$. The identity pins $\lambda^L = \lambda'^L$; cancelling the
    nonzero factor leaves $Z^{-1} W Z = Z'^{-1} W Z'$ for every $W$, so
    $Z' Z^{-1}$ commutes with the full matrix algebra, hence is a nonzero
    scalar.
\end{proof}

\subsection{\texorpdfstring{The optimal system size $n \ge 2L + 1$}{The optimal system size n >= 2L+1}}

The bound $n \ge 3L$ is not optimal. A bond of the closed chain is straddled by
the $L + 1$ overlapping windows of $L$ sites around it. For consecutive
windows, the deformed closed-chain coefficients have the form
\begin{align}
    \tr(B^{p}\, C_j(u_1 \cdots u_L)\, B^{u_{L+1}}\, B^{s})
    &=
    \tr(B^{p}\, B^{u_1}\, C_{j+1}(u_2 \cdots u_{L+1})\, B^{s}).
    \notag
\end{align}
Injectivity on the complementary arc strips these identities to
\begin{align}
    C_j(u_1 \cdots u_L)\,B^{u_{L+1}}
    &=
    B^{u_1}\,C_{j+1}(u_2 \cdots u_{L+1}),
    \notag
\end{align}
and chaining the relations across the $L + 1$ windows gives
$C_0(a)\,B^b = B^a\,C_L(b)$.
This is the overlap mechanism in \cite[Lemma~5]{Molnar2018NormalPEPS}.
The lemmas below assemble that argument: block injectivity extends to every
length, word products transport between equal-state chains, and the bond
operator is extracted and shown to act by conjugation, lowering the system size
to $n \ge 2L + 1$.

\begin{lemma}[Block injectivity extends to longer blocks]%
    \label{lem:isNBlkInjective_of_le}
    \lean{MPSTensor.isNBlkInjective_of_le}
    \leanok
    \uses{def:l_blk_injective}
    Let $A$ be a tensor of $D \times D$ matrices over physical dimension
    $d$, $L$-block injective with $0 < L$. Then $A$ is $m$-block injective
    for every $m \ge L$. This is the chain form of the statement that any
    region of at least the blocking size is injective, derived in the
    closing section of \cite{Molnar2018NormalPEPS} (``Normal MPS:
    alternative proof'') from its union lemma.
\end{lemma}
\begin{proof}\leanok
    \uses{def:eval_word}
    It suffices to pass from length $m$ to $m + 1$. Write the identity as
    a combination of length-$L$ word products and peel the leading letter
    of each: any matrix $P$ becomes a combination of terms
    $A^{i}\,(W P)$ with $W$ a word product of length $L - 1$. Each factor
    $W P$ is a combination of length-$m$ word products by the induction
    hypothesis, so each term is a combination of length-$(m+1)$ word
    products.
\end{proof}

\begin{lemma}[Word transport between same-state networks]%
    \label{lem:exists_mpvTransport}
    \lean{MPSTensor.exists_mpvTransport}
    \lean{MPSTensor.exists_evalWordTransport}
    \leanok
    \uses{def:l_blk_injective, def:mpv, def:eval_word}
    Let $A$ and $B$ be tensors of $D \times D$ matrices over physical
    dimension $d$, let $n = k + q$, and suppose the length-$k$ and
    length-$q$ word products of $A$ each span the full matrix algebra. If
    the matrix product vectors of $A$ and $B$ agree at system size $n$,
    then there is a linear map $\Lambda$ of the matrix algebra with
    $\Lambda(A^{x_1} \cdots A^{x_k}) = B^{x_1} \cdots B^{x_k}$ for every
    word $x$ of length $k$. The same holds when, instead of the traces,
    the word products of the two tensors are themselves equal at length
    $n$. This is the operator-transport mechanism underlying Lemma~5 of
    \cite{Molnar2018NormalPEPS}: an operator expanded over the window
    products of one network is re-read in the other.
\end{lemma}
\begin{proof}\leanok
    \uses{def:mpv, def:eval_word}
    Pair the matrix algebra against the complementary word products: let
    $\Psi_A(M)$ be the family of traces
    $\tr(M\,A^{y_1} \cdots A^{y_q})$ over the length-$q$ words $y$, and
    $\Psi_B$ its $B$-analog. $\Psi_A$ is injective because the length-$q$
    products span and the trace pairing is nondegenerate. On the spanning
    family of length-$k$ products the two pairings match:
    $\Psi_A(A^{x}) = \Psi_B(B^{x})$, since both compute closed-chain
    coefficients of total length $n$, equal by hypothesis. Hence the range
    of $\Psi_A$ lies in that of $\Psi_B$, forcing $\Psi_B$ to be injective
    by dimension count, and a left inverse of $\Psi_B$ composed with
    $\Psi_A$ is the transport. For the matrix-product version replace the
    traces by the products themselves.
\end{proof}

\begin{lemma}[Bond-operator extraction from overlapping windows]%
    \label{lem:overlapWindow_exists_bondOperator}
    \lean{MPSTensor.overlapWindow_exists_bondOperator}
    \lean{MPSTensor.bondOperator_unique}
    \leanok
    \uses{def:l_blk_injective, def:eval_word}
    Let $0 < L$ with $2L + 1 \le n$, let $B$ be a tensor of $D \times D$
    matrices over physical dimension $d$, $L$-block injective, and let
    $C_0, \ldots, C_L$ be deformed window tensors: each $C_j$ assigns a
    matrix to every word of length $L$, read as replacing the blocked
    tensor of the $L$-site window starting $j$ sites left of a fixed bond
    of the closed chain on $n$ sites. Suppose the deformed states agree
    for consecutive window positions: for every $j < L$, every prefix $p$
    of $j$ sites, every assignment $u$ on the $L + 1$ sites covered by the
    two windows, and every suffix $s$ on the remaining sites,
    \begin{align}
        \tr(B^{p}\, C_j(u_1 \cdots u_L)\, B^{u_{L+1}}\, B^{s})
        &=
        \tr\bigl(B^{p}\, B^{u_1}\, C_{j+1}(u_2 \cdots u_{L+1})
        B^{s}\bigr).
        \notag
    \end{align}
    Then the deformation is a virtual operation on the bond: there is a
    matrix $X$ such that, for all length-$L$ words $a$ and $b$,
    \begin{align}
        C_0(a) &= B^{a_1} \cdots B^{a_L}\, X,
        \notag\\
        C_L(b) &= X\, B^{b_1} \cdots B^{b_L},
        \notag
    \end{align}
    and $X$ is unique. This is Lemma~5
    of \cite{Molnar2018NormalPEPS} (its closing section, ``Normal MPS:
    alternative proof''), specialized to one site-independent tensor; the
    source states it for site-dependent tensors on five sites with $L = 2$,
    with the deformed window tensors arising as physical operators
    contracted with the blocked windows they act on; the site-dependent
    form, in the setting in which the source states the lemma, is
    Lemma~\ref{lem:chain_overlapWindow_exists_bondOperator}, and this
    statement is its constant specialization. The
    algebra-homomorphism clause of the source lemma is established with the
    insertion correspondence in
    Lemma~\ref{lem:exists_conjugation_of_mpv_eq}.
\end{lemma}
\begin{proof}\leanok
    \uses{lem:chain_overlapWindow_exists_bondOperator, def:eval_word}
    Place the tensor at every site of the site-dependent form
    (Lemma~\ref{lem:chain_overlapWindow_exists_bondOperator}): the arc
    products of the constant chain are the word products of the repeated
    tensor, and $L$-block injectivity of the tensor is window injectivity
    of the constant chain.
\end{proof}

\begin{lemma}[Conjugate word products at one length from Lemma~5]%
    \label{lem:exists_conjugation_of_mpv_eq}
    \lean{MPSTensor.exists_conjugation_of_mpv_eq}
    \leanok
    \uses{def:l_blk_injective, def:mpv, def:eval_word}
    Let $0 < L$ with $2L + 1 \le n$ and $0<D$, and let $A$ and $B$ be tensors of
    $D \times D$ matrices over physical dimension $d$, each $L$-block
    injective, whose matrix product vectors agree at system size $n$.
    Then there is an invertible matrix $Z$ such that, for every word $w$ of
    length $n$,
    \begin{align}
        B^{w_1} \cdots B^{w_n}
        &= Z\, A^{w_1} \cdots A^{w_n}\, Z^{-1}.
        \notag
    \end{align}
    This captures the insertion
    correspondence of Lemma~5 of \cite{Molnar2018NormalPEPS} and its
    algebra-homomorphism clause for one site-independent tensor: the map
    sending an insertion on one network to the corresponding insertion on
    the other is an automorphism of the matrix algebra, hence inner.
\end{lemma}
\begin{proof}\leanok
    \uses{cor:fundamentalTheorem_normalMPSChain_of_overlap,
        thm:mps_chain_gaugeEquiv_sameState}
    Regard $A$ and $B$ as constant site-dependent chains on $n$ sites.
    Window injectivity of a constant chain is block injectivity of the
    single tensor, and agreement of the matrix product vectors at size $n$
    is the same-state hypothesis, so the site-dependent overlapping-window
    capstone
    (Corollary~\ref{cor:fundamentalTheorem_normalMPSChain_of_overlap})
    supplies a per-bond gauge family relating the two chains. The arc
    products of a constant chain are the word products of the tensor, and
    the gauge conjugation of arc products
    (Theorem~\ref{thm:mps_chain_gaugeEquiv_sameState}), read at the full
    chain length $n$ where the end bond wraps to the starting bond, is
    conjugation of the length-$n$ word products by the gauge at the
    starting bond.
\end{proof}

\begin{lemma}[Rigidity of word products at one length]%
    \label{lem:proportional_of_evalWord_eq}
    \lean{MPSTensor.proportional_of_evalWord_eq}
    \leanok
    \uses{def:l_blk_injective, def:eval_word}
    Let $0 < L$ with $2L + 1 \le n$ and $0 < D$, and let $A$ and $C$ be
    tensors of $D \times D$ matrices over physical dimension $d$, each
    $L$-block injective, with equal word products at length $n$:
    $C^{w_1} \cdots C^{w_n} = A^{w_1} \cdots A^{w_n}$ for every word $w$ of
    length $n$. Then there is a constant $\lambda$ with $\lambda^n = 1$
    and $C^i = \lambda A^i$ for every $i$.
\end{lemma}
\begin{proof}\leanok
    \uses{lem:exists_mpvTransport, lem:isNBlkInjective_of_le,
        def:eval_word}
    Transport word products of $A$ to those of $C$ at the window length
    $L$, at $L + 1$, at the complementary length $n - 1 - L$, and at
    $n - L$ (Lemma~\ref{lem:exists_mpvTransport}, matrix-product version;
    the lengths lie between $L$ and $n - L$ by $2L + 1 \le n$). For a
    transport pair at lengths $k$ and $k + 1$, the transported identity
    $\Lambda_k(\Id)$ commutes with every letter of $C$ - its two one-letter
    extensions agree, $C^i\,\Lambda_k(\Id) = \Lambda_{k+1}(A^i) =
        \Lambda_k(\Id)\,C^i$ - hence with the full matrix algebra by block
    injectivity, so it is a scalar $c_k \Id$, nonzero because the transport
    is surjective onto the spanning word products, hence bijective.
    Peeling one letter off the full-length equality through the spanning
    windows of lengths $L$ and $n - 1 - L$ then leaves
    $C^i = (c_L\,c_{n-1-L})^{-1} A^i =: \lambda A^i$; evaluating any
    nonzero length-$n$ word product gives $\lambda^n = 1$.
\end{proof}

\begin{corollary}[Fundamental Theorem for translation-invariant normal MPS
        at $n \ge 2L + 1$, single-gauge form]%
    \label{cor:fundamentalTheorem_normalMPS_translationInvariant_of_overlap}
    \lean{TNLean.PEPS.fundamentalTheorem_normalMPS_translationInvariant_of_overlap}
    \leanok
    \uses{def:l_blk_injective, def:mpv}
    Let $0 < L$ with $2L + 1 \le n$ and $0<D$, and let $A$ and $B$ be tensors of
    $D \times D$ matrices over physical dimension $d$, each $L$-block
    injective, whose matrix product vectors agree at system size $n$:
    $V^{(n)}(A)_\sigma = V^{(n)}(B)_\sigma$ for every configuration
    $\sigma$. Then there are an invertible matrix $Z$ and a constant
    $\lambda$ with $\lambda^n = 1$ such that
    $B^i = \lambda\, Z^{-1} A^i Z$ for every $i$.
    This is the strengthened closed-chain corollary of
    \cite{Molnar2018NormalPEPS} (the corollary after Lemma~5 in its closing
    section ``Normal MPS: alternative proof''), proved through the
    overlapping windows at the optimal $n \ge 2L + 1$ rather than the
    $n \ge 3L$ of the blocking route; the single-gauge corollary
    \ref{cor:fundamentalTheorem_normalMPS_translationInvariant} delegates to
    it. The uniqueness clause of the source
    corollary - the gauge $Z$ is unique up to a multiplicative constant -
    needs no system size and is
    Corollary~\ref{cor:fundamentalTheorem_normalMPS_translationInvariant_gauge_unique}.
\end{corollary}
\begin{proof}\leanok
    \uses{lem:exists_conjugation_of_mpv_eq, lem:proportional_of_evalWord_eq}
    By Lemma~\ref{lem:exists_conjugation_of_mpv_eq} the word products of
    $B$ at length $n$ are those of $A$ conjugated by an invertible matrix.
    Absorbing the conjugation into $B$ - which preserves $L$-block
    injectivity - leaves two tensors with equal word products at length
    $n$, and Lemma~\ref{lem:proportional_of_evalWord_eq} makes them
    proportional letterwise by an $n$-th root of unity. Undoing the
    absorption gives $B^i = \lambda\, Z^{-1} A^i Z$.
\end{proof}

\begin{corollary}[Fundamental Theorem for translation-invariant normal MPS
        at $n \ge 2L + 1$, per-bond form]%
    \label{cor:fundamentalTheorem_normalMPS_of_overlap}
    \lean{TNLean.PEPS.fundamentalTheorem_normalMPS_of_overlap}
    \leanok
    \uses{def:l_blk_injective, def:mpv}
    Let $n$ be positive, let $0 < L$ with $2L + 1 \le n$ and $0<D$, and let
    $A$ and $B$ be tensors of $D \times D$ matrices over physical dimension $d$,
    each $L$-block injective, whose matrix product vectors agree at system
    size $n$. Then there are invertible matrices $Z_v$, one per bond of
    the closed chain ($v = 1, \ldots, n$, $n + 1 \equiv 1$), with
    $B^i = Z_v^{-1}\, A^i\, Z_{v+1}$ for every site $v$ and every $i$.
    This is the per-bond gauge family of the closed-chain corollary, proved
    at the source's strengthened bound $n \ge 2L + 1$
    (\cite{Molnar2018NormalPEPS}, the corollary after Lemma~5 in its
    closing section ``Normal MPS: alternative proof''). The named matrix-level
    Corollary~\ref{cor:fundamentalTheorem_normalMPS} delegates to it.
\end{corollary}
\begin{proof}\leanok
    \uses{cor:fundamentalTheorem_normalMPS_translationInvariant_of_overlap}
    The single-gauge form
    (Corollary~\ref{cor:fundamentalTheorem_normalMPS_translationInvariant_of_overlap})
    gives $Z$ and $\lambda$ with $\lambda^n = 1$ and
    $B^i = \lambda\,Z^{-1} A^i Z$. The family $Z_v := \lambda^{v} Z$
    dresses the gauge with powers of the root of unity: away from the seam
    the consecutive powers contribute the factor $\lambda$, and across the
    seam the wraparound contributes $\lambda^{1-n} = \lambda$ because
    $\lambda^n = 1$, so $B^i = Z_v^{-1} A^i Z_{v+1}$ at every site.
\end{proof}

\subsection{Site-dependent chains}

The source states its Lemma~5 for site-dependent tensors: five sites
carrying tensors $A_1, \ldots, A_5$ whose two-site windows are injective.
The entries above specialize to one repeated tensor; the remaining entries
of this section formalize the site-dependent setting itself.  A closed
chain carries one tensor per site, windows are arcs of consecutive sites,
and the closed-chain corollary delivers one gauge per bond.  Throughout,
all sites share one physical dimension $d$ and all bonds one bond
dimension $D$; the source poses no restriction on the local dimensions of
the site-dependent tensors, and the uniform-dimension restriction is
recorded in the route note
(\path{docs/paper-gaps/peps_normal_ft_section3_route.tex}).

\begin{definition}[Arc products and window injectivity of a site-dependent
        chain]%
    \label{def:chain_arc_window_injective}
    \lean{MPSChainTensor.arcEval}
    \lean{MPSChainTensor.IsWindowInjective}
    \leanok
    \uses{def:non_ti_chain}
    Let a closed chain on $n \ge 1$ sites carry one tensor of $D \times D$
    matrices over physical dimension $d$ per site, $A_0, \dots, A_{n-1}$.
    The \emph{arc product} of a word $x = x_1 \cdots x_\ell$ read from
    site $s$ is
    \[
        A_{(s)}^{x} = A_{s}^{x_1}\, A_{s+1}^{x_2} \cdots
        A_{s+\ell-1}^{x_\ell},
    \]
    site indices cyclic.  The chain is \emph{window injective at length
        $L$} when, for every site $s$, the arc products of the length-$L$
    words read from $s$ span the full matrix algebra.  This is the
    site-dependent setting of the closing section of
    \cite{Molnar2018NormalPEPS} (``Normal MPS: alternative proof''): an
    MPS on five sites in which the blocking of any two consecutive tensors
    is injective, for a general window length $L$.
\end{definition}

\begin{lemma}[Injectivity is window injectivity at length one]%
    \label{lem:isWindowInjective_one_of_isInjective}
    \lean{MPSChainTensor.isWindowInjective_one_of_isInjective}
    \leanok
    \uses{def:chain_arc_window_injective, def:injective}
    If every tensor in a closed chain is injective, then the chain is
    window injective at length \(1\).  Indeed, the length-one arc products
    at site \(s\) are precisely the matrices of the tensor at site \(s\).
\end{lemma}
\begin{proof}\leanok
    The family of length-one words is identified with the physical alphabet.
    Under this identification the length-one arc product is exactly the
    corresponding local matrix, so the two spanning conditions coincide.
\end{proof}

\begin{lemma}[Arc spanning from window injectivity]%
    \label{lem:isWindowInjective_arc_span}
    \lean{MPSChainTensor.IsWindowInjective.arc_span}
    \leanok
    \uses{def:chain_arc_window_injective}
    Let a site-dependent chain on $n \ge 1$ sites be window injective at
    length $L$ with $0 < L$.  Then for every $m \ge L$ and every site $s$,
    the arc products of the length-$m$ words read from $s$ span the full
    matrix algebra.  This is the site-dependent form of the remark that
    any region of at least the blocking size is injective
    (\cite{Molnar2018NormalPEPS}, closing section,
    cf.\ Lemma~\ref{lem:isNBlkInjective_of_le}).
\end{lemma}
\begin{proof}\leanok
    \uses{def:chain_arc_window_injective}
    Write the identity as a combination of window products at the starting
    site and peel the leading letter of each: any matrix becomes a
    combination of terms $A_s^{i}\,(W P)$ with $W$ an arc product of
    length $L - 1$ read from $s + 1$.  Each factor $W P$ is a combination
    of length-$m$ arc products read from $s + 1$ by induction, so each
    term is a combination of length-$(m+1)$ arc products read from $s$.
\end{proof}

\begin{lemma}[Bond-operator extraction from overlapping windows,
        site-dependent form]%
    \label{lem:chain_overlapWindow_exists_bondOperator}
    \lean{MPSChainTensor.overlapWindow_exists_bondOperator}
    \lean{MPSChainTensor.eq_of_span_mul_left}
    \leanok
    \uses{def:chain_arc_window_injective}
    Let $0 < L$ with $2L + 1 \le n$, let the closed chain on $n$ sites
    carry one tensor of $D \times D$ matrices per site, window injective
    at length $L$, and fix a site $r$.  Let $C_0, \ldots, C_L$ be deformed
    window tensors around the bond $(r+L-1,\, r+L)$: each $C_j$ assigns a
    matrix to every word of length $L$, read as replacing the blocked
    tensor of the $L$-site window starting at site $r + j$.  Suppose the
    deformed states agree for consecutive window positions: for every
    $j < L$, every prefix $p$ of $j$ sites read from $r$, every assignment
    $u$ on the $L + 1$ sites covered by the two windows, and every suffix
    $s$ on the remaining sites,
    \[
        \tr\bigl(A_{(r)}^{p}\, C_j(u_1 \cdots u_L)\,
        A_{r+j+L}^{u_{L+1}}\, A_{(r+j+L+1)}^{s}\bigr)
        = \tr\bigl(A_{(r)}^{p}\, A_{r+j}^{u_1}\,
        C_{j+1}(u_2 \cdots u_{L+1})\, A_{(r+j+L+1)}^{s}\bigr).
    \]
    Then the deformation is a virtual operation on the bond: there is a
    matrix $X$ with
    \[
        C_0(a) = A_{(r)}^{a}\, X
        \quad\text{and}\quad
        C_L(b) = X\, A_{(r+L)}^{b}
    \]
    for all length-$L$ words $a$, $b$, and $X$ is unique.  This is
    Lemma~5 of \cite{Molnar2018NormalPEPS} in the site-dependent setting
    in which the source states it, with the deformed window tensors
    arising as physical operators contracted with the blocked windows
    they act on; Lemma~\ref{lem:overlapWindow_exists_bondOperator} is its
    constant specialization.
\end{lemma}
\begin{proof}\leanok
    \uses{lem:isWindowInjective_arc_span, def:chain_arc_window_injective}
    Comparing the windows at $j$ and $j + 1$ through the complementary
    arc of $n - L - 1 \ge L$ sites, whose arc products span by
    Lemma~\ref{lem:isWindowInjective_arc_span}, strips the state equality
    to the letter level:
    $C_j(u_1 \cdots u_L)\,A_{r+j+L}^{u_{L+1}} =
        A_{r+j}^{u_1}\,C_{j+1}(u_2 \cdots u_{L+1})$ on every window of
    $L + 1$ sites.  Chaining the $L$ relations along the $2L$ sites
    flanking the bond turns the leftmost window into the rightmost one:
    $C_0(a)\,A_{(r+L)}^{b} = A_{(r)}^{a}\,C_L(b)$ for all length-$L$
    words.  Writing the identity as a combination
    $\sum_b \alpha_b A_{(r+L)}^{b} = I$, the matrix
    $X := \sum_b \alpha_b\,C_L(b)$ satisfies $A_{(r)}^{a} X = C_0(a)$;
    the mirrored combination of the $C_0$-family collapses onto the same
    matrix, giving $X\,A_{(r+L)}^{b} = C_L(b)$, and stripping the
    spanning factors makes $X$ unique.
\end{proof}

\begin{lemma}[Arc transport between same-state chains]%
    \label{lem:exists_arcTransport}
    \lean{MPSChainTensor.exists_arcTransport}
    \leanok
    \uses{def:chain_arc_window_injective, def:non_ti_chain}
    Let two site-dependent chains $A$, $B$ on $n$ sites generate the same
    closed-chain state, let $k + q = n$, and suppose the $A$-arc products
    of length $k$ read from a site $s$ and of length $q$ read from
    $s + k$ each span the full matrix algebra.  Then there is a linear
    map $\Lambda$ of the matrix algebra with
    $\Lambda(A_{(s)}^{x}) = B_{(s)}^{x}$ for every word $x$ of length
    $k$.
\end{lemma}
\begin{proof}\leanok
    \uses{lem:exists_mpvTransport, def:chain_arc_window_injective}
    As for one repeated tensor (Lemma~\ref{lem:exists_mpvTransport}):
    pair the matrix algebra against the complementary arc products.  The
    pairing is injective on the $A$-side by spanning and trace
    nondegeneracy, and the two pairings match on the spanning arc family
    because both compute closed-chain coefficients of total length $n$,
    equal after rotating the trace of the closed chain to the starting
    site $s$.
\end{proof}

\begin{lemma}[The insertion correspondence is an algebra homomorphism,
        site-dependent form]%
    \label{lem:chain_exists_insertionHom}
    \lean{MPSChainTensor.exists_insertionHom}
    \lean{MPSChainTensor.insertionHom_unique}
    \leanok
    \uses{def:chain_arc_window_injective, def:non_ti_chain}
    Let $0 < L$ with $2L + 1 \le n$, let $A$ and $B$ be site-dependent
    chains on $n$ sites, each window injective at length $L$, generating
    the same closed-chain state, and fix a site $r$.  Then there is a
    linear map $\Phi$ of the matrix algebra, unital and multiplicative,
    such that for every matrix $X$ and every word $w$ of length $n$,
    \[
        \tr\bigl(X\, A_{(r+L)}^{w}\bigr)
        = \tr\bigl(\Phi(X)\, B_{(r+L)}^{w}\bigr),
    \]
    and $\Phi$ is the unique linear map with this pairing property.  This
    is the closing clause of Lemma~5 of \cite{Molnar2018NormalPEPS} — the
    maps $O_1 \mapsto X$ and $O_3^T \mapsto X$ are uniquely defined and
    are algebra homomorphisms — in the site-dependent setting of the
    source, with the algebra structure carried by the insertions: $\Phi$
    sends an insertion on the bond $(r+L-1,\, r+L)$ of the first chain to
    the bond operator extracted on the second.
\end{lemma}
\begin{proof}\leanok
    \uses{lem:chain_overlapWindow_exists_bondOperator,
        lem:exists_arcTransport, lem:isWindowInjective_arc_span}
    A virtual insertion $X$ on the bond is realized on each of the
    $L + 1$ overlapping windows around it by deformed window tensors
    placing $X$ after the first $L - j$ letters of the window starting at
    site $r + j$.  Each deformation is transported to the second chain by
    the arc transport at its own starting site
    (Lemma~\ref{lem:exists_arcTransport}); the same-state hypothesis
    makes the transported deformations generate one common state, and the
    bond-operator extraction
    (Lemma~\ref{lem:chain_overlapWindow_exists_bondOperator}) produces
    $\Phi(X)$.  Uniqueness of the bond operator makes $\Phi$ linear,
    unital and multiplicative; the pairing against all closed-chain words
    follows by expanding the insertion over the window products, and pins
    $\Phi$ through the spanning arcs of length $n$
    (Lemma~\ref{lem:isWindowInjective_arc_span}) and nondegeneracy of the
    trace pairing.
\end{proof}

\begin{lemma}[Conjugate arc products at every bond, site-dependent form]%
    \label{lem:chain_exists_conjugation_of_sameState}
    \lean{MPSChainTensor.exists_conjugation_of_sameState}
    \leanok
    \uses{def:chain_arc_window_injective, def:non_ti_chain}
    Let $0 < L$ with $2L + 1 \le n$, and let $A$ and $B$ be
    site-dependent chains on $n$ sites with a common positive bond dimension,
    each window injective at length
    $L$, generating the same closed-chain state.  Then for every site $p$
    there is an invertible matrix $Z$ with
    \[
        B_{(p)}^{w} = Z\, A_{(p)}^{w}\, Z^{-1}
    \]
    for every word $w$ of length $n$.  The statement for one repeated
    tensor is Lemma~\ref{lem:exists_conjugation_of_mpv_eq}.
\end{lemma}
\begin{proof}\leanok
    \uses{lem:chain_exists_insertionHom, thm:skolem_noether}
    The insertion correspondence at the bond
    (Lemma~\ref{lem:chain_exists_insertionHom}) is unital and nonzero,
    hence an automorphism of the full matrix algebra, and by
    Skolem--Noether (Theorem~\ref{thm:skolem_noether}) it is conjugation
    by an invertible $Z$; pairing insertions against all closed-chain
    words read from $p$ and using nondegeneracy of the trace pairing
    transfers the conjugation to the arc products.
\end{proof}

\begin{corollary}[Fundamental Theorem for site-dependent normal closed
        chains at $n \ge 2L + 1$]%
    \label{cor:fundamentalTheorem_normalMPSChain_of_overlap}
    \lean{TNLean.PEPS.fundamentalTheorem_normalMPSChain_of_overlap}
    \leanok
    \uses{def:chain_arc_window_injective, def:non_ti_chain}
    Let $0 < L$ with $2L + 1 \le n$, and let $A$ and $B$ be
    site-dependent chains on $n$ sites with a common positive bond dimension,
    each window injective at length
    $L$, generating the same closed-chain state.  Then the chains are
    gauge equivalent: there are invertible matrices $Z_v$, one per bond
    of the closed chain, with
    \[
        B_v^i = Z_v\, A_v^i\, Z_{v+1}^{-1}
    \]
    for every site $v$ and every $i$, indices cyclic.  This is the
    closed-chain corollary of the closing section of
    \cite{Molnar2018NormalPEPS} assembled from its site-dependent
    Lemma~5 at every bond; the source displays the corollary for
    translation-invariant tensors
    (Corollary~\ref{cor:fundamentalTheorem_normalMPS_translationInvariant_of_overlap}),
    with the site-dependent Lemma~5 as its engine.
\end{corollary}
\begin{proof}\leanok
    \uses{lem:chain_exists_conjugation_of_sameState,
        lem:chain_overlapWindow_exists_bondOperator,
        lem:isWindowInjective_arc_span}
    The per-bond conjugations
    (Lemma~\ref{lem:chain_exists_conjugation_of_sameState}) give window
    covariance with nonzero scalars: on every arc of $m$ sites with
    $L \le m$ and $m + L \le n$, the conjugation at the near end
    intertwines the arc products of the two chains across the far bond,
    the bond-operator extraction
    (Lemma~\ref{lem:chain_overlapWindow_exists_bondOperator}) produces
    the connecting matrix, and the conjugation at the far end pins it to
    a nonzero scalar multiple of the gauge there through the centralizer
    of the full matrix algebra, so
    $B_{(p)}^{u} = c_{p,m} \cdot Z_p\, A_{(p)}^{u}\, Z_{p+m}^{-1}$ on all
    length-$m$ words $u$.  Comparing the covariance on windows of lengths
    $L + 1$ and $L$ through the spanning window products leaves the
    letterwise relation $B_v^i = \mu_v\, Z_v\, A_v^i\, Z_{v+1}^{-1}$ with
    nonzero scalars $\mu_v$; iterating it once around the closed chain
    against a nonzero arc product forces $\prod_v \mu_v = 1$, and
    dressing the gauges with the partial products of the scalars absorbs
    them, including across the seam.
\end{proof}

\begin{theorem}[Cyclic gauges preserve closed-chain states]%
    \label{thm:mps_chain_gaugeEquiv_sameState}
    \lean{MPSChainTensor.arcEval_eq_gauge_conj}
    \lean{MPSChainTensor.GaugeEquiv.sameState}
    \leanok
    \uses{def:chain_arc_window_injective, def:non_ti_chain}
    Let \(A_0,\ldots,A_{n-1}\) and \(B_0,\ldots,B_{n-1}\) be two
    site-dependent closed chains with a cyclic gauge relation
    \[
        B_v^i = Z_v\, A_v^i\, Z_{v+1}^{-1}
    \]
    at every site.  Then they generate the same closed-chain state.  More
    generally, the product along any arc is conjugated by the two gauges on
    the boundary bonds of that arc; for a full turn around the chain the
    two boundary bonds coincide, and cyclicity of trace removes the
    remaining conjugation.
\end{theorem}
\begin{proof}\leanok
    Iterate the sitewise gauge relation along the arc.  Consecutive
    boundary matrices cancel, leaving only the gauges at the two ends.
    For a word of length \(n\) read from the first site, the two ends are
    the same bond, so the closed trace is invariant under the resulting
    conjugation.
\end{proof}

\begin{corollary}[Fundamental Theorem for injective closed MPS chains]%
    \label{cor:fundamentalTheorem_injectiveMPSChain_of_sameState}
    \lean{TNLean.PEPS.fundamentalTheorem_injectiveMPSChain_of_sameState}
    \leanok
    \uses{cor:fundamentalTheorem_normalMPSChain_of_overlap,
        lem:isWindowInjective_one_of_isInjective}
    Let \(A_0,\ldots,A_{n-1}\) and \(B_0,\ldots,B_{n-1}\) be two
    site-dependent closed chains on \(n \ge 3\) sites, with common physical
    dimension and common positive bond dimension.  Suppose every local tensor in
    both chains is injective and the two closed-chain states are equal.
    Then there are invertible matrices \(Z_0,\ldots,Z_{n-1}\), one per
    bond of the closed chain, such that
    \[
        B_v^i = Z_v\, A_v^i\, Z_{v+1}^{-1}
    \]
    for every site \(v\) and every physical index \(i\), with indices read
    cyclically.  This is the \(L=1\), uniform-dimension instance of
    \cite[Theorem~\textup{thm:inj\_MPS}, lines~688--725]{Molnar2018NormalPEPS}.
\end{corollary}
\begin{proof}\leanok
    \uses{cor:fundamentalTheorem_normalMPSChain_of_overlap,
        lem:isWindowInjective_one_of_isInjective}
    Apply Corollary~\ref{cor:fundamentalTheorem_normalMPSChain_of_overlap}
    with \(L=1\).  The size condition becomes \(n \ge 3\), and
    Lemma~\ref{lem:isWindowInjective_one_of_isInjective} converts the
    sitewise injectivity hypotheses into the required window-injectivity
    hypotheses.
\end{proof}

\begin{definition}[Cyclic shift of a closed MPS chain]%
    \label{def:mps_chain_cyclic_shift}
    \lean{MPSChainTensor.cyclicShift}
    \lean{MPSChainTensor.IsCyclicShiftInvariantState}
    \leanok
    For a closed chain \(A_0,\ldots,A_{n-1}\), its cyclic shift is the
    chain whose \(v\)-th tensor is \(A_{v+1}\), with indices read modulo
    \(n\).  The generated state is cyclic-shift invariant when the
    closed-chain state of \(A_0,\ldots,A_{n-1}\) agrees with that of
    \(A_1,\ldots,A_{n-1},A_0\).
\end{definition}

\begin{corollary}[Gauge comparison with the cyclic shift]%
    \label{cor:fundamentalTheorem_injectiveMPSChain_cyclicShift}
    \lean{TNLean.PEPS.fundamentalTheorem_injectiveMPSChain_cyclicShift}
    \leanok
    \uses{cor:fundamentalTheorem_injectiveMPSChain_of_sameState,
        def:mps_chain_cyclic_shift}
    Let \(A_0,\ldots,A_{n-1}\) be an injective site-dependent closed chain
    with positive bond dimension on \(n\ge 3\) sites.  If its generated state
    is invariant under the
    cyclic shift \(A_v\mapsto A_{v+1}\), then there are invertible matrices
    \(Z_0,\ldots,Z_{n-1}\), one per bond, such that
    \[
        A_{v+1}^i = Z_v\, A_v^i\, Z_{v+1}^{-1}
    \]
    for every site \(v\) and physical index \(i\).  This is the first
    comparison step in the Applications-section translation-invariant
    description corollary of \cite[lines~1807--1824]{Molnar2018NormalPEPS}.
\end{corollary}
\begin{proof}\leanok
    Apply Corollary~\ref{cor:fundamentalTheorem_injectiveMPSChain_of_sameState}
    to the chain and its cyclic shift.  Injectivity is preserved by the
    shift, and the state equality is exactly the cyclic-shift invariance
    hypothesis.
\end{proof}

\begin{corollary}[Telescoping against the first tensor]%
    \label{cor:exists_gauge_to_first_of_cyclicShift_gaugeEquiv}
    \lean{MPSChainTensor.exists_gauge_to_first_of_cyclicShift_gaugeEquiv}
    \leanok
    \uses{def:mps_chain_cyclic_shift}
    Suppose a closed chain \(A_0,\ldots,A_{n-1}\) is gauge equivalent to its
    cyclic shift.  Then for every site \(v\) there are invertible matrices
    \(L_v\) and \(R_v\) such that
    \[
        A_v^i = L_v\, A_0^i\, R_v^{-1}
    \]
    for every physical index \(i\).  This is the telescoping step in
    \cite[lines~1828--1862]{Molnar2018NormalPEPS}, where the products denoted
    \(L_i\) and \(R_i\) express all tensors through the first one.
\end{corollary}
\begin{proof}\leanok
    Write the cyclic-shift gauge relation from one site to the next and
    iterate it along the chain.  At the first site one takes
    \(L_0=R_0=I\); passing from \(v\) to \(v+1\) multiplies the left and
    right virtual factors by the adjacent gauges.
\end{proof}

\begin{corollary}[Gauge-to-first description from a cyclic-invariant state]%
    \label{cor:exists_gauge_to_first_of_cyclicShiftInvariantState}
    \lean{TNLean.PEPS.exists_gauge_to_first_of_cyclicShiftInvariantState}
    \leanok
    \uses{cor:fundamentalTheorem_injectiveMPSChain_cyclicShift,
        cor:exists_gauge_to_first_of_cyclicShift_gaugeEquiv}
    Let \(A_0,\ldots,A_{n-1}\) be an injective site-dependent closed chain
    with positive bond dimension on \(n\ge 3\) sites.  If its generated state
    is invariant under the
    cyclic translation of the sites, then for every site \(v\) there are
    invertible matrices \(L_v\) and \(R_v\) such that
    \[
        A_v^i = L_v\, A_0^i\, R_v^{-1}
    \]
    for every physical index \(i\).  This is the uniform-bond-dimension form
    of the \(L_v,R_v\) step in the Applications-section proof of
    \cite[lines~1803--1890]{Molnar2018NormalPEPS}.
\end{corollary}
\begin{proof}\leanok
    Apply the injective chain theorem to the chain and its cyclic shift, then
    telescope the resulting bond gauges around the chain.
\end{proof}

\begin{corollary}[Translation-invariant description of an injective closed MPS chain,
        uniform bond dimension]%
    \label{cor:exists_constant_injectiveMPS_of_cyclicShiftInvariantState}
    \lean{TNLean.PEPS.exists_constant_injectiveMPS_of_cyclicShiftInvariantState}
    \leanok
    \uses{cor:fundamentalTheorem_injectiveMPSChain_cyclicShift,
        def:mps_chain_cyclic_shift}
    Let \(A_0,\ldots,A_{n-1}\) be an injective site-dependent closed chain on
    \(n\ge 3\) sites, with one physical dimension \(d\) and one positive bond
    dimension \(D\).  If the closed-chain state is invariant under the cyclic shift
    \(A_v\mapsto A_{v+1}\), then there is an injective tensor \(B\), with bond
    dimension \(D\), such that for every physical configuration \(\sigma\),
    \[
        \tr\!\left(
            A_0^{\sigma_0}\cdots A_{n-1}^{\sigma_{n-1}}\right)
        =
        \tr\!\left(
            B^{\sigma_0}\cdots B^{\sigma_{n-1}}\right).
    \]

    This is the uniform-bond-dimension target corresponding to the
    Applications-section corollary of
    \cite[lines~1803--1890]{Molnar2018NormalPEPS}.  The source also concludes
    that all bond dimensions of the original site-dependent representation are
    equal; in the present uniform-dimension statement that clause has already
    been built into the setting.
\end{corollary}
\begin{proof}\leanok
    Apply Corollary~\ref{cor:fundamentalTheorem_injectiveMPSChain_cyclicShift}
    to compare the chain with its cyclic shift, giving one invertible gauge
    $Z_v$ per bond with $A_{v+1}^i=Z_v\,A_v^i\,Z_{v+1}^{-1}$.  Telescoping
    the gauges into the running products $P_m=Z_{m-1}\cdots Z_0$ expresses
    every tensor through the first, $A_m^i=P_m\,(A_0^i Z_0)\,P_{m+1}^{-1}$, so
    the closed trace collapses to
    $\tr(B^{\sigma_0}\cdots B^{\sigma_{n-1}} P_n^{-1})$ with the
    tentative repeated tensor $B=A_0 Z_0$.  Reading the gauge relation across
    the seam $A_{n+1}\equiv A_0$ shows that the loop product $P_n$ commutes
    with the full matrix algebra spanned by $B$, hence is a scalar
    $\lambda\ne 0$.  Since $\mathbb{C}$ is algebraically closed there is an
    $n$-th root $c^n=\lambda^{-1}$; the rescaled tensor $c\,B$ is still
    injective and its constant chain generates the same closed state.
\end{proof}

\begin{corollary}[Translation-invariant description of an injective 2D PEPS,
        uniform bond dimension]%
    \label{cor:exists_constant_injectivePEPS_of_translationInvariantState}
    \notready
    \uses{cor:exists_constant_injectiveMPS_of_cyclicShiftInvariantState,
        thm:fundamentalTheorem_PEPS, def:peps_torus_translation_invariant,
        def:peps_torus_uniform_bond_dim,
        thm:peps_stateCoeff_translationInvariant}
    Let \(A\) be an injective PEPS on the discrete torus with orientation-uniform
    bond dimensions, horizontal dimension \(D_h\) and vertical dimension \(D_v\),
    and positive bond dimensions.  If the generated state is invariant under all
    lattice translations \(\tau_{a,b}\), that is, for every translation and every
    physical configuration \(\sigma\),
    \[
        \Coeff_{A}(\sigma\circ\tau_{a,b})
        =
        \Coeff_{A}(\sigma),
    \]
    then there is an injective torus tensor \(B\), with the same orientation-uniform
    bond dimensions \(D_h\) and \(D_v\), whose constant (site-independent) family
    generates the same state: for every physical configuration \(\sigma\),
    \[
        \Coeff_{A}(\sigma)
        =
        \Coeff_{B}(\sigma).
    \]

    This is the uniform-bond-dimension target corresponding to the
    Applications-section 2D corollary of
    \cite[lines~1892--1894]{Molnar2018NormalPEPS}.  The source additionally
    concludes that all vertical bond dimensions agree and all horizontal bond
    dimensions agree; in the present orientation-uniform setting those two
    clauses have already been built into the bond-dimension type, so they carry
    no content here.
\end{corollary}
\begin{proof}
    The source states this corollary without proof, remarking only that the
    one-dimensional argument extends; the route below is that intended
    extension.  Write \(A^{(j,k)}\) for the site tensor at lattice position
    \((j,k)\), with horizontal virtual legs of dimension \(D_h\) and vertical
    virtual legs of dimension \(D_v\).  Comparing the state with its
    \(\tau_{1,0}\)-translate through the injective PEPS Fundamental Theorem
    (Theorem~\ref{thm:fundamentalTheorem_PEPS}) produces, on each horizontal
    edge, invertible \(D_h\times D_h\) gauges \(X^{(j,k)}_h\) with
    \[
        A^{(j+1,k)} = \bigl(X^{(j,k)}_h\bigr)^{-1}\, A^{(j,k)}\, X^{(j+1,k)}_h,
    \]
    acting on the two horizontal virtual legs, and comparing with the
    \(\tau_{0,1}\)-translate produces, on each vertical edge, invertible
    \(D_v\times D_v\) gauges \(X^{(j,k)}_v\) with
    \[
        A^{(j,k+1)} = \bigl(X^{(j,k)}_v\bigr)^{-1}\, A^{(j,k)}\, X^{(j,k+1)}_v.
    \]
    Telescoping the horizontal gauges along a row reduces each site to the
    leftmost one,
    \[
        A^{(j,k)} = \bigl(L^{(j,k)}_h\bigr)^{-1}\, A^{(0,k)}\, R^{(j,k)}_h,
        \qquad
        L^{(j,k)}_h = X^{(0,k)}_h\cdots X^{(j-1,k)}_h,
        \quad
        R^{(j,k)}_h = X^{(1,k)}_h\cdots X^{(j,k)}_h,
    \]
    and the periodic boundary condition closes the row product into a single
    residual virtual operation, the 2D analogue of the one-dimensional identity
    \(R_vL_{v+1}^{-1}=Z_0^{-1}\):
    \[
        R^{(j,k)}_h\,\bigl(L^{(j+1,k)}_h\bigr)^{-1} = \bigl(X^{(0,k)}_h\bigr)^{-1}.
    \]
    The same telescoping closes each column, yielding one residual vertical
    operation.  Absorbing the two residual gauges into the corner site gives a
    single repeated tensor
    \[
        B = \bigl(X^{(0,0)}_v\bigr)^{-1}\,\bigl(X^{(0,0)}_h\bigr)^{-1}\, A^{(0,0)},
    \]
    whose constant family generates the same state.  The one-dimensional
    telescoping closure is stated as
    Corollary~\ref{cor:exists_constant_injectiveMPS_of_cyclicShiftInvariantState}.
    What remains is the two-directional PEPS assembly.  If
    \[
        H_{j,k}=R^{(j,k)}_h\bigl(L^{(j+1,k)}_h\bigr)^{-1},
        \qquad
        V_{j,k}=R^{(j,k)}_v\bigl(L^{(j,k+1)}_v\bigr)^{-1},
    \]
    are the residual horizontal and vertical operations, the horizontal residual
    acts on \(\mathbb C^{D_h}\) and the vertical residual acts on
    \(\mathbb C^{D_v}\).  The missing plaquette flatness identity must therefore
    be stated after placing them on the corner space
    \(\mathbb C^{D_h}\otimes\mathbb C^{D_v}\).  Equivalently, after the row and
    column telescopings make \(H_{j,k}=H_k\) and \(V_{j,k}=V_j\), one must prove
    \[
        (H_k\otimes I_{D_v})(I_{D_h}\otimes V_{j+1})
        =
        \lambda_{j,k}\,(I_{D_h}\otimes V_j)(H_{k+1}\otimes I_{D_v})
    \]
    for nonzero scalars \(\lambda_{j,k}\).  The scalar phases must also close
    around the two periodic seams,
    \[
        \prod_{k\in\mathbb Z/N_v\mathbb Z}\lambda_{j,k}=1,
        \qquad
        \prod_{j\in\mathbb Z/N_h\mathbb Z}\lambda_{j,k}=1.
    \]
    These identities are then absorbed into this single repeated 2D tensor.
\end{proof}

\begin{theorem}[Uniqueness of the injective closed-chain gauge]%
    \label{thm:fundamentalTheorem_injectiveMPSChain_gauge_unique}
    \lean{TNLean.PEPS.fundamentalTheorem_injectiveMPSChain_gauge_unique}
    \leanok
    \uses{cor:fundamentalTheorem_injectiveMPSChain_of_sameState,
        def:chain_arc_window_injective}
    Under the positive-bond-dimension hypotheses of
    Corollary~\ref{cor:fundamentalTheorem_injectiveMPSChain_of_sameState},
    suppose \(Z_v\) and \(Z'_v\) are two families of invertible matrices
    satisfying
    \[
        B_v^i = Z_v\, A_v^i\, Z_{v+1}^{-1}
        \quad\text{and}\quad
        B_v^i = Z'_v\, A_v^i\, (Z'_{v+1})^{-1}
    \]
    for every \(v,i\).  Then there is a single nonzero scalar \(c\) such
    that \(Z'_v = c Z_v\) for every bond \(v\).  This is the uniqueness
    clause in \cite[line~724]{Molnar2018NormalPEPS}.
\end{theorem}
\begin{proof}\leanok
    For each bond put \(C_v = Z_v^{-1}Z'_v\).  Comparing the two gauge
    relations gives
    \[
        C_v A_v^i = A_v^i C_{v+1}
    \]
    for every local matrix.  Since the matrices of \(A_v\) span the full
    matrix algebra, evaluating this identity on the identity matrix gives
    \(C_v=C_{v+1}\); hence all \(C_v\) are one common matrix.  The same
    relation then says that this common matrix commutes with the full
    matrix algebra, so it is a scalar matrix.  Its invertibility makes the
    scalar nonzero.
\end{proof}

Two remarks place these statements against \cite{Molnar2018NormalPEPS}.  The
injective matrix product states of the opening section of
\cite{Molnar2018NormalPEPS} are the case $L = 1$, in which a single tensor
already spans the matrix algebra; the corollaries above contain them at
$n \ge 3$, with no separate argument needed.  The site-dependent results are
stated for one physical dimension $d$ and one bond dimension $D$ around the
chain, while \cite{Molnar2018NormalPEPS} places no restriction on the local
dimensions; this uniform-dimension restriction is the only narrowing.
% Paper-gap note: docs/paper-gaps/peps_normal_ft_section3_route.tex

\section{The vertex-injective case and the residual edge scalars}

\begin{theorem}[Fundamental Theorem for normal PEPS, vertex-injective case]%
    \label{thm:fundamentalTheorem_normalPEPS_of_vertexInjective}
    \lean{TNLean.PEPS.fundamentalTheorem_normalPEPS_of_vertexInjective}
    \leanok
    \uses{thm:fundamentalTheorem_PEPS, def:peps_normalPEPSBlockingHypotheses,
        def:GaugeEquivPEPS, def:IsVertexInjective}
    Suppose PEPS tensors $A$ and $B$ are vertex injective with positive bond
    dimensions on a connected graph, generate the same state, and $A$ satisfies
    the normal blocking hypotheses. Then their defining tensors are related by
    a local gauge.
    This is the vertex-injective specialization of the general normal-PEPS
    theorem: when each single tensor is already injective the gauge follows
    from the injective Fundamental Theorem, so the region blocking hypotheses
    are not used. The general region-injective theorem, in which only blocked
    regions are injective, remains open; the missing step is the
    region-level insertion-algebra isomorphism that produces the per-edge gauge
    after blocking.
\end{theorem}
\begin{proof}\leanok
    Vertex injectivity, positive bond dimensions, same state, and connectivity
    are exactly the hypotheses of the injective Fundamental Theorem, which
    supplies the local gauge.
\end{proof}

\begin{definition}[Oriented endpoint scalar]%
    \label{def:edgeScalarAt}
    \lean{TNLean.PEPS.edgeScalarAt}
    \leanok
    \uses{def:edgeGaugeAt}
    An edge-scalar family $c=(c_e)_{e\in E}$ determines the endpoint scalar
    $c_{v,e}$ by $c_{v,e}=c_e$ at the lower endpoint of $e$ and
    $c_{v,e}=c_e^{-1}$ at the upper endpoint.
\end{definition}

\begin{definition}[Vertex-balanced edge scalars]%
    \label{def:IsVertexBalanced}
    \lean{TNLean.PEPS.IsVertexBalanced}
    \leanok
    \uses{def:edgeScalarAt}
    An edge-scalar family is vertex-balanced when
    $\prod_{e\ni v} c_{v,e}=1$ for every vertex $v$.
\end{definition}


\section{Balanced Edge Scalars}\label{sec:peps_balanced_edge_scalars}

The edge gauges are determined only up to a scalar on each edge. The residual
freedom is the balanced edge scalars: an assignment of a nonzero scalar to every
oriented edge whose product around each vertex is one.
\begin{center}
    \TNPEPSVertexScalarBalance
\end{center}
Two gauge families related by such an assignment give the same gauge
equivalence, which is the precise sense in which the gauges of the Fundamental
Theorem are unique up to a constant.

\begin{definition}[Gauge equivalence modulo balanced edge scalars]%
    \label{def:GaugeEquivModEdgeScalars}
    \lean{TNLean.PEPS.GaugeEquivModEdgeScalars}
    \leanok
    \uses{def:edgeGaugeAt, def:edgeScalarAt, def:IsVertexBalanced}
    Two gauge families are equivalent modulo balanced edge scalars if there is
    a nonzero scalar $c_e$ on each edge such that the corresponding endpoint
    actions satisfy $X_{v,e}=c_{v,e}Y_{v,e}$, and
    $\prod_{e\ni v}c_{v,e}=1$ for every vertex $v$.
\end{definition}

\begin{theorem}[Multiplication of oriented edge scalars]%
    \label{thm:edgeScalarAt_mul}
    \lean{TNLean.PEPS.edgeScalarAt_mul}
    \leanok
    \uses{def:edgeScalarAt}
    If two edge-scalar families are multiplied, then their endpoint scalars
    satisfy $(cd)_{v,e}=c_{v,e}d_{v,e}$ on every oriented half-edge.
    \begin{center}
        \TNPEPSEdgeGaugeOrientation
    \end{center}
\end{theorem}

\begin{proof}\leanok
    At the lower endpoint this is ordinary multiplication of the edge scalars.
    At the upper endpoint, $(c_ed_e)^{-1}=c_e^{-1}d_e^{-1}$ because the
    edge scalars lie in $\mathbb C^\times$.
\end{proof}

\begin{theorem}[Inversion of oriented edge scalars]%
    \label{thm:edgeScalarAt_inv}
    \lean{TNLean.PEPS.edgeScalarAt_inv}
    \leanok
    \uses{def:edgeScalarAt}
    Inverting an edge-scalar family gives
    $(c^{-1})_{v,e}=c_{v,e}^{-1}$ on every oriented endpoint.
\end{theorem}

\begin{proof}\leanok
    At the lower endpoint, $(c^{-1})_{v,e}=c_e^{-1}=c_{v,e}^{-1}$.
    At the upper endpoint,
    $(c^{-1})_{v,e}=(c_e^{-1})^{-1}=c_{v,e}^{-1}$.
\end{proof}

\begin{theorem}[Nonzero oriented edge scalars]%
    \label{thm:edgeScalarAt_ne_zero}
    \lean{TNLean.PEPS.edgeScalarAt_ne_zero}
    \leanok
    \uses{def:edgeScalarAt}
    Every oriented endpoint scalar satisfies $c_{v,e}\ne 0$.
\end{theorem}

\begin{proof}\leanok
    Edge scalars are units, and the inverse of a unit is again a unit.
\end{proof}

\begin{theorem}[Product of balanced edge scalars]%
    \label{thm:IsVertexBalanced_mul}
    \lean{TNLean.PEPS.IsVertexBalanced.mul}
    \leanok
    \uses{def:IsVertexBalanced, thm:edgeScalarAt_mul}
    The pointwise product of two vertex-balanced edge-scalar families is again
    vertex-balanced.
\end{theorem}

\begin{proof}\leanok
    For every vertex $v$,
    $\prod_{e\ni v}(cd)_{v,e}=
        (\prod_{e\ni v}c_{v,e})(\prod_{e\ni v}d_{v,e})=1\cdot 1=1$.
\end{proof}

\begin{theorem}[Inverse of balanced edge scalars]%
    \label{thm:IsVertexBalanced_inv}
    \lean{TNLean.PEPS.IsVertexBalanced.inv}
    \leanok
    \uses{def:IsVertexBalanced, thm:edgeScalarAt_inv}
    The pointwise inverse of a vertex-balanced edge-scalar family is again
    vertex-balanced.
\end{theorem}

\begin{proof}\leanok
    For every vertex $v$,
    $\prod_{e\ni v}(c^{-1})_{v,e}=
        (\prod_{e\ni v}c_{v,e})^{-1}=1^{-1}=1$.
\end{proof}

\begin{theorem}[Reflexivity of balanced edge-scalar equivalence]%
    \label{thm:GaugeEquivModEdgeScalars_refl}
    \lean{TNLean.PEPS.GaugeEquivModEdgeScalars.refl}
    \leanok
    \uses{def:GaugeEquivModEdgeScalars}
    Every edge-gauge family is equivalent to itself modulo balanced edge
    scalars, witnessed by $c_e=1$ for every edge $e$.
    \begin{center}
        \TNPEPSEdgeGaugeOrientation
    \end{center}
\end{theorem}

\begin{proof}\leanok
    If $c_e=1$, then $c_{v,e}=1$, $X_{v,e}=c_{v,e}X_{v,e}$, and
    $\prod_{e\ni v}c_{v,e}=1$.
\end{proof}

\begin{theorem}[Symmetry of balanced edge-scalar equivalence]%
    \label{thm:GaugeEquivModEdgeScalars_symm}
    \lean{TNLean.PEPS.GaugeEquivModEdgeScalars.symm}
    \leanok
    \uses{def:GaugeEquivModEdgeScalars, thm:IsVertexBalanced_inv,
        thm:edgeScalarAt_ne_zero}
    If one gauge family differs from another by balanced edge scalars, then the
    second differs from the first by the inverse balanced edge-scalar family.
    \begin{center}
        \TNPEPSEdgeGaugeOrientation
    \end{center}
\end{theorem}

\begin{proof}\leanok
    If $X_{v,e}=c_{v,e}Y_{v,e}$ with $c_{v,e}\ne 0$, then
    $Y_{v,e}=c_{v,e}^{-1}X_{v,e}$. The inverse scalar family is balanced by
    Theorem~\ref{thm:IsVertexBalanced_inv}.
\end{proof}

\begin{theorem}[Transitivity of balanced edge-scalar equivalence]%
    \label{thm:GaugeEquivModEdgeScalars_trans}
    \lean{TNLean.PEPS.GaugeEquivModEdgeScalars.trans}
    \leanok
    \uses{def:GaugeEquivModEdgeScalars, thm:IsVertexBalanced_mul,
        thm:edgeScalarAt_mul}
    If one gauge family differs from a second by balanced edge scalars, and the
    second differs from a third by balanced edge scalars, then the first differs
    from the third by the pointwise product of the two scalar families.
    \begin{center}
        \TNPEPSEdgeGaugeOrientation
    \end{center}
\end{theorem}

\begin{proof}\leanok
    From $X_{v,e}=c_{v,e}Y_{v,e}$ and $Y_{v,e}=d_{v,e}Z_{v,e}$ one obtains
    $X_{v,e}=(c_{v,e}d_{v,e})Z_{v,e}$. The product scalar family is balanced
    by Theorem~\ref{thm:IsVertexBalanced_mul}.
\end{proof}

\begin{theorem}[The balanced edge-scalar quotient is an equivalence relation]%
    \label{thm:GaugeEquivModEdgeScalars_equivalence}
    \lean{TNLean.PEPS.GaugeEquivModEdgeScalars.equivalence}
    \leanok
    \uses{thm:GaugeEquivModEdgeScalars_refl,
        thm:GaugeEquivModEdgeScalars_symm,
        thm:GaugeEquivModEdgeScalars_trans}
    Gauge equivalence modulo balanced edge scalars is reflexive, symmetric,
    and transitive.
\end{theorem}

\begin{proof}\leanok
    Reflexivity is witnessed by $1$, symmetry by $c^{-1}$, and transitivity
    by $cd$.
\end{proof}

\begin{definition}[The balanced edge-scalar quotient relation]%
    \label{def:GaugeEquivModEdgeScalars_setoid}
    \lean{TNLean.PEPS.GaugeEquivModEdgeScalars.setoid}
    \leanok
    \uses{def:GaugeEquivModEdgeScalars,
        thm:GaugeEquivModEdgeScalars_equivalence}
    The preceding equivalence relation is the quotient relation on edge-gauge
    families modulo vertex-balanced edge scalars.
\end{definition}

\begin{theorem}[Balanced edge scalars preserve the local gauge action]%
    \label{thm:GaugeEquivModEdgeScalars_gaugeVertex_eq}
    \lean{TNLean.PEPS.GaugeEquivModEdgeScalars.gaugeVertex_eq}
    \leanok
    \uses{def:GaugeEquivModEdgeScalars, def:gaugeVertex}
    If two gauge families differ by a balanced edge-scalar reweighting, then
    they produce the same gauged tensor at every vertex.
\end{theorem}

\begin{proof}\leanok\uses{def:GaugeEquivModEdgeScalars, def:gaugeVertex}
    In each summand of the gauged tensor at $v$, the edge scalars contribute
    the factor $\prod_{e\ni v}c_{v,e}=1$. Hence the summand, and therefore
    the whole local tensor, is unchanged.
\end{proof}

\begin{theorem}[Balanced edge scalars preserve the gauged tensor]%
    \label{thm:GaugeEquivModEdgeScalars_applyGauge_eq}
    \lean{TNLean.PEPS.GaugeEquivModEdgeScalars.applyGauge_eq}
    \leanok
    \uses{thm:GaugeEquivModEdgeScalars_gaugeVertex_eq, def:applyGauge}
    If two gauge families differ by a balanced edge-scalar reweighting, then
    applying them to the same PEPS tensor produces the same gauged tensor.
    \begin{center}
        \TNPEPSGaugeVertexAction
    \end{center}
\end{theorem}

\begin{proof}\leanok
    For every vertex $v$, virtual index $\eta$, and physical index $\sigma$,
    $(A^X)_v(\eta,\sigma)=(A^Y)_v(\eta,\sigma)$ for the two gauged tensors by
    Theorem~\ref{thm:GaugeEquivModEdgeScalars_gaugeVertex_eq}. Hence the
    two gauged tensors are equal.
\end{proof}

\begin{theorem}[Balanced edge scalars preserve the gauged state]%
    \label{thm:GaugeEquivModEdgeScalars_applyGauge_sameState}
    \lean{TNLean.PEPS.GaugeEquivModEdgeScalars.applyGauge_sameState}
    \leanok
    \uses{thm:GaugeEquivModEdgeScalars_applyGauge_eq, def:peps_same_state}
    If two gauge families differ by a balanced edge-scalar reweighting, then
    the PEPS tensors obtained by applying those gauge families represent the
    same state.
\end{theorem}

\begin{proof}\leanok
    The two gauged tensors are equal, hence all of their state coefficients are
    equal.
\end{proof}

\begin{theorem}[Gauge uniqueness modulo balanced edge scalars]%
    \label{thm:gauge_unique_mod_edge_scalars}
    \lean{TNLean.PEPS.gauge_unique_mod_edge_scalars}
    \leanok
    \uses{def:GaugeEquivModEdgeScalars, def:IsVertexInjective,
        thm:peps_piProductFormsScalar}
    Let every bond space be nonzero. If two gauge families relate the same
    vertex-injective PEPS tensor to the same target tensor, then their edge
    gauges differ by nonzero scalars $c_e$ whose oriented product at every
    vertex is $1$. Thus the correct graph quotient is by balanced edge scalars.
    % Correction documented in docs/paper-gaps/peps_gauge_edge_scalars.tex.
    % The nonzero-bond restriction is the standing normal-PEPS convention,
    % matching the existence direction; see the same paper-gap note.
\end{theorem}

\begin{proof}\leanok
    Combining the two gauge relations gives, at every vertex, equality of the
    gauged local tensors. Linear independence of the local tensor family
    promotes this to equality of the products of incident edge-gauge entries
    over every pair of virtual configurations. Each oriented endpoint gauge is
    invertible and, since every bond space is nonzero, indexed by a nonempty
    type; the many-leg scalar extraction then yields, at each vertex, nonzero
    proportionality constants between the two families whose product is $1$.
    Reading the lower-endpoint constant of each edge defines a single edge
    scalar $c_e$; the inverse relation transports it to the upper endpoint, and
    invertibility makes the proportionality constant unique, so the vertex
    products of the per-vertex constants identify with the oriented products of
    $c_e$. These are all $1$, so $c$ is vertex-balanced.
\end{proof}

\begin{remark}
    The local left inverse $\Psi_v$, the realization map $T \mapsto O_T$,
    the split at one incident edge, and the partition
    $V = \{u\} \sqcup (V \setminus \{u,v\}) \sqcup \{v\}$ isolate the local
    indices and maps needed for the edge-centered reduction of
    \cite[Section~3]{Molnar2018NormalPEPS}. The middle-region blocked coefficient is
    constructed by the preceding definitions, and the local-gauge step is the
    comparison with the corresponding three-site injective chain. The corrected
    uniqueness statement is a quotient by balanced edge scalars.
\end{remark}

